% Editable master: glossaries/english.md
% Complete export; generated by scripts/build_glossary_exports.py.
\newcommand{\runningtitle}{English–Old English}
% Layout inherited from the reviewed previous course; set a running title per unit.
% This is a small hand-maintained TeX layout, not a conversion pipeline.
\RequirePackage{fix-cm}
\documentclass[11pt,twoside]{article}
\usepackage{fontspec}
\setmainfont{Cambria}[
  Renderer=HarfBuzz,
  BoldFont=cambriab.ttf,
  ItalicFont=cambriai.ttf,
  BoldItalicFont=cambriaz.ttf
]
\setmonofont{lmmono10-regular.otf}[BoldFont=lmmonolt10-bold.otf]
\usepackage[paperwidth=165mm,paperheight=241mm,inner=21mm,outer=19mm,
  top=22mm,bottom=20mm,headheight=14pt,headsep=7mm,footskip=10mm]{geometry}
\usepackage{xcolor,array,booktabs,longtable,tabularx}
\usepackage{enumitem,fancyhdr}
\usepackage[unicode,hidelinks]{hyperref}
\definecolor{ink}{HTML}{000000}
\definecolor{muted}{HTML}{4C5558}
\setlength{\parindent}{0pt}
\setlength{\parskip}{3.5pt}
\AtBeginDocument{\fontsize{11}{14.3}\selectfont}
\setlength{\tabcolsep}{4pt}
\renewcommand{\arraystretch}{1.10}
\setlist[itemize]{leftmargin=1.3em,label=\textbullet,itemsep=2pt,topsep=3pt,parsep=0pt}
\setlist[enumerate]{leftmargin=1.8em,itemsep=3pt,topsep=4pt,parsep=0pt}
\widowpenalty=10000
\clubpenalty=10000
\emergencystretch=1em
\raggedbottom
\pagestyle{fancy}
\fancyhf{}
\providecommand{\runningtitle}{Old English}
\fancyhead[LE,RO]{\small\itshape \runningtitle}
\fancyhead[LO,RE]{\small Old English}
\fancyfoot[LE,RO]{\small\thepage}
\renewcommand{\headrulewidth}{0.3pt}
\newcommand{\booktitle}[2]{
  {\fontsize{9}{11}\selectfont\color{muted}#1\par}
  \vspace{3pt}
  {\fontsize{20}{24}\selectfont\color{ink}#2\par}
  \vspace{6pt}
}
\newcommand{\source}[1]{\par{\footnotesize\color{muted}#1\par}}
\newcommand{\readingtext}[1]{
  \par\addvspace{5pt}
  \begingroup\leftskip=1.5em
    \fontsize{12}{16}\selectfont #1\par
  \endgroup\addvspace{5pt}
}
\makeatletter
\renewcommand{\section}{\@startsection{section}{1}{\z@}
  {-2.4ex \@plus -0.6ex \@minus -0.2ex}{0.8ex \@plus 0.2ex}
  {\normalfont\fontsize{12}{15}\selectfont\bfseries\color{ink}}}
\renewcommand{\subsection}{\@startsection{subsection}{2}{\z@}
  {-1.8ex \@plus -0.5ex \@minus -0.2ex}{0.5ex \@plus 0.2ex}
  {\normalfont\fontsize{11}{14}\selectfont\bfseries\color{ink}}}
\makeatother

\usepackage{multicol}
\raggedcolumns
\newcommand{\needspace}[1]{\par\begingroup\skip0=\lastskip\vskip-\skip0\dimen0=#1\relax\advance\dimen0 by\skip0\vskip0pt plus\dimen0\penalty-100\vskip0pt plus-\dimen0\vskip\dimen0\penalty9999\vskip-\dimen0\vskip\skip0\endgroup}
\setlength{\columnsep}{6mm}
\hyphenpenalty=10000\exhyphenpenalty=10000
\emergencystretch=2em
\begin{document}
\addfontfeatures{Ligatures=TeXOff}
\raggedright
\pdfbookmark[0]{English–Old English exercise glossary}{outline-B00001}
% glossary-block:B00001:begin
\booktitle{GLOSSARY}{English–Old English exercise glossary}
% glossary-block:B00001:end
% glossary-block:B00002:begin
This glossary supplies English access to the vocabulary and expressions needed for the exercises in all eighteen units: Cook's items, original practice and reading work. The pronunciation practice supplies its own word meanings. English used only to give an instruction, such as “parse” or “give the genitive,” is not vocabulary to translate.\par
% glossary-block:B00002:end
% glossary-block:B00003:begin
Look up the intended meaning or the whole phrase. Choices under one English word are not always interchangeable. Read the qualifying meaning and case or construction note, then use the forms in the companion Old English–English glossary and the relevant unit. Exercise targets and supplied wording still apply; this reference does not turn supplied material into an independent-recall requirement.\par
% glossary-block:B00003:end
% glossary-block:B00004:begin
English articles, auxiliaries and dummy subjects often require a construction rather than a matching Old English word. Entries such as \textbf{a}, \textbf{do}, \textbf{will}, \textbf{there is}, and phrase entries explain these differences. A case label describes the Old English dependent: “genitive” requires a genitive form, not an extra word meaning “of.”\par
% glossary-block:B00004:end
% glossary-block:B00005:begin
Old English citation forms retain the book's spelling and length marks. Some source expressions are supplied only in their stated context. Inflected or conventional citation labels are identified where a normal infinitive or full paradigm is not established. Use the separate key to check a completed answer; this glossary is not a replacement key.\par
% glossary-block:B00005:end
\pdfbookmark[0]{Entries}{outline-B00006}
\clearpage
\fontsize{10.5}{13.2}\selectfont
\setlength{\parskip}{4pt}
\begin{multicols}{2}
% glossary-block:B00006:begin
\needspace{5\baselineskip}
\section*{Entries}
% glossary-block:B00006:end
\pdfbookmark[1]{A}{outline-B00007}
% glossary-block:B00007:begin
\needspace{5\baselineskip}
\subsection*{A}
% glossary-block:B00007:end
% glossary-block:B00008:begin
\textbf{a} — Ordinary English a/an usually has no separate Old English article. Ān is used in its taught a-certain-individual sense, not inserted into every indefinite noun phrase. No obligatory indefinite article: bare predicates/nouns often translate a/an. ān is used for the taught introducing-a-certain-individual sense, not inserted everywhere.\par
% glossary-block:B00008:end
% glossary-block:B00009:begin
\needspace{3\baselineskip}
\textbf{abandon} — \textbf{forlǣtan} (v.). abandon; forsake. Accusative person or thing left. Accusative thing/place left. Distinguish leave to someone and take leave of.\par
% glossary-block:B00009:end
% glossary-block:B00010:begin
\needspace{3\baselineskip}
\textbf{ability} — \textbf{miht} (n.). ability; power.\par
% glossary-block:B00010:end
% glossary-block:B00011:begin
\needspace{3\baselineskip}
\textbf{abode} — \textbf{wīc} (n.). dwelling; abode.\par
% glossary-block:B00011:end
% glossary-block:B00012:begin
\needspace{3\baselineskip}
\textbf{abounding in} — \textbf{welig} (adj.). rich in; abounding in. on + dative thing abounded in.\par
% glossary-block:B00012:end
% glossary-block:B00013:begin
\needspace{3\baselineskip}
\textbf{about} — \textbf{ābūtan} (prep.). around; about. With accusative. \textbf{be} (prep.). concerning; about. With dative. \textbf{ofer} (prep.). over; about. Object of rejoicing, with the assigned accusative frame. \textbf{swā-swā} (conj.). about; approximately. Approximation in the supplied source expression. \textbf{swilce} (adv.). about; approximately. \textbf{ymbe} (prep.). Approximate time, with accusative.\par
% glossary-block:B00013:end
% glossary-block:B00014:begin
\needspace{3\baselineskip}
\textbf{above} — \textbf{bufan} (prep.). With dative landmark in the taught spatial frame. \textbf{ofer} (prep.). over; across. Spatial frame; case varies with the specific construction. Do not impose a universal motion rule. Rank, with accusative. Rank, with explicitly supplied dative complement; retain this limit. above; more than. Preference in the supplied comparison. Spatial ofer can mean over/across; its case depends on the assigned spatial frame. Ofer in the authority sense over takes an accusative dependent. Rank and preference have further separately qualified senses.\par
% glossary-block:B00014:end
% glossary-block:B00015:begin
\needspace{3\baselineskip}
\textbf{above all} — \textbf{tōforan} (prep.). before; above all.\par
% glossary-block:B00015:end
% glossary-block:B00016:begin
\needspace{3\baselineskip}
\textbf{Abraham} — \textbf{Abraham} (proper name). Biblical patriarch; supplied name forms.\par
% glossary-block:B00016:end
% glossary-block:B00017:begin
\needspace{3\baselineskip}
\textbf{abundance} — \textbf{unrīm} (n.). abundance; uncounted quantity. Supplied quantity with genitive of the substance.\par
% glossary-block:B00017:end
% glossary-block:B00018:begin
\needspace{3\baselineskip}
\textbf{abyss} — \textbf{nēowolnes} (n.). abyss; the deep.\par
% glossary-block:B00018:end
% glossary-block:B00019:begin
\needspace{3\baselineskip}
\textbf{accept} — \textbf{onfōn} (v.). receive; accept. Dative thing received in this selected frame; the later accusative payment frame is separate. receive; accept. A person received into care in the supplied frame.\par
% glossary-block:B00019:end
% glossary-block:B00020:begin
\needspace{3\baselineskip}
\textbf{acceptable} — \textbf{gecwēme} (adj.). pleasing; acceptable.\par
% glossary-block:B00020:end
% glossary-block:B00021:begin
\needspace{3\baselineskip}
\textbf{accompanying} — \textbf{mid} (prep.; adv. in fixed expressions). with; accompanying. Dative person or cargo in the ordinary accompaniment frame.\par
% glossary-block:B00021:end
% glossary-block:B00022:begin
\needspace{3\baselineskip}
\textbf{accomplish} — \textbf{gefręmman} (v.). perform; accomplish. A work is the thing accomplished. \textbf{gefyllan} (v.). finish; accomplish. A work is the thing completed; distinct from fill. Gefyllan means finish/accomplish when a work is completed; the fill sense is separate.\par
% glossary-block:B00022:end
% glossary-block:B00023:begin
\needspace{3\baselineskip}
\textbf{accomplishment} — \textbf{cræft} (n.). skill; accomplishment.\par
% glossary-block:B00023:end
% glossary-block:B00024:begin
\needspace{3\baselineskip}
\textbf{accordingly} — \textbf{eornostlice} (adv.). accordingly; thus. \textbf{for-ðon-ðe} (conj.; adv. in for-ðon), \textbf{eornostlice} (adv.). The supplied consequence use of for-ðon-ðe means therefore/for that reason; eornostlice can mean accordingly/thus. Do not substitute the causal because sense.\par
% glossary-block:B00024:end
% glossary-block:B00025:begin
\needspace{3\baselineskip}
\textbf{according to} — \textbf{æfter} (prep.; conj. in fixed expressions). With dative. \textbf{be} (prep.). With dative. The supplied on + hīw expression means according to/in the forms of. Retain the whole supplied hīw phrase rather than treating on alone as a general according-to expression.\par
% glossary-block:B00025:end
% glossary-block:B00026:begin
\needspace{3\baselineskip}
\textbf{account} — \textbf{spell} (n.). account; discourse.\par
% glossary-block:B00026:end
% glossary-block:B00027:begin
\needspace{3\baselineskip}
\textbf{account heard} — \textbf{hlīsa} (n.). report; account heard.\par
% glossary-block:B00027:end
% glossary-block:B00028:begin
\needspace{3\baselineskip}
\textbf{a certain} — \textbf{ān} (numeral; adj.). Introduces an unnamed singular individual; not an obligatory indefinite article. \textbf{sum} (indefinite adjective/pronoun). a certain; some. \textbf{sum} (indefinite adjective/pronoun), \textbf{ān} (numeral; adj.). Sum or the taught introducing-an-individual ān can mean a certain. Ordinary English a/an does not require either word.\par
% glossary-block:B00028:end
% glossary-block:B00029:begin
\needspace{3\baselineskip}
\textbf{acknowledge} — \textbf{oncnāwan} (v.). recognize; acknowledge.\par
% glossary-block:B00029:end
% glossary-block:B00030:begin
\needspace{3\baselineskip}
\textbf{acorn} — \textbf{æcern} (n.).\par
% glossary-block:B00030:end
% glossary-block:B00031:begin
\needspace{3\baselineskip}
\textbf{across} — \textbf{ofer} (prep.). over; across. Spatial frame; case varies with the specific construction. Do not impose a universal motion rule.\par
% glossary-block:B00031:end
% glossary-block:B00032:begin
\needspace{3\baselineskip}
\textbf{act} — \textbf{plegian} (v.). act; perform. Performance context; ordinary play is a separate assigned sense.\par
% glossary-block:B00032:end
% glossary-block:B00033:begin
\needspace{3\baselineskip}
\textbf{add} — \textbf{tōgēiecan} (v.).  \textbf{underðēodan} (v.). add; subjoin. Supplied passive/participial addition to speech.\par
% glossary-block:B00033:end
% glossary-block:B00034:begin
\needspace{3\baselineskip}
\textbf{adder} — \textbf{næddre} (n.). snake; adder. Supplied weak feminine plural.\par
% glossary-block:B00034:end
% glossary-block:B00035:begin
\needspace{3\baselineskip}
\textbf{admonish} — \textbf{myngian} (v.). admonish; prompt. Supplied participial source use.\par
% glossary-block:B00035:end
% glossary-block:B00036:begin
\needspace{3\baselineskip}
\textbf{adorn} — \textbf{geglęngan} (v.). The supplied passive adornment expression uses mid + dative material. \textbf{geweorðian} (v.). distinguish; adorn. Accusative person adorned; mid + dative equipment or adornment.\par
% glossary-block:B00036:end
% glossary-block:B00037:begin
\needspace{3\baselineskip}
\textbf{adorned with} — \textbf{mid} (prep.; adv. in fixed expressions). with. Material or clothing in the supplied adornment phrase.\par
% glossary-block:B00037:end
% glossary-block:B00038:begin
\needspace{3\baselineskip}
\textbf{adversaries} — \textbf{wiðerwinna} (n.).\par
% glossary-block:B00038:end
% glossary-block:B00039:begin
\needspace{3\baselineskip}
\textbf{advise} — \textbf{lǣran} (v.). advise; instruct.\par
% glossary-block:B00039:end
% glossary-block:B00040:begin
\needspace{3\baselineskip}
\textbf{afar off} — \textbf{feorran} (adv.). feorran; the source supplies the from-afar paraphrase.\par
% glossary-block:B00040:end
% glossary-block:B00041:begin
\needspace{3\baselineskip}
\textbf{afflict} — \textbf{geðrǣstan} (v.). afflict; beat. Supplied participial source use. \textbf{ierman} (v.).\par
% glossary-block:B00041:end
% glossary-block:B00042:begin
\needspace{3\baselineskip}
\textbf{after} — \textbf{æfter} (prep.; conj. in fixed expressions). Temporal preposition with dative. In the supplied finite-clause conjunction æfter ðǣm ðe. after; following. Dative person followed, including motion after that person. after; in search of. Person sought or inquired after; keep this lexical frame distinct from temporal after. Temporal æfter + dative means after a person/event. The finite-clause conjunction is the complete supplied æfter ðǣm ðe. \textbf{siððan} (adv.; conj.). since; after; afterwards. Conjunction since/after versus adverb afterwards. Æfter ðǣm ðe is the complete supplied finite-clause conjunction after; ðe alone does not mean after.\par
% glossary-block:B00042:end
% glossary-block:B00043:begin
\needspace{3\baselineskip}
\textbf{after that} — Æfter ðon is the complete supplied temporal phrase after that.\par
% glossary-block:B00043:end
% glossary-block:B00044:begin
\needspace{3\baselineskip}
\textbf{afterwards} — \textbf{eft, ęft} (adv.). again; afterwards. \textbf{siððan} (adv.; conj.). since; after; afterwards. Conjunction since/after versus adverb afterwards.\par
% glossary-block:B00044:end
% glossary-block:B00045:begin
\needspace{3\baselineskip}
\textbf{again} — \textbf{eft, ęft} (adv.). again; also. The source note also permits the additive sense. again; back. again; afterwards. eft is the repeated-occurrence adverb. \textbf{mā} (adv.). more; again. In the negative phrase never more/never again. \textbf{ongēan} (adv.; prep.; conj.). back; again.\par
% glossary-block:B00045:end
% glossary-block:B00046:begin
\needspace{3\baselineskip}
\textbf{against} — \textbf{ongēan} (adv.; prep.; conj.). against; opposite. Accusative landmark or opponent. \textbf{wið} (prep.). against; opposed to. Accusative opponent. with; against. Dative fellow players in the exact supplied playing frame.\par
% glossary-block:B00046:end
% glossary-block:B00047:begin
\needspace{3\baselineskip}
\textbf{against that} — \textbf{ongēan} (adv.; prep.; conj.). by the time that. Complete supplied ongēan ðætte clause. Supplied ongēan ðætte + prospective subjunctive clause.\par
% glossary-block:B00047:end
% glossary-block:B00048:begin
\needspace{3\baselineskip}
\textbf{age} — \textbf{ieldu} (n.). age; old age.\par
% glossary-block:B00048:end
% glossary-block:B00049:begin
\needspace{3\baselineskip}
\textbf{agility} — \textbf{snelnes} (n.). skill; agility.\par
% glossary-block:B00049:end
% glossary-block:B00050:begin
\needspace{3\baselineskip}
\textbf{agreeably} — \textbf{gecwēmlīce} (adv.). pleasantly; agreeably.\par
% glossary-block:B00050:end
% glossary-block:B00051:begin
\needspace{3\baselineskip}
\textbf{agriculture} — \textbf{eorðtilð} (n.).\par
% glossary-block:B00051:end
% glossary-block:B00052:begin
\needspace{3\baselineskip}
\textbf{a hair} — \textbf{hǣr} (n.). hair; a hair.\par
% glossary-block:B00052:end
% glossary-block:B00053:begin
\needspace{3\baselineskip}
\textbf{ahead of} — \textbf{be-foran} (prep.). before; ahead of. Dative person, including going or sending ahead of that person.\par
% glossary-block:B00053:end
% glossary-block:B00054:begin
\needspace{3\baselineskip}
\textbf{a hundredfold} — \textbf{hundtēontig} (numeral). A hundred in return, in the supplied recompense expression.\par
% glossary-block:B00054:end
% glossary-block:B00055:begin
\needspace{3\baselineskip}
\textbf{aid} — \textbf{fultum} (n.). help; aid.\par
% glossary-block:B00055:end
% glossary-block:B00056:begin
\needspace{3\baselineskip}
\textbf{air} — \textbf{lyft} (n.).\par
% glossary-block:B00056:end
% glossary-block:B00057:begin
\needspace{3\baselineskip}
\textbf{a large number of} — \textbf{fela} (adj.). many; a large number of. Genitive of the things counted in the taught frame.\par
% glossary-block:B00057:end
% glossary-block:B00058:begin
\needspace{3\baselineskip}
\textbf{Alfred} — \textbf{Ælfred} (proper name).\par
% glossary-block:B00058:end
% glossary-block:B00059:begin
\needspace{3\baselineskip}
\textbf{a light} — \textbf{lēoht} (n.). light; luminary. Individual countable light, distinct from mass light. Mass light and an individual countable light/luminary are separate uses of lēoht. Use the count sense for a light or several lights.\par
% glossary-block:B00059:end
% glossary-block:B00060:begin
\needspace{3\baselineskip}
\textbf{alive} — \textbf{cucu} (adj.). alive; living. Supplied accusative neuter plural adjective. \textbf{cwic} (adj.). alive; living. \textbf{libban} (v.). living. Supplied weak participial adjective. The bracketed liveth yet is rendered by finite libban + gīet; an added copula is not required. The supplied whole phrase on līfe means alive; on alone does not.\par
% glossary-block:B00060:end
% glossary-block:B00061:begin
\needspace{3\baselineskip}
\textbf{all} — \textbf{eall} (adj.; substantive; adv.). all; whole. all people; everything. Substantival: plural for people, neuter singular for everything. Eall agrees with a noun in all/whole expressions. Substantival plural refers to all people; neuter singular can mean everything.\par
% glossary-block:B00061:end
% glossary-block:B00062:begin
\needspace{3\baselineskip}
\textbf{all around} — \textbf{ymbhwyrft} (n.). circuit. Whole on ymbhwyrfte expression means all around; supplied dependent landmark. Supplied on ymbhwyrfte + genitive landmark in this exercise.\par
% glossary-block:B00062:end
% glossary-block:B00063:begin
\needspace{3\baselineskip}
\textbf{allow} — \textbf{ālīefan} (v.). permit; allow. Impersonal permission frame, with dative permitted person; consult the complete supplied expression. \textbf{gelīefst} (v.; supplied pres. ind. 2 sg.). permit; allow. \textbf{līefan} (v.). allow; permit. \textbf{unnan} (v.). grant; allow. Dative person, genitive thing granted. Dative recipient and genitive thing granted; preserve the supported subjunctive command unne.\par
% glossary-block:B00063:end
% glossary-block:B00064:begin
\needspace{3\baselineskip}
\textbf{all people} — \textbf{eall} (adj.; substantive; adv.). all people; everything. Substantival: plural for people, neuter singular for everything.\par
% glossary-block:B00064:end
% glossary-block:B00065:begin
\needspace{3\baselineskip}
\textbf{almighty} — \textbf{ælmihtig} (adj.).\par
% glossary-block:B00065:end
% glossary-block:B00066:begin
\needspace{3\baselineskip}
\textbf{alone} — \textbf{ān} (numeral; adj.). Cardinal one/single and the agreeing adjective alone are distinct uses of ān. Preserve the supplied form where the exercise provides one.\par
% glossary-block:B00066:end
% glossary-block:B00067:begin
\needspace{3\baselineskip}
\textbf{along} — \textbf{geond} (prep.). along; through. With accusative.\par
% glossary-block:B00067:end
% glossary-block:B00068:begin
\needspace{3\baselineskip}
\textbf{along with} — \textbf{mid} (prep.; adv. in fixed expressions). with; along with. Supplied mid wes, be with them.\par
% glossary-block:B00068:end
% glossary-block:B00069:begin
\needspace{3\baselineskip}
\textbf{also} — \textbf{ēac} (adv.). Also in the combination ēac swilce. \textbf{eft, ęft} (adv.). again; also. The source note also permits the additive sense.\par
% glossary-block:B00069:end
% glossary-block:B00070:begin
\needspace{3\baselineskip}
\textbf{although} — \textbf{ðēah} (conj.; adv.). although; even though. ðēah with the taught concessive clause; a subjunctive need not make an admitted fact unreal. \textbf{ðēah ðe}. Supplied whole concessive conjunction in EX-XII-XIII-43-10.\par
% glossary-block:B00070:end
% glossary-block:B00071:begin
\needspace{3\baselineskip}
\textbf{always} — \textbf{ā} (adv.). ever; always. \textbf{ǣfre} (adv.). ever; always. \textbf{simle} (adv.).\par
% glossary-block:B00071:end
% glossary-block:B00072:begin
\needspace{3\baselineskip}
\textbf{am} — \textbf{wesan} (v.). be. Copular present/past identification or state; other auxiliary constructions are separate senses.\par
% glossary-block:B00072:end
% glossary-block:B00073:begin
\needspace{3\baselineskip}
\textbf{Amen} — \textbf{Amen} (interj.).\par
% glossary-block:B00073:end
% glossary-block:B00074:begin
\needspace{3\baselineskip}
\textbf{amid} — \textbf{on} (prep.). in; amid. Circumstances of the stated unrest.\par
% glossary-block:B00074:end
% glossary-block:B00075:begin
\needspace{3\baselineskip}
\textbf{among} — \textbf{betweox} (prep.). Accusative in the exact supplied phrases; not a universal case rule. among; between. \textbf{gemang} (prep.). Accusative group in the selected sending-among frame. \textbf{on} (prep.). among; in. Dative people among whom friends are obtained.\par
% glossary-block:B00075:end
% glossary-block:B00076:begin
\needspace{3\baselineskip}
\textbf{amount} — \textbf{micel} (adj.). much; an amount. Quantity sense. \textbf{dǣl} (n.). dǣl + genitive of the substance, as in sumne dǣl mętes; the amount is accusative object.\par
% glossary-block:B00076:end
% glossary-block:B00077:begin
\needspace{3\baselineskip}
\textbf{an} — Ordinary English a/an usually has no separate Old English article. Ān is used in its taught a-certain-individual sense, not inserted into every indefinite noun phrase. No obligatory indefinite article: bare predicates/nouns often translate a/an. ān is used for the taught introducing-a-certain-individual sense, not inserted everywhere.\par
% glossary-block:B00077:end
% glossary-block:B00078:begin
\needspace{3\baselineskip}
\textbf{and} — \textbf{and} (conj.). Simple coordination. \textbf{witodlice} (adv.). now; and. Discourse connective.\par
% glossary-block:B00078:end
% glossary-block:B00079:begin
\needspace{3\baselineskip}
\textbf{angel} — \textbf{ęngel} (n.).\par
% glossary-block:B00079:end
% glossary-block:B00080:begin
\needspace{3\baselineskip}
\textbf{animal} — \textbf{dēor} (n.). animal; beast. \textbf{nīeten} (n.). animal; beast; cattle. Cattle is the selected plural use; retain the early task distinction from dēor. \textbf{dēor} (n.), \textbf{nīeten} (n.). Dēor is the Unit 1 animal/beast target. Nīeten can also denote an animal in later contexts; use the noun specified by the exercise.\par
% glossary-block:B00080:end
% glossary-block:B00081:begin
\needspace{3\baselineskip}
\textbf{annihilate} — \textbf{fordilgian} (v.). destroy; annihilate.\par
% glossary-block:B00081:end
% glossary-block:B00082:begin
\needspace{3\baselineskip}
\textbf{announce} — \textbf{cȳðan} (v.). announce; make known. Dative person informed; accusative thing announced or a ðæt content clause.\par
% glossary-block:B00082:end
% glossary-block:B00083:begin
\needspace{3\baselineskip}
\textbf{anoint} — \textbf{gesmierwan} (v.). Supplied source form anointed. \textbf{smierwan} (v.). Accusative body part or person, mid + dative material.\par
% glossary-block:B00083:end
% glossary-block:B00084:begin
\needspace{3\baselineskip}
\textbf{another} — \textbf{ōðer} (adj.; ordinal). other; another.\par
% glossary-block:B00084:end
% glossary-block:B00085:begin
\needspace{3\baselineskip}
\textbf{answer} — \textbf{andswarian} (v.). answer; reply. Dative addressee; the reply may be given as direct speech.\par
% glossary-block:B00085:end
% glossary-block:B00086:begin
\needspace{3\baselineskip}
\textbf{anxiety} — \textbf{sorg} (n.). distress; anxiety; trouble; sorrow.\par
% glossary-block:B00086:end
% glossary-block:B00087:begin
\needspace{3\baselineskip}
\textbf{any} — \textbf{ǣnig} (indefinite adjective).  \textbf{hwilc} (interrogative/indefinite adj. or pron.). any; some. Indefinite adjective, distinct from an interrogative. \textbf{ǣnig} (indefinite adjective), \textbf{hwilc} (interrogative/indefinite adj. or pron.). Ǣnig means any; indefinite hwilc has a bounded any/some use. Hwā meaning anyone is a separate personal indefinite, not an adjective to insert before every noun.\par
% glossary-block:B00087:end
% glossary-block:B00088:begin
\needspace{3\baselineskip}
\textbf{any more} — \textbf{furðor} (adv.). further; any more. Adverb in the supplied no-spirit-more expression.\par
% glossary-block:B00088:end
% glossary-block:B00089:begin
\needspace{3\baselineskip}
\textbf{anyone} — \textbf{hwā, hwæt} (interrogative/indefinite pron.).\par
% glossary-block:B00089:end
% glossary-block:B00090:begin
\needspace{3\baselineskip}
\textbf{anything} — \textbf{āwiht} (indefinite substantive).\par
% glossary-block:B00090:end
% glossary-block:B00091:begin
\needspace{3\baselineskip}
\textbf{apart from} — \textbf{būtan} (prep.). besides; apart from. Dative excluded person or thing.\par
% glossary-block:B00091:end
% glossary-block:B00092:begin
\needspace{3\baselineskip}
\textbf{Apollo} — \textbf{Apollines} (proper name; source nom. sg.). The god, distinct from Apollonius.\par
% glossary-block:B00092:end
% glossary-block:B00093:begin
\needspace{3\baselineskip}
\textbf{Apollonius} — \textbf{Apollonius} (proper name). Supplied dative person in the assigned source expression.\par
% glossary-block:B00093:end
% glossary-block:B00094:begin
\needspace{3\baselineskip}
\textbf{appeal} — \textbf{bēn} (n.). petition; appeal.\par
% glossary-block:B00094:end
% glossary-block:B00095:begin
\needspace{3\baselineskip}
\textbf{appear} — \textbf{ætēowian} (v.).\par
% glossary-block:B00095:end
% glossary-block:B00096:begin
\needspace{3\baselineskip}
\textbf{appearance} — \textbf{ansīen} (n.).\par
% glossary-block:B00096:end
% glossary-block:B00097:begin
\needspace{3\baselineskip}
\textbf{appoint} — \textbf{gesęttan} (v.). appoint as lord. Supplied frame with person appointed, tō + dative role, and dative people governed. \textbf{sęttan} (v.). appoint; establish. With accusative entity/decree in the assigned frame.\par
% glossary-block:B00097:end
% glossary-block:B00098:begin
\needspace{3\baselineskip}
\textbf{appoint as} — \textbf{gesęttan} (v.). appoint as lord. Supplied frame with person appointed, tō + dative role, and dative people governed. Supplied appointment frame: accusative person, tō + dative role, and dative people governed.\par
% glossary-block:B00098:end
% glossary-block:B00099:begin
\needspace{3\baselineskip}
\textbf{approach} — \textbf{genēalǣcan} (v.). approach; draw near. Dative person, place or event approached. Supplied source form approached. \textbf{nēalǣcan} (v.). approach; draw near. Dative person or place approached when the goal is expressed. \textbf{nēalǣcan} (v.), \textbf{genēalǣcan} (v.). nēalǣcan and genēalǣcan mean approach/draw near; an expressed goal is dative. Preserve the verb selected by the exercise.\par
% glossary-block:B00099:end
% glossary-block:B00100:begin
\needspace{3\baselineskip}
\textbf{approximately} — \textbf{swā-swā} (conj.). about; approximately. Approximation in the supplied source expression. \textbf{swilce} (adv.). about; approximately.\par
% glossary-block:B00100:end
% glossary-block:B00101:begin
\needspace{3\baselineskip}
\textbf{Arcestrates} — \textbf{Arcestrates} (proper name). King of the region in the story; supplied name forms. The king; supplied name forms.\par
% glossary-block:B00101:end
% glossary-block:B00102:begin
\needspace{3\baselineskip}
\textbf{are} — \textbf{wesan} (v.). be. Copular present/past identification or state; other auxiliary constructions are separate senses. be. Supplied plural present copula.\par
% glossary-block:B00102:end
% glossary-block:B00103:begin
\needspace{3\baselineskip}
\textbf{arise} — \textbf{ārīsan} (v.). Intransitive arise; of + dative source in the bounded rise-from-death expression. \textbf{weorðan} (v.). arise; come about. Intransitive with a nominative subject, including the arising of a calm. \textbf{rīsan} (v.). Earlier taught rīsan supplies this command; prefixed ārīsan is separately used later.\par
% glossary-block:B00103:end
% glossary-block:B00104:begin
\needspace{3\baselineskip}
\textbf{arithmetic} — \textbf{gerīmcræft} (n.). arithmetic; computation.\par
% glossary-block:B00104:end
% glossary-block:B00105:begin
\needspace{3\baselineskip}
\textbf{army} — \textbf{hęre} (n.).  \textbf{weorod} (n.). host; army.\par
% glossary-block:B00105:end
% glossary-block:B00106:begin
\needspace{3\baselineskip}
\textbf{around} — \textbf{ābūtan} (prep.). around; about. With accusative.\par
% glossary-block:B00106:end
% glossary-block:B00107:begin
\needspace{3\baselineskip}
\textbf{arrange} — \textbf{gedōn} (v.). make; arrange. The present form/arrangement of the translation in the source.\par
% glossary-block:B00107:end
% glossary-block:B00108:begin
\needspace{3\baselineskip}
\textbf{array} — \textbf{frætwung} (n.).\par
% glossary-block:B00108:end
% glossary-block:B00109:begin
\needspace{3\baselineskip}
\textbf{arrive} — \textbf{becuman} (v.). come; arrive. \textbf{tōcuman} (v.). come; arrive.\par
% glossary-block:B00109:end
% glossary-block:B00110:begin
\needspace{3\baselineskip}
\textbf{as} — \textbf{swā} (adverb; conjunction). Only in sōna swā. as; in the same way as. as; like. Predicative comparison. Swā + adjective/adverb + swā frames equal degree, as … as: the first swā marks degree and the second introduces the standard. \textbf{ðæs-ðe} (conj.). Only in supplied ðæs-ðe mē ðyncð. Swā can introduce a predicative comparison, as/like; swā-swā means just as in its supplied comparison. Swilce means as if in the supplied apparent-rejuvenation comparison. Equal degree uses the separate paired swā … swā construction. \textbf{ðā} (conj.). Temporal as means when in the expressly supplied factual-past clause with ðā; comparison as uses a different construction.\par
% glossary-block:B00110:end
% glossary-block:B00111:begin
\needspace{3\baselineskip}
\textbf{as … as} — Swā + adjective/adverb + swā frames equal degree, as … as: the first swā marks degree and the second introduces the standard. Swā + adjective/adverb + swā frames equal degree, as … as. The first swā marks degree; the second introduces the standard. Comparative the … the uses a separate construction.\par
% glossary-block:B00111:end
% glossary-block:B00112:begin
\needspace{3\baselineskip}
\textbf{ascend} — \textbf{ūpāstīgan} (v.). rise up; ascend.\par
% glossary-block:B00112:end
% glossary-block:B00113:begin
\needspace{3\baselineskip}
\textbf{as far as} — \textbf{ōð} (prep.). as far as; up to.\par
% glossary-block:B00113:end
% glossary-block:B00114:begin
\needspace{3\baselineskip}
\textbf{as if} — \textbf{swilce} (adv.). Supplied apparent-rejuvenation comparison.\par
% glossary-block:B00114:end
% glossary-block:B00115:begin
\needspace{3\baselineskip}
\textbf{as it seems} — The supplied expression ðæs-ðe mē ðyncð means as it seems to me; ðæs-ðe alone introduces the as clause.\par
% glossary-block:B00115:end
% glossary-block:B00116:begin
\needspace{3\baselineskip}
\textbf{ask} — \textbf{āscian} (v.). ask about. Accusative person asked, genitive matter asked about; an indirect question may express the matter. \textbf{biddan} (v.). ask; request. Genitive thing requested; an accusative person may be added.\par
% glossary-block:B00116:end
% glossary-block:B00117:begin
\needspace{3\baselineskip}
\textbf{ask about} — \textbf{āscian} (v.). Accusative person asked, genitive matter asked about; an indirect question may express the matter. Supplied source form asked about.\par
% glossary-block:B00117:end
% glossary-block:B00118:begin
\needspace{3\baselineskip}
\textbf{ask for} — \textbf{giernan} (v.). desire; request; woo. Genitive desired thing/person in the selected frames, including the maiden wooed.\par
% glossary-block:B00118:end
% glossary-block:B00119:begin
\needspace{3\baselineskip}
\textbf{as long as} — The whole phrase swā lange swā means as long as and has inasmuch-as force in the assigned biblical context. Lange alone means long. Swā lange swā is the complete supplied phrase as long as; in the assigned biblical context it has inasmuch-as force. Swā lange swā is the complete supplied phrase as long as; in the assigned biblical context it has inasmuch-as force. Lange alone means long.\par
% glossary-block:B00119:end
% glossary-block:B00120:begin
\needspace{3\baselineskip}
\textbf{as lord} — \textbf{gesęttan} (v.). Supplied appointment frame: accusative person, tō + dative role, and dative people governed.\par
% glossary-block:B00120:end
% glossary-block:B00121:begin
\needspace{3\baselineskip}
\textbf{a smaller amount} — \textbf{lǣs} (degree word). less; a smaller amount.\par
% glossary-block:B00121:end
% glossary-block:B00122:begin
\needspace{3\baselineskip}
\textbf{ass} — \textbf{assa} (n.). ass; donkey.\par
% glossary-block:B00122:end
% glossary-block:B00123:begin
\needspace{3\baselineskip}
\textbf{assail} — \textbf{gestandan} (v.). assail; come upon. Supplied source sense. \textbf{onwinnan} (v.). assail; attack.\par
% glossary-block:B00123:end
% glossary-block:B00124:begin
\needspace{3\baselineskip}
\textbf{Asser} — \textbf{Asser} (proper name).\par
% glossary-block:B00124:end
% glossary-block:B00125:begin
\needspace{3\baselineskip}
\textbf{assist} — \textbf{gefultumian} (v.). help; assist.\par
% glossary-block:B00125:end
% glossary-block:B00126:begin
\needspace{3\baselineskip}
\textbf{as soon as} — \textbf{sōna} (adv.). In sōna swā. Sōna swā is the complete supplied expression as soon as; swā alone does not include the soon meaning. The complete supplied construction uses sōna with ðæs-ðe to mean as soon as. Retain both parts.\par
% glossary-block:B00126:end
% glossary-block:B00127:begin
\needspace{3\baselineskip}
\textbf{as wife} — Tō wīfe is the supplied role phrase as wife/to wife.\par
% glossary-block:B00127:end
% glossary-block:B00128:begin
\needspace{3\baselineskip}
\textbf{at} — \textbf{æt} (prep.).  \textbf{tō} (prep.). at; in response to. Only the exact at-the-prayers expression. at; to. Bounded supplied root expression with genitive dependent and dative root. \textbf{wið} (prep.). by; beside; at. Accusative landmark in the taught proximity phrase.\par
% glossary-block:B00128:end
% glossary-block:B00129:begin
\needspace{3\baselineskip}
\textbf{at first} — \textbf{ǣrest} (adv.). first; at first.\par
% glossary-block:B00129:end
% glossary-block:B00130:begin
\needspace{3\baselineskip}
\textbf{at hand} — \textbf{nēah} (adv.). near; at hand. nēah is the predicate in nigh at hand.\par
% glossary-block:B00130:end
% glossary-block:B00131:begin
\needspace{3\baselineskip}
\textbf{at length} — The supplied whole phrase æt nīehstan means at length. \textbf{nīehst} (adv.). Whole supplied æt nīehstan.\par
% glossary-block:B00131:end
% glossary-block:B00132:begin
\needspace{3\baselineskip}
\textbf{at noon} — The supplied whole phrase tō middes dæges means at noon/midday; middes and dæges are genitives. \textbf{middes} (adverb or particle). Whole supplied phrase tō middes dæges; preserve its genitives instead of applying ordinary tō + dative mechanically.\par
% glossary-block:B00132:end
% glossary-block:B00133:begin
\needspace{3\baselineskip}
\textbf{at once} — \textbf{ætgædere} (adv.). together; at once. Paired with samod in the supplied source. \textbf{sōna} (adv.). immediately; at once. sōna carries soon in this future-reference prompt.\par
% glossary-block:B00133:end
% glossary-block:B00134:begin
\needspace{3\baselineskip}
\textbf{at one's side} — \textbf{on} (prep.). Accusative hand in the bounded placement expression.\par
% glossary-block:B00134:end
% glossary-block:B00135:begin
\needspace{3\baselineskip}
\textbf{attack} — \textbf{onwinnan} (v.). assail; attack.\par
% glossary-block:B00135:end
% glossary-block:B00136:begin
\needspace{3\baselineskip}
\textbf{attend} — \textbf{ðēnian} (v.). serve; attend. Supplied preterite source form.\par
% glossary-block:B00136:end
% glossary-block:B00137:begin
\needspace{3\baselineskip}
\textbf{attendance} — \textbf{ðēnung} (n.). service; attendance.\par
% glossary-block:B00137:end
% glossary-block:B00138:begin
\needspace{3\baselineskip}
\textbf{attention} — \textbf{mōd} (n.). mind; attention.\par
% glossary-block:B00138:end
% glossary-block:B00139:begin
\needspace{3\baselineskip}
\textbf{attentive} — \textbf{behygdig} (adj.). attentive; thoughtful.\par
% glossary-block:B00139:end
% glossary-block:B00140:begin
\needspace{3\baselineskip}
\textbf{attentively} — \textbf{geornlice} (adv.). earnestly; attentively.\par
% glossary-block:B00140:end
% glossary-block:B00141:begin
\needspace{3\baselineskip}
\textbf{at that point} — \textbf{ðonne} (conj.; adv.). then; at that point.\par
% glossary-block:B00141:end
% glossary-block:B00142:begin
\needspace{3\baselineskip}
\textbf{at the left side} — \textbf{on} (prep.). at one's side. Accusative hand in the bounded placement expression.\par
% glossary-block:B00142:end
% glossary-block:B00143:begin
\needspace{3\baselineskip}
\textbf{at the outset} — \textbf{fyrmest} (adv.). first of all; at the outset.\par
% glossary-block:B00143:end
% glossary-block:B00144:begin
\needspace{3\baselineskip}
\textbf{at the right side} — \textbf{on} (prep.). at one's side. Accusative hand in the bounded placement expression.\par
% glossary-block:B00144:end
% glossary-block:B00145:begin
\needspace{3\baselineskip}
\textbf{at the root of} — The supplied phrase uses tō + dative root with a genitive dependent naming what has the root. Retain that full phrase; tō alone means at/to.\par
% glossary-block:B00145:end
% glossary-block:B00146:begin
\needspace{3\baselineskip}
\textbf{at work} — The supplied on + activity-noun phrases express at work/engaged in that activity. Use the whole activity phrase and its supplied form; on alone does not name the activity.\par
% glossary-block:B00146:end
% glossary-block:B00147:begin
\needspace{3\baselineskip}
\textbf{available time} — \textbf{ǣmetta} (n.). leisure; available time.\par
% glossary-block:B00147:end
% glossary-block:B00148:begin
\needspace{3\baselineskip}
\textbf{avoid} — \textbf{flēon} (v.). flee; flee from; avoid. Intransitive or with accusative person/thing avoided.\par
% glossary-block:B00148:end
% glossary-block:B00149:begin
\needspace{3\baselineskip}
\textbf{away} — \textbf{onwęg} (adv.).\par
% glossary-block:B00149:end
% glossary-block:B00150:begin
\needspace{3\baselineskip}
\textbf{a while ago} — The supplied whole phrase hwīlum ǣr means a while ago/formerly. Hwīlum alone has the separate earlier meaning sometimes. Hwīlum ǣr is the supplied whole expression a while ago. Hwīlum alone has the separate earlier meaning sometimes.\par
% glossary-block:B00150:end
% glossary-block:B00151:begin
\needspace{3\baselineskip}
\textbf{awl} — \textbf{ǣl} (n.).\par
% glossary-block:B00151:end
% glossary-block:B00152:begin
\needspace{3\baselineskip}
\textbf{ax} — \textbf{æx} (n.).\par
% glossary-block:B00152:end
\pdfbookmark[1]{B}{outline-B00153}
% glossary-block:B00153:begin
\needspace{5\baselineskip}
\subsection*{B}
% glossary-block:B00153:end
% glossary-block:B00154:begin
\textbf{back} — \textbf{eft, ęft} (adv.). again; back. \textbf{ongēan} (adv.; prep.; conj.). back; again.\par
% glossary-block:B00154:end
% glossary-block:B00155:begin
\needspace{3\baselineskip}
\textbf{bad} — \textbf{yfel} (adj.). bad; evil.\par
% glossary-block:B00155:end
% glossary-block:B00156:begin
\needspace{3\baselineskip}
\textbf{badly} — \textbf{yfele} (adv.). badly; wrongly. Adverb modifying doing; distinguish the adjective yfel.\par
% glossary-block:B00156:end
% glossary-block:B00157:begin
\needspace{3\baselineskip}
\textbf{ball} — \textbf{ðoðor} (n.).  \textbf{top} (n.). ball?. Tentative source interpretation; retain the question mark and uncertainty.\par
% glossary-block:B00157:end
% glossary-block:B00158:begin
\needspace{3\baselineskip}
\textbf{banquet} — \textbf{bēorscipe} (n.). feast; banquet. \textbf{gebēorscipe} (n.). feast; banquet. \textbf{gereord} (n.). meal; banquet.\par
% glossary-block:B00158:end
% glossary-block:B00159:begin
\needspace{3\baselineskip}
\textbf{Baptist} — \textbf{fuluhtere} (n.). baptizer; Baptist. Baptist as John's title in the assigned passage.\par
% glossary-block:B00159:end
% glossary-block:B00160:begin
\needspace{3\baselineskip}
\textbf{baptizer} — \textbf{fuluhtere} (n.). baptizer; Baptist. Baptist as John's title in the assigned passage.\par
% glossary-block:B00160:end
% glossary-block:B00161:begin
\needspace{3\baselineskip}
\textbf{barely} — \textbf{unēaðe} (adv.). with difficulty; barely.\par
% glossary-block:B00161:end
% glossary-block:B00162:begin
\needspace{3\baselineskip}
\textbf{bath} — \textbf{bæð} (n.).  \textbf{ðwēal} (n.). bath; bathing place.\par
% glossary-block:B00162:end
% glossary-block:B00163:begin
\needspace{3\baselineskip}
\textbf{bathe} — \textbf{baðian} (v.). Supplied source form bathed.\par
% glossary-block:B00163:end
% glossary-block:B00164:begin
\needspace{3\baselineskip}
\textbf{bathing place} — \textbf{bæðstęde} (n.). bathing place; gymnasium. Gymnasium is the contextual interpretation in the supplied reading. \textbf{ðwēal} (n.). bath; bathing place.\par
% glossary-block:B00164:end
% glossary-block:B00165:begin
\needspace{3\baselineskip}
\textbf{be} — Passive weorðan combines with a past participle for the assigned event or becoming. Lexical weorðan become and copular wesan be have separate uses. \textbf{wesan} (v.). Copular present/past identification or state; other auxiliary constructions are separate senses. Supplied plural present copula. Copular eom/wæs and bēo indicative series; later passive and compound uses require their own construction. Wesan + the assigned intransitive past participle supplies the taught compound construction. With cumen come it can render have/had come; English have is not an extra word to add to that construction. Passive wesan combines with the assigned past participle; be done includes the participial meaning as well as the finite be auxiliary. The supplied progressive construction is wesan + present participle. Ordinary English be doing can also be rendered by a simple finite Old English verb; this compound is not the universal default.\par
% glossary-block:B00165:end
% glossary-block:B00166:begin
\needspace{3\baselineskip}
\textbf{be able to} — \textbf{cunnan} (v.). know how; be able. Learned skill; plain infinitive complement. Mugan expresses capability, can/be able to, with a plain infinitive. Learned ability can instead use cunnan; neither construction takes an inserted tō before that infinitive. The supplied mugan form expresses capability, can/could/might be able to, with a plain infinitive. It is not the permission sense of English may/might. \textbf{mugan} (v.), \textbf{cunnan} (v.). Mugan expresses ability or capability. Cunnan + plain infinitive emphasizes learned ability. A modal complement is a plain infinitive, without tō.\par
% glossary-block:B00166:end
% glossary-block:B00167:begin
\needspace{3\baselineskip}
\textbf{be accustomed} — \textbf{gewunian} (v.). be accustomed; be wont.\par
% glossary-block:B00167:end
% glossary-block:B00168:begin
\needspace{3\baselineskip}
\textbf{be accustomed to} — \textbf{gewunian} (v.). gewunian + plain infinitive; no separate be auxiliary in the model.\par
% glossary-block:B00168:end
% glossary-block:B00169:begin
\needspace{3\baselineskip}
\textbf{be acquainted with} — \textbf{cunnan} (v.). know; be acquainted with. Accusative person known; distinct from learned ability.\par
% glossary-block:B00169:end
% glossary-block:B00170:begin
\needspace{3\baselineskip}
\textbf{be allowed to} — Mōtan expresses permission, may/be allowed to, with a plain infinitive. A wish beginning with English may instead uses the taught independent subjunctive. \textbf{mōtan} (v.). mōtan + plain infinitive; permission/prohibition differs from ability and independent wishes.\par
% glossary-block:B00170:end
% glossary-block:B00171:begin
\needspace{3\baselineskip}
\textbf{be amazed} — \textbf{wafian} (v.). Supplied source form.\par
% glossary-block:B00171:end
% glossary-block:B00172:begin
\needspace{3\baselineskip}
\textbf{bear} — \textbf{āberan} (v.). carry. Supplied carried form; retain the ā-prefix. \textbf{beran} (v.). carry. Accusative person or thing carried. bear; produce. Accusative produce; distinct from physical carrying. Accusative produce; distinguish physical carrying. \textbf{gefęrian} (v.). carry; bear. Supplied passive movement; was carried/borne.\par
% glossary-block:B00172:end
% glossary-block:B00173:begin
\needspace{3\baselineskip}
\textbf{bear fruit} — \textbf{beran} (v.). bear; produce. Accusative produce; distinct from physical carrying. Accusative produce; distinguish physical carrying.\par
% glossary-block:B00173:end
% glossary-block:B00174:begin
\needspace{3\baselineskip}
\textbf{bear in mind} — \textbf{geðęncean} (v.). remember; bear in mind.\par
% glossary-block:B00174:end
% glossary-block:B00175:begin
\needspace{3\baselineskip}
\textbf{bear rule} — \textbf{rīcsian} (v.). rule; have dominion. Used absolutely in the taught frame; no new object government. rīcsian can be used absolutely; English have/bear need no independent translation.\par
% glossary-block:B00175:end
% glossary-block:B00176:begin
\needspace{3\baselineskip}
\textbf{bear seed} — \textbf{beran} (v.). bear; produce. Accusative produce; distinct from physical carrying.\par
% glossary-block:B00176:end
% glossary-block:B00177:begin
\needspace{3\baselineskip}
\textbf{beast} — \textbf{dēor} (n.). animal; beast. \textbf{nīeten} (n.). animal; beast; cattle. Cattle is the selected plural use; retain the early task distinction from dēor. \textbf{dēor} (n.), \textbf{nīeten} (n.). Dēor is the Unit 1 animal/beast target. Nīeten can also denote an animal in later contexts; use the noun specified by the exercise.\par
% glossary-block:B00177:end
% glossary-block:B00178:begin
\needspace{3\baselineskip}
\textbf{be astonished} — \textbf{ofwundrian} (v.). Subject plus wesan and supplied participle, with degree/result tō ðan … ðæt.\par
% glossary-block:B00178:end
% glossary-block:B00179:begin
\needspace{3\baselineskip}
\textbf{beat} — \textbf{bēatan} (v.). beat; strike. Accusative person or thing struck. \textbf{geðrǣstan} (v.). afflict; beat. Supplied participial source use.\par
% glossary-block:B00179:end
% glossary-block:B00180:begin
\needspace{3\baselineskip}
\textbf{be at rest} — \textbf{licgan} (v.). lie; be at rest.\par
% glossary-block:B00180:end
% glossary-block:B00181:begin
\needspace{3\baselineskip}
\textbf{beautiful} — \textbf{fæger} (adj.). beautiful; fair.\par
% glossary-block:B00181:end
% glossary-block:B00182:begin
\needspace{3\baselineskip}
\textbf{be born} — \textbf{ācęnnan} (v.). Source supplies ācęnned in a passive expression; distinguish active beget/produce.\par
% glossary-block:B00182:end
% glossary-block:B00183:begin
\needspace{3\baselineskip}
\textbf{be called} — \textbf{hātan} (v.). be called; be named. Special hātte form in the assigned naming expression. The special form hātte means is/was called or named, according to context; do not substitute an ordinary active hātan form.\par
% glossary-block:B00183:end
% glossary-block:B00184:begin
\needspace{3\baselineskip}
\textbf{because} — \textbf{for-ðām} (conj.).  \textbf{for-ðām-ðe} (conj.).  \textbf{for-ðon-ðe} (conj.; adv. in for-ðon). Causal conjunction. \textbf{for-ðām} (conj.), \textbf{for-ðon-ðe} (conj.; adv. in for-ðon). For-ðām and for-ðon-ðe are causal conjunctions meaning because/for. They are distinct from prepositions meaning because of or for someone’s benefit.\par
% glossary-block:B00184:end
% glossary-block:B00185:begin
\needspace{3\baselineskip}
\textbf{because of} — \textbf{for} (prep.). Dative hindering circumstance. because of; on account of. for + dative cause in this supplied money expression; distinguish causal conjunction because.\par
% glossary-block:B00185:end
% glossary-block:B00186:begin
\needspace{3\baselineskip}
\textbf{become} — \textbf{geweorðan} (v.). come to be; become. become; come about. Passive weorðan combines with a past participle for the assigned event or becoming. Lexical weorðan become and copular wesan be have separate uses. \textbf{weorðan} (v.). Nominative subject and agreeing nominative predicate.\par
% glossary-block:B00186:end
% glossary-block:B00187:begin
\needspace{3\baselineskip}
\textbf{become calm} — \textbf{gesmyltan} (v.). Supplied preterite of the sea becoming calm.\par
% glossary-block:B00187:end
% glossary-block:B00188:begin
\needspace{3\baselineskip}
\textbf{become (passive)} — Passive wesan + past participle describes the assigned state/passive construction; passive weorðan + past participle describes becoming or the event. Simple lexical become is separate.\par
% glossary-block:B00188:end
% glossary-block:B00189:begin
\needspace{3\baselineskip}
\textbf{become still} — \textbf{gestillan} (v.). cease; become still. Intransitive in the source storm relative.\par
% glossary-block:B00189:end
% glossary-block:B00190:begin
\needspace{3\baselineskip}
\textbf{be doing} — The supplied progressive construction is wesan + present participle. Ordinary English be doing can also be rendered by a simple finite Old English verb; this compound is not the universal default.\par
% glossary-block:B00190:end
% glossary-block:B00191:begin
\needspace{3\baselineskip}
\textbf{be done} — Passive wesan combines with the assigned past participle; be done includes the participial meaning as well as the finite be auxiliary.\par
% glossary-block:B00191:end
% glossary-block:B00192:begin
\needspace{3\baselineskip}
\textbf{befall} — \textbf{becuman} (v.). come upon; befall. Retain the source's affected participant. \textbf{gelimpan} (v.). happen; befall. Supplied preterite subjunctive in might befall us. \textbf{geweorðan} (v.). come to pass; befall.\par
% glossary-block:B00192:end
% glossary-block:B00193:begin
\needspace{3\baselineskip}
\textbf{before} — \textbf{ǣr} (conj.). Ǣr introduces before; use the mood required by the taught factual or anticipated clause. \textbf{ǣr} (adv.). before; earlier. \textbf{ætforan} (prep.). before; in front of. Dative landmark/person in the supplied frame; retain ætforan rather than replacing it with beforan. \textbf{be-foran} (prep.). before; ahead of. Dative person, including going or sending ahead of that person. before; in front of. Dative landmark in the static frame. before; in front of. Accusative landmark in the positioning construal; static dative can be a different valid construal. \textbf{tōforan} (prep.). before; above all.\par
% glossary-block:B00193:end
% glossary-block:B00194:begin
\needspace{3\baselineskip}
\textbf{beget} — \textbf{ācęnnan} (v.). beget; bring forth.\par
% glossary-block:B00194:end
% glossary-block:B00195:begin
\needspace{3\baselineskip}
\textbf{begin} — \textbf{āginnan} (v.). Supplied source form began. \textbf{onginnan} (v.). Supplied preterite began. Supplied playing/praising action as complement.\par
% glossary-block:B00195:end
% glossary-block:B00196:begin
\needspace{3\baselineskip}
\textbf{beginning} — \textbf{anginn} (n.).\par
% glossary-block:B00196:end
% glossary-block:B00197:begin
\needspace{3\baselineskip}
\textbf{begotten of} — \textbf{of} (prep.). of; from. Origin of wisdom in the supplied begetting statement.\par
% glossary-block:B00197:end
% glossary-block:B00198:begin
\needspace{3\baselineskip}
\textbf{be gracious toward} — \textbf{glæd} (adj.). glæd in this contextual favour sense, wið + person; not merely emotionally glad.\par
% glossary-block:B00198:end
% glossary-block:B00199:begin
\needspace{3\baselineskip}
\textbf{be hale} — Gesund means well/in good health; hāl means hale/healthy. These adjectives describe the subject in be well/healthy; wel is the separate adverb for doing something well.\par
% glossary-block:B00199:end
% glossary-block:B00200:begin
\needspace{3\baselineskip}
\textbf{behave toward} — \textbf{dōn} (v.). deal with; behave toward. Assigned wið person frame; manner is separately expressed.\par
% glossary-block:B00200:end
% glossary-block:B00201:begin
\needspace{3\baselineskip}
\textbf{behead} — \textbf{behēafdian} (v.). Accusative person beheaded.\par
% glossary-block:B00201:end
% glossary-block:B00202:begin
\needspace{3\baselineskip}
\textbf{be healthy} — Gesund means well/in good health; hāl means hale/healthy. These adjectives describe the subject in be well/healthy; wel is the separate adverb for doing something well.\par
% glossary-block:B00202:end
% glossary-block:B00203:begin
\needspace{3\baselineskip}
\textbf{behold} — \textbf{behealdan} (v.). behold; look at. Accusative object seen. behealdan takes accusative objects; distinguish interjection efne behold. \textbf{efne} (interj.). behold!. Interjection, distinct from the verb behealdan look at.\par
% glossary-block:B00203:end
% glossary-block:B00204:begin
\needspace{3\baselineskip}
\textbf{be hungry} — \textbf{hyngran} (v.). Impersonal construction with accusative experiencer. Hyngran expresses be hungry impersonally, with an accusative experiencer. Alternatively, the adjective hungrig agrees with the subject of be. Use the construction named by the exercise.\par
% glossary-block:B00204:end
% glossary-block:B00205:begin
\needspace{3\baselineskip}
\textbf{be in astonishment} — \textbf{wafian} (v.). be amazed. Supplied source form.\par
% glossary-block:B00205:end
% glossary-block:B00206:begin
\needspace{3\baselineskip}
\textbf{be … -ing} — Simple OE finite forms can render English progressive aspect; wesan + present participle is only one bounded source construction, not a required default.\par
% glossary-block:B00206:end
% glossary-block:B00207:begin
\needspace{3\baselineskip}
\textbf{be it so} — The wish be it so uses independent third-singular subjunctive sīe with hit and swā. It does not use a permission verb.\par
% glossary-block:B00207:end
% glossary-block:B00208:begin
\needspace{3\baselineskip}
\textbf{believe} — \textbf{gelīefan} (v.).\par
% glossary-block:B00208:end
% glossary-block:B00209:begin
\needspace{3\baselineskip}
\textbf{believe in} — \textbf{on} (prep.). in. Accusative person believed/trusted in. \textbf{gelīefan} (v.). gelīefan on + accusative person in this assigned frame.\par
% glossary-block:B00209:end
% glossary-block:B00210:begin
\needspace{3\baselineskip}
\textbf{bellows} — \textbf{bęlg} (n.). Selected plural instrument use.\par
% glossary-block:B00210:end
% glossary-block:B00211:begin
\needspace{3\baselineskip}
\textbf{belong} — \textbf{belimpan} (v.). belong; relate. Supplied participial construction.\par
% glossary-block:B00211:end
% glossary-block:B00212:begin
\needspace{3\baselineskip}
\textbf{belong to} — Tō introduces the dependent in the supplied belonging construction; the verb supplies belong. Retain the complete supplied frame.\par
% glossary-block:B00212:end
% glossary-block:B00213:begin
\needspace{3\baselineskip}
\textbf{below} — \textbf{under} (prep.). under. Dative landmark in the taught static frame.\par
% glossary-block:B00213:end
% glossary-block:B00214:begin
\needspace{3\baselineskip}
\textbf{be lying} — \textbf{licgan} (v.). lie; be lying. Supplied participle describing the possessed persons, not a habban perfect.\par
% glossary-block:B00214:end
% glossary-block:B00215:begin
\needspace{3\baselineskip}
\textbf{be multiplied} — \textbf{gemanigfieldan} (v.). Assigned passive with agreeing gemanigfield/gemanigfielde; no independent active causative frame is assigned. Assigned passive use with wesan and agreeing gemanigfield/gemanigfielde; do not infer an unassigned active causative frame.\par
% glossary-block:B00215:end
% glossary-block:B00216:begin
\needspace{3\baselineskip}
\textbf{be named} — \textbf{hātan} (v.). be called; be named. Special hātte form in the assigned naming expression. The special form hātte means is/was called or named, according to context; do not substitute an ordinary active hātan form.\par
% glossary-block:B00216:end
% glossary-block:B00217:begin
\needspace{3\baselineskip}
\textbf{bend} — \textbf{bīegan} (v.). Supplied past bent; retain the source's unmarked ie in its actual form. \textbf{būgan} (v.).\par
% glossary-block:B00217:end
% glossary-block:B00218:begin
\needspace{3\baselineskip}
\textbf{benefit} — \textbf{fręmman} (v.). Dative beneficiary. Dative person benefited. \textbf{gōd} (n.). good; benefit; kindness. Nominal good in the repayment contrast. \textbf{ðearf} (n.). benefit; good; need. Source contextual nouns; retain the local expression.\par
% glossary-block:B00218:end
% glossary-block:B00219:begin
\needspace{3\baselineskip}
\textbf{benevolent} — \textbf{gōd} (adj.). kind; benevolent.\par
% glossary-block:B00219:end
% glossary-block:B00220:begin
\needspace{3\baselineskip}
\textbf{Benjamin} — \textbf{Beniamin} (proper name).\par
% glossary-block:B00220:end
% glossary-block:B00221:begin
\needspace{3\baselineskip}
\textbf{be obliged to} — Sculan expresses obligation, must/be obliged to, with a plain infinitive. Pure future reference does not automatically require sculan.\par
% glossary-block:B00221:end
% glossary-block:B00222:begin
\needspace{3\baselineskip}
\textbf{be (passive)} — Passive wesan + past participle describes the assigned state/passive construction; passive weorðan + past participle describes becoming or the event. Simple lexical become is separate.\par
% glossary-block:B00222:end
% glossary-block:B00223:begin
\needspace{3\baselineskip}
\textbf{be permitted} — \textbf{ālīefan} (v.). permit; allow. Impersonal permission frame, with dative permitted person; consult the complete supplied expression.\par
% glossary-block:B00223:end
% glossary-block:B00224:begin
\needspace{3\baselineskip}
\textbf{bereft of} — \textbf{berēafian} (v.). despoil; deprive of. Accusative person deprived, genitive thing lost; passive retains that genitive. Passive berēafod with genitive of what is lost; active person deprived is accusative.\par
% glossary-block:B00224:end
% glossary-block:B00225:begin
\needspace{3\baselineskip}
\textbf{beseech} — \textbf{biddan} (v.). beseech; entreat. Accusative person; the requested action can be a ðæt clause. Accusative person beseeched; ðæt + subjunctive gives the requested act. Prayer to God with gebiddan is a different entry.\par
% glossary-block:B00225:end
% glossary-block:B00226:begin
\needspace{3\baselineskip}
\textbf{beset} — \textbf{ymbsęllan} (v.). surround; beset. Supplied source form.\par
% glossary-block:B00226:end
% glossary-block:B00227:begin
\needspace{3\baselineskip}
\textbf{beside} — \textbf{wið} (prep.). by; beside; at. Accusative landmark in the taught proximity phrase.\par
% glossary-block:B00227:end
% glossary-block:B00228:begin
\needspace{3\baselineskip}
\textbf{besides} — \textbf{būtan} (prep.). besides; apart from. Dative excluded person or thing. \textbf{ðǣr-tō-ēacan} (adverbial expression). in addition to that; besides.\par
% glossary-block:B00228:end
% glossary-block:B00229:begin
\needspace{3\baselineskip}
\textbf{be silent} — \textbf{swigian} (v.). Supplied source sense.\par
% glossary-block:B00229:end
% glossary-block:B00230:begin
\needspace{3\baselineskip}
\textbf{best} — \textbf{bętst} (adj.).  Bętst is the superlative adverb best, the superlative of wel well. Its adjectival use is separately identified. For better, bętera is adjectival and bęt adverbial. For best, bętst is the corresponding adjectival or adverbial superlative in the taught uses. Choose degree and part of speech from the task.\par
% glossary-block:B00230:end
% glossary-block:B00231:begin
\needspace{3\baselineskip}
\textbf{bestowal} — \textbf{sęlen} (n.). bounty; bestowal.\par
% glossary-block:B00231:end
% glossary-block:B00232:begin
\needspace{3\baselineskip}
\textbf{be sufficient for} — \textbf{genyhtsumian} (v.). suffice; be sufficient for. Complete supplied construction with dative beneficiary. Supplied complete construction with dative beneficiary; plural books require plural subjunctive.\par
% glossary-block:B00232:end
% glossary-block:B00233:begin
\needspace{3\baselineskip}
\textbf{be to} — In the supplied interpretive question, sculan expresses what events should/were to mean. Retain that reported or prospective interpretation rather than substituting the ordinary obligation sense.\par
% glossary-block:B00233:end
% glossary-block:B00234:begin
\needspace{3\baselineskip}
\textbf{be to come} — \textbf{tō} (prep.). Whole supplied is tō cumanne; bounded be-to use, distinct from purpose after a verb of coming.\par
% glossary-block:B00234:end
% glossary-block:B00235:begin
\needspace{3\baselineskip}
\textbf{better} — \textbf{bętera} (adj.).  Bęt is the comparative adverb better, the comparative of wel well. Use bętera for an adjective describing a noun. For better, bętera is adjectival and bęt adverbial. For best, bętst is the corresponding adjectival or adverbial superlative in the taught uses. Choose degree and part of speech from the task.\par
% glossary-block:B00235:end
% glossary-block:B00236:begin
\needspace{3\baselineskip}
\textbf{between} — \textbf{betwēonan} (adverb or particle). between; to one another. Supplied reciprocal betwēonan him in 45.6. \textbf{betweox} (prep.). among; between.\par
% glossary-block:B00236:end
% glossary-block:B00237:begin
\needspace{3\baselineskip}
\textbf{be unwilling} — \textbf{willan} (v.). be unwilling; refuse. With ne or a licensed negative contraction; followed by plain infinitive. ne + willan, or the licensed contracted negative, with a plain infinitive.\par
% glossary-block:B00237:end
% glossary-block:B00238:begin
\needspace{3\baselineskip}
\textbf{be well} — Gesund means well/in good health; hāl means hale/healthy. These adjectives describe the subject in be well/healthy; wel is the separate adverb for doing something well.\par
% glossary-block:B00238:end
% glossary-block:B00239:begin
\needspace{3\baselineskip}
\textbf{be willing to} — \textbf{willan} (v.). willan + plain infinitive denotes volition; intention does not assert the action occurred.\par
% glossary-block:B00239:end
% glossary-block:B00240:begin
\needspace{3\baselineskip}
\textbf{be wondering} — \textbf{wundrian} (v.). wonder. Supplied source form.\par
% glossary-block:B00240:end
% glossary-block:B00241:begin
\needspace{3\baselineskip}
\textbf{be wont} — \textbf{gewunian} (v.). be accustomed; be wont.\par
% glossary-block:B00241:end
% glossary-block:B00242:begin
\needspace{3\baselineskip}
\textbf{be wont to} — \textbf{gewunian} (v.). gewunian + plain infinitive; no separate be auxiliary in the model.\par
% glossary-block:B00242:end
% glossary-block:B00243:begin
\needspace{3\baselineskip}
\textbf{be yet to come} — Supplied gīet sculon fīf means that five are yet to come, with be expressly omitted. This bounded ellipsis is not a general rule that English shall always takes sculan.\par
% glossary-block:B00243:end
% glossary-block:B00244:begin
\needspace{3\baselineskip}
\textbf{bid} — \textbf{bēodan} (v.). command; bid.\par
% glossary-block:B00244:end
% glossary-block:B00245:begin
\needspace{3\baselineskip}
\textbf{bid farewell} — \textbf{biddan} (v.). wish well; bid farewell. Supplied farewell formula. \textbf{grētan} (v.). greet; take leave of. Accusative person/place; greeting and farewell depend on context. grētan + accusative person/place; no separate take or bid is needed.\par
% glossary-block:B00245:end
% glossary-block:B00246:begin
\needspace{3\baselineskip}
\textbf{bind} — \textbf{bindan} (v.). bind; tie. Accusative person or thing bound.\par
% glossary-block:B00246:end
% glossary-block:B00247:begin
\needspace{3\baselineskip}
\textbf{bird} — \textbf{fugol} (n.). Fowl is collective/plural English here; use the required Old English plural.\par
% glossary-block:B00247:end
% glossary-block:B00248:begin
\needspace{3\baselineskip}
\textbf{birth} — \textbf{gebyrd} (n.).\par
% glossary-block:B00248:end
% glossary-block:B00249:begin
\needspace{3\baselineskip}
\textbf{bishop} — \textbf{bisceop} (n.).\par
% glossary-block:B00249:end
% glossary-block:B00250:begin
\needspace{3\baselineskip}
\textbf{black} — \textbf{blæc} (adj.).\par
% glossary-block:B00250:end
% glossary-block:B00251:begin
\needspace{3\baselineskip}
\textbf{blacksmith} — \textbf{smið} (n.).\par
% glossary-block:B00251:end
% glossary-block:B00252:begin
\needspace{3\baselineskip}
\textbf{blame} — \textbf{tǣlan} (v.). blame; censure. \textbf{wītan} (v.). blame; censure.\par
% glossary-block:B00252:end
% glossary-block:B00253:begin
\needspace{3\baselineskip}
\textbf{blameless} — \textbf{clǣne} (adj.). clean; pure; blameless. Source's blameless is contextual clǣne, not just physical clean.\par
% glossary-block:B00253:end
% glossary-block:B00254:begin
\needspace{3\baselineskip}
\textbf{bless} — \textbf{blētsian} (v.). Supplied third-singular preterite; God is the understood subject. \textbf{geblētsian} (v.).\par
% glossary-block:B00254:end
% glossary-block:B00255:begin
\needspace{3\baselineskip}
\textbf{blessed} — \textbf{ēadig} (adj.). happy; blessed.\par
% glossary-block:B00255:end
% glossary-block:B00256:begin
\needspace{3\baselineskip}
\textbf{blessing} — \textbf{blētsung} (n.).\par
% glossary-block:B00256:end
% glossary-block:B00257:begin
\needspace{3\baselineskip}
\textbf{bliss} — \textbf{drēam} (n.). joy; bliss.\par
% glossary-block:B00257:end
% glossary-block:B00258:begin
\needspace{3\baselineskip}
\textbf{bloom} — \textbf{blōstmian} (v.). blossom; bloom.\par
% glossary-block:B00258:end
% glossary-block:B00259:begin
\needspace{3\baselineskip}
\textbf{blossom} — \textbf{blōstma} (n.).  \textbf{blōstmian} (v.). blossom; bloom.\par
% glossary-block:B00259:end
% glossary-block:B00260:begin
\needspace{3\baselineskip}
\textbf{blow} — \textbf{blāwan} (v.). Intransitive in the assigned bellows sentence.\par
% glossary-block:B00260:end
% glossary-block:B00261:begin
\needspace{3\baselineskip}
\textbf{board} — \textbf{āstīgan} (v.). go aboard; board. Accusative ship boarded in the assigned frame. āstīgan takes the ship as direct accusative in these exercises.\par
% glossary-block:B00261:end
% glossary-block:B00262:begin
\needspace{3\baselineskip}
\textbf{bodily} — \textbf{līchamlic} (adj.).\par
% glossary-block:B00262:end
% glossary-block:B00263:begin
\needspace{3\baselineskip}
\textbf{body} — \textbf{līchama} (n.).\par
% glossary-block:B00263:end
% glossary-block:B00264:begin
\needspace{3\baselineskip}
\textbf{boil} — \textbf{sēoðan} (v.).  \textbf{weallan} (v.). surge; boil. Supplied present participle used adjectivally of the sea.\par
% glossary-block:B00264:end
% glossary-block:B00265:begin
\needspace{3\baselineskip}
\textbf{bold} — \textbf{beald} (adj.). bold; brave.\par
% glossary-block:B00265:end
% glossary-block:B00266:begin
\needspace{3\baselineskip}
\textbf{bone} — \textbf{bān} (n.). bone; ivory. Ivory is elpes bān, literally elephant's bone.\par
% glossary-block:B00266:end
% glossary-block:B00267:begin
\needspace{3\baselineskip}
\textbf{book} — \textbf{bōc} (n.).\par
% glossary-block:B00267:end
% glossary-block:B00268:begin
\needspace{3\baselineskip}
\textbf{book-Latin} — \textbf{Bōclæden} (n.). book-Latin; Latin.\par
% glossary-block:B00268:end
% glossary-block:B00269:begin
\needspace{3\baselineskip}
\textbf{border} — \textbf{gemǣre} (n.). boundary; border.\par
% glossary-block:B00269:end
% glossary-block:B00270:begin
\needspace{3\baselineskip}
\textbf{borders} — \textbf{landgemǣre} (n.). territory; borders.\par
% glossary-block:B00270:end
% glossary-block:B00271:begin
\needspace{3\baselineskip}
\textbf{both} — The assigned whole expression ǣgðer ge … ge means both … and, coordinating the two members. \textbf{bēgen} (numeral).\par
% glossary-block:B00271:end
% glossary-block:B00272:begin
\needspace{3\baselineskip}
\textbf{both ... and} — Ge … ge is the paired expression both … and; the two occurrences coordinate the two members.\par
% glossary-block:B00272:end
% glossary-block:B00273:begin
\needspace{3\baselineskip}
\textbf{both … and} — The assigned whole expression ǣgðer ge … ge means both … and, coordinating the two members. \textbf{ge} (conj.). Paired ge … ge. Ge … ge means both … and; ǣgðer ge … ge is the expanded paired expression. The adjective bēgen both is separate.\par
% glossary-block:B00273:end
% glossary-block:B00274:begin
\needspace{3\baselineskip}
\textbf{boundary} — \textbf{gemǣre} (n.). boundary; border.\par
% glossary-block:B00274:end
% glossary-block:B00275:begin
\needspace{3\baselineskip}
\textbf{bounty} — \textbf{sęlen} (n.). bounty; bestowal.\par
% glossary-block:B00275:end
% glossary-block:B00276:begin
\needspace{3\baselineskip}
\textbf{bow} — \textbf{ālūtan} (v.). bow; incline oneself. Reflexive meaning in the assigned use.\par
% glossary-block:B00276:end
% glossary-block:B00277:begin
\needspace{3\baselineskip}
\textbf{boy} — \textbf{cnapa} (n.). Masculine weak noun; the task's chosen word for lad.\par
% glossary-block:B00277:end
% glossary-block:B00278:begin
\needspace{3\baselineskip}
\textbf{brandish} — \textbf{bregdan} (v.). brandish; draw. \textbf{swingan} (v.). swing; brandish. Supplied past source form.\par
% glossary-block:B00278:end
% glossary-block:B00279:begin
\needspace{3\baselineskip}
\textbf{brave} — \textbf{beald} (adj.). bold; brave.\par
% glossary-block:B00279:end
% glossary-block:B00280:begin
\needspace{3\baselineskip}
\textbf{bread} — \textbf{hlāf} (n.). bread; loaf.\par
% glossary-block:B00280:end
% glossary-block:B00281:begin
\needspace{3\baselineskip}
\textbf{break} — \textbf{tōweorpan} (v.). break; overthrow. Break a commandment in the assigned use.\par
% glossary-block:B00281:end
% glossary-block:B00282:begin
\needspace{3\baselineskip}
\textbf{break up} — \textbf{tōberstan} (v.). break up; go to pieces.\par
% glossary-block:B00282:end
% glossary-block:B00283:begin
\needspace{3\baselineskip}
\textbf{brethren} — \textbf{brōðor} (n.). Brethren is an English plural alias, not another Old English lexeme.\par
% glossary-block:B00283:end
% glossary-block:B00284:begin
\needspace{3\baselineskip}
\textbf{bride} — \textbf{brȳd} (n.).\par
% glossary-block:B00284:end
% glossary-block:B00285:begin
\needspace{3\baselineskip}
\textbf{bridge} — \textbf{brycg} (n.).\par
% glossary-block:B00285:end
% glossary-block:B00286:begin
\needspace{3\baselineskip}
\textbf{brief} — \textbf{medmicel} (adj.). short; brief.\par
% glossary-block:B00286:end
% glossary-block:B00287:begin
\needspace{3\baselineskip}
\textbf{brightness} — \textbf{beorhtnes} (n.).\par
% glossary-block:B00287:end
% glossary-block:B00288:begin
\needspace{3\baselineskip}
\textbf{bring} — \textbf{bringan} (v.). Accusative thing or person brought; tō + dative destination/person, or the assigned bare dative recipient. Cook's explicit carry hint selects bringan with accusative presents and dative recipient. \textbf{gebringan} (v.). Supplied preterite with plural winds as subject; retain the prefix. \textbf{lǣdan} (v.). lead; bring; take. A person is the object in the assigned frame. \textbf{tōgelǣdan} (v.). Accusative thing brought, dative recipient.\par
% glossary-block:B00288:end
% glossary-block:B00289:begin
\needspace{3\baselineskip}
\textbf{bring about} — \textbf{dōn} (v.). cause; bring about. With a finite ðæt clause; its actor is nominative within that clause.\par
% glossary-block:B00289:end
% glossary-block:B00290:begin
\needspace{3\baselineskip}
\textbf{bring forth} — \textbf{ācęnnan} (v.). produce; bring forth. beget; bring forth. \textbf{forð} (adv.). forth. Adverbial particle with the bringing/producing verb. \textbf{tēon} (v.). bring forth; produce. Accusative thing produced with particle forð; use the assigned tēon forms. bring forth; produce. Supplied third-plural present subjunctive, with waters as subject.\par
% glossary-block:B00290:end
% glossary-block:B00291:begin
\needspace{3\baselineskip}
\textbf{Britain} — \textbf{Breoton} (proper name).\par
% glossary-block:B00291:end
% glossary-block:B00292:begin
\needspace{3\baselineskip}
\textbf{Britons} — \textbf{Bryttas} (ethnonym).\par
% glossary-block:B00292:end
% glossary-block:B00293:begin
\needspace{3\baselineskip}
\textbf{broad} — \textbf{brād} (adj.).\par
% glossary-block:B00293:end
% glossary-block:B00294:begin
\needspace{3\baselineskip}
\textbf{brother} — \textbf{brōðor} (n.). Brethren is an English plural alias, not another Old English lexeme.\par
% glossary-block:B00294:end
% glossary-block:B00295:begin
\needspace{3\baselineskip}
\textbf{buffet} — \textbf{gewǣgan} (v.). buffet; shake. Supplied participial source use.\par
% glossary-block:B00295:end
% glossary-block:B00296:begin
\needspace{3\baselineskip}
\textbf{build} — \textbf{gewyrcean} (v.). make; build. \textbf{wyrcean} (v.). make; build. Supplied unprefixed form; preserve explicit task controls. \textbf{wyrcean} (v.), \textbf{gewyrcean} (v.). Wyrcean and gewyrcean mean make/build in these senses. Preserve the prefixed or unprefixed verb selected by the exercise.\par
% glossary-block:B00296:end
% glossary-block:B00297:begin
\needspace{3\baselineskip}
\textbf{building} — \textbf{getimbrung} (n.). building; structure built.\par
% glossary-block:B00297:end
% glossary-block:B00298:begin
\needspace{3\baselineskip}
\textbf{burden} — \textbf{hęfignes} (n.).\par
% glossary-block:B00298:end
% glossary-block:B00299:begin
\needspace{3\baselineskip}
\textbf{burst} — \textbf{berstan} (v.).\par
% glossary-block:B00299:end
% glossary-block:B00300:begin
\needspace{3\baselineskip}
\textbf{but} — \textbf{ac} (conj.).  \textbf{sōðlice} (adv.). now; but. Discourse connective corresponding to source Latin autem, distinct from truly.\par
% glossary-block:B00300:end
% glossary-block:B00301:begin
\needspace{3\baselineskip}
\textbf{buy} — \textbf{bycgan} (v.). Accusative goods bought. \textbf{cēapian} (v.). trade; traffic. Traffic means trade; retain the supplied trading frame.\par
% glossary-block:B00301:end
% glossary-block:B00302:begin
\needspace{3\baselineskip}
\textbf{by} — \textbf{be} (prep.). Body part by which a person is led; supplied grasping phrase. \textbf{for} (prep.). by; for the sake of. Grounds of entreaty in the prayer; keep the supplied phrase. \textbf{ðurh, þurh} (prep.). through; by. Mediated action through a person in the supplied frame. Accusative oath standard. \textbf{wið} (prep.). by; beside; at. Accusative landmark in the taught proximity phrase.\par
% glossary-block:B00302:end
% glossary-block:B00303:begin
\needspace{3\baselineskip}
\textbf{by means of} — \textbf{mid} (prep.; adv. in fixed expressions). with; by means of. Dative instrument or means. with; by means of. \textbf{ðurh, þurh} (prep.). through; by means of. Supplied service/evidence expressions, including Ðurh hwæt.\par
% glossary-block:B00303:end
% glossary-block:B00304:begin
\needspace{3\baselineskip}
\textbf{by night} — On niht is the whole supplied temporal expression by night. Retain the exact phrase.\par
% glossary-block:B00304:end
% glossary-block:B00305:begin
\needspace{3\baselineskip}
\textbf{by permission} — \textbf{be} (prep.). Supplied permission expression.\par
% glossary-block:B00305:end
% glossary-block:B00306:begin
\needspace{3\baselineskip}
\textbf{by that} — \textbf{be} (prep.). by that; from that. Refers to the preceding proposition; supplied expression. by that; on that evidence. Evidence/inference frame, not an unrestricted source preposition.\par
% glossary-block:B00306:end
% glossary-block:B00307:begin
\needspace{3\baselineskip}
\textbf{by the same way} — Ðȳ ilcan wege is the supplied instrumental expression by the same way; no additional preposition is required under the exercise’s omission control.\par
% glossary-block:B00307:end
% glossary-block:B00308:begin
\needspace{3\baselineskip}
\textbf{by the time that} — \textbf{ongēan} (adv.; prep.; conj.). Complete supplied ongēan ðætte clause.\par
% glossary-block:B00308:end
% glossary-block:B00309:begin
\needspace{3\baselineskip}
\textbf{by what} — \textbf{ðurh, þurh} (prep.). through; by means of. Supplied service/evidence expressions, including Ðurh hwæt.\par
% glossary-block:B00309:end
\pdfbookmark[1]{C}{outline-B00310}
% glossary-block:B00310:begin
\needspace{5\baselineskip}
\subsection*{C}
% glossary-block:B00310:end
% glossary-block:B00311:begin
\textbf{call} — \textbf{cleopian} (v.). call; summon. Accusative person summoned; distinguish intransitive cry aloud. \textbf{hātan} (v.). call; name. \textbf{sęcgan} (v.). call; declare to be. Object with the stated agreeing predicate.\par
% glossary-block:B00311:end
% glossary-block:B00312:begin
\needspace{3\baselineskip}
\textbf{calm} — \textbf{smylte} (adj.). calm; smooth. \textbf{smyltnes} (n.). calm; calmness. \textbf{stilnes} (n.). stillness; calm. Supplied genitive singular.\par
% glossary-block:B00312:end
% glossary-block:B00313:begin
\needspace{3\baselineskip}
\textbf{calmest} — \textbf{smyltest} (adj.). very calm; calmest. Supplied superlative adjective describing the sea.\par
% glossary-block:B00313:end
% glossary-block:B00314:begin
\needspace{3\baselineskip}
\textbf{calmness} — \textbf{smyltnes} (n.). calm; calmness.\par
% glossary-block:B00314:end
% glossary-block:B00315:begin
\needspace{3\baselineskip}
\textbf{camel} — \textbf{olfend} (n.).\par
% glossary-block:B00315:end
% glossary-block:B00316:begin
\needspace{3\baselineskip}
\textbf{can} — Mugan expresses capability, can/be able to, with a plain infinitive. Learned ability can instead use cunnan; neither construction takes an inserted tō before that infinitive. The supplied mugan form expresses capability, can/could/might be able to, with a plain infinitive. It is not the permission sense of English may/might. \textbf{mugan} (v.), \textbf{cunnan} (v.). Mugan expresses ability or capability. Cunnan + plain infinitive emphasizes learned ability. A modal complement is a plain infinitive, without tō.\par
% glossary-block:B00316:end
% glossary-block:B00317:begin
\needspace{3\baselineskip}
\textbf{capacity} — \textbf{mǣð} (n.). measure; capacity.\par
% glossary-block:B00317:end
% glossary-block:B00318:begin
\needspace{3\baselineskip}
\textbf{cares} — \textbf{bisgu} (n.). occupations; cares. \textbf{woruldbisgu} (n.). worldly occupations; cares.\par
% glossary-block:B00318:end
% glossary-block:B00319:begin
\needspace{3\baselineskip}
\textbf{cargo} — \textbf{hlæst} (n.). wares; merchandise; cargo. Selected plural source use.\par
% glossary-block:B00319:end
% glossary-block:B00320:begin
\needspace{3\baselineskip}
\textbf{carpenter} — \textbf{trēowwyrhta} (n.).\par
% glossary-block:B00320:end
% glossary-block:B00321:begin
\needspace{3\baselineskip}
\textbf{carry} — \textbf{āberan} (v.). Supplied carried form; retain the ā-prefix. \textbf{beran} (v.). Accusative person or thing carried. \textbf{gefęrian} (v.). carry; bear. Supplied passive movement; was carried/borne. \textbf{lǣdan} (v.). carry; take. A person is conveyed in this supplied frame. \textbf{bringan} (v.). Cook's explicit carry hint selects bringan with accusative presents and dative recipient.\par
% glossary-block:B00321:end
% glossary-block:B00322:begin
\needspace{3\baselineskip}
\textbf{carry along} — \textbf{lǣdan} (v.). take; carry along.\par
% glossary-block:B00322:end
% glossary-block:B00323:begin
\needspace{3\baselineskip}
\textbf{carry out} — \textbf{dōn} (v.). fulfil; carry out. A word or command is fulfilled.\par
% glossary-block:B00323:end
% glossary-block:B00324:begin
\needspace{3\baselineskip}
\textbf{cast} — \textbf{weorpan} (v.). Accusative thing cast; a following destination has its own government.\par
% glossary-block:B00324:end
% glossary-block:B00325:begin
\needspace{3\baselineskip}
\textbf{cast a net} — \textbf{sęndan} (v.). sęndan takes the net as accusative object; on + accusative supplies the destination.\par
% glossary-block:B00325:end
% glossary-block:B00326:begin
\needspace{3\baselineskip}
\textbf{cast away} — \textbf{āweorpan} (v.). cast down; cast away.\par
% glossary-block:B00326:end
% glossary-block:B00327:begin
\needspace{3\baselineskip}
\textbf{cast down} — \textbf{āweorpan} (v.). cast down; cast away. Accusative people put down, of + dative point from which.\par
% glossary-block:B00327:end
% glossary-block:B00328:begin
\needspace{3\baselineskip}
\textbf{catch} — \textbf{fōn} (v.). seize; catch. Accusative patient in the selected frame. Accusative person overtaken in these tasks; the source's bracketed caught is the route. \textbf{gelǣccan} (v.). Supplied caught form.\par
% glossary-block:B00328:end
% glossary-block:B00329:begin
\needspace{3\baselineskip}
\textbf{cattle} — \textbf{nīeten} (n.). animal; beast; cattle. Cattle is the selected plural use; retain the early task distinction from dēor. The selected nīeten sense is plural cattle; dēor is the separate early animal/beast target.\par
% glossary-block:B00329:end
% glossary-block:B00330:begin
\needspace{3\baselineskip}
\textbf{cause} — \textbf{dōn} (v.). cause; bring about. With a finite ðæt clause; its actor is nominative within that clause. \textbf{intinga} (n.). cause; sake (in for … intingan). Within supplied for intingan expressions.\par
% glossary-block:B00330:end
% glossary-block:B00331:begin
\needspace{3\baselineskip}
\textbf{cause to be} — \textbf{gedōn} (v.). make; cause to be.\par
% glossary-block:B00331:end
% glossary-block:B00332:begin
\needspace{3\baselineskip}
\textbf{cause to stumble} — \textbf{ǣswician} (v.). offend; cause to stumble. Accusative person; ethical/metaphorical offence in the assigned use.\par
% glossary-block:B00332:end
% glossary-block:B00333:begin
\needspace{3\baselineskip}
\textbf{cease} — \textbf{blinnan} (v.). Supplied wæs blinnende periphrasis. \textbf{gestillan} (v.). cease; become still. Intransitive in the source storm relative. \textbf{geswīcan} (v.). Genitive thing ceased from. fail; cease. Intransitive without a genitive complement in the selected frame. Absolute command to the sea; distinct from cessation with a genitive complement.\par
% glossary-block:B00333:end
% glossary-block:B00334:begin
\needspace{3\baselineskip}
\textbf{cease from} — \textbf{geswīcan} (v.). cease. Genitive thing ceased from.\par
% glossary-block:B00334:end
% glossary-block:B00335:begin
\needspace{3\baselineskip}
\textbf{cease to hold} — \textbf{gewītan} (v.). pass away; cease to hold. Of a decree.\par
% glossary-block:B00335:end
% glossary-block:B00336:begin
\needspace{3\baselineskip}
\textbf{cedar} — \textbf{cēderbēam} (n.). cedar; cedar tree.\par
% glossary-block:B00336:end
% glossary-block:B00337:begin
\needspace{3\baselineskip}
\textbf{cedar tree} — \textbf{cēderbēam} (n.). cedar; cedar tree.\par
% glossary-block:B00337:end
% glossary-block:B00338:begin
\needspace{3\baselineskip}
\textbf{censure} — \textbf{tǣlan} (v.). blame; censure. \textbf{wītan} (v.). blame; censure.\par
% glossary-block:B00338:end
% glossary-block:B00339:begin
\needspace{3\baselineskip}
\textbf{chaff} — \textbf{ceaf} (n.).\par
% glossary-block:B00339:end
% glossary-block:B00340:begin
\needspace{3\baselineskip}
\textbf{chamber} — \textbf{būr} (n.). chamber; room.\par
% glossary-block:B00340:end
% glossary-block:B00341:begin
\needspace{3\baselineskip}
\textbf{chance} — \textbf{wēn} (n.). expectation; chance. wēn is ðæt + subjunctive supplies the taught perhaps construction.\par
% glossary-block:B00341:end
% glossary-block:B00342:begin
\needspace{3\baselineskip}
\textbf{chance upon} — \textbf{gefeallan} (v.).  \textbf{on} (prep.). upon; into. Supplied chancing-upon idiom.\par
% glossary-block:B00342:end
% glossary-block:B00343:begin
\needspace{3\baselineskip}
\textbf{change back} — \textbf{gewęndan} (v.). turn; change back. Supplied source form.\par
% glossary-block:B00343:end
% glossary-block:B00344:begin
\needspace{3\baselineskip}
\textbf{chastised} — \textbf{ðrēan} (v.). chastised; struck.\par
% glossary-block:B00344:end
% glossary-block:B00345:begin
\needspace{3\baselineskip}
\textbf{chief} — \textbf{ealdor} (n.). chief; king. This noun is distinct from the homonymous noun meaning life.\par
% glossary-block:B00345:end
% glossary-block:B00346:begin
\needspace{3\baselineskip}
\textbf{chief place} — \textbf{ealdordōm} (n.). chief place; primacy. \textbf{ealdorscipe} (n.). chief place; primacy.\par
% glossary-block:B00346:end
% glossary-block:B00347:begin
\needspace{3\baselineskip}
\textbf{child} — \textbf{bearn} (n.). son; child.\par
% glossary-block:B00347:end
% glossary-block:B00348:begin
\needspace{3\baselineskip}
\textbf{choose} — \textbf{cēosan} (v.). Accusative person or thing chosen. Accusative person selected.\par
% glossary-block:B00348:end
% glossary-block:B00349:begin
\needspace{3\baselineskip}
\textbf{Christ} — \textbf{Crīst} (n.).\par
% glossary-block:B00349:end
% glossary-block:B00350:begin
\needspace{3\baselineskip}
\textbf{church} — \textbf{cirice} (n.).\par
% glossary-block:B00350:end
% glossary-block:B00351:begin
\needspace{3\baselineskip}
\textbf{circuit} — \textbf{ymbhwyrft} (n.). Whole on ymbhwyrfte expression means all around; supplied dependent landmark.\par
% glossary-block:B00351:end
% glossary-block:B00352:begin
\needspace{3\baselineskip}
\textbf{citizens} — \textbf{ceasterwaran} (n.). townspeople; citizens.\par
% glossary-block:B00352:end
% glossary-block:B00353:begin
\needspace{3\baselineskip}
\textbf{city} — \textbf{burg} (n.). town; city. \textbf{ceaster} (n.). town; city.\par
% glossary-block:B00353:end
% glossary-block:B00354:begin
\needspace{3\baselineskip}
\textbf{clarion} — \textbf{biema} (n.). trumpet; clarion.\par
% glossary-block:B00354:end
% glossary-block:B00355:begin
\needspace{3\baselineskip}
\textbf{clean} — \textbf{clǣne} (adj.). pure; clean. clean; pure; blameless.\par
% glossary-block:B00355:end
% glossary-block:B00356:begin
\needspace{3\baselineskip}
\textbf{clearly} — Supplied superlative sweotolost means most clearly; do not use the positive sweotole for this degree. \textbf{sweotollice} (adv.).\par
% glossary-block:B00356:end
% glossary-block:B00357:begin
\needspace{3\baselineskip}
\textbf{cliff} — \textbf{clif} (n.).\par
% glossary-block:B00357:end
% glossary-block:B00358:begin
\needspace{3\baselineskip}
\textbf{cloak} — \textbf{sciccels} (n.).  \textbf{wǣfels} (n.).\par
% glossary-block:B00358:end
% glossary-block:B00359:begin
\needspace{3\baselineskip}
\textbf{close} — \textbf{geęndung} (n.). end; close.\par
% glossary-block:B00359:end
% glossary-block:B00360:begin
\needspace{3\baselineskip}
\textbf{clothe} — \textbf{begyrdan} (v.). gird; clothe. Supplied participial clothing expression. \textbf{bewǣfan} (v.). wrap; clothe. Supplied participial clothing expression. \textbf{gescrȳdan} (v.).  \textbf{scrȳdan} (v.). clothe; dress. Accusative person clothed in this introduced frame. \textbf{wrēon} (v.). cover; clothe.\par
% glossary-block:B00360:end
% glossary-block:B00361:begin
\needspace{3\baselineskip}
\textbf{clothed with} — \textbf{mid} (prep.; adv. in fixed expressions). with. Dative garment in the assigned clothing frame. with. Exact cloth/cloak complements of begyrd and bewǣfed.\par
% glossary-block:B00361:end
% glossary-block:B00362:begin
\needspace{3\baselineskip}
\textbf{clothing} — \textbf{gegierela} (n.). garment; clothing. \textbf{rēaf} (n.). raiment; clothing. clothing; garments. \textbf{rēaf} (n.), \textbf{gegierela} (n.). Wǣd denotes a garment, rēaf raiment/clothing, and gegierela a garment or clothing. Preserve any noun expressly supplied by the exercise.\par
% glossary-block:B00362:end
% glossary-block:B00363:begin
\needspace{3\baselineskip}
\textbf{cloud} — \textbf{wolcen} (n.).\par
% glossary-block:B00363:end
% glossary-block:B00364:begin
\needspace{3\baselineskip}
\textbf{coin} — \textbf{pęning} (n.). penny; coin.\par
% glossary-block:B00364:end
% glossary-block:B00365:begin
\needspace{3\baselineskip}
\textbf{collection} — \textbf{gegaderung} (n.). gathering; collection; congregation. Congregations of waters in the source.\par
% glossary-block:B00365:end
% glossary-block:B00366:begin
\needspace{3\baselineskip}
\textbf{come} — \textbf{becuman} (v.). come; arrive. \textbf{cuman} (v.).  \textbf{tōcuman} (v.). come; arrive.\par
% glossary-block:B00366:end
% glossary-block:B00367:begin
\needspace{3\baselineskip}
\textbf{come about} — \textbf{becuman} (v.). occur; come about. Of the calm in the source. \textbf{gelimpan} (v.). happen; come about. Supplied preterite source form. happen; come about. \textbf{geweorðan} (v.). become; come about. \textbf{weorðan} (v.). arise; come about. Intransitive with a nominative subject, including the arising of a calm. come about; result. Supplied resultative expression.\par
% glossary-block:B00367:end
% glossary-block:B00368:begin
\needspace{3\baselineskip}
\textbf{come nigh} — \textbf{nēalǣcan} (v.), \textbf{genēalǣcan} (v.). nēalǣcan and genēalǣcan mean approach/draw near; an expressed goal is dative. Preserve the verb selected by the exercise.\par
% glossary-block:B00368:end
% glossary-block:B00369:begin
\needspace{3\baselineskip}
\textbf{come out} — \textbf{ūtgangende} (v.). Supplied present participle; source permits its interpretation with earlier action.\par
% glossary-block:B00369:end
% glossary-block:B00370:begin
\needspace{3\baselineskip}
\textbf{come to an end} — \textbf{geęndian} (v.). end; come to an end.\par
% glossary-block:B00370:end
% glossary-block:B00371:begin
\needspace{3\baselineskip}
\textbf{come to be} — \textbf{geweorðan} (v.). come to be; become.\par
% glossary-block:B00371:end
% glossary-block:B00372:begin
\needspace{3\baselineskip}
\textbf{come to naught} — Weorðan tō nāhte is the supplied whole expression come to nothing/naught. Nāhte is the form in this result phrase; nāht alone means nothing.\par
% glossary-block:B00372:end
% glossary-block:B00373:begin
\needspace{3\baselineskip}
\textbf{come to nothing} — Weorðan tō nāhte is the supplied whole result expression come to nothing. Tō alone supplies only the result relation. Weorðan tō nāhte is the supplied whole expression come to nothing/naught. Nāhte is the form in this result phrase; nāht alone means nothing.\par
% glossary-block:B00373:end
% glossary-block:B00374:begin
\needspace{3\baselineskip}
\textbf{come to pass} — \textbf{geweorðan} (v.). come to pass; befall.\par
% glossary-block:B00374:end
% glossary-block:B00375:begin
\needspace{3\baselineskip}
\textbf{come up} — \textbf{ūpgān} (v.). go up; come up. From the water in the supplied expression.\par
% glossary-block:B00375:end
% glossary-block:B00376:begin
\needspace{3\baselineskip}
\textbf{come upon} — \textbf{becuman} (v.). come upon; befall. Retain the source's affected participant. \textbf{gestandan} (v.). assail; come upon. Supplied source sense.\par
% glossary-block:B00376:end
% glossary-block:B00377:begin
\needspace{3\baselineskip}
\textbf{comfort} — \textbf{frōfor} (n.). comfort; consolation.\par
% glossary-block:B00377:end
% glossary-block:B00378:begin
\needspace{3\baselineskip}
\textbf{command} — \textbf{bebēodan} (v.). command; instruct. Dative addressee and a ðæt clause with subjunctive for the commanded action. \textbf{bēodan} (v.). command; bid. \textbf{hǣs} (n.). command; order. \textbf{hātan} (v.). command; order. Accusative actor + plain infinitive for the commanded action; distinct from call/name.\par
% glossary-block:B00378:end
% glossary-block:B00379:begin
\needspace{3\baselineskip}
\textbf{commandment} — \textbf{bebod} (n.).\par
% glossary-block:B00379:end
% glossary-block:B00380:begin
\needspace{3\baselineskip}
\textbf{common} — \textbf{gemǣne} (adj.). common; shared.\par
% glossary-block:B00380:end
% glossary-block:B00381:begin
\needspace{3\baselineskip}
\textbf{companion} — \textbf{gefēra} (n.). companion; fellow.\par
% glossary-block:B00381:end
% glossary-block:B00382:begin
\needspace{3\baselineskip}
\textbf{companionship} — \textbf{gefērscipe} (n.).\par
% glossary-block:B00382:end
% glossary-block:B00383:begin
\needspace{3\baselineskip}
\textbf{company} — \textbf{flocc} (n.). group; company. \textbf{getæl} (n.). number; company. Singular collective getæl with the supplied dependent expression; retain source participial genitives. \textbf{męnigu} (n.). company; crowd. męnigu also denotes a multitude of things in the locally supplied spice expression.\par
% glossary-block:B00383:end
% glossary-block:B00384:begin
\needspace{3\baselineskip}
\textbf{compassion} — \textbf{mildheortnes} (n.).\par
% glossary-block:B00384:end
% glossary-block:B00385:begin
\needspace{3\baselineskip}
\textbf{compassionate} — \textbf{mildheort} (adj.). merciful; compassionate.\par
% glossary-block:B00385:end
% glossary-block:B00386:begin
\needspace{3\baselineskip}
\textbf{complete} — \textbf{fullfręmman} (v.). finish; complete. Supplied passive use.\par
% glossary-block:B00386:end
% glossary-block:B00387:begin
\needspace{3\baselineskip}
\textbf{completely} — Mid ealle is the fixed phrase completely; the individual adjective eall means all/whole. Mid ealle is the fixed phrase completely; mid alone does not have that meaning.\par
% glossary-block:B00387:end
% glossary-block:B00388:begin
\needspace{3\baselineskip}
\textbf{computation} — \textbf{gerīmcræft} (n.). arithmetic; computation.\par
% glossary-block:B00388:end
% glossary-block:B00389:begin
\needspace{3\baselineskip}
\textbf{conceal} — \textbf{diernan} (v.). conceal; hide. \textbf{forswīgian} (v.). keep silent about; conceal. \textbf{hȳdan} (v.). hide; conceal. Accusative thing hidden; wesan forms the assigned passive.\par
% glossary-block:B00389:end
% glossary-block:B00390:begin
\needspace{3\baselineskip}
\textbf{conceive} — \textbf{āsmēan} (v.). Exact supplied infinitive āsmēan in the heart-conceiving expression; not an independently assigned full paradigm. Supplied infinitive āsmēan; the heart is its subject.\par
% glossary-block:B00390:end
% glossary-block:B00391:begin
\needspace{3\baselineskip}
\textbf{concerning} — \textbf{be} (prep.). concerning; about. With dative.\par
% glossary-block:B00391:end
% glossary-block:B00392:begin
\needspace{3\baselineskip}
\textbf{condemn} — \textbf{dēman} (v.). Accusative person condemned.\par
% glossary-block:B00392:end
% glossary-block:B00393:begin
\needspace{3\baselineskip}
\textbf{confidence} — \textbf{trēownes} (n.). trust; confidence.\par
% glossary-block:B00393:end
% glossary-block:B00394:begin
\needspace{3\baselineskip}
\textbf{congregation} — \textbf{gegaderung} (n.). gathering; collection; congregation. Congregations of waters in the source.\par
% glossary-block:B00394:end
% glossary-block:B00395:begin
\needspace{3\baselineskip}
\textbf{conjure} — \textbf{hālsian} (v.). conjure; solemnly charge. Accusative person; ðurh + accusative oath standard; a requested ðæt clause may follow.\par
% glossary-block:B00395:end
% glossary-block:B00396:begin
\needspace{3\baselineskip}
\textbf{conjure by} — \textbf{ðurh, þurh} (prep.). by. Accusative oath standard.\par
% glossary-block:B00396:end
% glossary-block:B00397:begin
\needspace{3\baselineskip}
\textbf{conquer} — \textbf{oferwinnan} (v.). overcome; conquer.\par
% glossary-block:B00397:end
% glossary-block:B00398:begin
\needspace{3\baselineskip}
\textbf{consecrate} — \textbf{gehālgian} (v.). hallow; consecrate.\par
% glossary-block:B00398:end
% glossary-block:B00399:begin
\needspace{3\baselineskip}
\textbf{consider} — \textbf{smēagan} (v.). ponder; consider. Supplied preterite source form.\par
% glossary-block:B00399:end
% glossary-block:B00400:begin
\needspace{3\baselineskip}
\textbf{consolation} — \textbf{frōfor} (n.). comfort; consolation.\par
% glossary-block:B00400:end
% glossary-block:B00401:begin
\needspace{3\baselineskip}
\textbf{constant} — \textbf{singāl} (adj.). continual; constant.\par
% glossary-block:B00401:end
% glossary-block:B00402:begin
\needspace{3\baselineskip}
\textbf{container} — \textbf{fæt} (n.). vessel; container.\par
% glossary-block:B00402:end
% glossary-block:B00403:begin
\needspace{3\baselineskip}
\textbf{contend} — \textbf{feohtan} (v.). fight; contend. Intransitive; wið + accusative opponent.\par
% glossary-block:B00403:end
% glossary-block:B00404:begin
\needspace{3\baselineskip}
\textbf{continual} — \textbf{singāl} (adj.). continual; constant.\par
% glossary-block:B00404:end
% glossary-block:B00405:begin
\needspace{3\baselineskip}
\textbf{continue} — \textbf{ðurhwunian} (v.). Intransitive continuing of a building; distinguish the following dative-governing wiðstandan.\par
% glossary-block:B00405:end
% glossary-block:B00406:begin
\needspace{3\baselineskip}
\textbf{control} — \textbf{geweald} (n.). control; rule; dominion. Ordinary singular; on gewealde is the taught control phrase.\par
% glossary-block:B00406:end
% glossary-block:B00407:begin
\needspace{3\baselineskip}
\textbf{converse with} — \textbf{sprecan} (v.). speak with. mid + dative interlocutor in the taught production frame.\par
% glossary-block:B00407:end
% glossary-block:B00408:begin
\needspace{3\baselineskip}
\textbf{copper} — \textbf{ǣr} (n.).\par
% glossary-block:B00408:end
% glossary-block:B00409:begin
\needspace{3\baselineskip}
\textbf{correctly} — Supplied comparative rihtlicor means more correctly. Retain its source spelling without an inferred length mark.\par
% glossary-block:B00409:end
% glossary-block:B00410:begin
\needspace{3\baselineskip}
\textbf{costly} — \textbf{dēorwurð} (adj.). precious; costly.\par
% glossary-block:B00410:end
% glossary-block:B00411:begin
\needspace{3\baselineskip}
\textbf{could} — The supplied mugan form expresses capability, can/could/might be able to, with a plain infinitive. It is not the permission sense of English may/might. \textbf{mugan} (v.), \textbf{cunnan} (v.). Mugan expresses ability or capability. Cunnan + plain infinitive emphasizes learned ability. A modal complement is a plain infinitive, without tō.\par
% glossary-block:B00411:end
% glossary-block:B00412:begin
\needspace{3\baselineskip}
\textbf{counsellor} — \textbf{geðeahtend} (n.).\par
% glossary-block:B00412:end
% glossary-block:B00413:begin
\needspace{3\baselineskip}
\textbf{count} — \textbf{gerīman} (v.). number; count. Exact supplied infinitive gerīman; not an independently assigned full paradigm. Supplied infinitive gerīman, count; number.\par
% glossary-block:B00413:end
% glossary-block:B00414:begin
\needspace{3\baselineskip}
\textbf{countenance} — \textbf{andwlita} (n.). face; countenance.\par
% glossary-block:B00414:end
% glossary-block:B00415:begin
\needspace{3\baselineskip}
\textbf{country} — \textbf{land} (n.). land; country. \textbf{ðēod} (n.). region; country. \textbf{ēðel} (n.). ēðel in these home-country references; land remains a separate ordinary country word.\par
% glossary-block:B00415:end
% glossary-block:B00416:begin
\needspace{3\baselineskip}
\textbf{cover} — \textbf{oferwyrcean} (v.). overlay; cover. Accusative object, mid + dative material used to cover it. \textbf{wrēon} (v.). cover; clothe.\par
% glossary-block:B00416:end
% glossary-block:B00417:begin
\needspace{3\baselineskip}
\textbf{cowardly} — \textbf{earg} (adj.).\par
% glossary-block:B00417:end
% glossary-block:B00418:begin
\needspace{3\baselineskip}
\textbf{craft} — \textbf{cræft} (n.). trade; occupation; craft.\par
% glossary-block:B00418:end
% glossary-block:B00419:begin
\needspace{3\baselineskip}
\textbf{craftsman} — \textbf{wyrhta} (n.). workman; craftsman.\par
% glossary-block:B00419:end
% glossary-block:B00420:begin
\needspace{3\baselineskip}
\textbf{crag} — \textbf{torr} (n.). Probable contextual sense in Cook's note, distinct from earlier tower.\par
% glossary-block:B00420:end
% glossary-block:B00421:begin
\needspace{3\baselineskip}
\textbf{crash} — \textbf{gebrec} (n.). crash; noise; din. Supplied accusative singular.\par
% glossary-block:B00421:end
% glossary-block:B00422:begin
\needspace{3\baselineskip}
\textbf{create} — \textbf{gescieppan} (v.).  \textbf{scieppan} (v.). create; shape.\par
% glossary-block:B00422:end
% glossary-block:B00423:begin
\needspace{3\baselineskip}
\textbf{created being or thing} — \textbf{gesceaft} (n.). creature; created being or thing.\par
% glossary-block:B00423:end
% glossary-block:B00424:begin
\needspace{3\baselineskip}
\textbf{Creator} — \textbf{scieppend} (n.).  \textbf{wyrhta} (n.). Maker; Creator. Divine title, distinct from ordinary workman.\par
% glossary-block:B00424:end
% glossary-block:B00425:begin
\needspace{3\baselineskip}
\textbf{creature} — \textbf{gesceaft} (n.). creature; created being or thing.\par
% glossary-block:B00425:end
% glossary-block:B00426:begin
\needspace{3\baselineskip}
\textbf{crops} — \textbf{ierð} (n.). crops; produce of the earth.\par
% glossary-block:B00426:end
% glossary-block:B00427:begin
\needspace{3\baselineskip}
\textbf{cross} — \textbf{rōd} (n.).\par
% glossary-block:B00427:end
% glossary-block:B00428:begin
\needspace{3\baselineskip}
\textbf{crowd} — \textbf{męnigu} (n.). company; crowd. męnigu also denotes a multitude of things in the locally supplied spice expression.\par
% glossary-block:B00428:end
% glossary-block:B00429:begin
\needspace{3\baselineskip}
\textbf{crown} — \textbf{cynehelm} (n.).\par
% glossary-block:B00429:end
% glossary-block:B00430:begin
\needspace{3\baselineskip}
\textbf{cruelly} — \textbf{earmlice} (adverb or particle). cruelly; miserably.\par
% glossary-block:B00430:end
% glossary-block:B00431:begin
\needspace{3\baselineskip}
\textbf{cry} — \textbf{cleopian} (v.). cry aloud. Intransitive, with the words cried given as direct speech.\par
% glossary-block:B00431:end
% glossary-block:B00432:begin
\needspace{3\baselineskip}
\textbf{cry aloud} — \textbf{cleopian} (v.). Intransitive, with the words cried given as direct speech.\par
% glossary-block:B00432:end
% glossary-block:B00433:begin
\needspace{3\baselineskip}
\textbf{culpable} — \textbf{scyldig} (adj.). guilty; culpable.\par
% glossary-block:B00433:end
% glossary-block:B00434:begin
\needspace{3\baselineskip}
\textbf{cup} — \textbf{læfel} (n.).\par
% glossary-block:B00434:end
% glossary-block:B00435:begin
\needspace{3\baselineskip}
\textbf{custom} — \textbf{gewuna} (n.). custom; habit.\par
% glossary-block:B00435:end
% glossary-block:B00436:begin
\needspace{3\baselineskip}
\textbf{cut} — \textbf{snīðan} (v.).\par
% glossary-block:B00436:end
% glossary-block:B00437:begin
\needspace{3\baselineskip}
\textbf{cut down} — \textbf{forceorfan} (v.). cut down; cut off. Accusative tree or thing cut.\par
% glossary-block:B00437:end
% glossary-block:B00438:begin
\needspace{3\baselineskip}
\textbf{cut off} — \textbf{āceorfan} (v.). Accusative thing cut; of in the assigned expression is a separate off particle. \textbf{forceorfan} (v.). cut down; cut off. Accusative tree or thing cut.\par
% glossary-block:B00438:end
% glossary-block:B00439:begin
\needspace{3\baselineskip}
\textbf{Cyrenian} — \textbf{Cyrenisc} (adj.). Cyrenian; of Cyrene. Geographical adjective.\par
% glossary-block:B00439:end
\pdfbookmark[1]{D}{outline-B00440}
% glossary-block:B00440:begin
\needspace{5\baselineskip}
\subsection*{D}
% glossary-block:B00440:end
% glossary-block:B00441:begin
\textbf{daily} — \textbf{dæghwāmlīc} (adj.).\par
% glossary-block:B00441:end
% glossary-block:B00442:begin
\needspace{3\baselineskip}
\textbf{Danes} — \textbf{Dęne} (n.). Plural people-name.\par
% glossary-block:B00442:end
% glossary-block:B00443:begin
\needspace{3\baselineskip}
\textbf{danger} — \textbf{frēcnes} (n.). danger; peril.\par
% glossary-block:B00443:end
% glossary-block:B00444:begin
\needspace{3\baselineskip}
\textbf{dare to} — \textbf{durran} (v.). durran + plain infinitive; ne dearr denies daring.\par
% glossary-block:B00444:end
% glossary-block:B00445:begin
\needspace{3\baselineskip}
\textbf{darkness} — \textbf{ðīestru} (n.). Grammatically plural ðīestru; no singular is assigned.\par
% glossary-block:B00445:end
% glossary-block:B00446:begin
\needspace{3\baselineskip}
\textbf{daughter} — \textbf{dohtor} (n.).\par
% glossary-block:B00446:end
% glossary-block:B00447:begin
\needspace{3\baselineskip}
\textbf{day} — \textbf{dæg} (n.).\par
% glossary-block:B00447:end
% glossary-block:B00448:begin
\needspace{3\baselineskip}
\textbf{Day of Judgment} — \textbf{dōm} (n.). judgment. dōmes dæg means Day of Judgment.\par
% glossary-block:B00448:end
% glossary-block:B00449:begin
\needspace{3\baselineskip}
\textbf{dead} — \textbf{dēad} (adj.).\par
% glossary-block:B00449:end
% glossary-block:B00450:begin
\needspace{3\baselineskip}
\textbf{dead people} — \textbf{dēad} (adj.). Substantival adjective.\par
% glossary-block:B00450:end
% glossary-block:B00451:begin
\needspace{3\baselineskip}
\textbf{deal ill with} — Use the whole expression dōn earmlice wið + person, deal cruelly/ill with that person. Earmlice expresses the manner; it is not a second verb meaning deal.\par
% glossary-block:B00451:end
% glossary-block:B00452:begin
\needspace{3\baselineskip}
\textbf{deal with} — \textbf{dōn} (v.). deal with; behave toward. Assigned wið person frame; manner is separately expressed.\par
% glossary-block:B00452:end
% glossary-block:B00453:begin
\needspace{3\baselineskip}
\textbf{dear} — \textbf{lēof} (adj.).\par
% glossary-block:B00453:end
% glossary-block:B00454:begin
\needspace{3\baselineskip}
\textbf{dearer} — \textbf{lēofra} (adj.).\par
% glossary-block:B00454:end
% glossary-block:B00455:begin
\needspace{3\baselineskip}
\textbf{dearest} — \textbf{lēofost} (adj.).\par
% glossary-block:B00455:end
% glossary-block:B00456:begin
\needspace{3\baselineskip}
\textbf{death} — \textbf{dēað} (n.).\par
% glossary-block:B00456:end
% glossary-block:B00457:begin
\needspace{3\baselineskip}
\textbf{debt} — \textbf{gylt} (n.). Includes the prayer's moral debt.\par
% glossary-block:B00457:end
% glossary-block:B00458:begin
\needspace{3\baselineskip}
\textbf{debtor} — \textbf{gyltend} (n.).\par
% glossary-block:B00458:end
% glossary-block:B00459:begin
\needspace{3\baselineskip}
\textbf{declare to be} — \textbf{sęcgan} (v.). call; declare to be. Object with the stated agreeing predicate.\par
% glossary-block:B00459:end
% glossary-block:B00460:begin
\needspace{3\baselineskip}
\textbf{deed} — \textbf{dǣd} (n.).  \textbf{weorc} (n.). weorc includes deed/action in these prompts; dǣd is a separate lexical option elsewhere.\par
% glossary-block:B00460:end
% glossary-block:B00461:begin
\needspace{3\baselineskip}
\textbf{deeds} — \textbf{weorc} (n.). weorc includes deed/action in these prompts; dǣd is a separate lexical option elsewhere.\par
% glossary-block:B00461:end
% glossary-block:B00462:begin
\needspace{3\baselineskip}
\textbf{defence} — \textbf{gebeorg} (n.). protection; defence.\par
% glossary-block:B00462:end
% glossary-block:B00463:begin
\needspace{3\baselineskip}
\textbf{deity} — \textbf{God} (n.). god; deity. Common deity sense, distinct from the proper divine name.\par
% glossary-block:B00463:end
% glossary-block:B00464:begin
\needspace{3\baselineskip}
\textbf{delay} — \textbf{ęlcung} (n.).\par
% glossary-block:B00464:end
% glossary-block:B00465:begin
\needspace{3\baselineskip}
\textbf{deliver} — \textbf{ālīesan} (v.). deliver; release. Accusative person delivered or released; fram + dative marks separation.\par
% glossary-block:B00465:end
% glossary-block:B00466:begin
\needspace{3\baselineskip}
\textbf{demon} — \textbf{dēofol} (n.). devil; demon.\par
% glossary-block:B00466:end
% glossary-block:B00467:begin
\needspace{3\baselineskip}
\textbf{deny} — \textbf{forwiernan} (v.). deny; refuse. Dative person denied, genitive thing denied. Dative person denied, genitive thing denied.\par
% glossary-block:B00467:end
% glossary-block:B00468:begin
\needspace{3\baselineskip}
\textbf{depart} — \textbf{āgān} (v.). go away. Supplied gone-away form; distinct from āgan possess. \textbf{gewītan} (v.). depart; go away. Intransitive; fram + dative departure point is optional. The assigned frame requires no reflexive.\par
% glossary-block:B00468:end
% glossary-block:B00469:begin
\needspace{3\baselineskip}
\textbf{departure} — \textbf{forðfōr} (n.).\par
% glossary-block:B00469:end
% glossary-block:B00470:begin
\needspace{3\baselineskip}
\textbf{deprived of} — \textbf{berēafian} (v.). Passive berēafod with genitive of what is lost; active person deprived is accusative.\par
% glossary-block:B00470:end
% glossary-block:B00471:begin
\needspace{3\baselineskip}
\textbf{deprive of} — \textbf{berēafian} (v.). despoil; deprive of. Accusative person deprived, genitive thing lost; passive retains that genitive.\par
% glossary-block:B00471:end
% glossary-block:B00472:begin
\needspace{3\baselineskip}
\textbf{descend} — \textbf{āstīgan} (v.). Intransitive descend; of + dative departure point.\par
% glossary-block:B00472:end
% glossary-block:B00473:begin
\needspace{3\baselineskip}
\textbf{descendants} — \textbf{sǣd} (n.). seed; descendants. Collective descendants, distinct from literal seed.\par
% glossary-block:B00473:end
% glossary-block:B00474:begin
\needspace{3\baselineskip}
\textbf{deserve} — \textbf{āwyrcean} (v.). do; deserve. In the source, deserving through one's conduct; retain the complete local frame.\par
% glossary-block:B00474:end
% glossary-block:B00475:begin
\needspace{3\baselineskip}
\textbf{desire} — \textbf{gewilnian} (v.). desire; wish. \textbf{giernan} (v.). desire; request; woo. Genitive desired thing/person in the selected frames, including the maiden wooed. \textbf{lystan} (v.). please; be desired. Impersonal source expression; keep the complete supplied frame.\par
% glossary-block:B00475:end
% glossary-block:B00476:begin
\needspace{3\baselineskip}
\textbf{despair} — \textbf{ōrmōdnes} (n.). despair; hopelessness.\par
% glossary-block:B00476:end
% glossary-block:B00477:begin
\needspace{3\baselineskip}
\textbf{despise} — \textbf{rēcelēasian} (v.). neglect; despise.\par
% glossary-block:B00477:end
% glossary-block:B00478:begin
\needspace{3\baselineskip}
\textbf{despoil} — \textbf{berēafian} (v.). despoil; deprive of. Accusative person deprived, genitive thing lost; passive retains that genitive.\par
% glossary-block:B00478:end
% glossary-block:B00479:begin
\needspace{3\baselineskip}
\textbf{despoiled of} — \textbf{berēafian} (v.). Passive berēafod with genitive of what is lost; active person deprived is accusative.\par
% glossary-block:B00479:end
% glossary-block:B00480:begin
\needspace{3\baselineskip}
\textbf{destroy} — \textbf{fordilgian} (v.). destroy; annihilate. \textbf{fordōn} (v.).  \textbf{forniman} (v.). destroy; lay waste.\par
% glossary-block:B00480:end
% glossary-block:B00481:begin
\needspace{3\baselineskip}
\textbf{destruction} — \textbf{lor} (n.). destruction; ruin.\par
% glossary-block:B00481:end
% glossary-block:B00482:begin
\needspace{3\baselineskip}
\textbf{devil} — \textbf{dēofol} (n.). devil; demon.\par
% glossary-block:B00482:end
% glossary-block:B00483:begin
\needspace{3\baselineskip}
\textbf{did} — Dōn can replace an action already mentioned, as do/did in a comparison. This substitute verb is distinct from English auxiliary do in questions and negatives.\par
% glossary-block:B00483:end
% glossary-block:B00484:begin
\needspace{3\baselineskip}
\textbf{did (question or negative helper)} — No dummy dōn: inflect the lexical verb, with ne for negation and the taught question order. Lexical do/perform and substitute dōn are separate.\par
% glossary-block:B00484:end
% glossary-block:B00485:begin
\needspace{3\baselineskip}
\textbf{did (repeat an action)} — \textbf{dōn} (v.). dōn can replace the action already mentioned in a comparison; distinguish modern English dummy do.\par
% glossary-block:B00485:end
% glossary-block:B00486:begin
\needspace{3\baselineskip}
\textbf{dignity} — \textbf{horn} (n.). horn; strength; dignity. Strength/dignity is the figurative force in 37.14, not a replacement for the literal source horn. \textbf{weorðscipe} (n.). honour; dignity.\par
% glossary-block:B00486:end
% glossary-block:B00487:begin
\needspace{3\baselineskip}
\textbf{direct} — \textbf{āðęnian} (v.). direct; stretch toward. Directing the mind toward what is heard; supplied figurative phrase. \textbf{getǣcan} (v.). show; direct. Supplied source form. \textbf{gewissian} (v.). guide; direct.\par
% glossary-block:B00487:end
% glossary-block:B00488:begin
\needspace{3\baselineskip}
\textbf{dirty} — \textbf{horg} (adj.). dirty; filthy.\par
% glossary-block:B00488:end
% glossary-block:B00489:begin
\needspace{3\baselineskip}
\textbf{disciple} — \textbf{leornungcniht} (n.).\par
% glossary-block:B00489:end
% glossary-block:B00490:begin
\needspace{3\baselineskip}
\textbf{discourse} — \textbf{spell} (n.). account; discourse.\par
% glossary-block:B00490:end
% glossary-block:B00491:begin
\needspace{3\baselineskip}
\textbf{dish} — \textbf{disc} (n.).\par
% glossary-block:B00491:end
% glossary-block:B00492:begin
\needspace{3\baselineskip}
\textbf{dismiss} — \textbf{forlǣtan} (v.). let go; dismiss.\par
% glossary-block:B00492:end
% glossary-block:B00493:begin
\needspace{3\baselineskip}
\textbf{disorder} — \textbf{unstilnes} (n.). unrest; disorder.\par
% glossary-block:B00493:end
% glossary-block:B00494:begin
\needspace{3\baselineskip}
\textbf{display} — \textbf{ætīewan} (v.). show; display; reveal.\par
% glossary-block:B00494:end
% glossary-block:B00495:begin
\needspace{3\baselineskip}
\textbf{displease} — \textbf{ofðyncean} (v.). grieve; displease. Impersonal third singular with dative experiencer and a content clause.\par
% glossary-block:B00495:end
% glossary-block:B00496:begin
\needspace{3\baselineskip}
\textbf{disposition} — \textbf{willa} (n.). will; disposition.\par
% glossary-block:B00496:end
% glossary-block:B00497:begin
\needspace{3\baselineskip}
\textbf{dispute} — \textbf{geflīt} (n.). dispute; strife.\par
% glossary-block:B00497:end
% glossary-block:B00498:begin
\needspace{3\baselineskip}
\textbf{distinguish} — \textbf{geweorðian} (v.). distinguish; adorn. Accusative person adorned; mid + dative equipment or adornment. \textbf{tōdǣlen / tōdǣlden (source forms)} (verb; supplied inflected forms). separate; distinguish. Source forms tōdǣlen/tōdǣlden are retained; no silently invented citation spelling.\par
% glossary-block:B00498:end
% glossary-block:B00499:begin
\needspace{3\baselineskip}
\textbf{distress} — \textbf{gewinn} (n.). distress; struggle. \textbf{sorg} (n.). distress; anxiety; trouble; sorrow.\par
% glossary-block:B00499:end
% glossary-block:B00500:begin
\needspace{3\baselineskip}
\textbf{divide} — \textbf{gedǣlan} (v.). divide; separate. Preserve the assigned separation frame. \textbf{tōsceadan} (v.). separate; divide. \textbf{tōtwǣman} (v.). divide; separate. Waters divided from waters in the supplied frame.\par
% glossary-block:B00500:end
% glossary-block:B00501:begin
\needspace{3\baselineskip}
\textbf{divine favour} — \textbf{giefu} (n.). grace; divine favour.\par
% glossary-block:B00501:end
% glossary-block:B00502:begin
\needspace{3\baselineskip}
\textbf{do} — \textbf{āwyrcean} (v.). do; deserve. In the source, deserving through one's conduct; retain the complete local frame. Lexical dōn means do/perform. English auxiliary do/does/did in questions or negatives has no matching dōn: inflect the lexical verb and use the taught question order or ne. \textbf{dōn} (v.). do for; do to. Dative beneficiary/affected person; thing done is accusative when expressed. dōn can replace the action already mentioned in a comparison; distinguish modern English dummy do. Dōn can replace an action already mentioned, as do/did in a comparison. This substitute verb is distinct from English auxiliary do in questions and negatives. \textbf{gedōn} (v.). do; make. Supplied passive fulfilment formula. do; perform. No dummy dōn: inflect the lexical verb, with ne for negation and the taught question order. Lexical do/perform and substitute dōn are separate.\par
% glossary-block:B00502:end
% glossary-block:B00503:begin
\needspace{3\baselineskip}
\textbf{does} — No dummy dōn: inflect the lexical verb, with ne for negation and the taught question order. Lexical do/perform and substitute dōn are separate.\par
% glossary-block:B00503:end
% glossary-block:B00504:begin
\needspace{3\baselineskip}
\textbf{do for} — \textbf{dōn} (v.). do for; do to. Dative beneficiary/affected person; thing done is accusative when expressed.\par
% glossary-block:B00504:end
% glossary-block:B00505:begin
\needspace{3\baselineskip}
\textbf{dog} — \textbf{hund} (n.).\par
% glossary-block:B00505:end
% glossary-block:B00506:begin
\needspace{3\baselineskip}
\textbf{do good to} — \textbf{fręmman} (v.). benefit. Dative beneficiary.\par
% glossary-block:B00506:end
% glossary-block:B00507:begin
\needspace{3\baselineskip}
\textbf{dolphin} — \textbf{męreswīn} (n.).\par
% glossary-block:B00507:end
% glossary-block:B00508:begin
\needspace{3\baselineskip}
\textbf{dominion} — \textbf{geweald} (n.). control; rule; dominion. Ordinary singular; on gewealde is the taught control phrase. \textbf{rīce} (n.). Genitive rīces after geswīcan in ceased to bear rule. Do not insert a literal bear verb.\par
% glossary-block:B00508:end
% glossary-block:B00509:begin
\needspace{3\baselineskip}
\textbf{donkey} — \textbf{assa} (n.). ass; donkey.\par
% glossary-block:B00509:end
% glossary-block:B00510:begin
\needspace{3\baselineskip}
\textbf{doom} — \textbf{dōm} (n.). doom; fate.\par
% glossary-block:B00510:end
% glossary-block:B00511:begin
\needspace{3\baselineskip}
\textbf{do to} — \textbf{dōn} (v.). do for; do to. Dative beneficiary/affected person; thing done is accusative when expressed.\par
% glossary-block:B00511:end
% glossary-block:B00512:begin
\needspace{3\baselineskip}
\textbf{double} — \textbf{twīfeald} (adj.). double; twofold.\par
% glossary-block:B00512:end
% glossary-block:B00513:begin
\needspace{3\baselineskip}
\textbf{dragon} — \textbf{draca} (n.).\par
% glossary-block:B00513:end
% glossary-block:B00514:begin
\needspace{3\baselineskip}
\textbf{draw} — \textbf{bregdan} (v.). brandish; draw. \textbf{tēon} (v.). draw; pull.\par
% glossary-block:B00514:end
% glossary-block:B00515:begin
\needspace{3\baselineskip}
\textbf{draw near} — \textbf{genēalǣcan} (v.). approach; draw near. Dative person, place or event approached. approach. Supplied source form approached. \textbf{nēalǣcan} (v.). approach; draw near. Dative person or place approached when the goal is expressed. \textbf{nēalǣcan} (v.), \textbf{genēalǣcan} (v.). nēalǣcan and genēalǣcan mean approach/draw near; an expressed goal is dative. Preserve the verb selected by the exercise.\par
% glossary-block:B00515:end
% glossary-block:B00516:begin
\needspace{3\baselineskip}
\textbf{draw nigh} — \textbf{nēalǣcan} (v.), \textbf{genēalǣcan} (v.). nēalǣcan and genēalǣcan mean approach/draw near; an expressed goal is dative. Preserve the verb selected by the exercise.\par
% glossary-block:B00516:end
% glossary-block:B00517:begin
\needspace{3\baselineskip}
\textbf{dreadful} — \textbf{ęgeslic} (adj.). dreadful; terrible.\par
% glossary-block:B00517:end
% glossary-block:B00518:begin
\needspace{3\baselineskip}
\textbf{dress} — \textbf{scrȳdan} (v.). clothe; dress. Accusative person clothed in this introduced frame.\par
% glossary-block:B00518:end
% glossary-block:B00519:begin
\needspace{3\baselineskip}
\textbf{dress oneself} — \textbf{scrȳdan} (v.). Supplied reflexive use.\par
% glossary-block:B00519:end
% glossary-block:B00520:begin
\needspace{3\baselineskip}
\textbf{drink} — \textbf{dręnc} (n.).  \textbf{drincan} (v.). Accusative liquid drunk.\par
% glossary-block:B00520:end
% glossary-block:B00521:begin
\needspace{3\baselineskip}
\textbf{drive} — \textbf{drīfan} (v.). Accusative person or thing driven.\par
% glossary-block:B00521:end
% glossary-block:B00522:begin
\needspace{3\baselineskip}
\textbf{drive out} — \textbf{ādrīfan} (v.). drive out; expel. Accusative person expelled; fram + dative may mark the point of separation.\par
% glossary-block:B00522:end
% glossary-block:B00523:begin
\needspace{3\baselineskip}
\textbf{dry land} — \textbf{drȳgnes} (n.).\par
% glossary-block:B00523:end
% glossary-block:B00524:begin
\needspace{3\baselineskip}
\textbf{during} — \textbf{on} (prep.). during; in. Supplied lifetime expression. Dative time expression.\par
% glossary-block:B00524:end
% glossary-block:B00525:begin
\needspace{3\baselineskip}
\textbf{dust} — \textbf{dūst} (n.).\par
% glossary-block:B00525:end
% glossary-block:B00526:begin
\needspace{3\baselineskip}
\textbf{dwell} — \textbf{wunian} (v.). dwell; remain.\par
% glossary-block:B00526:end
% glossary-block:B00527:begin
\needspace{3\baselineskip}
\textbf{dwelling} — \textbf{hām} (n.). home; dwelling. \textbf{hof} (n.).  \textbf{wīc} (n.). dwelling; abode.\par
% glossary-block:B00527:end
\pdfbookmark[1]{E}{outline-B00528}
% glossary-block:B00528:begin
\needspace{5\baselineskip}
\subsection*{E}
% glossary-block:B00528:end
% glossary-block:B00529:begin
\textbf{each} — \textbf{ǣghwilc} (indefinite adjective). every; each. each; every. \textbf{ǣlc} (adj.). each; every. With singular noun; substantival each follows the case required by its clause. \textbf{gehwilc} (indefinite adjective). each; every. Ǣlc and gehwilc mean each/every with a singular noun. Personal gehwā means each person/everyone; neuter gehwæt refers to things. Select case from the clause role.\par
% glossary-block:B00529:end
% glossary-block:B00530:begin
\needspace{3\baselineskip}
\textbf{each one} — \textbf{gehwā / gehwæt} (pron.). each one; everyone; everything. Personal gehwā and neuter gehwæt; the early reference status does not by itself establish production.\par
% glossary-block:B00530:end
% glossary-block:B00531:begin
\needspace{3\baselineskip}
\textbf{ear} — \textbf{ēare} (n.).\par
% glossary-block:B00531:end
% glossary-block:B00532:begin
\needspace{3\baselineskip}
\textbf{earlier} — \textbf{ǣr} (adv.). before; earlier. \textbf{ǣrra} (adj.). former; earlier.\par
% glossary-block:B00532:end
% glossary-block:B00533:begin
\needspace{3\baselineskip}
\textbf{earnestly} — \textbf{geornlice} (adv.). earnestly; attentively.\par
% glossary-block:B00533:end
% glossary-block:B00534:begin
\needspace{3\baselineskip}
\textbf{earth} — \textbf{eorðe} (n.).  \textbf{middangeard} (n.). world; earth.\par
% glossary-block:B00534:end
% glossary-block:B00535:begin
\needspace{3\baselineskip}
\textbf{earthly} — \textbf{eorðlic} (adj.).\par
% glossary-block:B00535:end
% glossary-block:B00536:begin
\needspace{3\baselineskip}
\textbf{East Anglians} — \textbf{Eastęngle} (proper name).\par
% glossary-block:B00536:end
% glossary-block:B00537:begin
\needspace{3\baselineskip}
\textbf{eat} — \textbf{etan} (v.).\par
% glossary-block:B00537:end
% glossary-block:B00538:begin
\needspace{3\baselineskip}
\textbf{educate} — \textbf{getēon} (v.). train; educate.\par
% glossary-block:B00538:end
% glossary-block:B00539:begin
\needspace{3\baselineskip}
\textbf{educated} — \textbf{gelǣred} (adj.). educated; trained; skilled.\par
% glossary-block:B00539:end
% glossary-block:B00540:begin
\needspace{3\baselineskip}
\textbf{Egypt} — \textbf{Egypta} (proper name). Egyptians. Egypta land denotes Egypt, literally land of the Egyptians; supplied name forms.\par
% glossary-block:B00540:end
% glossary-block:B00541:begin
\needspace{3\baselineskip}
\textbf{Egyptians} — \textbf{Egypta} (proper name). Egypta land denotes Egypt, literally land of the Egyptians; supplied name forms. \textbf{Egypte} (ethnonym).\par
% glossary-block:B00541:end
% glossary-block:B00542:begin
\needspace{3\baselineskip}
\textbf{eight} — \textbf{eahta} (cardinal numeral). Cardinal 8.\par
% glossary-block:B00542:end
% glossary-block:B00543:begin
\needspace{3\baselineskip}
\textbf{eldest} — \textbf{ieldest} (adj.). eldest; oldest. Substantival person/eldest brother in this source context.\par
% glossary-block:B00543:end
% glossary-block:B00544:begin
\needspace{3\baselineskip}
\textbf{elephant} — \textbf{elp} (n.).\par
% glossary-block:B00544:end
% glossary-block:B00545:begin
\needspace{3\baselineskip}
\textbf{eleven} — \textbf{endlefan} (cardinal numeral). Cardinal 11. \textbf{endleofan} (numeral).\par
% glossary-block:B00545:end
% glossary-block:B00546:begin
\needspace{3\baselineskip}
\textbf{eleventh} — \textbf{endlefta} (ordinal).\par
% glossary-block:B00546:end
% glossary-block:B00547:begin
\needspace{3\baselineskip}
\textbf{empty} — \textbf{ǣmtig} (adj.). empty; void. \textbf{īdel} (adj.). empty; void.\par
% glossary-block:B00547:end
% glossary-block:B00548:begin
\needspace{3\baselineskip}
\textbf{end} — \textbf{ęnde} (n.).  \textbf{geęndian} (v.). end; come to an end. \textbf{geęndung} (n.). end; close.\par
% glossary-block:B00548:end
% glossary-block:B00549:begin
\needspace{3\baselineskip}
\textbf{endure} — \textbf{drēogan} (v.).  \textbf{ðolian} (v.). suffer; endure. Accusative hardship endured. \textbf{ðurhwunian} (v.). remain; endure; last. Intransitive continuing of a building; distinguish the following dative-governing wiðstandan.\par
% glossary-block:B00549:end
% glossary-block:B00550:begin
\needspace{3\baselineskip}
\textbf{enemy} — \textbf{fēond} (n.). enemy; foe. \textbf{gewinna} (n.). enemy; opponent. Supplied plural source forms.\par
% glossary-block:B00550:end
% glossary-block:B00551:begin
\needspace{3\baselineskip}
\textbf{engaged in} — The supplied on + activity-noun phrases express at work/engaged in that activity. Use the whole activity phrase and its supplied form; on alone does not name the activity.\par
% glossary-block:B00551:end
% glossary-block:B00552:begin
\needspace{3\baselineskip}
\textbf{English} — \textbf{Englisc} (adjective used substantivally). English, the language. Language noun. \textbf{Englisc} (adj.).\par
% glossary-block:B00552:end
% glossary-block:B00553:begin
\needspace{3\baselineskip}
\textbf{English people} — \textbf{Engle} (n.). the English. People, distinct from the language/adjective.\par
% glossary-block:B00553:end
% glossary-block:B00554:begin
\needspace{3\baselineskip}
\textbf{enslave} — \textbf{ðēowian} (v.). enslave; reduce to servitude.\par
% glossary-block:B00554:end
% glossary-block:B00555:begin
\needspace{3\baselineskip}
\textbf{enslaved men} — \textbf{ðēowa manna} (noun phrase). enslaved men; servants. Whole supplied phrase ðēowa manna.\par
% glossary-block:B00555:end
% glossary-block:B00556:begin
\needspace{3\baselineskip}
\textbf{enslaved people} — \textbf{ðēow} (adj.). the unfree; enslaved people. Substantival adjective in the source comparison.\par
% glossary-block:B00556:end
% glossary-block:B00557:begin
\needspace{3\baselineskip}
\textbf{enslaved person} — \textbf{ðēow} (n.). servant; enslaved person.\par
% glossary-block:B00557:end
% glossary-block:B00558:begin
\needspace{3\baselineskip}
\textbf{enter} — \textbf{gangan} (v.). go; enter. Supplied circumstantial participial phrase. \textbf{ingān} (v.). enter; go in. enter; go in. Supplied past source form. \textbf{ingangan} (v.). enter; go in. Supplied present-participle phrase.\par
% glossary-block:B00558:end
% glossary-block:B00559:begin
\needspace{3\baselineskip}
\textbf{entreat} — \textbf{biddan} (v.). beseech; entreat. Accusative person; the requested action can be a ðæt clause. \textbf{hālsian} (v.). entreat; implore.\par
% glossary-block:B00559:end
% glossary-block:B00560:begin
\needspace{3\baselineskip}
\textbf{envelop} — \textbf{ymbsęllan} (v.). surround; envelop. Supplied source form.\par
% glossary-block:B00560:end
% glossary-block:B00561:begin
\needspace{3\baselineskip}
\textbf{envoys} — \textbf{ǣrendwreca} (n.). messengers; envoys.\par
% glossary-block:B00561:end
% glossary-block:B00562:begin
\needspace{3\baselineskip}
\textbf{equal} — \textbf{gelīc} (adj.). equal; like. Dative standard of comparison. \textbf{gelīca} (n.). equal; peer.\par
% glossary-block:B00562:end
% glossary-block:B00563:begin
\needspace{3\baselineskip}
\textbf{escape} — \textbf{ætberstan} (v.).\par
% glossary-block:B00563:end
% glossary-block:B00564:begin
\needspace{3\baselineskip}
\textbf{establish} — \textbf{gesęttan} (v.). set; place; establish. \textbf{gestaðelian} (v.). establish; make steadfast. \textbf{sęttan} (v.). appoint; establish. With accusative entity/decree in the assigned frame.\par
% glossary-block:B00564:end
% glossary-block:B00565:begin
\needspace{3\baselineskip}
\textbf{eternal} — \textbf{ēce} (adj.). Adjective ēce, agreeing with the noun it qualifies.\par
% glossary-block:B00565:end
% glossary-block:B00566:begin
\needspace{3\baselineskip}
\textbf{Etheldred} — \textbf{Æðelðrȳð} (proper name). Etheldred; Æthelthryth. Proper name.\par
% glossary-block:B00566:end
% glossary-block:B00567:begin
\needspace{3\baselineskip}
\textbf{even} — \textbf{efne} (interj.). just; precisely; even.\par
% glossary-block:B00567:end
% glossary-block:B00568:begin
\needspace{3\baselineskip}
\textbf{evening} — \textbf{ǣfen} (n.).\par
% glossary-block:B00568:end
% glossary-block:B00569:begin
\needspace{3\baselineskip}
\textbf{even though} — \textbf{ðēah} (conj.; adv.). although; even though. ðēah with the taught concessive clause; a subjunctive need not make an admitted fact unreal. \textbf{ðēah ðe}. Supplied whole concessive conjunction in EX-XII-XIII-43-10.\par
% glossary-block:B00569:end
% glossary-block:B00570:begin
\needspace{3\baselineskip}
\textbf{ever} — \textbf{ā} (adv.). ever; always. \textbf{ǣfre} (adv.). ever; always.\par
% glossary-block:B00570:end
% glossary-block:B00571:begin
\needspace{3\baselineskip}
\textbf{everlasting} — \textbf{ēce} (adj.). Adjective ēce, agreeing with the noun it qualifies.\par
% glossary-block:B00571:end
% glossary-block:B00572:begin
\needspace{3\baselineskip}
\textbf{every} — \textbf{ǣghwilc} (indefinite adjective). every; each. each; every. \textbf{ǣlc} (adj.). each; every. With singular noun; substantival each follows the case required by its clause. \textbf{gehwilc} (indefinite adjective). each; every. Ǣlc and gehwilc mean each/every with a singular noun. Personal gehwā means each person/everyone; neuter gehwæt refers to things. Select case from the clause role.\par
% glossary-block:B00572:end
% glossary-block:B00573:begin
\needspace{3\baselineskip}
\textbf{everyone} — \textbf{gehwā / gehwæt} (pron.). each one; everyone; everything. Personal gehwā and neuter gehwæt; the early reference status does not by itself establish production. Ǣlc and gehwilc mean each/every with a singular noun. Personal gehwā means each person/everyone; neuter gehwæt refers to things. Select case from the clause role.\par
% glossary-block:B00573:end
% glossary-block:B00574:begin
\needspace{3\baselineskip}
\textbf{everything} — \textbf{eall} (adj.; substantive; adv.). all people; everything. Substantival: plural for people, neuter singular for everything. Eall agrees with a noun in all/whole expressions. Substantival plural refers to all people; neuter singular can mean everything. \textbf{gehwā / gehwæt} (pron.). each one; everyone; everything. Personal gehwā and neuter gehwæt; the early reference status does not by itself establish production.\par
% glossary-block:B00574:end
% glossary-block:B00575:begin
\needspace{3\baselineskip}
\textbf{evil} — \textbf{yfel} (adj.). bad; evil. Substantive noun in fram yfele. evil; evil people. Substantival adjective in the assigned use.\par
% glossary-block:B00575:end
% glossary-block:B00576:begin
\needspace{3\baselineskip}
\textbf{evil people} — \textbf{yfel} (adj.). evil; evil people. Substantival adjective in the assigned use.\par
% glossary-block:B00576:end
% glossary-block:B00577:begin
\needspace{3\baselineskip}
\textbf{exactly} — \textbf{eall} (adj.; substantive; adv.). just; exactly. Degree emphasis in eall swā.\par
% glossary-block:B00577:end
% glossary-block:B00578:begin
\needspace{3\baselineskip}
\textbf{exalt} — \textbf{ūp-ā-hębban} (v.). raise; exalt. Accusative people raised; the prefixed verb supplies up.\par
% glossary-block:B00578:end
% glossary-block:B00579:begin
\needspace{3\baselineskip}
\textbf{exalted} — \textbf{hēah} (adj.). high in rank; exalted. Rank/status sense in the supplied source.\par
% glossary-block:B00579:end
% glossary-block:B00580:begin
\needspace{3\baselineskip}
\textbf{example} — \textbf{bȳsen} (n.). example; illustration.\par
% glossary-block:B00580:end
% glossary-block:B00581:begin
\needspace{3\baselineskip}
\textbf{exceedingly} — \textbf{ðearle} (adv.). greatly; exceedingly.\par
% glossary-block:B00581:end
% glossary-block:B00582:begin
\needspace{3\baselineskip}
\textbf{except} — \textbf{būtan} (conj.). Būtan introducing the supplied conditional exception; distinguish a prepositional complement.\par
% glossary-block:B00582:end
% glossary-block:B00583:begin
\needspace{3\baselineskip}
\textbf{except if} — \textbf{būtan} (conj.). unless; except if.\par
% glossary-block:B00583:end
% glossary-block:B00584:begin
\needspace{3\baselineskip}
\textbf{expect} — \textbf{wēnan} (v.). suppose; expect; think. Genitive thing expected in the selected nominal frame; content-clause uses are separately taught.\par
% glossary-block:B00584:end
% glossary-block:B00585:begin
\needspace{3\baselineskip}
\textbf{expectation} — \textbf{wēn} (n.). expectation; chance. wēn is ðæt + subjunctive supplies the taught perhaps construction.\par
% glossary-block:B00585:end
% glossary-block:B00586:begin
\needspace{3\baselineskip}
\textbf{expel} — \textbf{ādrīfan} (v.). drive out; expel. Accusative person expelled; fram + dative may mark the point of separation. \textbf{āflīeman} (v.). put to flight; expel.\par
% glossary-block:B00586:end
% glossary-block:B00587:begin
\needspace{3\baselineskip}
\textbf{explain} — \textbf{geręccan} (v.). explain; render.\par
% glossary-block:B00587:end
% glossary-block:B00588:begin
\needspace{3\baselineskip}
\textbf{extend} — \textbf{brǣdan} (v.). spread; extend. Supplied source form spreads.\par
% glossary-block:B00588:end
% glossary-block:B00589:begin
\needspace{3\baselineskip}
\textbf{eye} — \textbf{ēage} (n.).\par
% glossary-block:B00589:end
\pdfbookmark[1]{F}{outline-B00590}
% glossary-block:B00590:begin
\needspace{5\baselineskip}
\subsection*{F}
% glossary-block:B00590:end
% glossary-block:B00591:begin
\textbf{face} — \textbf{andwlita} (n.). face; countenance. \textbf{brādnes} (n.). surface; face. Of the deep in the assigned expression.\par
% glossary-block:B00591:end
% glossary-block:B00592:begin
\needspace{3\baselineskip}
\textbf{fact} — \textbf{ðing} (n.). thing; fact; matter. Matter known/discussed, distinct from a generic physical object.\par
% glossary-block:B00592:end
% glossary-block:B00593:begin
\needspace{3\baselineskip}
\textbf{fail} — \textbf{geswīcan} (v.). fail; cease. Intransitive without a genitive complement in the selected frame.\par
% glossary-block:B00593:end
% glossary-block:B00594:begin
\needspace{3\baselineskip}
\textbf{fair} — \textbf{fæger} (adj.). beautiful; fair. \textbf{gesyndig} (adj.). favourable; fair. Of winds.\par
% glossary-block:B00594:end
% glossary-block:B00595:begin
\needspace{3\baselineskip}
\textbf{fairer} — \textbf{fægerra} (adj.). more beautiful; fairer.\par
% glossary-block:B00595:end
% glossary-block:B00596:begin
\needspace{3\baselineskip}
\textbf{faith} — \textbf{gelēafa} (n.).\par
% glossary-block:B00596:end
% glossary-block:B00597:begin
\needspace{3\baselineskip}
\textbf{fall} — \textbf{āhrēosan} (v.). fall; fall upon. \textbf{feallan} (v.).  \textbf{hrēosan} (v.). Supplied present participle falling.\par
% glossary-block:B00597:end
% glossary-block:B00598:begin
\needspace{3\baselineskip}
\textbf{fall in love} — \textbf{gefeallan} (v.). Figurative falling of the heart; supplied idiom differs in prefix from befeallan in 42.1.\par
% glossary-block:B00598:end
% glossary-block:B00599:begin
\needspace{3\baselineskip}
\textbf{fall in love with} — \textbf{befeallan} (v.). fall in love. Supplied idiom makes lufu love the subject, its possessor the lover, and on + accusative the beloved. \textbf{on} (prep.). upon. Accusative beloved in the supplied love-falls-upon idiom. In the supplied idiom, the lover’s lufu love is the subject of befeallan, and on + accusative identifies the beloved. The whole expression means fall in love with; the person is not the literal subject of fall.\par
% glossary-block:B00599:end
% glossary-block:B00600:begin
\needspace{3\baselineskip}
\textbf{fall upon} — \textbf{āhrēosan} (v.). fall; fall upon.\par
% glossary-block:B00600:end
% glossary-block:B00601:begin
\needspace{3\baselineskip}
\textbf{fame} — \textbf{hlīsa} (n.).\par
% glossary-block:B00601:end
% glossary-block:B00602:begin
\needspace{3\baselineskip}
\textbf{famine} — \textbf{hunger} (n.). famine; hunger.\par
% glossary-block:B00602:end
% glossary-block:B00603:begin
\needspace{3\baselineskip}
\textbf{far} — \textbf{feorr} (adv.). far; far off.\par
% glossary-block:B00603:end
% glossary-block:B00604:begin
\needspace{3\baselineskip}
\textbf{far and wide} — \textbf{wīde} (adv.). widely; far and wide.\par
% glossary-block:B00604:end
% glossary-block:B00605:begin
\needspace{3\baselineskip}
\textbf{farmer} — \textbf{ierðling} (n.). farmer; plowman.\par
% glossary-block:B00605:end
% glossary-block:B00606:begin
\needspace{3\baselineskip}
\textbf{Farne} — \textbf{Farne} (proper name). Island; supplied proper-name form after in.\par
% glossary-block:B00606:end
% glossary-block:B00607:begin
\needspace{3\baselineskip}
\textbf{far off} — \textbf{feorr} (adv.). far; far off.\par
% glossary-block:B00607:end
% glossary-block:B00608:begin
\needspace{3\baselineskip}
\textbf{fate} — \textbf{dōm} (n.). doom; fate.\par
% glossary-block:B00608:end
% glossary-block:B00609:begin
\needspace{3\baselineskip}
\textbf{father} — \textbf{fæder} (n.). Includes the spiritual father relationship in the assigned prompts.\par
% glossary-block:B00609:end
% glossary-block:B00610:begin
\needspace{3\baselineskip}
\textbf{favorable} — \textbf{gesyndig} (adj.). favourable; fair. Of winds.\par
% glossary-block:B00610:end
% glossary-block:B00611:begin
\needspace{3\baselineskip}
\textbf{favorably disposed} — \textbf{glæd} (adj.). gracious; favourably disposed. With wið + person in the supplied gracious-toward frame.\par
% glossary-block:B00611:end
% glossary-block:B00612:begin
\needspace{3\baselineskip}
\textbf{favourable} — \textbf{gesyndig} (adj.). favourable; fair. Of winds.\par
% glossary-block:B00612:end
% glossary-block:B00613:begin
\needspace{3\baselineskip}
\textbf{favourably disposed} — \textbf{glæd} (adj.). gracious; favourably disposed. With wið + person in the supplied gracious-toward frame.\par
% glossary-block:B00613:end
% glossary-block:B00614:begin
\needspace{3\baselineskip}
\textbf{fear} — \textbf{ondrǣdan} (v.). Accusative person feared in 24.6; the supplied reflexive construction in 48.6 has dative self and genitive thing feared. Keep the frames separate.\par
% glossary-block:B00614:end
% glossary-block:B00615:begin
\needspace{3\baselineskip}
\textbf{feast} — \textbf{bēorscipe} (n.). feast; banquet. \textbf{gebēorscipe} (n.). feast; banquet.\par
% glossary-block:B00615:end
% glossary-block:B00616:begin
\needspace{3\baselineskip}
\textbf{feathered} — \textbf{gefiðred} (adj.). Adjective agreeing with the fowl.\par
% glossary-block:B00616:end
% glossary-block:B00617:begin
\needspace{3\baselineskip}
\textbf{feed} — \textbf{fēdan} (v.). feed; nourish. Accusative being fed; food may be mid + dative or a bare instrumental in the assigned expressions.\par
% glossary-block:B00617:end
% glossary-block:B00618:begin
\needspace{3\baselineskip}
\textbf{feet} — \textbf{fōt} (n.). foot. Supplied body-part form feet; distinct from the length-measure sense. foot. Unit of length, distinct from the body-part sense.\par
% glossary-block:B00618:end
% glossary-block:B00619:begin
\needspace{3\baselineskip}
\textbf{fellow} — \textbf{gefēra} (n.). companion; fellow.\par
% glossary-block:B00619:end
% glossary-block:B00620:begin
\needspace{3\baselineskip}
\textbf{fervently} — \textbf{inweardlīce} (adv.). fervently; wholeheartedly.\par
% glossary-block:B00620:end
% glossary-block:B00621:begin
\needspace{3\baselineskip}
\textbf{fetch} — \textbf{feccean} (v.).\par
% glossary-block:B00621:end
% glossary-block:B00622:begin
\needspace{3\baselineskip}
\textbf{few} — \textbf{fēawe} (adj.). few people. Substantival use in the assigned expression.\par
% glossary-block:B00622:end
% glossary-block:B00623:begin
\needspace{3\baselineskip}
\textbf{few people} — \textbf{fēawe} (adj.). Substantival use in the assigned expression.\par
% glossary-block:B00623:end
% glossary-block:B00624:begin
\needspace{3\baselineskip}
\textbf{field} — \textbf{æcer} (n.).\par
% glossary-block:B00624:end
% glossary-block:B00625:begin
\needspace{3\baselineskip}
\textbf{fierce} — \textbf{angrīslic} (adj.). fierce; raging.\par
% glossary-block:B00625:end
% glossary-block:B00626:begin
\needspace{3\baselineskip}
\textbf{fiercely} — \textbf{ðearlice} (adv.). violently; fiercely.\par
% glossary-block:B00626:end
% glossary-block:B00627:begin
\needspace{3\baselineskip}
\textbf{fierceness} — \textbf{rēðnes} (n.). fury; fierceness.\par
% glossary-block:B00627:end
% glossary-block:B00628:begin
\needspace{3\baselineskip}
\textbf{fiercer} — \textbf{swīðra} (adj.). stronger; fiercer. Comparative degree, distinct from right-side sense.\par
% glossary-block:B00628:end
% glossary-block:B00629:begin
\needspace{3\baselineskip}
\textbf{fifth} — \textbf{fīfta} (ordinal numeral).\par
% glossary-block:B00629:end
% glossary-block:B00630:begin
\needspace{3\baselineskip}
\textbf{fifty} — \textbf{fīftig} (numeral).\par
% glossary-block:B00630:end
% glossary-block:B00631:begin
\needspace{3\baselineskip}
\textbf{fight} — \textbf{campian} (v.).  \textbf{feohtan} (v.). fight; contend. Intransitive; wið + accusative opponent.\par
% glossary-block:B00631:end
% glossary-block:B00632:begin
\needspace{3\baselineskip}
\textbf{fill} — \textbf{āfyllan} (v.).  \textbf{gefyllan} (v.). Gefyllan means fill when something is filled; the finish/accomplish sense is separate.\par
% glossary-block:B00632:end
% glossary-block:B00633:begin
\needspace{3\baselineskip}
\textbf{filthy} — \textbf{horg} (adj.). dirty; filthy.\par
% glossary-block:B00633:end
% glossary-block:B00634:begin
\needspace{3\baselineskip}
\textbf{find} — \textbf{findan} (v.). Accusative person or thing found. \textbf{gemētan} (v.). find; meet. Accusative person or thing found/met. Supplied past participle gemētt in a passive. \textbf{mētan} (v.). meet; find.\par
% glossary-block:B00634:end
% glossary-block:B00635:begin
\needspace{3\baselineskip}
\textbf{find out} — \textbf{gewitan} (v.). find out; learn. Knowledge verb, distinct from gewītan depart.\par
% glossary-block:B00635:end
% glossary-block:B00636:begin
\needspace{3\baselineskip}
\textbf{finger} — \textbf{finger} (n.).\par
% glossary-block:B00636:end
% glossary-block:B00637:begin
\needspace{3\baselineskip}
\textbf{finish} — \textbf{fullfręmman} (v.). finish; complete. Supplied passive use. \textbf{gefyllan} (v.). finish; accomplish. A work is the thing completed; distinct from fill. Gefyllan means finish/accomplish when a work is completed; the fill sense is separate.\par
% glossary-block:B00637:end
% glossary-block:B00638:begin
\needspace{3\baselineskip}
\textbf{fire} — \textbf{fȳr} (n.).\par
% glossary-block:B00638:end
% glossary-block:B00639:begin
\needspace{3\baselineskip}
\textbf{fire from heaven} — \textbf{heofonfȳr} (n.). fire from heaven; heavenly fire.\par
% glossary-block:B00639:end
% glossary-block:B00640:begin
\needspace{3\baselineskip}
\textbf{firmament} — \textbf{fæstnes} (n.).\par
% glossary-block:B00640:end
% glossary-block:B00641:begin
\needspace{3\baselineskip}
\textbf{first} — \textbf{ǣrest} (adv.). first; at first. \textbf{ǣresta} (ordinal numeral).  \textbf{fyrmest} (adv.).\par
% glossary-block:B00641:end
% glossary-block:B00642:begin
\needspace{3\baselineskip}
\textbf{first of all} — \textbf{fyrmest} (adv.). first of all; at the outset.\par
% glossary-block:B00642:end
% glossary-block:B00643:begin
\needspace{3\baselineskip}
\textbf{fish} — \textbf{fisc} (n.).\par
% glossary-block:B00643:end
% glossary-block:B00644:begin
\needspace{3\baselineskip}
\textbf{fisher} — \textbf{fiscere} (n.). Masculine noun; fisherman is only one English surface form of this lookup.\par
% glossary-block:B00644:end
% glossary-block:B00645:begin
\needspace{3\baselineskip}
\textbf{fisherman} — \textbf{fiscere} (n.). Masculine noun; fisherman is only one English surface form of this lookup.\par
% glossary-block:B00645:end
% glossary-block:B00646:begin
\needspace{3\baselineskip}
\textbf{fishermen} — \textbf{fiscere} (n.). Masculine noun; fisherman is only one English surface form of this lookup.\par
% glossary-block:B00646:end
% glossary-block:B00647:begin
\needspace{3\baselineskip}
\textbf{fishers} — \textbf{fiscere} (n.). Masculine noun; fisherman is only one English surface form of this lookup.\par
% glossary-block:B00647:end
% glossary-block:B00648:begin
\needspace{3\baselineskip}
\textbf{fish-hook} — \textbf{angel} (n.). Compound English spelling and optional fish- lead to the same fishing noun; distinguish ęngel angel.\par
% glossary-block:B00648:end
% glossary-block:B00649:begin
\needspace{3\baselineskip}
\textbf{fishhook} — \textbf{angel} (n.). Compound English spelling and optional fish- lead to the same fishing noun; distinguish ęngel angel.\par
% glossary-block:B00649:end
% glossary-block:B00650:begin
\needspace{3\baselineskip}
\textbf{fishing} — \textbf{fiscnoð} (n.).  Gān on fiscnoð means go fishing; fiscnoð is the noun fishing and on is part of this supplied phrase.\par
% glossary-block:B00650:end
% glossary-block:B00651:begin
\needspace{3\baselineskip}
\textbf{fish kind} — \textbf{fisccynn} (n.). kinds of fish. Supplied accusative neuter plural.\par
% glossary-block:B00651:end
% glossary-block:B00652:begin
\needspace{3\baselineskip}
\textbf{five} — \textbf{fīf} (cardinal numeral). Cardinal 5.\par
% glossary-block:B00652:end
% glossary-block:B00653:begin
\needspace{3\baselineskip}
\textbf{flee} — \textbf{flēon} (v.). flee; flee from; avoid. Intransitive or with accusative person/thing avoided. Supplied third-plural present; distinct from flēogan fly.\par
% glossary-block:B00653:end
% glossary-block:B00654:begin
\needspace{3\baselineskip}
\textbf{flee from} — \textbf{flēon} (v.). flee; flee from; avoid. Intransitive or with accusative person/thing avoided.\par
% glossary-block:B00654:end
% glossary-block:B00655:begin
\needspace{3\baselineskip}
\textbf{fleet} — \textbf{sciphęre} (n.). fleet; naval force.\par
% glossary-block:B00655:end
% glossary-block:B00656:begin
\needspace{3\baselineskip}
\textbf{flesh} — \textbf{flǣsc} (n.).\par
% glossary-block:B00656:end
% glossary-block:B00657:begin
\needspace{3\baselineskip}
\textbf{flood} — \textbf{flōd} (n.). flood; stream of water.\par
% glossary-block:B00657:end
% glossary-block:B00658:begin
\needspace{3\baselineskip}
\textbf{flow} — \textbf{flōwan} (v.).\par
% glossary-block:B00658:end
% glossary-block:B00659:begin
\needspace{3\baselineskip}
\textbf{fly} — \textbf{flēogan} (v.).\par
% glossary-block:B00659:end
% glossary-block:B00660:begin
\needspace{3\baselineskip}
\textbf{fodder} — \textbf{fōdor} (n.). Provender is routed by Cook's bracketed fodder to fōdor.\par
% glossary-block:B00660:end
% glossary-block:B00661:begin
\needspace{3\baselineskip}
\textbf{foe} — \textbf{fēond} (n.). enemy; foe.\par
% glossary-block:B00661:end
% glossary-block:B00662:begin
\needspace{3\baselineskip}
\textbf{follow} — \textbf{æfter} (prep.; conj. in fixed expressions). after; following. Dative person followed, including motion after that person.\par
% glossary-block:B00662:end
% glossary-block:B00663:begin
\needspace{3\baselineskip}
\textbf{food} — \textbf{ēst} (n.). provisions; food. \textbf{męte} (n.).\par
% glossary-block:B00663:end
% glossary-block:B00664:begin
\needspace{3\baselineskip}
\textbf{foot} — \textbf{fōt} (n.). Supplied body-part form feet; distinct from the length-measure sense. Unit of length, distinct from the body-part sense.\par
% glossary-block:B00664:end
% glossary-block:B00665:begin
\needspace{3\baselineskip}
\textbf{footstool} — \textbf{fōtsceamul} (n.).\par
% glossary-block:B00665:end
% glossary-block:B00666:begin
\needspace{3\baselineskip}
\textbf{for} — \textbf{for} (prep.). for; on behalf of. Dative beneficiary. for; for the sake of. Case follows each complete supplied accusative/dative expression. in exchange for. With dative. for; to the benefit of. \textbf{tō} (prep.). for; to serve as. tō + dative function; a separate dative beneficiary may accompany it. for; to the benefit of. Care of a body part in the supplied expression. \textbf{wið} (prep.). in return for. With dative; distinct from opponent/hostility.\par
% glossary-block:B00666:end
% glossary-block:B00667:begin
\needspace{3\baselineskip}
\textbf{for aid} — Tō + the supplied aid/protection noun forms the whole purpose expression for aid/for protection; the noun supplies the aid/protection meaning.\par
% glossary-block:B00667:end
% glossary-block:B00668:begin
\needspace{3\baselineskip}
\textbf{for a sign} — Tō + a dative noun names the assigned function, for food or for a sign. The noun supplies food/sign; a separate dative beneficiary may also occur. \textbf{tō} (prep.). tō + dative function, with an optional separate dative beneficiary.\par
% glossary-block:B00668:end
% glossary-block:B00669:begin
\needspace{3\baselineskip}
\textbf{for (causal)} — \textbf{for-ðām} (conj.), \textbf{for-ðon-ðe} (conj.; adv. in for-ðon). For-ðām and for-ðon-ðe are causal conjunctions meaning because/for. They are distinct from prepositions meaning because of or for someone’s benefit.\par
% glossary-block:B00669:end
% glossary-block:B00670:begin
\needspace{3\baselineskip}
\textbf{foreign} — \textbf{ęlðēodig} (adj.).\par
% glossary-block:B00670:end
% glossary-block:B00671:begin
\needspace{3\baselineskip}
\textbf{foreigners} — \textbf{ęlðēodig} (adj.). foreigners; strangers. Substantival adjective.\par
% glossary-block:B00671:end
% glossary-block:B00672:begin
\needspace{3\baselineskip}
\textbf{foremost} — \textbf{furðra} (adj.). foremost; prior.\par
% glossary-block:B00672:end
% glossary-block:B00673:begin
\needspace{3\baselineskip}
\textbf{forest} — \textbf{wudu} (n.). wood; forest. Collective vegetation, not a new material sense. \textbf{wudubearu} (n.). forest; grove.\par
% glossary-block:B00673:end
% glossary-block:B00674:begin
\needspace{3\baselineskip}
\textbf{for ever} — Tō worulde is the whole fixed phrase for ever/forever.\par
% glossary-block:B00674:end
% glossary-block:B00675:begin
\needspace{3\baselineskip}
\textbf{forever} — Tō worulde is the whole fixed phrase for ever/forever.\par
% glossary-block:B00675:end
% glossary-block:B00676:begin
\needspace{3\baselineskip}
\textbf{for ever and ever} — \textbf{ā} (adv.). Preserve the source-supplied repetition ā ā ā where given; this is not three separately chosen words.\par
% glossary-block:B00676:end
% glossary-block:B00677:begin
\needspace{3\baselineskip}
\textbf{for food} — Tō + a dative noun names the assigned function, for food or for a sign. The noun supplies food/sign; a separate dative beneficiary may also occur. \textbf{tō} (prep.). tō + dative function, with an optional separate dative beneficiary.\par
% glossary-block:B00677:end
% glossary-block:B00678:begin
\needspace{3\baselineskip}
\textbf{forgive} — \textbf{forgiefan} (v.). Accusative sin, dative person forgiven; in a passive the sin becomes nominative and the person remains dative.\par
% glossary-block:B00678:end
% glossary-block:B00679:begin
\needspace{3\baselineskip}
\textbf{for its part} — \textbf{sōðlice} (adv.). now; but. Discourse connective corresponding to source Latin autem, distinct from truly.\par
% glossary-block:B00679:end
% glossary-block:B00680:begin
\needspace{3\baselineskip}
\textbf{form} — \textbf{hīw} (n.). form; kind. Supplied dative neuter plural.\par
% glossary-block:B00680:end
% glossary-block:B00681:begin
\needspace{3\baselineskip}
\textbf{former} — \textbf{ǣrra} (adj.). former; earlier.\par
% glossary-block:B00681:end
% glossary-block:B00682:begin
\needspace{3\baselineskip}
\textbf{formerly} — The supplied whole phrase hwīlum ǣr means a while ago/formerly. Hwīlum alone has the separate earlier meaning sometimes.\par
% glossary-block:B00682:end
% glossary-block:B00683:begin
\needspace{3\baselineskip}
\textbf{for protection} — Tō + the supplied aid/protection noun forms the whole purpose expression for aid/for protection; the noun supplies the aid/protection meaning.\par
% glossary-block:B00683:end
% glossary-block:B00684:begin
\needspace{3\baselineskip}
\textbf{for (recipient or beneficiary)} — A dative noun/pronoun often expresses the recipient/beneficiary without an added preposition. Distinguish destinations and lexical government.\par
% glossary-block:B00684:end
% glossary-block:B00685:begin
\needspace{3\baselineskip}
\textbf{forsake} — \textbf{forlǣtan} (v.). abandon; forsake. Accusative person or thing left.\par
% glossary-block:B00685:end
% glossary-block:B00686:begin
\needspace{3\baselineskip}
\textbf{forth} — \textbf{forð} (adv.). Adverbial particle with the bringing/producing verb.\par
% glossary-block:B00686:end
% glossary-block:B00687:begin
\needspace{3\baselineskip}
\textbf{for that reason} — \textbf{for-ðon-ðe} (conj.; adv. in for-ðon). therefore; for that reason. Adverbial consequence, distinct from because.\par
% glossary-block:B00687:end
% glossary-block:B00688:begin
\needspace{3\baselineskip}
\textbf{for the benefit of} — \textbf{for} (prep.). Supplied for + dative benefit phrase; its possessed noun controls possessive agreement.\par
% glossary-block:B00688:end
% glossary-block:B00689:begin
\needspace{3\baselineskip}
\textbf{for the sake of} — \textbf{for} (prep.). for; for the sake of. Case follows each complete supplied accusative/dative expression. by; for the sake of. Grounds of entreaty in the prayer; keep the supplied phrase. Intinga means cause/sake. The supplied for intingan expression renders for the sake of; retain its full supplied dependent phrase.\par
% glossary-block:B00689:end
% glossary-block:B00690:begin
\needspace{3\baselineskip}
\textbf{fortification} — \textbf{fæsten} (n.). fortification; stronghold.\par
% glossary-block:B00690:end
% glossary-block:B00691:begin
\needspace{3\baselineskip}
\textbf{forty} — \textbf{fēowertig} (numeral).\par
% glossary-block:B00691:end
% glossary-block:B00692:begin
\needspace{3\baselineskip}
\textbf{for what purpose} — Tō hwon is the whole expression why/for what purpose.\par
% glossary-block:B00692:end
% glossary-block:B00693:begin
\needspace{3\baselineskip}
\textbf{for what reason} — \textbf{for} (prep.). why; for what reason. In for hwȳ or for hwon.\par
% glossary-block:B00693:end
% glossary-block:B00694:begin
\needspace{3\baselineskip}
\textbf{foul} — \textbf{fūl} (adj.). foul; vile.\par
% glossary-block:B00694:end
% glossary-block:B00695:begin
\needspace{3\baselineskip}
\textbf{foundation} — \textbf{grundweall} (n.).\par
% glossary-block:B00695:end
% glossary-block:B00696:begin
\needspace{3\baselineskip}
\textbf{four} — \textbf{fēower} (cardinal numeral). Cardinal 4.\par
% glossary-block:B00696:end
% glossary-block:B00697:begin
\needspace{3\baselineskip}
\textbf{fourth} — \textbf{fēorða} (ordinal numeral).\par
% glossary-block:B00697:end
% glossary-block:B00698:begin
\needspace{3\baselineskip}
\textbf{fowl} — \textbf{fugol} (n.). Fowl is collective/plural English here; use the required Old English plural.\par
% glossary-block:B00698:end
% glossary-block:B00699:begin
\needspace{3\baselineskip}
\textbf{fox} — \textbf{fox} (n.).\par
% glossary-block:B00699:end
% glossary-block:B00700:begin
\needspace{3\baselineskip}
\textbf{free} — \textbf{frēo} (adj.).  Frēo is the adjective free. In go free it describes the subject with appropriate agreement; it is not itself a verb meaning go.\par
% glossary-block:B00700:end
% glossary-block:B00701:begin
\needspace{3\baselineskip}
\textbf{frequently} — \textbf{gelōmlice} (adv.). frequently; regularly. \textbf{genēahhe} (adv.). often; frequently.\par
% glossary-block:B00701:end
% glossary-block:B00702:begin
\needspace{3\baselineskip}
\textbf{friend} — \textbf{frēond} (n.).\par
% glossary-block:B00702:end
% glossary-block:B00703:begin
\needspace{3\baselineskip}
\textbf{frighten} — \textbf{āfǣran} (v.).\par
% glossary-block:B00703:end
% glossary-block:B00704:begin
\needspace{3\baselineskip}
\textbf{from} — \textbf{æt} (prep.). æt + dative person with learn. from; out of. Proceeds of work, in the assigned expression. \textbf{fram} (prep.). Separation, with dative in the taught frame. Only the supplied rest-from-work construction. \textbf{of} (prep.). from; out of. With dative. of; from. Portion from a whole in the supplied partitive phrases; dative whole. Starting condition of old age in the supplied rejuvenation comparison. of; from. Origin of wisdom in the supplied begetting statement. Starting generation in the supplied temporal range.\par
% glossary-block:B00704:end
% glossary-block:B00705:begin
\needspace{3\baselineskip}
\textbf{from afar} — \textbf{feorran} (adv.). feorran; the source supplies the from-afar paraphrase.\par
% glossary-block:B00705:end
% glossary-block:B00706:begin
\needspace{3\baselineskip}
\textbf{from generation to generation} — \textbf{of} (prep.). from. Starting generation in the supplied temporal range. On introduces the succeeding generation in the supplied from-generation-to-generation range. Retain the full supplied temporal phrase; on alone does not mean from generation to generation.\par
% glossary-block:B00706:end
% glossary-block:B00707:begin
\needspace{3\baselineskip}
\textbf{from heaven} — \textbf{heofonlice} (adv.). from heaven; heavenly. Source expression at Cook 143.17.\par
% glossary-block:B00707:end
% glossary-block:B00708:begin
\needspace{3\baselineskip}
\textbf{from here} — \textbf{heonan} (adv.). hence; from here.\par
% glossary-block:B00708:end
% glossary-block:B00709:begin
\needspace{3\baselineskip}
\textbf{from that} — \textbf{be} (prep.). by that; from that. Refers to the preceding proposition; supplied expression. by that; on that evidence. Evidence/inference frame, not an unrestricted source preposition.\par
% glossary-block:B00709:end
% glossary-block:B00710:begin
\needspace{3\baselineskip}
\textbf{from the north} — \textbf{norðan} (adv.).\par
% glossary-block:B00710:end
% glossary-block:B00711:begin
\needspace{3\baselineskip}
\textbf{from there} — \textbf{ðanan} (adv.). thence; from there.\par
% glossary-block:B00711:end
% glossary-block:B00712:begin
\needspace{3\baselineskip}
\textbf{from the south} — \textbf{sūðan} (adv.). in the south; from the south. Retain the supplied spatial expression.\par
% glossary-block:B00712:end
% glossary-block:B00713:begin
\needspace{3\baselineskip}
\textbf{from the west} — \textbf{westan} (adv.).\par
% glossary-block:B00713:end
% glossary-block:B00714:begin
\needspace{3\baselineskip}
\textbf{from where} — \textbf{hwanon} (adv.). whence; from where.\par
% glossary-block:B00714:end
% glossary-block:B00715:begin
\needspace{3\baselineskip}
\textbf{fruit} — \textbf{wæstm} (n.). fruit; produce. Corrected short æ in wæstm under C01, despite the retained erroneous source hint wǣstm.\par
% glossary-block:B00715:end
% glossary-block:B00716:begin
\needspace{3\baselineskip}
\textbf{fruit-bearing} — \textbf{æppelbǣre} (adj.).\par
% glossary-block:B00716:end
% glossary-block:B00717:begin
\needspace{3\baselineskip}
\textbf{fruitful} — \textbf{wæstmbǣre} (adj.).\par
% glossary-block:B00717:end
% glossary-block:B00718:begin
\needspace{3\baselineskip}
\textbf{fulfil} — \textbf{dōn} (v.). fulfil; carry out. A word or command is fulfilled. \textbf{gefyllan} (v.). Accusative prediction fulfilled; passive subjunctive in the assigned purpose construction.\par
% glossary-block:B00718:end
% glossary-block:B00719:begin
\needspace{3\baselineskip}
\textbf{fulfill} — \textbf{dōn} (v.). fulfil; carry out. A word or command is fulfilled. \textbf{gefyllan} (v.). fulfil. Accusative prediction fulfilled; passive subjunctive in the assigned purpose construction.\par
% glossary-block:B00719:end
% glossary-block:B00720:begin
\needspace{3\baselineskip}
\textbf{furlong} — \textbf{furlang} (n.).\par
% glossary-block:B00720:end
% glossary-block:B00721:begin
\needspace{3\baselineskip}
\textbf{further} — \textbf{furðor} (adv.). further; any more. Adverb in the supplied no-spirit-more expression. \textbf{gīet} (adv.). still; further.\par
% glossary-block:B00721:end
% glossary-block:B00722:begin
\needspace{3\baselineskip}
\textbf{fury} — \textbf{rēðnes} (n.). fury; fierceness.\par
% glossary-block:B00722:end
\pdfbookmark[1]{G}{outline-B00723}
% glossary-block:B00723:begin
\needspace{5\baselineskip}
\subsection*{G}
% glossary-block:B00723:end
% glossary-block:B00724:begin
\textbf{gain} — \textbf{gestrēon} (n.). gain; profit.\par
% glossary-block:B00724:end
% glossary-block:B00725:begin
\needspace{3\baselineskip}
\textbf{game} — \textbf{plega} (n.). play; game.\par
% glossary-block:B00725:end
% glossary-block:B00726:begin
\needspace{3\baselineskip}
\textbf{garment} — \textbf{gegierela} (n.). garment; clothing. \textbf{wǣd} (n.).  \textbf{wǣd} (n.), \textbf{gegierela} (n.). Wǣd denotes a garment, rēaf raiment/clothing, and gegierela a garment or clothing. Preserve any noun expressly supplied by the exercise.\par
% glossary-block:B00726:end
% glossary-block:B00727:begin
\needspace{3\baselineskip}
\textbf{garments} — \textbf{rēaf} (n.). clothing; garments.\par
% glossary-block:B00727:end
% glossary-block:B00728:begin
\needspace{3\baselineskip}
\textbf{gate} — \textbf{geat} (n.).\par
% glossary-block:B00728:end
% glossary-block:B00729:begin
\needspace{3\baselineskip}
\textbf{gather} — \textbf{gegaderian} (v.). Accusative people or things gathered.\par
% glossary-block:B00729:end
% glossary-block:B00730:begin
\needspace{3\baselineskip}
\textbf{gathering} — \textbf{gegaderung} (n.). gathering; collection; congregation. Congregations of waters in the source.\par
% glossary-block:B00730:end
% glossary-block:B00731:begin
\needspace{3\baselineskip}
\textbf{gaze} — \textbf{gesihð} (n.). sight; gaze. Supplied accusative singular.\par
% glossary-block:B00731:end
% glossary-block:B00732:begin
\needspace{3\baselineskip}
\textbf{gem} — \textbf{gimm} (n.). gem; precious stone.\par
% glossary-block:B00732:end
% glossary-block:B00733:begin
\needspace{3\baselineskip}
\textbf{generation} — \textbf{cnēoris} (n.).\par
% glossary-block:B00733:end
% glossary-block:B00734:begin
\needspace{3\baselineskip}
\textbf{gift} — \textbf{giefu} (n.). Supplied plural gifts. Divine grace in Cook's bracketed gift route; not a license to merge all secular present/gift senses. \textbf{lāc} (n.). gift; present.\par
% glossary-block:B00734:end
% glossary-block:B00735:begin
\needspace{3\baselineskip}
\textbf{gird} — \textbf{begyrdan} (v.). gird; clothe. Supplied participial clothing expression.\par
% glossary-block:B00735:end
% glossary-block:B00736:begin
\needspace{3\baselineskip}
\textbf{girl} — \textbf{mǣden} (n.). girl; maiden.\par
% glossary-block:B00736:end
% glossary-block:B00737:begin
\needspace{3\baselineskip}
\textbf{give} — \textbf{forgiefan} (v.). give; grant. Supplied granted sense at Cook 143.18; not forgive here. give; grant. Accusative thing given and dative recipient; the thing may be understood from context. \textbf{gesęllan} (v.).  \textbf{giefan} (v.).  \textbf{sęllan} (v.). Sęllan means give in the giving sense; an expressed recipient is dative. Its sell sense is separate.\par
% glossary-block:B00737:end
% glossary-block:B00738:begin
\needspace{3\baselineskip}
\textbf{give back} — \textbf{āgiefan} (v.). give back; restore.\par
% glossary-block:B00738:end
% glossary-block:B00739:begin
\needspace{3\baselineskip}
\textbf{give light to} — \textbf{ālīehtan} (v.). illuminate; give light to. Accusative person or thing illuminated. ālīehtan takes an accusative object, not a dative inferred from English to.\par
% glossary-block:B00739:end
% glossary-block:B00740:begin
\needspace{3\baselineskip}
\textbf{give to drink} — \textbf{scęncan} (v.). pour out; give to drink. Accusative drink, dative recipient. \textbf{sęllan} (v.). Supplied dative person + plain drincan, without tō. Supplied sęllan + dative person + plain drincan; no added tō.\par
% glossary-block:B00740:end
% glossary-block:B00741:begin
\needspace{3\baselineskip}
\textbf{glad} — \textbf{blīðe} (adj.). merry; glad. \textbf{glæd} (adj.).\par
% glossary-block:B00741:end
% glossary-block:B00742:begin
\needspace{3\baselineskip}
\textbf{gladly} — \textbf{blīðe} (adv.). joyfully; gladly.\par
% glossary-block:B00742:end
% glossary-block:B00743:begin
\needspace{3\baselineskip}
\textbf{glorify} — \textbf{mǣrsian} (v.). make great; magnify; glorify. \textbf{wuldrian} (v.). Accusative person; an accusative reflexive expresses glorifying oneself.\par
% glossary-block:B00743:end
% glossary-block:B00744:begin
\needspace{3\baselineskip}
\textbf{glory} — \textbf{mægenðrymm} (n.). majesty; glory. \textbf{mægenðrymnes} (n.). glory; majesty. \textbf{wuldor} (n.).\par
% glossary-block:B00744:end
% glossary-block:B00745:begin
\needspace{3\baselineskip}
\textbf{go} — \textbf{faran} (v.). go; travel. \textbf{gān} (v.). go; walk. Supplied imperative gang. \textbf{gangan} (v.). go; enter. Supplied circumstantial participial phrase.\par
% glossary-block:B00745:end
% glossary-block:B00746:begin
\needspace{3\baselineskip}
\textbf{go aboard} — \textbf{āstīgan} (v.). go aboard; board. Accusative ship boarded in the assigned frame. āstīgan takes the ship as direct accusative in these exercises.\par
% glossary-block:B00746:end
% glossary-block:B00747:begin
\needspace{3\baselineskip}
\textbf{goat} — \textbf{gāt} (n.).\par
% glossary-block:B00747:end
% glossary-block:B00748:begin
\needspace{3\baselineskip}
\textbf{go away} — \textbf{āgān} (v.). Supplied gone-away form; distinct from āgan possess. \textbf{gewītan} (v.). depart; go away. Intransitive; fram + dative departure point is optional. The assigned frame requires no reflexive.\par
% glossary-block:B00748:end
% glossary-block:B00749:begin
\needspace{3\baselineskip}
\textbf{god} — \textbf{God} (n.). god; deity. Common deity sense, distinct from the proper divine name. A god or deity; this common-noun sense is distinct from the proper name God.\par
% glossary-block:B00749:end
% glossary-block:B00750:begin
\needspace{3\baselineskip}
\textbf{God} — \textbf{God} (n.). Proper divine name; no productive plural sense is assigned. The proper divine name God; no productive plural is assigned.\par
% glossary-block:B00750:end
% glossary-block:B00751:begin
\needspace{3\baselineskip}
\textbf{go fishing} — Gān on fiscnoð means go fishing; fiscnoð is the noun fishing and on is part of this supplied phrase. Gān on fiscnoð is the supplied whole expression go fishing. Fiscnoð is the noun fishing; on is part of this phrase, not a verb meaning go.\par
% glossary-block:B00751:end
% glossary-block:B00752:begin
\needspace{3\baselineskip}
\textbf{go free} — Frēo is the adjective free. In go free it describes the subject with appropriate agreement; it is not itself a verb meaning go.\par
% glossary-block:B00752:end
% glossary-block:B00753:begin
\needspace{3\baselineskip}
\textbf{go in} — \textbf{ingān} (v.). enter; go in. enter; go in. Supplied past source form. \textbf{ingangan} (v.). enter; go in. Supplied present-participle phrase.\par
% glossary-block:B00753:end
% glossary-block:B00754:begin
\needspace{3\baselineskip}
\textbf{gold} — \textbf{gold} (n.).\par
% glossary-block:B00754:end
% glossary-block:B00755:begin
\needspace{3\baselineskip}
\textbf{good} — \textbf{gōd} (adj.). Ordinary gōd good. The person-directed kind/benevolent use is a separate sense. \textbf{gōd} (n.). good; benefit; kindness. Nominal good in the repayment contrast. \textbf{ðearf} (n.). benefit; good; need. Source contextual nouns; retain the local expression.\par
% glossary-block:B00755:end
% glossary-block:B00756:begin
\needspace{3\baselineskip}
\textbf{goods} — \textbf{ǣht} (n.). possessions; goods. \textbf{gōd} (n.). goods; property. Plural lexical noun.\par
% glossary-block:B00756:end
% glossary-block:B00757:begin
\needspace{3\baselineskip}
\textbf{good service} — \textbf{earnung} (n.). merit; good service. Source plural merits/services.\par
% glossary-block:B00757:end
% glossary-block:B00758:begin
\needspace{3\baselineskip}
\textbf{good things} — \textbf{gōd} (n.). goods; property. Plural lexical noun.\par
% glossary-block:B00758:end
% glossary-block:B00759:begin
\needspace{3\baselineskip}
\textbf{go out} — \textbf{ūtgān} (v.). Supplied past source form.\par
% glossary-block:B00759:end
% glossary-block:B00760:begin
\needspace{3\baselineskip}
\textbf{Gospel} — \textbf{godspell} (n.).\par
% glossary-block:B00760:end
% glossary-block:B00761:begin
\needspace{3\baselineskip}
\textbf{go to pieces} — \textbf{tōberstan} (v.). break up; go to pieces.\par
% glossary-block:B00761:end
% glossary-block:B00762:begin
\needspace{3\baselineskip}
\textbf{go up} — \textbf{ūpgān} (v.). go up; come up. From the water in the supplied expression.\par
% glossary-block:B00762:end
% glossary-block:B00763:begin
\needspace{3\baselineskip}
\textbf{grace} — \textbf{giefu} (n.). grace; divine favour. Divine grace in Cook's bracketed gift route; not a license to merge all secular present/gift senses.\par
% glossary-block:B00763:end
% glossary-block:B00764:begin
\needspace{3\baselineskip}
\textbf{gracious} — \textbf{glæd} (adj.). gracious; favourably disposed. With wið + person in the supplied gracious-toward frame.\par
% glossary-block:B00764:end
% glossary-block:B00765:begin
\needspace{3\baselineskip}
\textbf{grant} — \textbf{forgiefan} (v.). give; grant. Supplied granted sense at Cook 143.18; not forgive here. give; grant. Accusative thing given and dative recipient; the thing may be understood from context. \textbf{unnan} (v.). grant; allow. Dative person, genitive thing granted. Dative recipient and genitive thing granted; preserve the supported subjunctive command unne.\par
% glossary-block:B00765:end
% glossary-block:B00766:begin
\needspace{3\baselineskip}
\textbf{grass} — \textbf{gærs} (n.).\par
% glossary-block:B00766:end
% glossary-block:B00767:begin
\needspace{3\baselineskip}
\textbf{great} — \textbf{mǣre} (adj.). great; renowned. Distinct from simple physical size where the task requires the renowned sense. \textbf{micel} (adj.). great; large. great; loud. Loud is the contextual force with a voice.\par
% glossary-block:B00767:end
% glossary-block:B00768:begin
\needspace{3\baselineskip}
\textbf{greater} — \textbf{māra} (adj.). greater; larger. Comparative adjective. greater; more important. Value/importance in the life/raiment comparison. Mā expresses more in amount or number in its separate uses; māra expresses a greater amount. Comparative swīðor expresses more intensely or preferential rather. Use an agreeing comparative adjective for greater quality/value; these uses are not interchangeable.\par
% glossary-block:B00768:end
% glossary-block:B00769:begin
\needspace{3\baselineskip}
\textbf{greater amount} — \textbf{māra} (adj.). more; a greater amount. Nominal/attributive quantity, distinct from more in number.\par
% glossary-block:B00769:end
% glossary-block:B00770:begin
\needspace{3\baselineskip}
\textbf{greatest} — \textbf{mǣst} (adj.). greatest; largest. Superlative adjective.\par
% glossary-block:B00770:end
% glossary-block:B00771:begin
\needspace{3\baselineskip}
\textbf{greatly} — \textbf{ðearle} (adv.). greatly; exceedingly. \textbf{swīðe} (adv.). Degree adverb swīðe; retain the separate great/large adjective micel.\par
% glossary-block:B00771:end
% glossary-block:B00772:begin
\needspace{3\baselineskip}
\textbf{greet} — \textbf{grētan} (v.). greet; take leave of. Accusative person/place; greeting and farewell depend on context. Greeting with accusative person; the same verb also has farewell/leave-taking uses.\par
% glossary-block:B00772:end
% glossary-block:B00773:begin
\needspace{3\baselineskip}
\textbf{grief} — \textbf{sārnes} (n.).\par
% glossary-block:B00773:end
% glossary-block:B00774:begin
\needspace{3\baselineskip}
\textbf{grieve} — \textbf{ofðyncean} (v.). grieve; displease. Impersonal third singular with dative experiencer and a content clause. Impersonal third-singular predicate, dative experiencer, and ðæt clause for the cause.\par
% glossary-block:B00774:end
% glossary-block:B00775:begin
\needspace{3\baselineskip}
\textbf{grievous} — \textbf{hęfig} (adj.).\par
% glossary-block:B00775:end
% glossary-block:B00776:begin
\needspace{3\baselineskip}
\textbf{group} — \textbf{flocc} (n.). group; company.\par
% glossary-block:B00776:end
% glossary-block:B00777:begin
\needspace{3\baselineskip}
\textbf{grove} — \textbf{wudubearu} (n.). forest; grove.\par
% glossary-block:B00777:end
% glossary-block:B00778:begin
\needspace{3\baselineskip}
\textbf{grow} — \textbf{weaxan} (v.). Intransitive increase in these commands/wishes; not the transitive īecan increase something.\par
% glossary-block:B00778:end
% glossary-block:B00779:begin
\needspace{3\baselineskip}
\textbf{guest} — \textbf{cuma} (n.). stranger; guest.\par
% glossary-block:B00779:end
% glossary-block:B00780:begin
\needspace{3\baselineskip}
\textbf{guest-house} — \textbf{giesthūs} (n.). inn; guest-house.\par
% glossary-block:B00780:end
% glossary-block:B00781:begin
\needspace{3\baselineskip}
\textbf{guide} — \textbf{gewissian} (v.). guide; direct.\par
% glossary-block:B00781:end
% glossary-block:B00782:begin
\needspace{3\baselineskip}
\textbf{guilty} — \textbf{scyldig} (adj.). guilty; culpable.\par
% glossary-block:B00782:end
% glossary-block:B00783:begin
\needspace{3\baselineskip}
\textbf{gymnasium} — \textbf{bæðstęde} (n.). bathing place; gymnasium. Gymnasium is the contextual interpretation in the supplied reading.\par
% glossary-block:B00783:end
\pdfbookmark[1]{H}{outline-B00784}
% glossary-block:B00784:begin
\needspace{5\baselineskip}
\subsection*{H}
% glossary-block:B00784:end
% glossary-block:B00785:begin
\textbf{habit} — \textbf{gewuna} (n.). custom; habit.\par
% glossary-block:B00785:end
% glossary-block:B00786:begin
\needspace{3\baselineskip}
\textbf{had} — Auxiliary habban combines with a transitive past participle in the taught completed-action/result construction. English have/has/had does not require habban in every perfect; a simple Old English tense may express the intended time relation. \textbf{habban} (v.). Possessive habban takes an accusative possession. Auxiliary habban combines with a transitive past participle in the taught perfect construction; it is not an automatic translation of every English have.\par
% glossary-block:B00786:end
% glossary-block:B00787:begin
\needspace{3\baselineskip}
\textbf{had …} — Preserve the action's time relation. Simple preterite, sometimes with ǣr, often renders English perfect/pluperfect; use an analytic construction only where taught/requested.\par
% glossary-block:B00787:end
% glossary-block:B00788:begin
\needspace{3\baselineskip}
\textbf{had come} — Wesan + the assigned intransitive past participle supplies the taught compound construction. With cumen come it can render have/had come; English have is not an extra word to add to that construction.\par
% glossary-block:B00788:end
% glossary-block:B00789:begin
\needspace{3\baselineskip}
\textbf{had rather} — \textbf{lēofra} (adj.). preferable; dearer. Dative experiencer + is lēofre + infinitive in the taught preference frame. The preference construction is dative experiencer + is lēofre + plain infinitive, is preferable to someone. It renders had/would rather without requiring past time. Wician in the supplied example is the action preferred, sojourn/live with, not a word meaning prefer.\par
% glossary-block:B00789:end
% glossary-block:B00790:begin
\needspace{3\baselineskip}
\textbf{had to} — \textbf{sculan} (v.). sculan + plain infinitive for obligation; English have to is not possessive habban. Supplied have something to offer in Unit 6 is a different construction.\par
% glossary-block:B00790:end
% glossary-block:B00791:begin
\needspace{3\baselineskip}
\textbf{hail} — \textbf{hægel} (n.).\par
% glossary-block:B00791:end
% glossary-block:B00792:begin
\needspace{3\baselineskip}
\textbf{hair} — \textbf{hǣr} (n.). hair; a hair.\par
% glossary-block:B00792:end
% glossary-block:B00793:begin
\needspace{3\baselineskip}
\textbf{hale} — \textbf{hāl} (adj.). hale; well. Supplied greeting. hale; healthy.\par
% glossary-block:B00793:end
% glossary-block:B00794:begin
\needspace{3\baselineskip}
\textbf{half} — \textbf{healf} (adj.).\par
% glossary-block:B00794:end
% glossary-block:B00795:begin
\needspace{3\baselineskip}
\textbf{hall} — \textbf{heall} (n.).\par
% glossary-block:B00795:end
% glossary-block:B00796:begin
\needspace{3\baselineskip}
\textbf{hallow} — \textbf{gehālgian} (v.). hallow; consecrate.\par
% glossary-block:B00796:end
% glossary-block:B00797:begin
\needspace{3\baselineskip}
\textbf{halt} — \textbf{healt} (adj.). lame; halt. Adjective describing impaired walking; not the command halt/stop.\par
% glossary-block:B00797:end
% glossary-block:B00798:begin
\needspace{3\baselineskip}
\textbf{hammer} — \textbf{slęcg} (n.).\par
% glossary-block:B00798:end
% glossary-block:B00799:begin
\needspace{3\baselineskip}
\textbf{hand} — \textbf{hand} (n.). hand; side. Side in right-/left-hand expressions.\par
% glossary-block:B00799:end
% glossary-block:B00800:begin
\needspace{3\baselineskip}
\textbf{handmaiden} — \textbf{ðīnen} (n.).\par
% glossary-block:B00800:end
% glossary-block:B00801:begin
\needspace{3\baselineskip}
\textbf{happen} — \textbf{gelimpan} (v.). happen; befall. Supplied preterite subjunctive in might befall us. happen; come about. Supplied preterite source form. happen; come about.\par
% glossary-block:B00801:end
% glossary-block:B00802:begin
\needspace{3\baselineskip}
\textbf{happy} — \textbf{ēadig} (adj.). happy; blessed.\par
% glossary-block:B00802:end
% glossary-block:B00803:begin
\needspace{3\baselineskip}
\textbf{hard to count} — \textbf{earfoðrīme} (adj.). innumerable; hard to count. Assigned source adjective.\par
% glossary-block:B00803:end
% glossary-block:B00804:begin
\needspace{3\baselineskip}
\textbf{harp} — \textbf{hearpe} (n.).  \textbf{hearpian} (v.). Supplied genitive harper/harping construction, with instrument; do not reconstruct it from the infinitive alone. Preserve the supplied genitive harper/harping expression; headword alone does not reconstruct its syntax.\par
% glossary-block:B00804:end
% glossary-block:B00805:begin
\needspace{3\baselineskip}
\textbf{harper} — \textbf{hearpere} (n.).\par
% glossary-block:B00805:end
% glossary-block:B00806:begin
\needspace{3\baselineskip}
\textbf{harp pick} — \textbf{hearpenægl} (n.). plectrum; harp pick.\par
% glossary-block:B00806:end
% glossary-block:B00807:begin
\needspace{3\baselineskip}
\textbf{harp strings} — \textbf{hearpestręng} (n.).\par
% glossary-block:B00807:end
% glossary-block:B00808:begin
\needspace{3\baselineskip}
\textbf{harvest} — \textbf{hærfest} (n.). Masculine noun; rich in uses welig on + dative fruits.\par
% glossary-block:B00808:end
% glossary-block:B00809:begin
\needspace{3\baselineskip}
\textbf{has} — Auxiliary habban combines with a transitive past participle in the taught completed-action/result construction. English have/has/had does not require habban in every perfect; a simple Old English tense may express the intended time relation. \textbf{habban} (v.). Possessive habban takes an accusative possession. Auxiliary habban combines with a transitive past participle in the taught perfect construction; it is not an automatic translation of every English have.\par
% glossary-block:B00809:end
% glossary-block:B00810:begin
\needspace{3\baselineskip}
\textbf{hatchet} — \textbf{adesa} (n.).\par
% glossary-block:B00810:end
% glossary-block:B00811:begin
\needspace{3\baselineskip}
\textbf{hate} — \textbf{hatian} (v.).\par
% glossary-block:B00811:end
% glossary-block:B00812:begin
\needspace{3\baselineskip}
\textbf{have} — \textbf{habban} (v.). have; possess. Accusative possession; distinguish the perfect auxiliary. Possessive habban takes an accusative possession. Auxiliary habban combines with a transitive past participle in the taught perfect construction; it is not an automatic translation of every English have. Auxiliary habban combines with a transitive past participle in the taught completed-action/result construction. English have/has/had does not require habban in every perfect; a simple Old English tense may express the intended time relation. \textbf{āgan} (v.). have; possess. Accusative possession; positive present and preterite indicative in the taught scope.\par
% glossary-block:B00812:end
% glossary-block:B00813:begin
\needspace{3\baselineskip}
\textbf{have as wife} — \textbf{tō} (prep.). habban + accusative woman + tō wīfe in the supplied role frame.\par
% glossary-block:B00813:end
% glossary-block:B00814:begin
\needspace{3\baselineskip}
\textbf{have come} — Wesan + the assigned intransitive past participle supplies the taught compound construction. With cumen come it can render have/had come; English have is not an extra word to add to that construction.\par
% glossary-block:B00814:end
% glossary-block:B00815:begin
\needspace{3\baselineskip}
\textbf{have compassion on} — \textbf{gemiltsian} (v.). pity; have compassion on. Dative person. gemiltsian + dative person. A wish uses the appropriate subjunctive, not a separate have verb.\par
% glossary-block:B00815:end
% glossary-block:B00816:begin
\needspace{3\baselineskip}
\textbf{have … (completed action)} — Preserve the action's time relation. Simple preterite, sometimes with ǣr, often renders English perfect/pluperfect; use an analytic construction only where taught/requested.\par
% glossary-block:B00816:end
% glossary-block:B00817:begin
\needspace{3\baselineskip}
\textbf{have dominion} — \textbf{rīcsian} (v.). rule; have dominion. Used absolutely in the taught frame; no new object government. rīcsian can be used absolutely; English have/bear need no independent translation.\par
% glossary-block:B00817:end
% glossary-block:B00818:begin
\needspace{3\baselineskip}
\textbf{have dominion over} — \textbf{ofer} (prep.). over. Accusative dependent in the taught authority frame.\par
% glossary-block:B00818:end
% glossary-block:B00819:begin
\needspace{3\baselineskip}
\textbf{have mercy on} — \textbf{gemiltsian} (v.). pity; have compassion on. Dative person. gemiltsian + dative person. A wish uses the appropriate subjunctive, not a separate have verb.\par
% glossary-block:B00819:end
% glossary-block:B00820:begin
\needspace{3\baselineskip}
\textbf{have not} — \textbf{nabban} (v.).\par
% glossary-block:B00820:end
% glossary-block:B00821:begin
\needspace{3\baselineskip}
\textbf{have on} — \textbf{on} (adv.). on. Adverb in the cloak-wearing relative clause.\par
% glossary-block:B00821:end
% glossary-block:B00822:begin
\needspace{3\baselineskip}
\textbf{have speech with} — \textbf{wið} (prep.). with. Supplied conversational counterpart; do not infer hostility.\par
% glossary-block:B00822:end
% glossary-block:B00823:begin
\needspace{3\baselineskip}
\textbf{have to} — \textbf{sculan} (v.). sculan + plain infinitive for obligation; English have to is not possessive habban. Supplied have something to offer in Unit 6 is a different construction.\par
% glossary-block:B00823:end
% glossary-block:B00824:begin
\needspace{3\baselineskip}
\textbf{he} — Hē is masculine singular nominative. Third-person pronouns follow the grammatical gender of their referent. Grammatical gender and antecedent determine reference; it may be represented by masculine/feminine forms for a gendered inanimate noun.\par
% glossary-block:B00824:end
% glossary-block:B00825:begin
\needspace{3\baselineskip}
\textbf{head} — \textbf{hēafod} (n.).\par
% glossary-block:B00825:end
% glossary-block:B00826:begin
\needspace{3\baselineskip}
\textbf{healthy} — \textbf{hāl} (adj.). hale; healthy.\par
% glossary-block:B00826:end
% glossary-block:B00827:begin
\needspace{3\baselineskip}
\textbf{hear} — \textbf{gehīeran} (v.). Supplied prefixed third-singular preterite. hear; listen. \textbf{hīeran} (v.). hear an action. Accusative perceived actor plus plain infinitive in the taught perception construction. \textbf{hlystan} (v.). listen to; heed. Genitive person or thing listened to; hear is the contextual prompt wording. hlystan governs genitive of the speech heard; do not copy hīeran's ordinary accusative frame.\par
% glossary-block:B00827:end
% glossary-block:B00828:begin
\needspace{3\baselineskip}
\textbf{hear someone do something} — \textbf{hīeran} (v.). hear an action. Accusative perceived actor plus plain infinitive in the taught perception construction.\par
% glossary-block:B00828:end
% glossary-block:B00829:begin
\needspace{3\baselineskip}
\textbf{heart} — \textbf{heorte} (n.).  \textbf{mōd} (n.). heart; mind.\par
% glossary-block:B00829:end
% glossary-block:B00830:begin
\needspace{3\baselineskip}
\textbf{heat} — \textbf{hǣtu} (n.).\par
% glossary-block:B00830:end
% glossary-block:B00831:begin
\needspace{3\baselineskip}
\textbf{heathen people} — \textbf{hǣðen} (adj.).\par
% glossary-block:B00831:end
% glossary-block:B00832:begin
\needspace{3\baselineskip}
\textbf{heaven} — \textbf{heofon} (n.). heaven; sky. \textbf{heofone} (n.). heaven; sky. Weak feminine lexical alternative in the source.\par
% glossary-block:B00832:end
% glossary-block:B00833:begin
\needspace{3\baselineskip}
\textbf{heavenly} — \textbf{heofonlic} (adj.).  \textbf{heofonlice} (adv.). from heaven; heavenly. Source expression at Cook 143.17.\par
% glossary-block:B00833:end
% glossary-block:B00834:begin
\needspace{3\baselineskip}
\textbf{heavenly fire} — \textbf{heofonfȳr} (n.). fire from heaven; heavenly fire.\par
% glossary-block:B00834:end
% glossary-block:B00835:begin
\needspace{3\baselineskip}
\textbf{heavenly powers} — \textbf{mægen} (n.). power; strength. Heavenly powers in 37.2 is the contextual plural use.\par
% glossary-block:B00835:end
% glossary-block:B00836:begin
\needspace{3\baselineskip}
\textbf{heed} — \textbf{hlystan} (v.). listen to; heed. Genitive person or thing listened to; hear is the contextual prompt wording. hlystan governs genitive of the speech heard; do not copy hīeran's ordinary accusative frame.\par
% glossary-block:B00836:end
% glossary-block:B00837:begin
\needspace{3\baselineskip}
\textbf{height} — \textbf{hēanes} (n.). height; highest point.\par
% glossary-block:B00837:end
% glossary-block:B00838:begin
\needspace{3\baselineskip}
\textbf{heights} — \textbf{hēanes} (n.). heights; highest places.\par
% glossary-block:B00838:end
% glossary-block:B00839:begin
\needspace{3\baselineskip}
\textbf{hell} — \textbf{hęll} (n.).\par
% glossary-block:B00839:end
% glossary-block:B00840:begin
\needspace{3\baselineskip}
\textbf{help} — \textbf{fultum} (n.). help; aid. help; support. \textbf{gefultumian} (v.). help; assist. \textbf{helpan} (v.). Dative person helped.\par
% glossary-block:B00840:end
% glossary-block:B00841:begin
\needspace{3\baselineskip}
\textbf{hence} — \textbf{heonan} (adv.). hence; from here.\par
% glossary-block:B00841:end
% glossary-block:B00842:begin
\needspace{3\baselineskip}
\textbf{hence forth} — \textbf{heonanforð} (adv.). heonanforð as a whole adverb.\par
% glossary-block:B00842:end
% glossary-block:B00843:begin
\needspace{3\baselineskip}
\textbf{henceforth} — \textbf{heonanforð} (adv.). heonanforð as a whole adverb.\par
% glossary-block:B00843:end
% glossary-block:B00844:begin
\needspace{3\baselineskip}
\textbf{her} — Hīe is feminine singular accusative; hiere is dative or genitive. Possessive hiere remains invariant before the possessed noun. \textbf{hē} (personal pron.). his; her; its; their. Invariant genitive possessors, unlike agreeing possessive adjectives. Invariant personal genitive possessor; distinguish the same surface used as a governed personal case. Grammatical gender and antecedent determine reference; it may be represented by masculine/feminine forms for a gendered inanimate noun.\par
% glossary-block:B00844:end
% glossary-block:B00845:begin
\needspace{3\baselineskip}
\textbf{herb} — \textbf{wyrt} (n.). herb; plant. Cook specifically routes these spices to wyrt; wyrtgemang is the separate earlier spice word.\par
% glossary-block:B00845:end
% glossary-block:B00846:begin
\needspace{3\baselineskip}
\textbf{here} — \textbf{hēr} (adv.).  \textbf{hider} (adv.). hither; here. Motion toward the speaker, distinct from location hēr.\par
% glossary-block:B00846:end
% glossary-block:B00847:begin
\needspace{3\baselineskip}
\textbf{herself} — An ordinary reflexive uses the personal pronoun in the case governed by its verb or preposition. Agreeing self adds the taught emphasis; English -self does not automatically require a separate self. Ordinary reflexive personal pronouns take the governed case; agreeing self adds the taught emphasis. English -self does not automatically require self.\par
% glossary-block:B00847:end
% glossary-block:B00848:begin
\needspace{3\baselineskip}
\textbf{hew} — \textbf{hēawan} (v.).\par
% glossary-block:B00848:end
% glossary-block:B00849:begin
\needspace{3\baselineskip}
\textbf{hide} — \textbf{diernan} (v.). conceal; hide. \textbf{hȳdan} (v.). hide; conceal. Accusative thing hidden; wesan forms the assigned passive.\par
% glossary-block:B00849:end
% glossary-block:B00850:begin
\needspace{3\baselineskip}
\textbf{Hierusalem} — \textbf{Hierusalem} (proper name).\par
% glossary-block:B00850:end
% glossary-block:B00851:begin
\needspace{3\baselineskip}
\textbf{high} — \textbf{hēah} (adj.). high in rank; exalted. Rank/status sense in the supplied source. \textbf{hēah} (adv.). Adverbial height, distinct from an agreeing adjective.\par
% glossary-block:B00851:end
% glossary-block:B00852:begin
\needspace{3\baselineskip}
\textbf{high-born} — \textbf{æðelboren} (adj.).\par
% glossary-block:B00852:end
% glossary-block:B00853:begin
\needspace{3\baselineskip}
\textbf{high cliff} — \textbf{hēaclif} (n.).\par
% glossary-block:B00853:end
% glossary-block:B00854:begin
\needspace{3\baselineskip}
\textbf{Highest} — \textbf{hēah} (adj.). the Highest. Divine title.\par
% glossary-block:B00854:end
% glossary-block:B00855:begin
\needspace{3\baselineskip}
\textbf{highest places} — \textbf{hēanes} (n.). heights; highest places.\par
% glossary-block:B00855:end
% glossary-block:B00856:begin
\needspace{3\baselineskip}
\textbf{highest point} — \textbf{hēanes} (n.). height; highest point.\par
% glossary-block:B00856:end
% glossary-block:B00857:begin
\needspace{3\baselineskip}
\textbf{high office} — \textbf{ār} (n.). honour; high position. Contextual rank or office, not a new general meaning of every ār form.\par
% glossary-block:B00857:end
% glossary-block:B00858:begin
\needspace{3\baselineskip}
\textbf{high position} — \textbf{ār} (n.). honour; high position. Contextual rank or office, not a new general meaning of every ār form.\par
% glossary-block:B00858:end
% glossary-block:B00859:begin
\needspace{3\baselineskip}
\textbf{hight} — \textbf{hātan} (v.). be called; be named. Special hātte form in the assigned naming expression. The special form hātte means is/was called or named, according to context; do not substitute an ordinary active hātan form.\par
% glossary-block:B00859:end
% glossary-block:B00860:begin
\needspace{3\baselineskip}
\textbf{hill} — \textbf{dūn} (n.). mountain; hill. \textbf{hyll} (n.).\par
% glossary-block:B00860:end
% glossary-block:B00861:begin
\needspace{3\baselineskip}
\textbf{him} — Hine is masculine singular accusative; him is masculine/neuter singular dative. Choose the governed case. Grammatical gender and antecedent determine reference; it may be represented by masculine/feminine forms for a gendered inanimate noun.\par
% glossary-block:B00861:end
% glossary-block:B00862:begin
\needspace{3\baselineskip}
\textbf{himself} — An ordinary reflexive uses the personal pronoun in the case governed by its verb or preposition. Agreeing self adds the taught emphasis; English -self does not automatically require a separate self. Ordinary reflexive personal pronouns take the governed case; agreeing self adds the taught emphasis. English -self does not automatically require self.\par
% glossary-block:B00862:end
% glossary-block:B00863:begin
\needspace{3\baselineskip}
\textbf{hire} — \textbf{hȳrian} (v.).  \textbf{mēd} (n.). hire; wages paid for work.\par
% glossary-block:B00863:end
% glossary-block:B00864:begin
\needspace{3\baselineskip}
\textbf{his} — \textbf{hē} (personal pron.). his; her; its; their. Invariant genitive possessors, unlike agreeing possessive adjectives. Invariant personal genitive possessor; distinguish the same surface used as a governed personal case.\par
% glossary-block:B00864:end
% glossary-block:B00865:begin
\needspace{3\baselineskip}
\textbf{hither} — \textbf{hider} (adv.). hither; here. Motion toward the speaker, distinct from location hēr.\par
% glossary-block:B00865:end
% glossary-block:B00866:begin
\needspace{3\baselineskip}
\textbf{hoard} — \textbf{hord} (n.). hoard; treasure.\par
% glossary-block:B00866:end
% glossary-block:B00867:begin
\needspace{3\baselineskip}
\textbf{hoard of gold} — \textbf{goldhord} (n.). treasure; hoard of gold.\par
% glossary-block:B00867:end
% glossary-block:B00868:begin
\needspace{3\baselineskip}
\textbf{hold} — \textbf{healdan} (v.).\par
% glossary-block:B00868:end
% glossary-block:B00869:begin
\needspace{3\baselineskip}
\textbf{hold chief place} — Healdan ealdorscipe means hold chief place or primacy: healdan takes accusative ealdorscipe. Retain the exercise’s omission of the article.\par
% glossary-block:B00869:end
% glossary-block:B00870:begin
\needspace{3\baselineskip}
\textbf{hole} — \textbf{hol} (n.).  \textbf{ðyrel} (n.).\par
% glossary-block:B00870:end
% glossary-block:B00871:begin
\needspace{3\baselineskip}
\textbf{holy} — \textbf{hālig} (adj.).  \textbf{Sanctus} (foreign adjective). Saint; holy. Latin word in the source title.\par
% glossary-block:B00871:end
% glossary-block:B00872:begin
\needspace{3\baselineskip}
\textbf{holy one} — \textbf{hālga} (n.). saint; holy one.\par
% glossary-block:B00872:end
% glossary-block:B00873:begin
\needspace{3\baselineskip}
\textbf{holy ones} — \textbf{hālig} (adj.). holy ones; saints. Substantival adjective in the source.\par
% glossary-block:B00873:end
% glossary-block:B00874:begin
\needspace{3\baselineskip}
\textbf{home} — \textbf{hām} (adv.). home; homeward. Directional adverb; distinct from the noun home. \textbf{hām} (n.). home; dwelling.\par
% glossary-block:B00874:end
% glossary-block:B00875:begin
\needspace{3\baselineskip}
\textbf{homeland} — \textbf{ēðel} (n.). native land; homeland. ēðel in these home-country references; land remains a separate ordinary country word.\par
% glossary-block:B00875:end
% glossary-block:B00876:begin
\needspace{3\baselineskip}
\textbf{homeward} — \textbf{hām} (adv.). home; homeward. Directional adverb; distinct from the noun home. \textbf{hāmweard} (adverb or particle). Supplied hāmweard with wesan; distinguish noun hām and directional adverb hām.\par
% glossary-block:B00876:end
% glossary-block:B00877:begin
\needspace{3\baselineskip}
\textbf{homewards} — \textbf{hām} (adv.). home; homeward. Directional adverb; distinct from the noun home.\par
% glossary-block:B00877:end
% glossary-block:B00878:begin
\needspace{3\baselineskip}
\textbf{honey} — \textbf{hunig} (n.).\par
% glossary-block:B00878:end
% glossary-block:B00879:begin
\needspace{3\baselineskip}
\textbf{honor} — \textbf{ār} (n.). honour; high position. Contextual rank or office, not a new general meaning of every ār form.\par
% glossary-block:B00879:end
% glossary-block:B00880:begin
\needspace{3\baselineskip}
\textbf{honour} — \textbf{ār} (n.). honour; high position. Contextual rank or office, not a new general meaning of every ār form. \textbf{weorðscipe} (n.). honour; dignity.\par
% glossary-block:B00880:end
% glossary-block:B00881:begin
\needspace{3\baselineskip}
\textbf{hook} — \textbf{angel} (n.). Compound English spelling and optional fish- lead to the same fishing noun; distinguish ęngel angel.\par
% glossary-block:B00881:end
% glossary-block:B00882:begin
\needspace{3\baselineskip}
\textbf{hope} — \textbf{tōhopa} (n.).\par
% glossary-block:B00882:end
% glossary-block:B00883:begin
\needspace{3\baselineskip}
\textbf{hopelessness} — \textbf{ōrmōdnes} (n.). despair; hopelessness.\par
% glossary-block:B00883:end
% glossary-block:B00884:begin
\needspace{3\baselineskip}
\textbf{horn} — \textbf{horn} (n.). horn; strength; dignity. Strength/dignity is the figurative force in 37.14, not a replacement for the literal source horn.\par
% glossary-block:B00884:end
% glossary-block:B00885:begin
\needspace{3\baselineskip}
\textbf{horse} — \textbf{hors} (n.).\par
% glossary-block:B00885:end
% glossary-block:B00886:begin
\needspace{3\baselineskip}
\textbf{host} — \textbf{weorod} (n.). host; army.\par
% glossary-block:B00886:end
% glossary-block:B00887:begin
\needspace{3\baselineskip}
\textbf{hot} — \textbf{hāt} (adj.).\par
% glossary-block:B00887:end
% glossary-block:B00888:begin
\needspace{3\baselineskip}
\textbf{hour} — \textbf{tīd} (n.). Time measure in the supplied intervals.\par
% glossary-block:B00888:end
% glossary-block:B00889:begin
\needspace{3\baselineskip}
\textbf{house} — \textbf{hūs} (n.).\par
% glossary-block:B00889:end
% glossary-block:B00890:begin
\needspace{3\baselineskip}
\textbf{how} — \textbf{hū} (adv.). how?.\par
% glossary-block:B00890:end
% glossary-block:B00891:begin
\needspace{3\baselineskip}
\textbf{human being} — \textbf{mann} (n.; impersonal pronoun). man; human being.\par
% glossary-block:B00891:end
% glossary-block:B00892:begin
\needspace{3\baselineskip}
\textbf{humble} — \textbf{ēaðmōd} (adj.). lowly; humble. \textbf{hīenan} (v.). oppress; humble.\par
% glossary-block:B00892:end
% glossary-block:B00893:begin
\needspace{3\baselineskip}
\textbf{humble people} — \textbf{ēaðmōd} (adj.). humble people; the lowly. Substantival adjective.\par
% glossary-block:B00893:end
% glossary-block:B00894:begin
\needspace{3\baselineskip}
\textbf{humility} — \textbf{ēaðmōdnes} (n.).\par
% glossary-block:B00894:end
% glossary-block:B00895:begin
\needspace{3\baselineskip}
\textbf{hundred} — \textbf{hund} (cardinal). Only the complete supplied hundred-sheep expression at its first appearance; later numerical frames remain locally specified. \textbf{hundtēontig} (numeral). Within the complete supplied 144,000 expression; do not infer an independent high-number system. High quantities retain the complete local numerical expressions; a glossary headword does not replace the supplied genitive/counting construction.\par
% glossary-block:B00895:end
% glossary-block:B00896:begin
\needspace{3\baselineskip}
\textbf{hundredfold} — \textbf{hundtēontig} (numeral). a hundredfold. A hundred in return, in the supplied recompense expression.\par
% glossary-block:B00896:end
% glossary-block:B00897:begin
\needspace{3\baselineskip}
\textbf{hunger} — \textbf{hunger} (n.). famine; hunger. \textbf{hyngran} (v.). be hungry. Impersonal construction with accusative experiencer.\par
% glossary-block:B00897:end
% glossary-block:B00898:begin
\needspace{3\baselineskip}
\textbf{hungering} — \textbf{hyngran} (v.). hungry; hungering. Present participle used substantivally in the assigned frame.\par
% glossary-block:B00898:end
% glossary-block:B00899:begin
\needspace{3\baselineskip}
\textbf{hungry} — \textbf{hungrig} (adj.).  \textbf{hyngran} (v.). hungry; hungering. Present participle used substantivally in the assigned frame.\par
% glossary-block:B00899:end
% glossary-block:B00900:begin
\needspace{3\baselineskip}
\textbf{hungry people} — \textbf{hungrig} (adj.). Substantival adjective.\par
% glossary-block:B00900:end
% glossary-block:B00901:begin
\needspace{3\baselineskip}
\textbf{hurt} — \textbf{scęððan} (v.). hurt; injure. Dative person harmed.\par
% glossary-block:B00901:end
% glossary-block:B00902:begin
\needspace{3\baselineskip}
\textbf{husband} — \textbf{wer} (n.).\par
% glossary-block:B00902:end
\pdfbookmark[1]{I}{outline-B00903}
% glossary-block:B00903:begin
\needspace{5\baselineskip}
\subsection*{I}
% glossary-block:B00903:end
% glossary-block:B00904:begin
\textbf{I} — Ic is the nominative singular subject I. The governed object form me is mē in both dative and accusative. Select case from the verb/preposition and preserve dual versus plural; dual subjects use plural verb forms.\par
% glossary-block:B00904:end
% glossary-block:B00905:begin
\needspace{3\baselineskip}
\textbf{ice} — \textbf{īs} (n.).\par
% glossary-block:B00905:end
% glossary-block:B00906:begin
\needspace{3\baselineskip}
\textbf{idle} — \textbf{īdel} (adj.). Distinct from empty/void.\par
% glossary-block:B00906:end
% glossary-block:B00907:begin
\needspace{3\baselineskip}
\textbf{idleness} — \textbf{īdelnes} (n.).\par
% glossary-block:B00907:end
% glossary-block:B00908:begin
\needspace{3\baselineskip}
\textbf{if} — \textbf{gif} (conj.). Mood is determined by the open, tentative or counterfactual condition, not by gif alone. gif does not by itself choose a mood; open, tentative and counterfactual conditions have distinct taught frames.\par
% glossary-block:B00908:end
% glossary-block:B00909:begin
\needspace{3\baselineskip}
\textbf{ill} — \textbf{yfele} (adv.). Adverb modifying doing; distinguish the adjective yfel.\par
% glossary-block:B00909:end
% glossary-block:B00910:begin
\needspace{3\baselineskip}
\textbf{illuminate} — \textbf{ālīehtan} (v.). illuminate; give light to. Accusative person or thing illuminated. ālīehtan takes an accusative object, not a dative inferred from English to.\par
% glossary-block:B00910:end
% glossary-block:B00911:begin
\needspace{3\baselineskip}
\textbf{illumination} — \textbf{līehting} (n.).\par
% glossary-block:B00911:end
% glossary-block:B00912:begin
\needspace{3\baselineskip}
\textbf{illustration} — \textbf{bȳsen} (n.). example; illustration.\par
% glossary-block:B00912:end
% glossary-block:B00913:begin
\needspace{3\baselineskip}
\textbf{illustrious} — \textbf{æðele} (adj.). noble; illustrious.\par
% glossary-block:B00913:end
% glossary-block:B00914:begin
\needspace{3\baselineskip}
\textbf{image} — \textbf{andlīcnes} (n.).\par
% glossary-block:B00914:end
% glossary-block:B00915:begin
\needspace{3\baselineskip}
\textbf{imagine} — \textbf{āsmēan} (v.). conceive. Exact supplied infinitive āsmēan in the heart-conceiving expression; not an independently assigned full paradigm. Supplied infinitive āsmēan; the heart is its subject.\par
% glossary-block:B00915:end
% glossary-block:B00916:begin
\needspace{3\baselineskip}
\textbf{immediately} — \textbf{sōna} (adv.). immediately; at once. sōna carries soon in this future-reference prompt.\par
% glossary-block:B00916:end
% glossary-block:B00917:begin
\needspace{3\baselineskip}
\textbf{implore} — \textbf{hālsian} (v.). entreat; implore.\par
% glossary-block:B00917:end
% glossary-block:B00918:begin
\needspace{3\baselineskip}
\textbf{in} — \textbf{in} (prep.). in; on. With dative in the exact locations/circumstances. in; inward. Adverbial direction. \textbf{on} (prep.). in; on. Dative for the taught ordinary location construction. Supported expression in his majesty/glory. in; with respect to. Quality or domain of description. Supplied circumstance of the storm. Accusative person believed/trusted in. among; in. Dative people among whom friends are obtained. Circumstance of the day's heat in the supplied source. during. Dative time expression. in; wearing. Supplied on scēapa gegierelum, in sheep's garments. Source/object of joy in the supplied frame. in; amid. Circumstances of the stated unrest.\par
% glossary-block:B00918:end
% glossary-block:B00919:begin
\needspace{3\baselineskip}
\textbf{in accordance with} — \textbf{tō} (prep.). A will as the goal of guidance in the supplied prayer.\par
% glossary-block:B00919:end
% glossary-block:B00920:begin
\needspace{3\baselineskip}
\textbf{in addition to that} — \textbf{ðǣr-tō-ēacan} (adverbial expression). in addition to that; besides.\par
% glossary-block:B00920:end
% glossary-block:B00921:begin
\needspace{3\baselineskip}
\textbf{inasmuch as} — Swā lange swā is the complete supplied phrase as long as; in the assigned biblical context it has inasmuch-as force. Swā lange swā is the complete supplied phrase as long as; in the assigned biblical context it has inasmuch-as force. Lange alone means long.\par
% glossary-block:B00921:end
% glossary-block:B00922:begin
\needspace{3\baselineskip}
\textbf{in a storm} — \textbf{on} (prep.). in. Supplied circumstance of the storm.\par
% glossary-block:B00922:end
% glossary-block:B00923:begin
\needspace{3\baselineskip}
\textbf{in circuit of} — \textbf{ymbhwyrft} (n.). Supplied on ymbhwyrfte + genitive landmark in this exercise.\par
% glossary-block:B00923:end
% glossary-block:B00924:begin
\needspace{3\baselineskip}
\textbf{incline} — \textbf{āhyldan} (v.). lay; incline.\par
% glossary-block:B00924:end
% glossary-block:B00925:begin
\needspace{3\baselineskip}
\textbf{incline oneself} — \textbf{ālūtan} (v.). bow; incline oneself. Reflexive meaning in the assigned use.\par
% glossary-block:B00925:end
% glossary-block:B00926:begin
\needspace{3\baselineskip}
\textbf{in control} — \textbf{on} (prep.). in one's control. on + dative possessed geweald.\par
% glossary-block:B00926:end
% glossary-block:B00927:begin
\needspace{3\baselineskip}
\textbf{in control of} — On + dative possessed geweald expresses in/under someone’s control. In the assigned habban construction, the things controlled are its accusative object; the controller is expressed within the geweald phrase.\par
% glossary-block:B00927:end
% glossary-block:B00928:begin
\needspace{3\baselineskip}
\textbf{increase} — \textbf{īecan} (v.).  \textbf{weaxan} (v.). Intransitive increase in these commands/wishes; not the transitive īecan increase something.\par
% glossary-block:B00928:end
% glossary-block:B00929:begin
\needspace{3\baselineskip}
\textbf{indeed} — \textbf{sōðlice} (adv.). truly; indeed. Adverbial confirmation; distinct from discourse connective sōðlice where the key uses but/now.\par
% glossary-block:B00929:end
% glossary-block:B00930:begin
\needspace{3\baselineskip}
\textbf{indescribable} — \textbf{unāsęcgende} (adj.).\par
% glossary-block:B00930:end
% glossary-block:B00931:begin
\needspace{3\baselineskip}
\textbf{in exchange for} — \textbf{for} (prep.). With dative. \textbf{wið} (prep.). in return for. With dative; distinct from opponent/hostility. wið + dative exchange, distinct from wið + accusative opponent.\par
% glossary-block:B00931:end
% glossary-block:B00932:begin
\needspace{3\baselineskip}
\textbf{in fact} — \textbf{sōðlice} (adv.). Adverbial confirmation; distinct from discourse connective sōðlice where the key uses but/now.\par
% glossary-block:B00932:end
% glossary-block:B00933:begin
\needspace{3\baselineskip}
\textbf{infirmity} — \textbf{untrymnes} (n.).\par
% glossary-block:B00933:end
% glossary-block:B00934:begin
\needspace{3\baselineskip}
\textbf{inflict slaughter} — \textbf{ongeslēan} (v.). inflict slaughter; slay.\par
% glossary-block:B00934:end
% glossary-block:B00935:begin
\needspace{3\baselineskip}
\textbf{in forms of} — The supplied on + hīw expression means according to/in the forms of. Retain the whole supplied hīw phrase rather than treating on alone as a general according-to expression.\par
% glossary-block:B00935:end
% glossary-block:B00936:begin
\needspace{3\baselineskip}
\textbf{in front of} — \textbf{ætforan} (prep.). before; in front of. Dative landmark/person in the supplied frame; retain ætforan rather than replacing it with beforan. \textbf{be-foran} (prep.). before; in front of. Dative landmark in the static frame. before; in front of. Accusative landmark in the positioning construal; static dative can be a different valid construal.\par
% glossary-block:B00936:end
% glossary-block:B00937:begin
\needspace{3\baselineskip}
\textbf{in glory} — \textbf{on} (prep.). in. Supported expression in his majesty/glory.\par
% glossary-block:B00937:end
% glossary-block:B00938:begin
\needspace{3\baselineskip}
\textbf{in good health} — \textbf{gesund} (adj.). well; in good health.\par
% glossary-block:B00938:end
% glossary-block:B00939:begin
\needspace{3\baselineskip}
\textbf{ingot} — \textbf{węcg} (n.). metal; ingot. Supplied genitive plural dependent on veins.\par
% glossary-block:B00939:end
% glossary-block:B00940:begin
\needspace{3\baselineskip}
\textbf{in it} — \textbf{ðǣron} (adverb or particle). therein; in it. Supplied ðǣron points back to the sacks' location.\par
% glossary-block:B00940:end
% glossary-block:B00941:begin
\needspace{3\baselineskip}
\textbf{injure} — \textbf{scęððan} (v.). hurt; injure. Dative person harmed.\par
% glossary-block:B00941:end
% glossary-block:B00942:begin
\needspace{3\baselineskip}
\textbf{inn} — \textbf{giesthūs} (n.). inn; guest-house.\par
% glossary-block:B00942:end
% glossary-block:B00943:begin
\needspace{3\baselineskip}
\textbf{innocent} — \textbf{clǣne} (adj.). Source's blameless is contextual clǣne, not just physical clean.\par
% glossary-block:B00943:end
% glossary-block:B00944:begin
\needspace{3\baselineskip}
\textbf{innumerable} — \textbf{earfoðrīme} (adj.). innumerable; hard to count. Assigned source adjective.\par
% glossary-block:B00944:end
% glossary-block:B00945:begin
\needspace{3\baselineskip}
\textbf{in one's life} — \textbf{on} (prep.). during; in. Supplied lifetime expression.\par
% glossary-block:B00945:end
% glossary-block:B00946:begin
\needspace{3\baselineskip}
\textbf{in one's name} — On + dative naman with a genitive possessor means in someone’s name. The whole name phrase, not on alone, supplies this meaning.\par
% glossary-block:B00946:end
% glossary-block:B00947:begin
\needspace{3\baselineskip}
\textbf{in order that} — \textbf{ðæt} (conj.). so that; in order that. Purpose conjunction with the assigned subjunctive construction.\par
% glossary-block:B00947:end
% glossary-block:B00948:begin
\needspace{3\baselineskip}
\textbf{in order to} — Tō + inflected infinitive expresses the taught purpose, to/in order to. Plain infinitives after modals do not take an inserted tō; recipient to may instead be expressed by a dative.\par
% glossary-block:B00948:end
% glossary-block:B00949:begin
\needspace{3\baselineskip}
\textbf{in place of} — \textbf{for} (prep.). instead of; in place of. Accusative in the supplied substitution frame. for + accusative in this supplied substitution frame.\par
% glossary-block:B00949:end
% glossary-block:B00950:begin
\needspace{3\baselineskip}
\textbf{inquire after} — \textbf{æfter} (prep.; conj. in fixed expressions). after; in search of. Person sought or inquired after; keep this lexical frame distinct from temporal after.\par
% glossary-block:B00950:end
% glossary-block:B00951:begin
\needspace{3\baselineskip}
\textbf{in readiness for} — \textbf{ongēan} (adv.; prep.; conj.). by the time that. Complete supplied ongēan ðætte clause. Supplied ongēan ðætte + prospective subjunctive clause.\par
% glossary-block:B00951:end
% glossary-block:B00952:begin
\needspace{3\baselineskip}
\textbf{in response to} — \textbf{tō} (prep.). at; in response to. Only the exact at-the-prayers expression.\par
% glossary-block:B00952:end
% glossary-block:B00953:begin
\needspace{3\baselineskip}
\textbf{in return for} — \textbf{wið} (prep.). With dative; distinct from opponent/hostility. wið + dative exchange, distinct from wið + accusative opponent.\par
% glossary-block:B00953:end
% glossary-block:B00954:begin
\needspace{3\baselineskip}
\textbf{in search of} — \textbf{æfter} (prep.; conj. in fixed expressions). after; in search of. Person sought or inquired after; keep this lexical frame distinct from temporal after.\par
% glossary-block:B00954:end
% glossary-block:B00955:begin
\needspace{3\baselineskip}
\textbf{inside} — \textbf{innan} (adv.). within; inside.\par
% glossary-block:B00955:end
% glossary-block:B00956:begin
\needspace{3\baselineskip}
\textbf{instead of} — \textbf{for} (prep.). instead of; in place of. Accusative in the supplied substitution frame. for + accusative in this supplied substitution frame.\par
% glossary-block:B00956:end
% glossary-block:B00957:begin
\needspace{3\baselineskip}
\textbf{instruct} — \textbf{bebēodan} (v.). command; instruct. Dative addressee and a ðæt clause with subjunctive for the commanded action. \textbf{lǣran} (v.). advise; instruct. teach; instruct.\par
% glossary-block:B00957:end
% glossary-block:B00958:begin
\needspace{3\baselineskip}
\textbf{intelligibly} — Supplied superlative andgietfullicost means most intelligibly; do not use the positive andgietfullice for this degree.\par
% glossary-block:B00958:end
% glossary-block:B00959:begin
\needspace{3\baselineskip}
\textbf{intend} — \textbf{willan} (v.). Willan + plain infinitive expresses the reported intention in the assigned task. Distinguish intention from a mechanically inserted future auxiliary.\par
% glossary-block:B00959:end
% glossary-block:B00960:begin
\needspace{3\baselineskip}
\textbf{intend to} — \textbf{willan} (v.). Willan + plain infinitive expresses the reported intention in the assigned task. Distinguish intention from a mechanically inserted future auxiliary.\par
% glossary-block:B00960:end
% glossary-block:B00961:begin
\needspace{3\baselineskip}
\textbf{intent} — \textbf{geornfull} (adj.). occupied; intent. Source mental-attention context.\par
% glossary-block:B00961:end
% glossary-block:B00962:begin
\needspace{3\baselineskip}
\textbf{interpreter} — \textbf{wealhstōd} (n.). interpreter; translator.\par
% glossary-block:B00962:end
% glossary-block:B00963:begin
\needspace{3\baselineskip}
\textbf{interval} — \textbf{fæc} (n.). interval; period of time; while. Noun in the supplied time expression. \textbf{fierst} (n.). period; interval. Masculine noun in the source time expression.\par
% glossary-block:B00963:end
% glossary-block:B00964:begin
\needspace{3\baselineskip}
\textbf{in the beginning} — The supplied whole phrase on anginne means in the beginning.\par
% glossary-block:B00964:end
% glossary-block:B00965:begin
\needspace{3\baselineskip}
\textbf{in the heat} — \textbf{on} (prep.). in. Circumstance of the day's heat in the supplied source.\par
% glossary-block:B00965:end
% glossary-block:B00966:begin
\needspace{3\baselineskip}
\textbf{in the midst of} — \textbf{tōmiddes} (prep.). With dative landmark. tōmiddes + dative landmark; English of does not require genitive here.\par
% glossary-block:B00966:end
% glossary-block:B00967:begin
\needspace{3\baselineskip}
\textbf{in the morning} — The supplied whole phrase on morgen means in the morning; morgen is accusative singular here. On morgen is the supplied expression in the morning; morgen is accusative singular in this expression.\par
% glossary-block:B00967:end
% glossary-block:B00968:begin
\needspace{3\baselineskip}
\textbf{in the name of} — On + dative naman with a genitive possessor means in someone’s name. The whole name phrase, not on alone, supplies this meaning.\par
% glossary-block:B00968:end
% glossary-block:B00969:begin
\needspace{3\baselineskip}
\textbf{in the same way as} — \textbf{swā} (adverb; conjunction). as; in the same way as.\par
% glossary-block:B00969:end
% glossary-block:B00970:begin
\needspace{3\baselineskip}
\textbf{in the south} — The supplied whole phrase be sūðan means in the south. Be alone does not mean south. \textbf{sūðan} (adv.). in the south; from the south. Retain the supplied spatial expression.\par
% glossary-block:B00970:end
% glossary-block:B00971:begin
\needspace{3\baselineskip}
\textbf{into} — \textbf{in-tō} (prep.). into; destination expressed with dative after intō. \textbf{on} (prep.). into; in two. Supplied division expression. into; onto. Accusative goal of placement in the taught frame. upon; into. Supplied chancing-upon idiom. Destination of sale in the supplied expression. \textbf{tō} (prep.). to; into. Result in supplied tō lore, tō sorge or tō nāhte. into; to. Resulting language/literary form.\par
% glossary-block:B00971:end
% glossary-block:B00972:begin
\needspace{3\baselineskip}
\textbf{in two} — \textbf{on} (prep.). into; in two. Supplied division expression. \textbf{twēgen} (cardinal numeral). two; in two. Two parts in the supplied division expression.\par
% glossary-block:B00972:end
% glossary-block:B00973:begin
\needspace{3\baselineskip}
\textbf{in vain} — \textbf{hōlunga} (adv.). hōlunga; not at all in vain uses nālæs hōlunga with the intended negative scope.\par
% glossary-block:B00973:end
% glossary-block:B00974:begin
\needspace{3\baselineskip}
\textbf{invisible} — \textbf{ungesewenlic} (adj.).\par
% glossary-block:B00974:end
% glossary-block:B00975:begin
\needspace{3\baselineskip}
\textbf{inward} — \textbf{in} (prep.). in; inward. Adverbial direction.\par
% glossary-block:B00975:end
% glossary-block:B00976:begin
\needspace{3\baselineskip}
\textbf{in worldly life} — \textbf{on} (prep.). Supplied temporal/situational expression.\par
% glossary-block:B00976:end
% glossary-block:B00977:begin
\needspace{3\baselineskip}
\textbf{iron} — \textbf{īsern} (n.). Metal noun. \textbf{īsern} (adj.). iron; made of iron. Adjective of material.\par
% glossary-block:B00977:end
% glossary-block:B00978:begin
\needspace{3\baselineskip}
\textbf{is} — \textbf{wesan} (v.). be. Copular present/past identification or state; other auxiliary constructions are separate senses. Copular eom/wæs and bēo indicative series; later passive and compound uses require their own construction.\par
% glossary-block:B00978:end
% glossary-block:B00979:begin
\needspace{3\baselineskip}
\textbf{island} — \textbf{īgland} (n.).\par
% glossary-block:B00979:end
% glossary-block:B00980:begin
\needspace{3\baselineskip}
\textbf{Israel} — \textbf{Israel} (proper name).\par
% glossary-block:B00980:end
% glossary-block:B00981:begin
\needspace{3\baselineskip}
\textbf{Israelites} — \textbf{Israela} (proper name). Supplied genitive plural name form.\par
% glossary-block:B00981:end
% glossary-block:B00982:begin
\needspace{3\baselineskip}
\textbf{is to come} — \textbf{tō} (prep.). Whole supplied is tō cumanne; bounded be-to use, distinct from purpose after a verb of coming.\par
% glossary-block:B00982:end
% glossary-block:B00983:begin
\needspace{3\baselineskip}
\textbf{it} — Hit is neuter singular nominative/accusative; dative is him and genitive his. An inanimate referent of masculine or feminine grammatical gender takes the corresponding gendered pronoun instead. Grammatical gender and antecedent determine reference; it may be represented by masculine/feminine forms for a gendered inanimate noun.\par
% glossary-block:B00983:end
% glossary-block:B00984:begin
\needspace{3\baselineskip}
\textbf{it (anticipatory or impersonal)} — Hit occurs in some taught impersonal/anticipatory frames; other supplied passives expressly omit it. Follow that whole frame.\par
% glossary-block:B00984:end
% glossary-block:B00985:begin
\needspace{3\baselineskip}
\textbf{its} — \textbf{hē} (personal pron.). his; her; its; their. Invariant genitive possessors, unlike agreeing possessive adjectives. Invariant personal genitive possessor; distinguish the same surface used as a governed personal case.\par
% glossary-block:B00985:end
% glossary-block:B00986:begin
\needspace{3\baselineskip}
\textbf{itself} — Emphatic self agrees with the person or thing emphasized: supplied dēaðes selfes means death itself. This is emphasis on an inanimate thing, not a reflexive object.\par
% glossary-block:B00986:end
% glossary-block:B00987:begin
\needspace{3\baselineskip}
\textbf{ivory} — \textbf{bān} (n.). bone; ivory. Ivory is elpes bān, literally elephant's bone. \textbf{elpenbǣnen} (adj.). ivory; made of ivory. Adjective of material.\par
% glossary-block:B00987:end
\pdfbookmark[1]{J}{outline-B00988}
% glossary-block:B00988:begin
\needspace{5\baselineskip}
\subsection*{J}
% glossary-block:B00988:end
% glossary-block:B00989:begin
\textbf{Jacob} — \textbf{Iacob} (proper name).\par
% glossary-block:B00989:end
% glossary-block:B00990:begin
\needspace{3\baselineskip}
\textbf{Jerusalem} — \textbf{Hierusalem} (proper name). English spelling Jerusalem points to the retained Old English name Hierusalem.\par
% glossary-block:B00990:end
% glossary-block:B00991:begin
\needspace{3\baselineskip}
\textbf{Jesus} — \textbf{Hǣlend} (n.). Saviour; Jesus. Proper title.\par
% glossary-block:B00991:end
% glossary-block:B00992:begin
\needspace{3\baselineskip}
\textbf{jet} — \textbf{gagates} (n.). jet (Latin loan-name). The stone; Latin loan-name in the source.\par
% glossary-block:B00992:end
% glossary-block:B00993:begin
\needspace{3\baselineskip}
\textbf{Jews} — \textbf{Iudea} (proper name). Supplied genitive plural name form.\par
% glossary-block:B00993:end
% glossary-block:B00994:begin
\needspace{3\baselineskip}
\textbf{John} — \textbf{Johannes} (proper name).\par
% glossary-block:B00994:end
% glossary-block:B00995:begin
\needspace{3\baselineskip}
\textbf{join} — \textbf{gemęngan} (v.). join; mingle. Supplied source form.\par
% glossary-block:B00995:end
% glossary-block:B00996:begin
\needspace{3\baselineskip}
\textbf{Joseph} — \textbf{Iosep} (proper name).\par
% glossary-block:B00996:end
% glossary-block:B00997:begin
\needspace{3\baselineskip}
\textbf{journey} — \textbf{sīðfæt} (n.). journey; passage. \textbf{faran} (v.). Verb faran; distinct from noun journey sīðfæt.\par
% glossary-block:B00997:end
% glossary-block:B00998:begin
\needspace{3\baselineskip}
\textbf{joy} — \textbf{drēam} (n.). joy; bliss.\par
% glossary-block:B00998:end
% glossary-block:B00999:begin
\needspace{3\baselineskip}
\textbf{joyfully} — \textbf{blīðe} (adv.). joyfully; gladly.\par
% glossary-block:B00999:end
% glossary-block:B01000:begin
\needspace{3\baselineskip}
\textbf{Judah} — \textbf{Iudas} (proper name).\par
% glossary-block:B01000:end
% glossary-block:B01001:begin
\needspace{3\baselineskip}
\textbf{judge} — \textbf{dēma} (n.).  \textbf{dēman} (v.).\par
% glossary-block:B01001:end
% glossary-block:B01002:begin
\needspace{3\baselineskip}
\textbf{judgment} — \textbf{dōm} (n.). dōmes dæg means Day of Judgment. A judgment made in judging, distinct from doom/fate.\par
% glossary-block:B01002:end
% glossary-block:B01003:begin
\needspace{3\baselineskip}
\textbf{just} — \textbf{eall} (adj.; substantive; adv.). just; exactly. Degree emphasis in eall swā. \textbf{efne} (interj.). just; precisely; even. \textbf{riht} (adj.). right; just.\par
% glossary-block:B01003:end
% glossary-block:B01004:begin
\needspace{3\baselineskip}
\textbf{just as} — \textbf{swā-swā} (conj.).\par
% glossary-block:B01004:end
% glossary-block:B01005:begin
\needspace{3\baselineskip}
\textbf{just as … as} — Swā + adjective/adverb + swā frames equal degree, as … as. The first swā marks degree; the second introduces the standard. Comparative the … the uses a separate construction.\par
% glossary-block:B01005:end
\pdfbookmark[1]{K}{outline-B01006}
% glossary-block:B01006:begin
\needspace{5\baselineskip}
\subsection*{K}
% glossary-block:B01006:end
% glossary-block:B01007:begin
\textbf{keep} — \textbf{healdan} (v.). keep; observe. Accusative saying or command observed.\par
% glossary-block:B01007:end
% glossary-block:B01008:begin
\needspace{3\baselineskip}
\textbf{keep silent} — \textbf{swigian} (v.). be silent. Supplied source sense.\par
% glossary-block:B01008:end
% glossary-block:B01009:begin
\needspace{3\baselineskip}
\textbf{keep silent about} — \textbf{forswīgian} (v.). keep silent about; conceal.\par
% glossary-block:B01009:end
% glossary-block:B01010:begin
\needspace{3\baselineskip}
\textbf{kill} — \textbf{cwellan} (v.).  \textbf{ofslēan} (v.). slay; kill. \textbf{slēan} (v.). slay; kill. Accusative person killed. Accusative person; the bounded negative command uses selected short-vowel imperative sleah.\par
% glossary-block:B01010:end
% glossary-block:B01011:begin
\needspace{3\baselineskip}
\textbf{kind} — \textbf{cynn} (n.).  \textbf{gōd} (adj.). kind; benevolent. Gōd means kind or benevolent in the taught person-directed frame. \textbf{hīw} (n.). form; kind. Supplied dative neuter plural.\par
% glossary-block:B01011:end
% glossary-block:B01012:begin
\needspace{3\baselineskip}
\textbf{kindness} — \textbf{ār-fæstnes} (n.).  \textbf{gōd} (n.). good; benefit; kindness. Nominal good in the repayment contrast.\par
% glossary-block:B01012:end
% glossary-block:B01013:begin
\needspace{3\baselineskip}
\textbf{kinds of fish} — \textbf{fisccynn} (n.). Supplied accusative neuter plural.\par
% glossary-block:B01013:end
% glossary-block:B01014:begin
\needspace{3\baselineskip}
\textbf{kind toward} — \textbf{ofer} (prep.). toward. In gōd ofer + accusative people.\par
% glossary-block:B01014:end
% glossary-block:B01015:begin
\needspace{3\baselineskip}
\textbf{king} — \textbf{cyning} (n.).  \textbf{ealdor} (n.). chief; king. This noun is distinct from the homonymous noun meaning life.\par
% glossary-block:B01015:end
% glossary-block:B01016:begin
\needspace{3\baselineskip}
\textbf{kingdom} — \textbf{rīce} (n.).\par
% glossary-block:B01016:end
% glossary-block:B01017:begin
\needspace{3\baselineskip}
\textbf{kiss} — \textbf{coss} (n.).  \textbf{cyssan} (v.).\par
% glossary-block:B01017:end
% glossary-block:B01018:begin
\needspace{3\baselineskip}
\textbf{knee} — \textbf{cnēo} (n.). Supplied accusative plural knees.\par
% glossary-block:B01018:end
% glossary-block:B01019:begin
\needspace{3\baselineskip}
\textbf{knock} — \textbf{cnucian} (v.).\par
% glossary-block:B01019:end
% glossary-block:B01020:begin
\needspace{3\baselineskip}
\textbf{know} — \textbf{cnāwan} (v.).  \textbf{cunnan} (v.). know; be acquainted with. Accusative person known; distinct from learned ability. \textbf{witan} (v.). A stated fact or the answer to a question as complement. Accusative person in this taught frame; cunnan has a separate acquaintance sense. \textbf{cunnan} (v.), \textbf{witan} (v.). Cunnan + accusative person expresses acquaintance. Witan can take an accusative person in its taught personal sense, or a fact/question as content. Learned ability is the separate cunnan + plain infinitive construction.\par
% glossary-block:B01020:end
% glossary-block:B01021:begin
\needspace{3\baselineskip}
\textbf{know a person} — \textbf{witan} (v.). know. Accusative person in this taught frame; cunnan has a separate acquaintance sense.\par
% glossary-block:B01021:end
% glossary-block:B01022:begin
\needspace{3\baselineskip}
\textbf{know by recognition} — \textbf{oncnāwan} (v.). recognize; know by recognition.\par
% glossary-block:B01022:end
% glossary-block:B01023:begin
\needspace{3\baselineskip}
\textbf{know how} — \textbf{cunnan} (v.). know how; be able. Learned skill; plain infinitive complement.\par
% glossary-block:B01023:end
% glossary-block:B01024:begin
\needspace{3\baselineskip}
\textbf{knowledge} — \textbf{andgiet} (n.). andgiet denotes the knowledge possessed and desired in this exercise.\par
% glossary-block:B01024:end
% glossary-block:B01025:begin
\needspace{3\baselineskip}
\textbf{know that} — \textbf{witan} (v.). know. A stated fact or the answer to a question as complement.\par
% glossary-block:B01025:end
% glossary-block:B01026:begin
\needspace{3\baselineskip}
\textbf{know whether} — \textbf{witan} (v.). know. A stated fact or the answer to a question as complement.\par
% glossary-block:B01026:end
\pdfbookmark[1]{L}{outline-B01027}
% glossary-block:B01027:begin
\needspace{5\baselineskip}
\subsection*{L}
% glossary-block:B01027:end
% glossary-block:B01028:begin
\textbf{labor} — \textbf{gedeorf} (n.). labor; toil. \textbf{gewinn} (n.). labour; work. Noun, distinct from labour as a verb. Noun in the expression from his labour. \textbf{winnan} (v.). Verb; the subject labours with the hands.\par
% glossary-block:B01028:end
% glossary-block:B01029:begin
\needspace{3\baselineskip}
\textbf{labour} — \textbf{gewinn} (n.). labour; work. Noun, distinct from labour as a verb. Noun in the expression from his labour. \textbf{winnan} (v.). Verb; the subject labours with the hands.\par
% glossary-block:B01029:end
% glossary-block:B01030:begin
\needspace{3\baselineskip}
\textbf{lad} — \textbf{cnapa} (n.). Masculine weak noun; the task's chosen word for lad.\par
% glossary-block:B01030:end
% glossary-block:B01031:begin
\needspace{3\baselineskip}
\textbf{lady} — \textbf{hlǣfdige} (n.).\par
% glossary-block:B01031:end
% glossary-block:B01032:begin
\needspace{3\baselineskip}
\textbf{lame} — \textbf{healt} (adj.). lame; halt. Adjective describing impaired walking; not the command halt/stop.\par
% glossary-block:B01032:end
% glossary-block:B01033:begin
\needspace{3\baselineskip}
\textbf{land} — \textbf{land} (n.). land; country.\par
% glossary-block:B01033:end
% glossary-block:B01034:begin
\needspace{3\baselineskip}
\textbf{large} — \textbf{micel} (adj.). great; large.\par
% glossary-block:B01034:end
% glossary-block:B01035:begin
\needspace{3\baselineskip}
\textbf{larger} — \textbf{māra} (adj.). greater; larger. Comparative adjective.\par
% glossary-block:B01035:end
% glossary-block:B01036:begin
\needspace{3\baselineskip}
\textbf{largest} — \textbf{mǣst} (adj.). greatest; largest. Superlative adjective.\par
% glossary-block:B01036:end
% glossary-block:B01037:begin
\needspace{3\baselineskip}
\textbf{last} — \textbf{ðurhwunian} (v.). remain; endure; last.\par
% glossary-block:B01037:end
% glossary-block:B01038:begin
\needspace{3\baselineskip}
\textbf{Latin} — \textbf{Bōclæden} (n.). book-Latin; Latin. \textbf{Læden} (n.). The language.\par
% glossary-block:B01038:end
% glossary-block:B01039:begin
\needspace{3\baselineskip}
\textbf{law} — \textbf{ǣ} (n.).\par
% glossary-block:B01039:end
% glossary-block:B01040:begin
\needspace{3\baselineskip}
\textbf{lay} — \textbf{āhyldan} (v.). lay; incline. \textbf{lęcgean} (v.). lay; place. Supplied source form. lay; place.\par
% glossary-block:B01040:end
% glossary-block:B01041:begin
\needspace{3\baselineskip}
\textbf{layman} — \textbf{ceorl} (n.).\par
% glossary-block:B01041:end
% glossary-block:B01042:begin
\needspace{3\baselineskip}
\textbf{lay out} — \textbf{lęcgean} (v.). lęcgean + accusative thing + forð.\par
% glossary-block:B01042:end
% glossary-block:B01043:begin
\needspace{3\baselineskip}
\textbf{lay waste} — \textbf{forniman} (v.). destroy; lay waste.\par
% glossary-block:B01043:end
% glossary-block:B01044:begin
\needspace{3\baselineskip}
\textbf{Lazarus} — \textbf{Lazarus} (proper name). The parable's poor man; supplied plural denotes people like him.\par
% glossary-block:B01044:end
% glossary-block:B01045:begin
\needspace{3\baselineskip}
\textbf{Lazaruses} — \textbf{Lazarus} (proper name). Lazarus. The parable's poor man; supplied plural denotes people like him.\par
% glossary-block:B01045:end
% glossary-block:B01046:begin
\needspace{3\baselineskip}
\textbf{lead} — \textbf{lǣdan} (v.). lead; bring; take. A person is the object in the assigned frame. \textbf{lēad} (n.). Metal noun, not the verb lead.\par
% glossary-block:B01046:end
% glossary-block:B01047:begin
\needspace{3\baselineskip}
\textbf{leader} — \textbf{ealdormann} (n.). prince; leader; noble.\par
% glossary-block:B01047:end
% glossary-block:B01048:begin
\needspace{3\baselineskip}
\textbf{learn} — \textbf{geleornian} (v.). learn; study. learn; master. Learn a skill or performance in the source. \textbf{gewitan} (v.). find out; learn. Knowledge verb, distinct from gewītan depart. \textbf{leornian} (v.). Accusative thing learned; æt + dative person/source.\par
% glossary-block:B01048:end
% glossary-block:B01049:begin
\needspace{3\baselineskip}
\textbf{learn from} — \textbf{æt} (prep.). from. æt + dative person with learn. leornian æt + dative person in the assigned learning frame.\par
% glossary-block:B01049:end
% glossary-block:B01050:begin
\needspace{3\baselineskip}
\textbf{learning} — \textbf{lār} (n.). learning; study. \textbf{leornung} (n.). study; learning.\par
% glossary-block:B01050:end
% glossary-block:B01051:begin
\needspace{3\baselineskip}
\textbf{learn of} — \textbf{æt} (prep.). from. æt + dative person with learn. leornian æt + dative person in the assigned learning frame.\par
% glossary-block:B01051:end
% glossary-block:B01052:begin
\needspace{3\baselineskip}
\textbf{least} — \textbf{lǣst} (adj.). In rank or importance.\par
% glossary-block:B01052:end
% glossary-block:B01053:begin
\needspace{3\baselineskip}
\textbf{leave} — \textbf{forlǣtan} (v.). leave to. Thing left and recipient retain their taught roles. Accusative thing/place left. Distinguish leave to someone and take leave of.\par
% glossary-block:B01053:end
% glossary-block:B01054:begin
\needspace{3\baselineskip}
\textbf{leave to} — \textbf{forlǣtan} (v.). Thing left and recipient retain their taught roles.\par
% glossary-block:B01054:end
% glossary-block:B01055:begin
\needspace{3\baselineskip}
\textbf{left} — \textbf{winstre} (adj.). Supplied side phrase.\par
% glossary-block:B01055:end
% glossary-block:B01056:begin
\needspace{3\baselineskip}
\textbf{left-hand side} — \textbf{winstre} (adj.). left. Supplied side phrase.\par
% glossary-block:B01056:end
% glossary-block:B01057:begin
\needspace{3\baselineskip}
\textbf{leisure} — \textbf{ǣmetta} (n.). leisure; available time.\par
% glossary-block:B01057:end
% glossary-block:B01058:begin
\needspace{3\baselineskip}
\textbf{less} — \textbf{lǣs} (degree word). less; a smaller amount.\par
% glossary-block:B01058:end
% glossary-block:B01059:begin
\needspace{3\baselineskip}
\textbf{lest} — \textbf{ðȳ-lǣs-ðe} (conj.). lest; so that not. ðȳ-lǣs-ðe with the assigned subjunctive negative-purpose clause.\par
% glossary-block:B01059:end
% glossary-block:B01060:begin
\needspace{3\baselineskip}
\textbf{let} — Lǣtan is lexical let/allow. Wishes introduced by English let use the taught independent subjunctive; they do not require an inserted lǣtan. An independent wish/hortatory uses a subjunctive; let does not automatically require lǣtan. Literal permission to go uses lǣtan + accusative actor + infinitive.\par
% glossary-block:B01060:end
% glossary-block:B01061:begin
\needspace{3\baselineskip}
\textbf{let go} — \textbf{forlǣtan} (v.). let go; dismiss.\par
% glossary-block:B01061:end
% glossary-block:B01062:begin
\needspace{3\baselineskip}
\textbf{letter} — \textbf{gewrit} (n.). letter; written document.\par
% glossary-block:B01062:end
% glossary-block:B01063:begin
\needspace{3\baselineskip}
\textbf{levin} — \textbf{līeg} (n.). thunderbolt; lightning. Source levin sense.\par
% glossary-block:B01063:end
% glossary-block:B01064:begin
\needspace{3\baselineskip}
\textbf{lick} — \textbf{liccian} (v.).\par
% glossary-block:B01064:end
% glossary-block:B01065:begin
\needspace{3\baselineskip}
\textbf{lie} — \textbf{licgan} (v.). lie; be at rest. lie; be lying. Supplied participle describing the possessed persons, not a habban perfect.\par
% glossary-block:B01065:end
% glossary-block:B01066:begin
\needspace{3\baselineskip}
\textbf{life} — \textbf{līf} (n.). An individual's lived life, including cognate accusative with libban.\par
% glossary-block:B01066:end
% glossary-block:B01067:begin
\needspace{3\baselineskip}
\textbf{lift} — \textbf{āhębban} (v.). raise; lift. \textbf{hębban} (v.). raise; lift.\par
% glossary-block:B01067:end
% glossary-block:B01068:begin
\needspace{3\baselineskip}
\textbf{lift up} — \textbf{ūp-ā-hębban} (v.). Supplied source form. Supplied first-plural preterite. Accusative people raised; the prefixed verb supplies up.\par
% glossary-block:B01068:end
% glossary-block:B01069:begin
\needspace{3\baselineskip}
\textbf{light} — \textbf{lēoht} (n.). light; luminary. Individual countable light, distinct from mass light. Mass light and an individual countable light/luminary are separate uses of lēoht. Use the count sense for a light or several lights.\par
% glossary-block:B01069:end
% glossary-block:B01070:begin
\needspace{3\baselineskip}
\textbf{lightning} — \textbf{līeg} (n.). thunderbolt; lightning. Source levin sense.\par
% glossary-block:B01070:end
% glossary-block:B01071:begin
\needspace{3\baselineskip}
\textbf{like} — \textbf{gelīc} (adj.). equal; like. Dative standard of comparison. \textbf{swā} (adverb; conjunction). as; like. Predicative comparison. Swā can introduce a predicative comparison, as/like; swā-swā means just as in its supplied comparison. Swilce means as if in the supplied apparent-rejuvenation comparison. Equal degree uses the separate paired swā … swā construction.\par
% glossary-block:B01071:end
% glossary-block:B01072:begin
\needspace{3\baselineskip}
\textbf{likeness} — \textbf{gelīcnes} (n.).\par
% glossary-block:B01072:end
% glossary-block:B01073:begin
\needspace{3\baselineskip}
\textbf{likewise} — \textbf{gelīce} (adv.). likewise; similarly. \textbf{swilce} (adv.). Adverb, distinct from inflected swilc such.\par
% glossary-block:B01073:end
% glossary-block:B01074:begin
\needspace{3\baselineskip}
\textbf{linen cloth} — \textbf{scīete} (n.).\par
% glossary-block:B01074:end
% glossary-block:B01075:begin
\needspace{3\baselineskip}
\textbf{lion} — \textbf{lēo} (n.). lion (Latin loan-form).\par
% glossary-block:B01075:end
% glossary-block:B01076:begin
\needspace{3\baselineskip}
\textbf{listen} — \textbf{gehīeran} (v.). hear; listen.\par
% glossary-block:B01076:end
% glossary-block:B01077:begin
\needspace{3\baselineskip}
\textbf{listen to} — \textbf{hlystan} (v.). listen to; heed. Genitive person or thing listened to; hear is the contextual prompt wording. hlystan governs genitive of the speech heard; do not copy hīeran's ordinary accusative frame.\par
% glossary-block:B01077:end
% glossary-block:B01078:begin
\needspace{3\baselineskip}
\textbf{little} — \textbf{lȳtel} (adj.). little; small.\par
% glossary-block:B01078:end
% glossary-block:B01079:begin
\needspace{3\baselineskip}
\textbf{little by little} — \textbf{styccemǣlum} (adv.). Single adverb styccemǣlum.\par
% glossary-block:B01079:end
% glossary-block:B01080:begin
\needspace{3\baselineskip}
\textbf{live} — \textbf{libban} (v.). The bracketed liveth yet is rendered by finite libban + gīet; an added copula is not required.\par
% glossary-block:B01080:end
% glossary-block:B01081:begin
\needspace{3\baselineskip}
\textbf{lives} — \textbf{līf} (n.). life. An individual's lived life, including cognate accusative with libban.\par
% glossary-block:B01081:end
% glossary-block:B01082:begin
\needspace{3\baselineskip}
\textbf{living} — \textbf{cucu} (adj.). alive; living. Supplied accusative neuter plural adjective. \textbf{cwic} (adj.). alive; living. \textbf{libban} (v.). Supplied weak participial adjective.\par
% glossary-block:B01082:end
% glossary-block:B01083:begin
\needspace{3\baselineskip}
\textbf{living people} — \textbf{cwic} (adj.). Substantival adjective.\par
% glossary-block:B01083:end
% glossary-block:B01084:begin
\needspace{3\baselineskip}
\textbf{lodge} — \textbf{wician} (v.). sojourn; lodge. Intransitive; mid + dative person with whom one lodges.\par
% glossary-block:B01084:end
% glossary-block:B01085:begin
\needspace{3\baselineskip}
\textbf{long} — \textbf{lang} (adj.). Adjective agreeing with the time interval fæc. \textbf{lange} (adv.). In supplied swā lange swā; its biblical inasmuch-as force is explained locally.\par
% glossary-block:B01085:end
% glossary-block:B01086:begin
\needspace{3\baselineskip}
\textbf{longer} — \textbf{lęng} (adverb or particle). lęng as the comparative adverb; clause negation supplies no longer.\par
% glossary-block:B01086:end
% glossary-block:B01087:begin
\needspace{3\baselineskip}
\textbf{look} — \textbf{besēon} (v.). look; turn one's gaze. Supplied source form and frame.\par
% glossary-block:B01087:end
% glossary-block:B01088:begin
\needspace{3\baselineskip}
\textbf{look at} — \textbf{behealdan} (v.). behold; look at. Accusative object seen. observe. Supplied source form observed. behealdan takes accusative objects; distinguish interjection efne behold.\par
% glossary-block:B01088:end
% glossary-block:B01089:begin
\needspace{3\baselineskip}
\textbf{look for} — \textbf{sēcan} (v.). sēcan with accusative sought thing, no extra word for English for/out.\par
% glossary-block:B01089:end
% glossary-block:B01090:begin
\needspace{3\baselineskip}
\textbf{look toward} — \textbf{besēon} (v.). Target of the look follows the supplied phrase.\par
% glossary-block:B01090:end
% glossary-block:B01091:begin
\needspace{3\baselineskip}
\textbf{lord} — \textbf{dryhten} (n.). lord; the Lord. \textbf{hlāford} (n.). lord; master.\par
% glossary-block:B01091:end
% glossary-block:B01092:begin
\needspace{3\baselineskip}
\textbf{lose} — \textbf{forlēosan} (v.). Accusative thing lost.\par
% glossary-block:B01092:end
% glossary-block:B01093:begin
\needspace{3\baselineskip}
\textbf{loss} — \textbf{lyre} (n.).\par
% glossary-block:B01093:end
% glossary-block:B01094:begin
\needspace{3\baselineskip}
\textbf{loud} — \textbf{hlūd} (adj.).  \textbf{micel} (adj.). great; loud. Loud is the contextual force with a voice.\par
% glossary-block:B01094:end
% glossary-block:B01095:begin
\needspace{3\baselineskip}
\textbf{love} — \textbf{lufe} (n.).  \textbf{lufian} (v.).  \textbf{lufu} (n.).\par
% glossary-block:B01095:end
% glossary-block:B01096:begin
\needspace{3\baselineskip}
\textbf{lover} — \textbf{lufiend} (n.).\par
% glossary-block:B01096:end
% glossary-block:B01097:begin
\needspace{3\baselineskip}
\textbf{lowly} — \textbf{ēaðmōd} (adj.). lowly; humble.\par
% glossary-block:B01097:end
% glossary-block:B01098:begin
\needspace{3\baselineskip}
\textbf{luminary} — \textbf{lēoht} (n.). light; luminary. Individual countable light, distinct from mass light. Mass light and an individual countable light/luminary are separate uses of lēoht. Use the count sense for a light or several lights.\par
% glossary-block:B01098:end
% glossary-block:B01099:begin
\needspace{3\baselineskip}
\textbf{lust} — \textbf{gālnes} (n.). Assigned source sense.\par
% glossary-block:B01099:end
\pdfbookmark[1]{M}{outline-B01100}
% glossary-block:B01100:begin
\needspace{5\baselineskip}
\subsection*{M}
% glossary-block:B01100:end
% glossary-block:B01101:begin
\textbf{made of iron} — \textbf{īsern} (adj.). iron; made of iron. Adjective of material.\par
% glossary-block:B01101:end
% glossary-block:B01102:begin
\needspace{3\baselineskip}
\textbf{made of ivory} — \textbf{elpenbǣnen} (adj.). ivory; made of ivory. Adjective of material.\par
% glossary-block:B01102:end
% glossary-block:B01103:begin
\needspace{3\baselineskip}
\textbf{made of silver} — \textbf{seolfren} (adj.). silver; made of silver. Material adjective, distinct from the noun seolfor.\par
% glossary-block:B01103:end
% glossary-block:B01104:begin
\needspace{3\baselineskip}
\textbf{made of stone} — \textbf{stānen} (adj.). stone; made of stone. Material adjective.\par
% glossary-block:B01104:end
% glossary-block:B01105:begin
\needspace{3\baselineskip}
\textbf{magnify} — \textbf{mǣrsian} (v.). make great; magnify; glorify.\par
% glossary-block:B01105:end
% glossary-block:B01106:begin
\needspace{3\baselineskip}
\textbf{maiden} — \textbf{fǣmne} (n.). maiden; young woman. \textbf{mǣden} (n.). girl; maiden.\par
% glossary-block:B01106:end
% glossary-block:B01107:begin
\needspace{3\baselineskip}
\textbf{maidenhood} — \textbf{mægðhād} (n.). maidenhood; virginity.\par
% glossary-block:B01107:end
% glossary-block:B01108:begin
\needspace{3\baselineskip}
\textbf{maintain} — \textbf{healdan} (v.). maintain; sustain.\par
% glossary-block:B01108:end
% glossary-block:B01109:begin
\needspace{3\baselineskip}
\textbf{majesty} — \textbf{mægenðrymm} (n.). majesty; glory. \textbf{mægenðrymnes} (n.). glory; majesty.\par
% glossary-block:B01109:end
% glossary-block:B01110:begin
\needspace{3\baselineskip}
\textbf{make} — \textbf{dōn} (v.). make; render. Accusative object with agreeing object predicate. \textbf{gedōn} (v.). do; make. Supplied passive fulfilment formula. make; arrange. The present form/arrangement of the translation in the source. make; cause to be. \textbf{gewyrcean} (v.). make; build. make; render. Accusative object with an agreeing predicate noun/adjective. \textbf{wyrcean} (v.). make; build. Supplied unprefixed form; preserve explicit task controls. \textbf{wyrcean} (v.), \textbf{gewyrcean} (v.). Wyrcean and gewyrcean mean make/build in these senses. Preserve the prefixed or unprefixed verb selected by the exercise.\par
% glossary-block:B01110:end
% glossary-block:B01111:begin
\needspace{3\baselineskip}
\textbf{make a decree} — \textbf{sęttan} (v.). appoint; establish. With accusative entity/decree in the assigned frame.\par
% glossary-block:B01111:end
% glossary-block:B01112:begin
\needspace{3\baselineskip}
\textbf{make a signal} — The supplied expression slōg tācen means made a signal; tācen is its object. Slēan alone has the separate strike sense.\par
% glossary-block:B01112:end
% glossary-block:B01113:begin
\needspace{3\baselineskip}
\textbf{make clear} — \textbf{gesweotolian} (v.). reveal; make clear. Supplied source form.\par
% glossary-block:B01113:end
% glossary-block:B01114:begin
\needspace{3\baselineskip}
\textbf{make great} — \textbf{mǣrsian} (v.). make great; magnify; glorify.\par
% glossary-block:B01114:end
% glossary-block:B01115:begin
\needspace{3\baselineskip}
\textbf{make known} — \textbf{cȳðan} (v.). announce; make known. Dative person informed; accusative thing announced or a ðæt content clause.\par
% glossary-block:B01115:end
% glossary-block:B01116:begin
\needspace{3\baselineskip}
\textbf{make manifest} — \textbf{gesweotolian} (v.). make manifest; reveal.\par
% glossary-block:B01116:end
% glossary-block:B01117:begin
\needspace{3\baselineskip}
\textbf{Maker} — \textbf{wyrhta} (n.). Maker; Creator. Divine title, distinct from ordinary workman.\par
% glossary-block:B01117:end
% glossary-block:B01118:begin
\needspace{3\baselineskip}
\textbf{make steadfast} — \textbf{gestaðelian} (v.). establish; make steadfast.\par
% glossary-block:B01118:end
% glossary-block:B01119:begin
\needspace{3\baselineskip}
\textbf{man} — \textbf{mann} (n.; impersonal pronoun). man; human being. \textbf{wer} (n.). Supplied genitive singular.\par
% glossary-block:B01119:end
% glossary-block:B01120:begin
\needspace{3\baselineskip}
\textbf{manifold} — \textbf{manigfeald} (adj.). manifold; many.\par
% glossary-block:B01120:end
% glossary-block:B01121:begin
\needspace{3\baselineskip}
\textbf{many} — \textbf{fela} (adj.). many; a large number of. Genitive of the things counted in the taught frame. \textbf{manig} (adj.).  \textbf{manigfeald} (adj.). manifold; many.\par
% glossary-block:B01121:end
% glossary-block:B01122:begin
\needspace{3\baselineskip}
\textbf{marine} — \textbf{sǣlic} (adj.). marine; of the sea.\par
% glossary-block:B01122:end
% glossary-block:B01123:begin
\needspace{3\baselineskip}
\textbf{Mary} — \textbf{Maria} (proper name).\par
% glossary-block:B01123:end
% glossary-block:B01124:begin
\needspace{3\baselineskip}
\textbf{master} — \textbf{geleornian} (v.). learn; master. Learn a skill or performance in the source. \textbf{hlāford} (n.). lord; master. \textbf{lārēow} (n.). master; teacher.\par
% glossary-block:B01124:end
% glossary-block:B01125:begin
\needspace{3\baselineskip}
\textbf{matter} — \textbf{ðing} (n.). thing; fact; matter. Matter known/discussed, distinct from a generic physical object.\par
% glossary-block:B01125:end
% glossary-block:B01126:begin
\needspace{3\baselineskip}
\textbf{may} — Mōtan expresses permission, may/be allowed to, with a plain infinitive. A wish beginning with English may instead uses the taught independent subjunctive.\par
% glossary-block:B01126:end
% glossary-block:B01127:begin
\needspace{3\baselineskip}
\textbf{may (permission)} — \textbf{mōtan} (v.). mōtan + plain infinitive; permission/prohibition differs from ability and independent wishes.\par
% glossary-block:B01127:end
% glossary-block:B01128:begin
\needspace{3\baselineskip}
\textbf{may (wish)} — An independent wish/hortatory uses a subjunctive; let does not automatically require lǣtan. Literal permission to go uses lǣtan + accusative actor + infinitive.\par
% glossary-block:B01128:end
% glossary-block:B01129:begin
\needspace{3\baselineskip}
\textbf{me} — Mē is first-person singular dative or accusative me; choose its case from the verb or preposition. Select case from the verb/preposition and preserve dual versus plural; dual subjects use plural verb forms.\par
% glossary-block:B01129:end
% glossary-block:B01130:begin
\needspace{3\baselineskip}
\textbf{meal} — \textbf{gereord} (n.). meal; banquet.\par
% glossary-block:B01130:end
% glossary-block:B01131:begin
\needspace{3\baselineskip}
\textbf{meaning} — \textbf{andgiet} (n.). meaning; sense.\par
% glossary-block:B01131:end
% glossary-block:B01132:begin
\needspace{3\baselineskip}
\textbf{means of access} — \textbf{weg} (n.). way; means of access. Figurative access to the Father in the declaration; distinct from physical road.\par
% glossary-block:B01132:end
% glossary-block:B01133:begin
\needspace{3\baselineskip}
\textbf{measure} — \textbf{mǣð} (n.). measure; capacity.\par
% glossary-block:B01133:end
% glossary-block:B01134:begin
\needspace{3\baselineskip}
\textbf{meet} — \textbf{gemētan} (v.). find; meet. Accusative person or thing found/met. \textbf{mētan} (v.). meet; find.\par
% glossary-block:B01134:end
% glossary-block:B01135:begin
\needspace{3\baselineskip}
\textbf{mention before} — \textbf{foresęcgean} (v.). mention before; tell of beforehand.\par
% glossary-block:B01135:end
% glossary-block:B01136:begin
\needspace{3\baselineskip}
\textbf{merchandise} — \textbf{hlæst} (n.). wares; merchandise; cargo. Selected plural source use.\par
% glossary-block:B01136:end
% glossary-block:B01137:begin
\needspace{3\baselineskip}
\textbf{merchant} — \textbf{mangere} (n.).\par
% glossary-block:B01137:end
% glossary-block:B01138:begin
\needspace{3\baselineskip}
\textbf{mercies} — \textbf{mildheortnes} (n.). mercy; mercies.\par
% glossary-block:B01138:end
% glossary-block:B01139:begin
\needspace{3\baselineskip}
\textbf{merciful} — \textbf{mildheort} (adj.). merciful; compassionate.\par
% glossary-block:B01139:end
% glossary-block:B01140:begin
\needspace{3\baselineskip}
\textbf{mercy} — \textbf{mildheortnes} (n.). mercy; mercies.\par
% glossary-block:B01140:end
% glossary-block:B01141:begin
\needspace{3\baselineskip}
\textbf{merit} — \textbf{earnung} (n.). merit; good service. Source plural merits/services.\par
% glossary-block:B01141:end
% glossary-block:B01142:begin
\needspace{3\baselineskip}
\textbf{merry} — \textbf{blīðe} (adj.). merry; glad.\par
% glossary-block:B01142:end
% glossary-block:B01143:begin
\needspace{3\baselineskip}
\textbf{messengers} — \textbf{ǣrendwreca} (n.). messengers; envoys.\par
% glossary-block:B01143:end
% glossary-block:B01144:begin
\needspace{3\baselineskip}
\textbf{metal} — \textbf{węcg} (n.). metal; ingot. Supplied genitive plural dependent on veins.\par
% glossary-block:B01144:end
% glossary-block:B01145:begin
\needspace{3\baselineskip}
\textbf{Michael} — \textbf{Michael} (proper name).\par
% glossary-block:B01145:end
% glossary-block:B01146:begin
\needspace{3\baselineskip}
\textbf{midday} — The supplied whole phrase tō middes dæges means at noon/midday; middes and dæges are genitives.\par
% glossary-block:B01146:end
% glossary-block:B01147:begin
\needspace{3\baselineskip}
\textbf{might} — The supplied mugan form expresses capability, can/could/might be able to, with a plain infinitive. It is not the permission sense of English may/might. In the exact supplied purpose construction, willan accompanies the watching/seeing infinitives and is rendered might. This does not make willan a general equivalent of English might or pure future will.\par
% glossary-block:B01147:end
% glossary-block:B01148:begin
\needspace{3\baselineskip}
\textbf{mighty} — \textbf{mihtig} (adj.). mighty; powerful.\par
% glossary-block:B01148:end
% glossary-block:B01149:begin
\needspace{3\baselineskip}
\textbf{mind} — \textbf{geðanc} (n.). thought; mind. \textbf{mōd} (n.). mind; attention. heart; mind.\par
% glossary-block:B01149:end
% glossary-block:B01150:begin
\needspace{3\baselineskip}
\textbf{mine} — \textbf{mīn} (possessive adjective). my; mine. Mīn my/mine: one first-person possessor; the possessive agrees with the possessed noun.\par
% glossary-block:B01150:end
% glossary-block:B01151:begin
\needspace{3\baselineskip}
\textbf{mingle} — \textbf{gemęngan} (v.). join; mingle. Supplied source form. mix. Mix sound with song in the supplied performance expression.\par
% glossary-block:B01151:end
% glossary-block:B01152:begin
\needspace{3\baselineskip}
\textbf{miserably} — \textbf{earmlice} (adverb or particle). cruelly; miserably.\par
% glossary-block:B01152:end
% glossary-block:B01153:begin
\needspace{3\baselineskip}
\textbf{mix} — \textbf{gemęngan} (v.). Mix sound with song in the supplied performance expression.\par
% glossary-block:B01153:end
% glossary-block:B01154:begin
\needspace{3\baselineskip}
\textbf{mixture of spices} — \textbf{wyrtgemang} (n.). spice; mixture of spices.\par
% glossary-block:B01154:end
% glossary-block:B01155:begin
\needspace{3\baselineskip}
\textbf{monastery} — \textbf{mynster} (n.).\par
% glossary-block:B01155:end
% glossary-block:B01156:begin
\needspace{3\baselineskip}
\textbf{money} — \textbf{feoh} (n.).\par
% glossary-block:B01156:end
% glossary-block:B01157:begin
\needspace{3\baselineskip}
\textbf{monk} — \textbf{munuc} (n.).\par
% glossary-block:B01157:end
% glossary-block:B01158:begin
\needspace{3\baselineskip}
\textbf{month} — \textbf{mōnað} (n.). Source plural months.\par
% glossary-block:B01158:end
% glossary-block:B01159:begin
\needspace{3\baselineskip}
\textbf{moon} — \textbf{mōna} (n.).\par
% glossary-block:B01159:end
% glossary-block:B01160:begin
\needspace{3\baselineskip}
\textbf{more} — \textbf{mā} (adv.). Adverb of amount. more; again. In the negative phrase never more/never again. More in number, distinct from degree or greater value. \textbf{māra} (adj.). greater; more important. Value/importance in the life/raiment comparison. more; a greater amount. Nominal/attributive quantity, distinct from more in number. Comparative swīðor expresses rather/more in the preference construction; the positive adverb swīðe means greatly/very. Comparative swīðor expresses more in the degree of an action or feeling, not a count of objects. Mā expresses more in amount or number in its separate uses; māra expresses a greater amount. Comparative swīðor expresses more intensely or preferential rather. Use an agreeing comparative adjective for greater quality/value; these uses are not interchangeable.\par
% glossary-block:B01160:end
% glossary-block:B01161:begin
\needspace{3\baselineskip}
\textbf{more beautiful} — \textbf{fægerra} (adj.). more beautiful; fairer.\par
% glossary-block:B01161:end
% glossary-block:B01162:begin
\needspace{3\baselineskip}
\textbf{more correctly} — Supplied comparative rihtlicor means more correctly. Retain its source spelling without an inferred length mark.\par
% glossary-block:B01162:end
% glossary-block:B01163:begin
\needspace{3\baselineskip}
\textbf{more important} — \textbf{māra} (adj.). greater; more important. Value/importance in the life/raiment comparison.\par
% glossary-block:B01163:end
% glossary-block:B01164:begin
\needspace{3\baselineskip}
\textbf{more precious} — \textbf{dēorwierðra} (adj.). Comparative adjective; distinct from an adverb modifying an action.\par
% glossary-block:B01164:end
% glossary-block:B01165:begin
\needspace{3\baselineskip}
\textbf{more severely} — Supplied comparative swīðlicor means more violently/more severely; it is not the positive swīðlice.\par
% glossary-block:B01165:end
% glossary-block:B01166:begin
\needspace{3\baselineskip}
\textbf{more than} — \textbf{ofer} (prep.). above; more than. Preference in the supplied comparison.\par
% glossary-block:B01166:end
% glossary-block:B01167:begin
\needspace{3\baselineskip}
\textbf{more violently} — Supplied comparative swīðlicor means more violently/more severely; it is not the positive swīðlice.\par
% glossary-block:B01167:end
% glossary-block:B01168:begin
\needspace{3\baselineskip}
\textbf{more widely} — Wīdor is the comparative adverb more widely; wīde is the positive widely.\par
% glossary-block:B01168:end
% glossary-block:B01169:begin
\needspace{3\baselineskip}
\textbf{morning} — \textbf{morgen} (n.).\par
% glossary-block:B01169:end
% glossary-block:B01170:begin
\needspace{3\baselineskip}
\textbf{most} — Swīðost means most/very greatly. It is the recognition-only superlative reference form of swīðe.\par
% glossary-block:B01170:end
% glossary-block:B01171:begin
\needspace{3\baselineskip}
\textbf{most beautiful} — \textbf{fægerest} (adj.).\par
% glossary-block:B01171:end
% glossary-block:B01172:begin
\needspace{3\baselineskip}
\textbf{most clearly} — Supplied superlative sweotolost means most clearly; do not use the positive sweotole for this degree.\par
% glossary-block:B01172:end
% glossary-block:B01173:begin
\needspace{3\baselineskip}
\textbf{most intelligibly} — Supplied superlative andgietfullicost means most intelligibly; do not use the positive andgietfullice for this degree.\par
% glossary-block:B01173:end
% glossary-block:B01174:begin
\needspace{3\baselineskip}
\textbf{most widely} — Wīdost is the superlative adverb most widely; wīde is the positive widely.\par
% glossary-block:B01174:end
% glossary-block:B01175:begin
\needspace{3\baselineskip}
\textbf{mother} — \textbf{mōdor} (n.).\par
% glossary-block:B01175:end
% glossary-block:B01176:begin
\needspace{3\baselineskip}
\textbf{mountain} — \textbf{dūn} (n.). mountain; hill. \textbf{munt} (n.).\par
% glossary-block:B01176:end
% glossary-block:B01177:begin
\needspace{3\baselineskip}
\textbf{mouth} — \textbf{mūð} (n.).\par
% glossary-block:B01177:end
% glossary-block:B01178:begin
\needspace{3\baselineskip}
\textbf{move} — \textbf{astyrian} (v.). move; set in motion. Supplied source sense.\par
% glossary-block:B01178:end
% glossary-block:B01179:begin
\needspace{3\baselineskip}
\textbf{moving} — \textbf{styriendlic} (adj.).\par
% glossary-block:B01179:end
% glossary-block:B01180:begin
\needspace{3\baselineskip}
\textbf{much} — \textbf{micel} (adj.). much; an amount. Quantity sense. Swīðe expresses much as degree, as in love very much; swā swīðe expresses so much. This differs from micel expressing quantity.\par
% glossary-block:B01180:end
% glossary-block:B01181:begin
\needspace{3\baselineskip}
\textbf{multiply} — \textbf{gemanigfieldan} (v.). be multiplied. Assigned passive with agreeing gemanigfield/gemanigfielde; no independent active causative frame is assigned. Assigned passive use with wesan and agreeing gemanigfield/gemanigfielde; do not infer an unassigned active causative frame.\par
% glossary-block:B01181:end
% glossary-block:B01182:begin
\needspace{3\baselineskip}
\textbf{multitude} — \textbf{getæl} (n.). Singular collective getæl with the supplied dependent expression; retain source participial genitives. \textbf{męnigu} (n.). męnigu also denotes a multitude of things in the locally supplied spice expression.\par
% glossary-block:B01182:end
% glossary-block:B01183:begin
\needspace{3\baselineskip}
\textbf{music} — \textbf{swēg} (n.). music; sound.\par
% glossary-block:B01183:end
% glossary-block:B01184:begin
\needspace{3\baselineskip}
\textbf{musical skill} — \textbf{swēgcræft} (n.).\par
% glossary-block:B01184:end
% glossary-block:B01185:begin
\needspace{3\baselineskip}
\textbf{must} — Sculan expresses obligation, must/be obliged to, with a plain infinitive. Pure future reference does not automatically require sculan. \textbf{sculan} (v.). sculan + plain infinitive for obligation; English have to is not possessive habban. Supplied have something to offer in Unit 6 is a different construction.\par
% glossary-block:B01185:end
% glossary-block:B01186:begin
\needspace{3\baselineskip}
\textbf{my} — \textbf{mīn} (possessive adjective). my; mine. Mīn my/mine: one first-person possessor; the possessive agrees with the possessed noun.\par
% glossary-block:B01186:end
% glossary-block:B01187:begin
\needspace{3\baselineskip}
\textbf{myself} — An ordinary reflexive uses the personal pronoun in the case governed by its verb or preposition. Agreeing self adds the taught emphasis; English -self does not automatically require a separate self. Ordinary reflexive personal pronouns take the governed case; agreeing self adds the taught emphasis. English -self does not automatically require self.\par
% glossary-block:B01187:end
\pdfbookmark[1]{N}{outline-B01188}
% glossary-block:B01188:begin
\needspace{5\baselineskip}
\subsection*{N}
% glossary-block:B01188:end
% glossary-block:B01189:begin
\textbf{naked} — \textbf{nacod} (adj.).\par
% glossary-block:B01189:end
% glossary-block:B01190:begin
\needspace{3\baselineskip}
\textbf{naked people} — \textbf{nacod} (adj.). Substantival adjective.\par
% glossary-block:B01190:end
% glossary-block:B01191:begin
\needspace{3\baselineskip}
\textbf{name} — \textbf{hātan} (v.). call; name. Active hātan call/name; the supplied naming expression may have an accusative object and accusative name. \textbf{nama} (n.).\par
% glossary-block:B01191:end
% glossary-block:B01192:begin
\needspace{3\baselineskip}
\textbf{narrow} — \textbf{nearu} (adj.).\par
% glossary-block:B01192:end
% glossary-block:B01193:begin
\needspace{3\baselineskip}
\textbf{nation} — \textbf{mǣgð} (n.). tribe; nation. \textbf{ðēod} (n.). people; nation.\par
% glossary-block:B01193:end
% glossary-block:B01194:begin
\needspace{3\baselineskip}
\textbf{nations} — \textbf{folc} (n.). peoples. Distinct peoples or national groups; plural sense.\par
% glossary-block:B01194:end
% glossary-block:B01195:begin
\needspace{3\baselineskip}
\textbf{native land} — \textbf{ēðel} (n.). native land; homeland. ēðel in these home-country references; land remains a separate ordinary country word.\par
% glossary-block:B01195:end
% glossary-block:B01196:begin
\needspace{3\baselineskip}
\textbf{naught} — \textbf{nāht} (substantive). nothing; naught. \textbf{nāwiht} (indefinite pronoun/substantive). nothing; naught.\par
% glossary-block:B01196:end
% glossary-block:B01197:begin
\needspace{3\baselineskip}
\textbf{naval force} — \textbf{sciphęre} (n.). fleet; naval force.\par
% glossary-block:B01197:end
% glossary-block:B01198:begin
\needspace{3\baselineskip}
\textbf{near} — \textbf{nēah} (adv.). near; at hand. nēah is the predicate in nigh at hand.\par
% glossary-block:B01198:end
% glossary-block:B01199:begin
\needspace{3\baselineskip}
\textbf{necessarily} — \textbf{nīede} (adverb or particle). necessarily; of necessity. Archaic English needs is an adverb here; do not use the verb need.\par
% glossary-block:B01199:end
% glossary-block:B01200:begin
\needspace{3\baselineskip}
\textbf{need} — \textbf{beðurfan} (v.). need; require. Genitive thing needed. Genitive thing needed; a relative particle may represent that governed relation. \textbf{ðearf} (n.). benefit; good; need. Source contextual nouns; retain the local expression.\par
% glossary-block:B01200:end
% glossary-block:B01201:begin
\needspace{3\baselineskip}
\textbf{needle} — \textbf{nǣdl} (n.).\par
% glossary-block:B01201:end
% glossary-block:B01202:begin
\needspace{3\baselineskip}
\textbf{needs} — \textbf{nīede} (adverb or particle). Archaic English needs is an adverb here; do not use the verb need.\par
% glossary-block:B01202:end
% glossary-block:B01203:begin
\needspace{3\baselineskip}
\textbf{need to} — \textbf{ðurfan} (v.). ðurfan + plain infinitive; distinguish beðurfan + genitive thing needed.\par
% glossary-block:B01203:end
% glossary-block:B01204:begin
\needspace{3\baselineskip}
\textbf{needy person} — \textbf{ðearfa} (n.). poor person; needy person. Noun referring to a person.\par
% glossary-block:B01204:end
% glossary-block:B01205:begin
\needspace{3\baselineskip}
\textbf{neglect} — \textbf{rēcelēasian} (v.). neglect; despise.\par
% glossary-block:B01205:end
% glossary-block:B01206:begin
\needspace{3\baselineskip}
\textbf{neighbor} — \textbf{nēahgebūr} (n.).\par
% glossary-block:B01206:end
% glossary-block:B01207:begin
\needspace{3\baselineskip}
\textbf{neighbour} — \textbf{nēahgebūr} (n.).\par
% glossary-block:B01207:end
% glossary-block:B01208:begin
\needspace{3\baselineskip}
\textbf{neither} — \textbf{nāhwæðer} (adverb or particle).\par
% glossary-block:B01208:end
% glossary-block:B01209:begin
\needspace{3\baselineskip}
\textbf{neither … nor} — \textbf{nāhwæðer} (adverb or particle). Whole supplied nāhwæðer ne … ne with finite negative scope; neither excludes both alternatives.\par
% glossary-block:B01209:end
% glossary-block:B01210:begin
\needspace{3\baselineskip}
\textbf{nest} — \textbf{nest} (n.).\par
% glossary-block:B01210:end
% glossary-block:B01211:begin
\needspace{3\baselineskip}
\textbf{net} — \textbf{nętt} (n.).\par
% glossary-block:B01211:end
% glossary-block:B01212:begin
\needspace{3\baselineskip}
\textbf{never} — \textbf{nǣfre} (adv.). Nǣfre means never; retain finite ne or the licensed negative contraction in negative concord.\par
% glossary-block:B01212:end
% glossary-block:B01213:begin
\needspace{3\baselineskip}
\textbf{never more} — Mā means more/again within the negative expression never more/never again; the clause must supply the negation.\par
% glossary-block:B01213:end
% glossary-block:B01214:begin
\needspace{3\baselineskip}
\textbf{nevertheless} — \textbf{ðā-hwæðre} (conj.). nevertheless; yet. Whole supplied connective. \textbf{ðēah} (conj.; adv.). nevertheless; yet. Independent adverb, distinct from concessive conjunction although.\par
% glossary-block:B01214:end
% glossary-block:B01215:begin
\needspace{3\baselineskip}
\textbf{new} — \textbf{nīwe} (adj.).\par
% glossary-block:B01215:end
% glossary-block:B01216:begin
\needspace{3\baselineskip}
\textbf{nigh} — \textbf{nēah} (adv.). nēah is the predicate in nigh at hand.\par
% glossary-block:B01216:end
% glossary-block:B01217:begin
\needspace{3\baselineskip}
\textbf{night} — \textbf{niht} (n.).\par
% glossary-block:B01217:end
% glossary-block:B01218:begin
\needspace{3\baselineskip}
\textbf{nine} — \textbf{nigon} (cardinal numeral). Cardinal 9.\par
% glossary-block:B01218:end
% glossary-block:B01219:begin
\needspace{3\baselineskip}
\textbf{ninety} — \textbf{hundnigontig} (numeral).\par
% glossary-block:B01219:end
% glossary-block:B01220:begin
\needspace{3\baselineskip}
\textbf{ninth} — \textbf{nigoða} (ordinal).\par
% glossary-block:B01220:end
% glossary-block:B01221:begin
\needspace{3\baselineskip}
\textbf{no} — \textbf{nān} (indefinite). no; no one.\par
% glossary-block:B01221:end
% glossary-block:B01222:begin
\needspace{3\baselineskip}
\textbf{noble} — \textbf{æðele} (adj.). noble; illustrious. \textbf{ealdormann} (n.). prince; leader; noble.\par
% glossary-block:B01222:end
% glossary-block:B01223:begin
\needspace{3\baselineskip}
\textbf{noble people} — \textbf{æðele} (adj.). Substantival adjective in the source.\par
% glossary-block:B01223:end
% glossary-block:B01224:begin
\needspace{3\baselineskip}
\textbf{nobody} — \textbf{nān} (indefinite). Substantival nān means no one; retain finite ne or the licensed negative contraction in negative concord.\par
% glossary-block:B01224:end
% glossary-block:B01225:begin
\needspace{3\baselineskip}
\textbf{noise} — \textbf{gebrec} (n.). crash; noise; din. Supplied accusative singular.\par
% glossary-block:B01225:end
% glossary-block:B01226:begin
\needspace{3\baselineskip}
\textbf{no longer} — \textbf{lęng} (adverb or particle). lęng as the comparative adverb; clause negation supplies no longer.\par
% glossary-block:B01226:end
% glossary-block:B01227:begin
\needspace{3\baselineskip}
\textbf{no man} — \textbf{nān} (indefinite). Substantival nān can render no man/no one where the task omits a separate man; retain negative concord.\par
% glossary-block:B01227:end
% glossary-block:B01228:begin
\needspace{3\baselineskip}
\textbf{noon} — The supplied whole phrase tō middes dæges means at noon/midday; middes and dæges are genitives. \textbf{middes} (adverb or particle). Whole supplied phrase tō middes dæges; preserve its genitives instead of applying ordinary tō + dative mechanically.\par
% glossary-block:B01228:end
% glossary-block:B01229:begin
\needspace{3\baselineskip}
\textbf{no one} — \textbf{nān} (indefinite). no; no one. Substantival nān means no one; retain finite ne or the licensed negative contraction in negative concord.\par
% glossary-block:B01229:end
% glossary-block:B01230:begin
\needspace{3\baselineskip}
\textbf{northeasterly} — \textbf{ēastnorðerne} (adj.).\par
% glossary-block:B01230:end
% glossary-block:B01231:begin
\needspace{3\baselineskip}
\textbf{not} — \textbf{nā} (adv.). Negates a contrasted constituent, distinct from finite ne. Ne is the clause negative not, placed with the finite verb in the taught construction. English do not/did not requires no extra Old English auxiliary do.\par
% glossary-block:B01231:end
% glossary-block:B01232:begin
\needspace{3\baselineskip}
\textbf{not at all} — \textbf{nālæs} (adv.).\par
% glossary-block:B01232:end
% glossary-block:B01233:begin
\needspace{3\baselineskip}
\textbf{not be} — \textbf{wesan} (v.). Supplied negative nǣre in reported supposition.\par
% glossary-block:B01233:end
% glossary-block:B01234:begin
\needspace{3\baselineskip}
\textbf{not have} — \textbf{nāgan} (v.). not have; not possess. Negative possessive verb in its supplied forms.\par
% glossary-block:B01234:end
% glossary-block:B01235:begin
\needspace{3\baselineskip}
\textbf{nothing} — \textbf{nāht} (substantive). nothing; naught. \textbf{nāwiht} (indefinite pronoun/substantive). nothing; naught. Nāwiht means nothing in the bounded accusative construction; retain finite ne or the licensed negative contraction.\par
% glossary-block:B01235:end
% glossary-block:B01236:begin
\needspace{3\baselineskip}
\textbf{not know} — \textbf{nytan} (v.). Supplied negative knowing form.\par
% glossary-block:B01236:end
% glossary-block:B01237:begin
\needspace{3\baselineskip}
\textbf{not possess} — \textbf{nāgan} (v.). not have; not possess. Negative possessive verb in its supplied forms.\par
% glossary-block:B01237:end
% glossary-block:B01238:begin
\needspace{3\baselineskip}
\textbf{not working} — \textbf{īdel} (adj.). idle. Distinct from empty/void.\par
% glossary-block:B01238:end
% glossary-block:B01239:begin
\needspace{3\baselineskip}
\textbf{nourish} — \textbf{āfēdan} (v.). support; nourish. Accusative person supported. \textbf{fēdan} (v.). feed; nourish. Accusative being fed; food may be mid + dative or a bare instrumental in the assigned expressions.\par
% glossary-block:B01239:end
% glossary-block:B01240:begin
\needspace{3\baselineskip}
\textbf{now} — \textbf{nū} (adv.; conj.).  \textbf{sōðlice} (adv.). now; but. Discourse connective corresponding to source Latin autem, distinct from truly. \textbf{witodlice} (adv.). now; and. Discourse connective.\par
% glossary-block:B01240:end
% glossary-block:B01241:begin
\needspace{3\baselineskip}
\textbf{number} — \textbf{gerīman} (v.). number; count. Exact supplied infinitive gerīman; not an independently assigned full paradigm. Supplied infinitive gerīman, count; number. \textbf{getæl} (n.). number; company. Singular collective getæl with the supplied dependent expression; retain source participial genitives.\par
% glossary-block:B01241:end
\pdfbookmark[1]{O}{outline-B01242}
% glossary-block:B01242:begin
\needspace{5\baselineskip}
\subsection*{O}
% glossary-block:B01242:end
% glossary-block:B01243:begin
\textbf{O} — \textbf{ēalā} (interj.). O!. Interjection in address, exclamation or wish.\par
% glossary-block:B01243:end
% glossary-block:B01244:begin
\needspace{3\baselineskip}
\textbf{oath} — \textbf{āð} (n.).\par
% glossary-block:B01244:end
% glossary-block:B01245:begin
\needspace{3\baselineskip}
\textbf{obedience} — \textbf{gehīersumnes} (n.).  \textbf{hīernes} (n.).\par
% glossary-block:B01245:end
% glossary-block:B01246:begin
\needspace{3\baselineskip}
\textbf{obey} — \textbf{gehīersumian} (v.). serve; obey. Supplied source form and government. \textbf{hīeran} (v.). Dative person obeyed. Dative person obeyed; distinct from hearing a sound/object.\par
% glossary-block:B01246:end
% glossary-block:B01247:begin
\needspace{3\baselineskip}
\textbf{observe} — \textbf{behealdan} (v.). Supplied source form observed. \textbf{healdan} (v.). keep; observe. Accusative saying or command observed. \textbf{scēawian} (v.). watch; observe. Accusative person or event watched, including a journey.\par
% glossary-block:B01247:end
% glossary-block:B01248:begin
\needspace{3\baselineskip}
\textbf{obtain} — \textbf{tilian} (v.). obtain; provide. Genitive thing procured; dative reflexive beneficiary in the assigned frame.\par
% glossary-block:B01248:end
% glossary-block:B01249:begin
\needspace{3\baselineskip}
\textbf{obtain for oneself} — \textbf{tilian} (v.). tilian + dative reflexive beneficiary + genitive thing obtained.\par
% glossary-block:B01249:end
% glossary-block:B01250:begin
\needspace{3\baselineskip}
\textbf{occasion} — \textbf{sīð} (n.). time; occasion.\par
% glossary-block:B01250:end
% glossary-block:B01251:begin
\needspace{3\baselineskip}
\textbf{occupation} — \textbf{cræft} (n.). trade; occupation; craft.\par
% glossary-block:B01251:end
% glossary-block:B01252:begin
\needspace{3\baselineskip}
\textbf{occupations} — \textbf{bisgu} (n.). occupations; cares.\par
% glossary-block:B01252:end
% glossary-block:B01253:begin
\needspace{3\baselineskip}
\textbf{occupied} — \textbf{geornfull} (adj.). occupied; intent. Source mental-attention context.\par
% glossary-block:B01253:end
% glossary-block:B01254:begin
\needspace{3\baselineskip}
\textbf{occupy} — \textbf{bisgian} (v.). occupy; trouble.\par
% glossary-block:B01254:end
% glossary-block:B01255:begin
\needspace{3\baselineskip}
\textbf{occur} — \textbf{becuman} (v.). occur; come about. Of the calm in the source.\par
% glossary-block:B01255:end
% glossary-block:B01256:begin
\needspace{3\baselineskip}
\textbf{of} — \textbf{of} (prep.). of; from. Portion from a whole in the supplied partitive phrases; dative whole. of; from. Origin of wisdom in the supplied begetting statement. Possessive/relational English of or 's often becomes a dependent genitive, not preposition of. Partitive, source and lexical government require their own frames.\par
% glossary-block:B01256:end
% glossary-block:B01257:begin
\needspace{3\baselineskip}
\textbf{of Cyrene} — \textbf{Cyrenisc} (adj.). Cyrenian; of Cyrene. Geographical adjective.\par
% glossary-block:B01257:end
% glossary-block:B01258:begin
\needspace{3\baselineskip}
\textbf{off} — \textbf{of} (adv.).\par
% glossary-block:B01258:end
% glossary-block:B01259:begin
\needspace{3\baselineskip}
\textbf{offend} — \textbf{ǣswician} (v.). offend; cause to stumble. Accusative person; ethical/metaphorical offence in the assigned use.\par
% glossary-block:B01259:end
% glossary-block:B01260:begin
\needspace{3\baselineskip}
\textbf{offer} — \textbf{bēodan} (v.).\par
% glossary-block:B01260:end
% glossary-block:B01261:begin
\needspace{3\baselineskip}
\textbf{offspring} — \textbf{sǣd} (n.). seed; descendants. Collective descendants, distinct from literal seed.\par
% glossary-block:B01261:end
% glossary-block:B01262:begin
\needspace{3\baselineskip}
\textbf{of me} — Mīn is the first-person singular personal genitive, of me. Do not add the preposition of; attributive my is a separately taught agreeing possessive.\par
% glossary-block:B01262:end
% glossary-block:B01263:begin
\needspace{3\baselineskip}
\textbf{of necessity} — \textbf{nīede} (adverb or particle). necessarily; of necessity. Archaic English needs is an adverb here; do not use the verb need.\par
% glossary-block:B01263:end
% glossary-block:B01264:begin
\needspace{3\baselineskip}
\textbf{often} — \textbf{genēahhe} (adv.). often; frequently. \textbf{oft} (adv.).\par
% glossary-block:B01264:end
% glossary-block:B01265:begin
\needspace{3\baselineskip}
\textbf{of the sea} — \textbf{sǣlic} (adj.). marine; of the sea.\par
% glossary-block:B01265:end
% glossary-block:B01266:begin
\needspace{3\baselineskip}
\textbf{of us} — Uncer is the personal genitive of us two; ūre is plural of us. Do not add the preposition of; attributive our uses separately taught agreeing possessives.\par
% glossary-block:B01266:end
% glossary-block:B01267:begin
\needspace{3\baselineskip}
\textbf{of us two} — \textbf{ic} (personal pron.). Personal genitive dual after a selection word.\par
% glossary-block:B01267:end
% glossary-block:B01268:begin
\needspace{3\baselineskip}
\textbf{of what sort} — \textbf{hūlic} (interrogative adjective). of what sort?.\par
% glossary-block:B01268:end
% glossary-block:B01269:begin
\needspace{3\baselineskip}
\textbf{of winter} — \textbf{winterlic} (adj.). winter; of winter.\par
% glossary-block:B01269:end
% glossary-block:B01270:begin
\needspace{3\baselineskip}
\textbf{of you} — Personal genitive ðīn refers to one addressee, incer to two, ēower to plural addressees. English of is expressed by the genitive, not an added preposition; attributive your is separately taught.\par
% glossary-block:B01270:end
% glossary-block:B01271:begin
\needspace{3\baselineskip}
\textbf{oil} — \textbf{ęle} (n.).\par
% glossary-block:B01271:end
% glossary-block:B01272:begin
\needspace{3\baselineskip}
\textbf{ointment} — \textbf{sealf} (n.).\par
% glossary-block:B01272:end
% glossary-block:B01273:begin
\needspace{3\baselineskip}
\textbf{old} — \textbf{eald} (adj.).\par
% glossary-block:B01273:end
% glossary-block:B01274:begin
\needspace{3\baselineskip}
\textbf{old age} — \textbf{ieldu} (n.). age; old age.\par
% glossary-block:B01274:end
% glossary-block:B01275:begin
\needspace{3\baselineskip}
\textbf{older} — \textbf{ieldra} (adj.).\par
% glossary-block:B01275:end
% glossary-block:B01276:begin
\needspace{3\baselineskip}
\textbf{oldest} — \textbf{ieldest} (adj.). eldest; oldest. Substantival person/eldest brother in this source context.\par
% glossary-block:B01276:end
% glossary-block:B01277:begin
\needspace{3\baselineskip}
\textbf{old people} — \textbf{eald} (adj.). old people; the elderly. Substantival adjective.\par
% glossary-block:B01277:end
% glossary-block:B01278:begin
\needspace{3\baselineskip}
\textbf{on} — \textbf{in} (prep.). in; on. With dative in the exact locations/circumstances. \textbf{on} (prep.). in; on. Dative for the taught ordinary location construction. Time preposition in the stated-day expression. \textbf{on} (adv.). Adverb in the cloak-wearing relative clause. \textbf{upon} (prep.). on; onto.\par
% glossary-block:B01278:end
% glossary-block:B01279:begin
\needspace{3\baselineskip}
\textbf{on account of} — \textbf{for} (prep.). because of; on account of. for + dative cause in this supplied money expression; distinguish causal conjunction because.\par
% glossary-block:B01279:end
% glossary-block:B01280:begin
\needspace{3\baselineskip}
\textbf{on behalf of} — \textbf{for} (prep.). for; on behalf of. Dative beneficiary.\par
% glossary-block:B01280:end
% glossary-block:B01281:begin
\needspace{3\baselineskip}
\textbf{once more} — \textbf{eft, ęft} (adv.). eft is the repeated-occurrence adverb.\par
% glossary-block:B01281:end
% glossary-block:B01282:begin
\needspace{3\baselineskip}
\textbf{one} — \textbf{ān} (numeral; adj.). one; single. Cardinal one/single and the agreeing adjective alone are distinct uses of ān. Preserve the supplied form where the exercise provides one. \textbf{mann} (n.; impersonal pronoun). one; people in general. Impersonal man with singular finite verb.\par
% glossary-block:B01282:end
% glossary-block:B01283:begin
\needspace{3\baselineskip}
\textbf{one of} — The supplied partitive expression combines the required word for one with of + dative whole. Of alone means of/from, not one of.\par
% glossary-block:B01283:end
% glossary-block:B01284:begin
\needspace{3\baselineskip}
\textbf{one (people in general)} — Impersonal man has a singular finite verb; distinguish numeral ān and one selected member of a group.\par
% glossary-block:B01284:end
% glossary-block:B01285:begin
\needspace{3\baselineskip}
\textbf{one who delivers} — \textbf{ālīesend} (n.). Redeemer; one who delivers.\par
% glossary-block:B01285:end
% glossary-block:B01286:begin
\needspace{3\baselineskip}
\textbf{one who shows mercy} — \textbf{gemiltsiend} (n.). one who shows mercy; pitier.\par
% glossary-block:B01286:end
% glossary-block:B01287:begin
\needspace{3\baselineskip}
\textbf{only-begotten} — \textbf{āncęnned} (adj.).\par
% glossary-block:B01287:end
% glossary-block:B01288:begin
\needspace{3\baselineskip}
\textbf{on one's side} — The bounded on … hand phrase means on someone’s side; hand is accusative even with standan. Retain the supplied possessor and side expression.\par
% glossary-block:B01288:end
% glossary-block:B01289:begin
\needspace{3\baselineskip}
\textbf{on that evidence} — \textbf{be} (prep.). by that; from that. Refers to the preceding proposition; supplied expression. by that; on that evidence. Evidence/inference frame, not an unrestricted source preposition.\par
% glossary-block:B01289:end
% glossary-block:B01290:begin
\needspace{3\baselineskip}
\textbf{on the side of} — The bounded on … hand phrase means on someone’s side; hand is accusative even with standan. Retain the supplied possessor and side expression.\par
% glossary-block:B01290:end
% glossary-block:B01291:begin
\needspace{3\baselineskip}
\textbf{on the way home} — \textbf{hāmweard} (adverb or particle). Supplied hāmweard with wesan; distinguish noun hām and directional adverb hām.\par
% glossary-block:B01291:end
% glossary-block:B01292:begin
\needspace{3\baselineskip}
\textbf{onto} — \textbf{in} (prep.). toward; onto. Supplied phrase directing the mind toward these things. \textbf{on} (prep.). into; onto. Accusative goal of placement in the taught frame. \textbf{upon} (prep.). on; onto.\par
% glossary-block:B01292:end
% glossary-block:B01293:begin
\needspace{3\baselineskip}
\textbf{on top of} — \textbf{uppan} (prep.). upon; on top of.\par
% glossary-block:B01293:end
% glossary-block:B01294:begin
\needspace{3\baselineskip}
\textbf{open} — \textbf{geopenian} (v.). Accusative thing opened, including mouth or eyes. \textbf{ontȳnan} (v.). Supplied preterite opened. \textbf{open} (adj.).  \textbf{undōn} (v.). undo; open.\par
% glossary-block:B01294:end
% glossary-block:B01295:begin
\needspace{3\baselineskip}
\textbf{opponent} — \textbf{gewinna} (n.). enemy; opponent. Supplied plural source forms.\par
% glossary-block:B01295:end
% glossary-block:B01296:begin
\needspace{3\baselineskip}
\textbf{oppose} — \textbf{standan} (v.). oppose; stand against. Supplied standan ongēan frame, distinct from simple stand.\par
% glossary-block:B01296:end
% glossary-block:B01297:begin
\needspace{3\baselineskip}
\textbf{opposed to} — \textbf{wið} (prep.). against; opposed to. Accusative opponent.\par
% glossary-block:B01297:end
% glossary-block:B01298:begin
\needspace{3\baselineskip}
\textbf{opposite} — \textbf{ongēan} (adv.; prep.; conj.). against; opposite. Accusative landmark or opponent.\par
% glossary-block:B01298:end
% glossary-block:B01299:begin
\needspace{3\baselineskip}
\textbf{oppress} — \textbf{hīenan} (v.). oppress; humble.\par
% glossary-block:B01299:end
% glossary-block:B01300:begin
\needspace{3\baselineskip}
\textbf{or} — \textbf{oððe} (conj.).\par
% glossary-block:B01300:end
% glossary-block:B01301:begin
\needspace{3\baselineskip}
\textbf{oratory} — \textbf{cirice} (n.). Chad's separate place of prayer; distinct from the ordinary church sense.\par
% glossary-block:B01301:end
% glossary-block:B01302:begin
\needspace{3\baselineskip}
\textbf{order} — \textbf{hǣs} (n.). command; order. \textbf{hātan} (v.). command; order. Accusative actor + plain infinitive for the commanded action; distinct from call/name.\par
% glossary-block:B01302:end
% glossary-block:B01303:begin
\needspace{3\baselineskip}
\textbf{order someone to} — \textbf{hātan} (v.). hātan + accusative actor + plain infinitive.\par
% glossary-block:B01303:end
% glossary-block:B01304:begin
\needspace{3\baselineskip}
\textbf{ore} — \textbf{ōra} (n.). vein of metal; ore (sense uncertain). Supplied dative plural after on.\par
% glossary-block:B01304:end
% glossary-block:B01305:begin
\needspace{3\baselineskip}
\textbf{Osburg} — \textbf{Osburg} (proper name).\par
% glossary-block:B01305:end
% glossary-block:B01306:begin
\needspace{3\baselineskip}
\textbf{other} — \textbf{ōðer} (adj.; ordinal). other; another.\par
% glossary-block:B01306:end
% glossary-block:B01307:begin
\needspace{3\baselineskip}
\textbf{ought to} — \textbf{sculan} (v.). sculan + plain infinitive for obligation; English have to is not possessive habban. Supplied have something to offer in Unit 6 is a different construction.\par
% glossary-block:B01307:end
% glossary-block:B01308:begin
\needspace{3\baselineskip}
\textbf{our} — \textbf{ic} (personal pron.). Personal genitive plural used attributively in the exact ūre hǣle expression. \textbf{uncer} (possessive adjective). our; ours. Two possessors; agreeing possessive. \textbf{ūre} (possessive adjective). our; ours. Plural possessors; agreeing possessive. \textbf{uncer} (possessive adjective), \textbf{ūre} (possessive adjective). Uncer means belonging to us two; ūre belonging to us in the plural. The possessive agrees with the possessed noun.\par
% glossary-block:B01308:end
% glossary-block:B01309:begin
\needspace{3\baselineskip}
\textbf{ours} — \textbf{uncer} (possessive adjective). our; ours. Two possessors; agreeing possessive. \textbf{ūre} (possessive adjective). our; ours. Plural possessors; agreeing possessive. \textbf{uncer} (possessive adjective), \textbf{ūre} (possessive adjective). Uncer means belonging to us two; ūre belonging to us in the plural. The possessive agrees with the possessed noun.\par
% glossary-block:B01309:end
% glossary-block:B01310:begin
\needspace{3\baselineskip}
\textbf{ourselves} — An ordinary reflexive uses the personal pronoun in the case governed by its verb or preposition. Agreeing self adds the taught emphasis; English -self does not automatically require a separate self. Ordinary reflexive personal pronouns take the governed case; agreeing self adds the taught emphasis. English -self does not automatically require self.\par
% glossary-block:B01310:end
% glossary-block:B01311:begin
\needspace{3\baselineskip}
\textbf{out} — \textbf{ūt} (adv.).\par
% glossary-block:B01311:end
% glossary-block:B01312:begin
\needspace{3\baselineskip}
\textbf{out of} — \textbf{æt} (prep.). from; out of. Proceeds of work, in the assigned expression. \textbf{of} (prep.). from; out of. With dative.\par
% glossary-block:B01312:end
% glossary-block:B01313:begin
\needspace{3\baselineskip}
\textbf{outside} — \textbf{būtan} (prep.). Būtan + dative means without a companion or outside a city in the respective sense. Conditional unless is separate. \textbf{ūte} (adv.).\par
% glossary-block:B01313:end
% glossary-block:B01314:begin
\needspace{3\baselineskip}
\textbf{oven} — \textbf{ofen} (n.).\par
% glossary-block:B01314:end
% glossary-block:B01315:begin
\needspace{3\baselineskip}
\textbf{over} — \textbf{ofer} (prep.). over; across. Spatial frame; case varies with the specific construction. Do not impose a universal motion rule. Accusative dependent in the taught authority frame. over; about. Object of rejoicing, with the assigned accusative frame. Spatial ofer can mean over/across; its case depends on the assigned spatial frame. Ofer in the authority sense over takes an accusative dependent. Rank and preference have further separately qualified senses.\par
% glossary-block:B01315:end
% glossary-block:B01316:begin
\needspace{3\baselineskip}
\textbf{overcome} — \textbf{oferwinnan} (v.). overcome; conquer.\par
% glossary-block:B01316:end
% glossary-block:B01317:begin
\needspace{3\baselineskip}
\textbf{overlay} — \textbf{oferwyrcean} (v.). overlay; cover. Accusative object, mid + dative material used to cover it.\par
% glossary-block:B01317:end
% glossary-block:B01318:begin
\needspace{3\baselineskip}
\textbf{overtake} — \textbf{fōn} (v.). Accusative person overtaken in these tasks; the source's bracketed caught is the route.\par
% glossary-block:B01318:end
% glossary-block:B01319:begin
\needspace{3\baselineskip}
\textbf{overthrow} — \textbf{tōweorpan} (v.). break; overthrow. Break a commandment in the assigned use.\par
% glossary-block:B01319:end
% glossary-block:B01320:begin
\needspace{3\baselineskip}
\textbf{own} — \textbf{āgen} (adj.).  \textbf{āgan} (v.). have; possess. Accusative possession; positive present and preterite indicative in the taught scope.\par
% glossary-block:B01320:end
\pdfbookmark[1]{P}{outline-B01321}
% glossary-block:B01321:begin
\needspace{5\baselineskip}
\subsection*{P}
% glossary-block:B01321:end
% glossary-block:B01322:begin
\textbf{palm branch} — \textbf{palmtwīg} (n.). palm twig; palm branch.\par
% glossary-block:B01322:end
% glossary-block:B01323:begin
\needspace{3\baselineskip}
\textbf{palm twig} — \textbf{palmtwīg} (n.). palm twig; palm branch.\par
% glossary-block:B01323:end
% glossary-block:B01324:begin
\needspace{3\baselineskip}
\textbf{parch} — \textbf{forbærnan} (v.). parch; scorch.\par
% glossary-block:B01324:end
% glossary-block:B01325:begin
\needspace{3\baselineskip}
\textbf{part} — \textbf{dǣl} (n.). dǣl + genitive of the substance, as in sumne dǣl mętes; the amount is accusative object.\par
% glossary-block:B01325:end
% glossary-block:B01326:begin
\needspace{3\baselineskip}
\textbf{pass} — \textbf{gewītan} (v.). pass away; cease to hold. Of a decree.\par
% glossary-block:B01326:end
% glossary-block:B01327:begin
\needspace{3\baselineskip}
\textbf{passage} — \textbf{sīðfæt} (n.). journey; passage.\par
% glossary-block:B01327:end
% glossary-block:B01328:begin
\needspace{3\baselineskip}
\textbf{pass away} — \textbf{gewītan} (v.). pass away; cease to hold. Of a decree.\par
% glossary-block:B01328:end
% glossary-block:B01329:begin
\needspace{3\baselineskip}
\textbf{path} — \textbf{weg} (n.). way; path. Weg means road, path or way.\par
% glossary-block:B01329:end
% glossary-block:B01330:begin
\needspace{3\baselineskip}
\textbf{pay} — \textbf{forgieldan} (v.). pay; repay; requite. Dative recipient, accusative thing repaid in the selected frame.\par
% glossary-block:B01330:end
% glossary-block:B01331:begin
\needspace{3\baselineskip}
\textbf{peace} — \textbf{sibb} (n.).\par
% glossary-block:B01331:end
% glossary-block:B01332:begin
\needspace{3\baselineskip}
\textbf{pearl} — \textbf{męregrēote} (n.).\par
% glossary-block:B01332:end
% glossary-block:B01333:begin
\needspace{3\baselineskip}
\textbf{peer} — \textbf{gelīca} (n.). equal; peer.\par
% glossary-block:B01333:end
% glossary-block:B01334:begin
\needspace{3\baselineskip}
\textbf{penalty} — \textbf{wīte} (n.). punishment; penalty.\par
% glossary-block:B01334:end
% glossary-block:B01335:begin
\needspace{3\baselineskip}
\textbf{penetrate} — \textbf{fēolan} (v.). Preserve the supplied/taught lexical frame; not a synonym for every kind of entering.\par
% glossary-block:B01335:end
% glossary-block:B01336:begin
\needspace{3\baselineskip}
\textbf{penny} — \textbf{pęning} (n.). penny; coin.\par
% glossary-block:B01336:end
% glossary-block:B01337:begin
\needspace{3\baselineskip}
\textbf{Pentapolis} — \textbf{Pentapolis} (proper place name). Cyrenian region in this story.\par
% glossary-block:B01337:end
% glossary-block:B01338:begin
\needspace{3\baselineskip}
\textbf{people} — \textbf{folc} (n.). Collective singular. \textbf{ðēod} (n.). people; nation.\par
% glossary-block:B01338:end
% glossary-block:B01339:begin
\needspace{3\baselineskip}
\textbf{people in general} — \textbf{mann} (n.; impersonal pronoun). one; people in general. Impersonal man with singular finite verb.\par
% glossary-block:B01339:end
% glossary-block:B01340:begin
\needspace{3\baselineskip}
\textbf{people not noble} — \textbf{unæðele} (adj.). Substantival adjective in the supplied contrast.\par
% glossary-block:B01340:end
% glossary-block:B01341:begin
\needspace{3\baselineskip}
\textbf{peoples} — \textbf{folc} (n.). Distinct peoples or national groups; plural sense.\par
% glossary-block:B01341:end
% glossary-block:B01342:begin
\needspace{3\baselineskip}
\textbf{perceive} — \textbf{oncnāwan} (v.). recognize; perceive. Supplied past source form. \textbf{ongietan} (v.). understand; perceive. Supplied infinitive after a modal. perceive; realize. understand; perceive.\par
% glossary-block:B01342:end
% glossary-block:B01343:begin
\needspace{3\baselineskip}
\textbf{perform} — \textbf{āwyrcean} (v.). do; deserve. In the source, deserving through one's conduct; retain the complete local frame. \textbf{dōn} (v.). do; perform. \textbf{ęfnan} (v.).  \textbf{forðtēon} (v.). present; perform. Source performance sense; retain the actual form's spelling. \textbf{gedōn} (v.). do; perform. \textbf{gefręmman} (v.). perform; accomplish. A work is the thing accomplished. \textbf{plegian} (v.). act; perform. Performance context; ordinary play is a separate assigned sense. \textbf{wyrcean} (v.). work; perform. Accusative work or deed.\par
% glossary-block:B01343:end
% glossary-block:B01344:begin
\needspace{3\baselineskip}
\textbf{perhaps} — \textbf{wēn} (n.). expectation; chance. wēn is ðæt + subjunctive supplies the taught perhaps construction. wēn is ðæt + present subjunctive in the taught possibility frame.\par
% glossary-block:B01344:end
% glossary-block:B01345:begin
\needspace{3\baselineskip}
\textbf{peril} — \textbf{frēcnes} (n.). danger; peril. \textbf{pliht} (n.). peril; risk.\par
% glossary-block:B01345:end
% glossary-block:B01346:begin
\needspace{3\baselineskip}
\textbf{period} — \textbf{fierst} (n.). period; interval. Masculine noun in the source time expression.\par
% glossary-block:B01346:end
% glossary-block:B01347:begin
\needspace{3\baselineskip}
\textbf{perish} — \textbf{forweorðan} (v.).\par
% glossary-block:B01347:end
% glossary-block:B01348:begin
\needspace{3\baselineskip}
\textbf{permission} — \textbf{lēaf} (n.).\par
% glossary-block:B01348:end
% glossary-block:B01349:begin
\needspace{3\baselineskip}
\textbf{permit} — \textbf{ālīefan} (v.). permit; allow. Impersonal permission frame, with dative permitted person; consult the complete supplied expression. \textbf{gelīefst} (v.; supplied pres. ind. 2 sg.). permit; allow. \textbf{līefan} (v.). allow; permit.\par
% glossary-block:B01349:end
% glossary-block:B01350:begin
\needspace{3\baselineskip}
\textbf{persecute} — \textbf{ēhtan} (v.). pursue; persecute. Accusative person in the selected production frame; a separately supported genitive alternative is not the selected target.\par
% glossary-block:B01350:end
% glossary-block:B01351:begin
\needspace{3\baselineskip}
\textbf{petition} — \textbf{bēn} (n.). prayer; petition. Source dative plural after tō. petition; appeal.\par
% glossary-block:B01351:end
% glossary-block:B01352:begin
\needspace{3\baselineskip}
\textbf{Picts} — \textbf{Piht} (ethnonym). Supplied people-name forms.\par
% glossary-block:B01352:end
% glossary-block:B01353:begin
\needspace{3\baselineskip}
\textbf{pitier} — \textbf{gemiltsiend} (n.). one who shows mercy; pitier.\par
% glossary-block:B01353:end
% glossary-block:B01354:begin
\needspace{3\baselineskip}
\textbf{pity} — \textbf{gemiltsian} (v.). pity; have compassion on. Dative person. gemiltsian + dative person. A wish uses the appropriate subjunctive, not a separate have verb.\par
% glossary-block:B01354:end
% glossary-block:B01355:begin
\needspace{3\baselineskip}
\textbf{place} — \textbf{āsęttan} (v.). set; place. \textbf{dōn} (v.). put; place. Accusative object; on + accusative goal in the assigned frame. \textbf{gelōgian} (v.). place; set. \textbf{gesęttan} (v.). set; place. Supplied gesętte. set; place; establish. \textbf{lęcgean} (v.). lay; place. Supplied source form. lay; place. \textbf{sęttan} (v.). set; place. Place an accusative object; destination is a separate phrase.\par
% glossary-block:B01355:end
% glossary-block:B01356:begin
\needspace{3\baselineskip}
\textbf{placed} — \textbf{gesęttan} (v.). placed; situated. Supplied participial complement of understood us; not the identical-looking finite preterite.\par
% glossary-block:B01356:end
% glossary-block:B01357:begin
\needspace{3\baselineskip}
\textbf{place of prayer} — \textbf{cirice} (n.). oratory. Chad's separate place of prayer; distinct from the ordinary church sense.\par
% glossary-block:B01357:end
% glossary-block:B01358:begin
\needspace{3\baselineskip}
\textbf{plait} — \textbf{windan} (v.). wind; plait. Accusative crown; material follows of. Cook's platted means plaited.\par
% glossary-block:B01358:end
% glossary-block:B01359:begin
\needspace{3\baselineskip}
\textbf{plant} — \textbf{wyrt} (n.). herb; plant.\par
% glossary-block:B01359:end
% glossary-block:B01360:begin
\needspace{3\baselineskip}
\textbf{plat} — \textbf{windan} (v.). Accusative crown; material follows of. Cook's platted means plaited.\par
% glossary-block:B01360:end
% glossary-block:B01361:begin
\needspace{3\baselineskip}
\textbf{platted} — \textbf{windan} (v.). Accusative crown; material follows of. Cook's platted means plaited.\par
% glossary-block:B01361:end
% glossary-block:B01362:begin
\needspace{3\baselineskip}
\textbf{play} — \textbf{plega} (n.). play; game. \textbf{plegian} (v.). Intransitive; the reading supplies local wið participant and mid instrument phrases.\par
% glossary-block:B01362:end
% glossary-block:B01363:begin
\needspace{3\baselineskip}
\textbf{play harp} — \textbf{hearpian} (v.).\par
% glossary-block:B01363:end
% glossary-block:B01364:begin
\needspace{3\baselineskip}
\textbf{play the harp} — \textbf{hearpian} (v.). harp. Supplied genitive harper/harping construction, with instrument; do not reconstruct it from the infinitive alone. Preserve the supplied genitive harper/harping expression; headword alone does not reconstruct its syntax.\par
% glossary-block:B01364:end
% glossary-block:B01365:begin
\needspace{3\baselineskip}
\textbf{play with} — \textbf{wið} (prep.). with; against. Dative fellow players in the exact supplied playing frame.\par
% glossary-block:B01365:end
% glossary-block:B01366:begin
\needspace{3\baselineskip}
\textbf{pleasant} — \textbf{wynsum} (adj.). pleasant; winsome.\par
% glossary-block:B01366:end
% glossary-block:B01367:begin
\needspace{3\baselineskip}
\textbf{pleasantly} — \textbf{gecwēmlīce} (adv.). pleasantly; agreeably.\par
% glossary-block:B01367:end
% glossary-block:B01368:begin
\needspace{3\baselineskip}
\textbf{please} — \textbf{lician} (v.). The thing pleasing is the subject; the persons pleased are dative in the supplied reading. \textbf{lystan} (v.). please; be desired. Impersonal source expression; keep the complete supplied frame.\par
% glossary-block:B01368:end
% glossary-block:B01369:begin
\needspace{3\baselineskip}
\textbf{pleasing} — \textbf{gecwēme} (adj.). pleasing; acceptable.\par
% glossary-block:B01369:end
% glossary-block:B01370:begin
\needspace{3\baselineskip}
\textbf{plectrum} — \textbf{hearpenægl} (n.). plectrum; harp pick.\par
% glossary-block:B01370:end
% glossary-block:B01371:begin
\needspace{3\baselineskip}
\textbf{plow} — \textbf{ęrian} (v.).\par
% glossary-block:B01371:end
% glossary-block:B01372:begin
\needspace{3\baselineskip}
\textbf{plowman} — \textbf{ierðling} (n.). farmer; plowman.\par
% glossary-block:B01372:end
% glossary-block:B01373:begin
\needspace{3\baselineskip}
\textbf{plowshare} — \textbf{sulhscear} (n.).\par
% glossary-block:B01373:end
% glossary-block:B01374:begin
\needspace{3\baselineskip}
\textbf{plunder} — \textbf{hęrgian} (v.). plunder; ravage.\par
% glossary-block:B01374:end
% glossary-block:B01375:begin
\needspace{3\baselineskip}
\textbf{ply} — \textbf{begān} (v.). ply; practise. Accusative trade or occupation.\par
% glossary-block:B01375:end
% glossary-block:B01376:begin
\needspace{3\baselineskip}
\textbf{poem} — \textbf{lēoð} (n.). poem; poetry; verse.\par
% glossary-block:B01376:end
% glossary-block:B01377:begin
\needspace{3\baselineskip}
\textbf{poet} — \textbf{lēoðwyrhta} (n.).\par
% glossary-block:B01377:end
% glossary-block:B01378:begin
\needspace{3\baselineskip}
\textbf{poetry} — \textbf{lēoð} (n.). poem; poetry; verse.\par
% glossary-block:B01378:end
% glossary-block:B01379:begin
\needspace{3\baselineskip}
\textbf{ponder} — \textbf{smēagan} (v.). ponder; consider. Supplied preterite source form.\par
% glossary-block:B01379:end
% glossary-block:B01380:begin
\needspace{3\baselineskip}
\textbf{poor man} — \textbf{wǣdla} (n.), \textbf{ðearfa} (n.). Wǣdla denotes a poor or destitute man; ðearfa a poor or needy person. Use the noun specified by the exercise.\par
% glossary-block:B01380:end
% glossary-block:B01381:begin
\needspace{3\baselineskip}
\textbf{poor or destitute man} — \textbf{wǣdla} (n.).\par
% glossary-block:B01381:end
% glossary-block:B01382:begin
\needspace{3\baselineskip}
\textbf{poor person} — \textbf{ðearfa} (n.). poor person; needy person. Noun referring to a person. \textbf{wǣdla} (n.), \textbf{ðearfa} (n.). Wǣdla denotes a poor or destitute man; ðearfa a poor or needy person. Use the noun specified by the exercise.\par
% glossary-block:B01382:end
% glossary-block:B01383:begin
\needspace{3\baselineskip}
\textbf{possess} — \textbf{habban} (v.). have; possess. Accusative possession; distinguish the perfect auxiliary. \textbf{āgan} (v.). have; possess. Accusative possession; positive present and preterite indicative in the taught scope.\par
% glossary-block:B01383:end
% glossary-block:B01384:begin
\needspace{3\baselineskip}
\textbf{possessions} — \textbf{ǣht} (n.). possessions; goods.\par
% glossary-block:B01384:end
% glossary-block:B01385:begin
\needspace{3\baselineskip}
\textbf{pound} — \textbf{pund} (n.). Weight measure in the supplied gold quantity.\par
% glossary-block:B01385:end
% glossary-block:B01386:begin
\needspace{3\baselineskip}
\textbf{pour} — \textbf{scęncan} (v.). pour out; give to drink. Accusative drink, dative recipient.\par
% glossary-block:B01386:end
% glossary-block:B01387:begin
\needspace{3\baselineskip}
\textbf{pour out} — \textbf{scęncan} (v.). pour out; give to drink. Accusative drink, dative recipient.\par
% glossary-block:B01387:end
% glossary-block:B01388:begin
\needspace{3\baselineskip}
\textbf{power} — \textbf{mægen} (n.). power; strength. Heavenly powers in 37.2 is the contextual plural use. \textbf{miht} (n.). ability; power.\par
% glossary-block:B01388:end
% glossary-block:B01389:begin
\needspace{3\baselineskip}
\textbf{powerful} — \textbf{mihtig} (adj.). mighty; powerful.\par
% glossary-block:B01389:end
% glossary-block:B01390:begin
\needspace{3\baselineskip}
\textbf{powerful people} — \textbf{miht} (n.). powers; powerful people. Those holding high station in the source's concluding human application.\par
% glossary-block:B01390:end
% glossary-block:B01391:begin
\needspace{3\baselineskip}
\textbf{powers} — \textbf{miht} (n.). powers; powerful people. Those holding high station in the source's concluding human application.\par
% glossary-block:B01391:end
% glossary-block:B01392:begin
\needspace{3\baselineskip}
\textbf{practice} — \textbf{begān} (v.). ply; practise. Accusative trade or occupation. pursue; practise. Study is the thing pursued.\par
% glossary-block:B01392:end
% glossary-block:B01393:begin
\needspace{3\baselineskip}
\textbf{practise} — \textbf{begān} (v.). ply; practise. Accusative trade or occupation. pursue; practise. Study is the thing pursued.\par
% glossary-block:B01393:end
% glossary-block:B01394:begin
\needspace{3\baselineskip}
\textbf{praise} — \textbf{hęrian} (v.). Genitive plural present-participle use, of people praising. Supplied source form praised. \textbf{lof} (n.).\par
% glossary-block:B01394:end
% glossary-block:B01395:begin
\needspace{3\baselineskip}
\textbf{pray} — \textbf{biddan} (v.). beseech; entreat. Accusative person; the requested action can be a ðæt clause. Accusative person beseeched; ðæt + subjunctive gives the requested act. Prayer to God with gebiddan is a different entry. \textbf{gebiddan} (v.). Supplied wæs gebiddende, rendered prayed in this source. With an accusative reflexive pronoun in the assigned prayer frame. For prayer, gebiddan takes the assigned accusative reflexive pronoun. Older English pray meaning beseech/entreat a person uses biddan + accusative person; a requested action can be a ðæt clause.\par
% glossary-block:B01395:end
% glossary-block:B01396:begin
\needspace{3\baselineskip}
\textbf{prayer} — \textbf{bēn} (n.). prayer; petition. Source dative plural after tō. \textbf{gebed} (n.).\par
% glossary-block:B01396:end
% glossary-block:B01397:begin
\needspace{3\baselineskip}
\textbf{precious} — \textbf{dēorwierðe} (adj.). precious; valuable. \textbf{dēorwurð} (adj.). precious; costly.\par
% glossary-block:B01397:end
% glossary-block:B01398:begin
\needspace{3\baselineskip}
\textbf{precious stone} — \textbf{gimm} (n.). gem; precious stone.\par
% glossary-block:B01398:end
% glossary-block:B01399:begin
\needspace{3\baselineskip}
\textbf{precisely} — \textbf{efne} (interj.). just; precisely; even.\par
% glossary-block:B01399:end
% glossary-block:B01400:begin
\needspace{3\baselineskip}
\textbf{prefer} — \textbf{lēofra} (adj.). preferable; dearer. Dative experiencer + is lēofre + infinitive in the taught preference frame. The preference construction is dative experiencer + is lēofre + plain infinitive, is preferable to someone. It renders had/would rather without requiring past time. Wician in the supplied example is the action preferred, sojourn/live with, not a word meaning prefer.\par
% glossary-block:B01400:end
% glossary-block:B01401:begin
\needspace{3\baselineskip}
\textbf{preferable} — \textbf{lēofra} (adj.). preferable; dearer. Dative experiencer + is lēofre + infinitive in the taught preference frame.\par
% glossary-block:B01401:end
% glossary-block:B01402:begin
\needspace{3\baselineskip}
\textbf{prepare} — \textbf{gearcian} (v.).  \textbf{gierwan} (v.). Accusative thing prepared.\par
% glossary-block:B01402:end
% glossary-block:B01403:begin
\needspace{3\baselineskip}
\textbf{present} — \textbf{andweard} (adj.).  \textbf{forðtēon} (v.). present; perform. Source performance sense; retain the actual form's spelling. \textbf{giefu} (n.). gift. Supplied plural gifts. \textbf{lāc} (n.). gift; present.\par
% glossary-block:B01403:end
% glossary-block:B01404:begin
\needspace{3\baselineskip}
\textbf{price} — \textbf{weorð} (n.). price; value.\par
% glossary-block:B01404:end
% glossary-block:B01405:begin
\needspace{3\baselineskip}
\textbf{priest} — \textbf{mæsseprēost} (n.).\par
% glossary-block:B01405:end
% glossary-block:B01406:begin
\needspace{3\baselineskip}
\textbf{primacy} — \textbf{ealdordōm} (n.). chief place; primacy. \textbf{ealdorscipe} (n.). chief place; primacy.\par
% glossary-block:B01406:end
% glossary-block:B01407:begin
\needspace{3\baselineskip}
\textbf{prince} — \textbf{ealdormann} (n.). prince; leader; noble.\par
% glossary-block:B01407:end
% glossary-block:B01408:begin
\needspace{3\baselineskip}
\textbf{princess} — \textbf{cwēn} (n.). The context distinguishes princess from queen. The selected accusative singular cwēn is endingless.\par
% glossary-block:B01408:end
% glossary-block:B01409:begin
\needspace{3\baselineskip}
\textbf{prior} — \textbf{furðra} (adj.). foremost; prior.\par
% glossary-block:B01409:end
% glossary-block:B01410:begin
\needspace{3\baselineskip}
\textbf{prison} — \textbf{cweartern} (n.).\par
% glossary-block:B01410:end
% glossary-block:B01411:begin
\needspace{3\baselineskip}
\textbf{produce} — \textbf{ācęnnan} (v.). produce; bring forth. \textbf{beran} (v.). bear; produce. Accusative produce; distinct from physical carrying. \textbf{tēon} (v.). bring forth; produce. Accusative thing produced with particle forð; use the assigned tēon forms. bring forth; produce. Supplied third-plural present subjunctive, with waters as subject. \textbf{wæstm} (n.). fruit; produce.\par
% glossary-block:B01411:end
% glossary-block:B01412:begin
\needspace{3\baselineskip}
\textbf{produce fruit} — \textbf{wyrcean} (v.). yield; produce. Accusative fruit in the taught/source expression. Accusative fruit; birð wæstm is a recognized lexical alternative in the key.\par
% glossary-block:B01412:end
% glossary-block:B01413:begin
\needspace{3\baselineskip}
\textbf{produce of the earth} — \textbf{ierð} (n.). crops; produce of the earth.\par
% glossary-block:B01413:end
% glossary-block:B01414:begin
\needspace{3\baselineskip}
\textbf{produce seed} — \textbf{beran} (v.). Accusative produce; distinguish physical carrying.\par
% glossary-block:B01414:end
% glossary-block:B01415:begin
\needspace{3\baselineskip}
\textbf{productive of} — \textbf{beran} (v.). rich in; productive of. Supplied berende on + dative ore expression; participial adjective.\par
% glossary-block:B01415:end
% glossary-block:B01416:begin
\needspace{3\baselineskip}
\textbf{profit} — \textbf{gestrēon} (n.). gain; profit.\par
% glossary-block:B01416:end
% glossary-block:B01417:begin
\needspace{3\baselineskip}
\textbf{promise} — \textbf{behātan} (v.).  \textbf{gehātan} (v.). Accusative thing promised; a dative recipient may be added.\par
% glossary-block:B01417:end
% glossary-block:B01418:begin
\needspace{3\baselineskip}
\textbf{prompt} — \textbf{myngian} (v.). admonish; prompt. Supplied participial source use.\par
% glossary-block:B01418:end
% glossary-block:B01419:begin
\needspace{3\baselineskip}
\textbf{property} — \textbf{gōd} (n.). goods; property. Plural lexical noun.\par
% glossary-block:B01419:end
% glossary-block:B01420:begin
\needspace{3\baselineskip}
\textbf{prophet} — \textbf{wītga} (n.).\par
% glossary-block:B01420:end
% glossary-block:B01421:begin
\needspace{3\baselineskip}
\textbf{prostrate oneself} — \textbf{āstręccean} (v.). Reflexive accusative; the before phrase retains its own case/construal.\par
% glossary-block:B01421:end
% glossary-block:B01422:begin
\needspace{3\baselineskip}
\textbf{protect} — \textbf{gescieldan} (v.). shield; protect.\par
% glossary-block:B01422:end
% glossary-block:B01423:begin
\needspace{3\baselineskip}
\textbf{protection} — \textbf{gebeorg} (n.). protection; defence.\par
% glossary-block:B01423:end
% glossary-block:B01424:begin
\needspace{3\baselineskip}
\textbf{proud} — \textbf{wlanc} (adj.). proud; wealthy. Contextually the rich man addressed.\par
% glossary-block:B01424:end
% glossary-block:B01425:begin
\needspace{3\baselineskip}
\textbf{provender} — \textbf{fōdor} (n.). Provender is routed by Cook's bracketed fodder to fōdor.\par
% glossary-block:B01425:end
% glossary-block:B01426:begin
\needspace{3\baselineskip}
\textbf{provide} — \textbf{tilian} (v.). obtain; provide. Genitive thing procured; dative reflexive beneficiary in the assigned frame.\par
% glossary-block:B01426:end
% glossary-block:B01427:begin
\needspace{3\baselineskip}
\textbf{provide for oneself} — \textbf{tilian} (v.). obtain; provide. Genitive thing procured; dative reflexive beneficiary in the assigned frame. tilian + dative reflexive beneficiary + genitive thing obtained.\par
% glossary-block:B01427:end
% glossary-block:B01428:begin
\needspace{3\baselineskip}
\textbf{provisions} — \textbf{ēst} (n.). provisions; food.\par
% glossary-block:B01428:end
% glossary-block:B01429:begin
\needspace{3\baselineskip}
\textbf{public place} — \textbf{strǣt} (n.). Street used as a gathering place in the assigned prompt.\par
% glossary-block:B01429:end
% glossary-block:B01430:begin
\needspace{3\baselineskip}
\textbf{pull} — \textbf{tēon} (v.). draw; pull.\par
% glossary-block:B01430:end
% glossary-block:B01431:begin
\needspace{3\baselineskip}
\textbf{punishment} — \textbf{wīte} (n.). punishment; penalty.\par
% glossary-block:B01431:end
% glossary-block:B01432:begin
\needspace{3\baselineskip}
\textbf{pure} — \textbf{clǣne} (adj.). pure; clean. clean; pure; blameless.\par
% glossary-block:B01432:end
% glossary-block:B01433:begin
\needspace{3\baselineskip}
\textbf{purple cloth} — \textbf{gōdwębb} (n.).\par
% glossary-block:B01433:end
% glossary-block:B01434:begin
\needspace{3\baselineskip}
\textbf{pursue} — \textbf{begān} (v.). pursue; practise. Study is the thing pursued. \textbf{ēhtan} (v.). pursue; persecute. Accusative person in the selected production frame; a separately supported genitive alternative is not the selected target.\par
% glossary-block:B01434:end
% glossary-block:B01435:begin
\needspace{3\baselineskip}
\textbf{put} — \textbf{dōn} (v.). put; place. Accusative object; on + accusative goal in the assigned frame. \textbf{sęttan} (v.). Place an accusative object; destination is a separate phrase.\par
% glossary-block:B01435:end
% glossary-block:B01436:begin
\needspace{3\baselineskip}
\textbf{put down} — \textbf{āweorpan} (v.). Accusative people put down, of + dative point from which. \textbf{sęttan} (v.). Sęttan means put down/place an accusative thing, such as seed. This placement sense is distinct from āweorpan cast down/depose.\par
% glossary-block:B01436:end
% glossary-block:B01437:begin
\needspace{3\baselineskip}
\textbf{put to flight} — \textbf{āflīeman} (v.). put to flight; expel.\par
% glossary-block:B01437:end
\pdfbookmark[1]{Q}{outline-B01438}
% glossary-block:B01438:begin
\needspace{5\baselineskip}
\subsection*{Q}
% glossary-block:B01438:end
% glossary-block:B01439:begin
\textbf{quantity} — \textbf{unrīm} (n.). abundance; uncounted quantity. Supplied quantity with genitive of the substance.\par
% glossary-block:B01439:end
% glossary-block:B01440:begin
\needspace{3\baselineskip}
\textbf{queen} — \textbf{cwēn} (n.). The context distinguishes princess from queen. The selected accusative singular cwēn is endingless.\par
% glossary-block:B01440:end
% glossary-block:B01441:begin
\needspace{3\baselineskip}
\textbf{quell} — \textbf{gestillan} (v.). quell; settle. Accusative disturbance settled; distinguish intransitive become still.\par
% glossary-block:B01441:end
% glossary-block:B01442:begin
\needspace{3\baselineskip}
\textbf{quickly} — \textbf{hrædlīce} (adv.).  \textbf{hraðe} (adv.).\par
% glossary-block:B01442:end
% glossary-block:B01443:begin
\needspace{3\baselineskip}
\textbf{quickness} — \textbf{hrædnes} (n.). speed; quickness.\par
% glossary-block:B01443:end
% glossary-block:B01444:begin
\needspace{3\baselineskip}
\textbf{quiet} — \textbf{stille} (adj.). still; quiet.\par
% glossary-block:B01444:end
% glossary-block:B01445:begin
\needspace{3\baselineskip}
\textbf{quietness} — \textbf{stilnes} (n.). stillness; quietness.\par
% glossary-block:B01445:end
\pdfbookmark[1]{R}{outline-B01446}
% glossary-block:B01446:begin
\needspace{5\baselineskip}
\subsection*{R}
% glossary-block:B01446:end
% glossary-block:B01447:begin
\textbf{raging} — \textbf{angrīslic} (adj.). fierce; raging.\par
% glossary-block:B01447:end
% glossary-block:B01448:begin
\needspace{3\baselineskip}
\textbf{raiment} — \textbf{rēaf} (n.). raiment; clothing. Wǣd denotes a garment, rēaf raiment/clothing, and gegierela a garment or clothing. Preserve any noun expressly supplied by the exercise.\par
% glossary-block:B01448:end
% glossary-block:B01449:begin
\needspace{3\baselineskip}
\textbf{rain} — \textbf{regn} (n.).  \textbf{rīnan} (v.). Supplied impersonal source clause with past rīnde.\par
% glossary-block:B01449:end
% glossary-block:B01450:begin
\needspace{3\baselineskip}
\textbf{raise} — \textbf{āhębban} (v.). raise; lift. \textbf{hębban} (v.). raise; lift. \textbf{ūp-ā-hębban} (v.). lift up. Supplied source form. lift up. Supplied first-plural preterite. raise; exalt. Accusative people raised; the prefixed verb supplies up.\par
% glossary-block:B01450:end
% glossary-block:B01451:begin
\needspace{3\baselineskip}
\textbf{rather} — Comparative swīðor expresses rather/more in the preference construction; the positive adverb swīðe means greatly/very.\par
% glossary-block:B01451:end
% glossary-block:B01452:begin
\needspace{3\baselineskip}
\textbf{ravage} — \textbf{hęrgian} (v.). plunder; ravage.\par
% glossary-block:B01452:end
% glossary-block:B01453:begin
\needspace{3\baselineskip}
\textbf{read} — \textbf{rǣdan} (v.). Accusative text, or absolute with no expressed object.\par
% glossary-block:B01453:end
% glossary-block:B01454:begin
\needspace{3\baselineskip}
\textbf{reading books} — \textbf{bēcrǣding} (n.).\par
% glossary-block:B01454:end
% glossary-block:B01455:begin
\needspace{3\baselineskip}
\textbf{ready} — \textbf{gearu} (adj.).\par
% glossary-block:B01455:end
% glossary-block:B01456:begin
\needspace{3\baselineskip}
\textbf{realize} — \textbf{ongietan} (v.). perceive; realize.\par
% glossary-block:B01456:end
% glossary-block:B01457:begin
\needspace{3\baselineskip}
\textbf{reap} — \textbf{rīpan} (v.).\par
% glossary-block:B01457:end
% glossary-block:B01458:begin
\needspace{3\baselineskip}
\textbf{receive} — \textbf{onfōn} (v.). receive; accept. Dative thing received in this selected frame; the later accusative payment frame is separate. receive; accept. A person received into care in the supplied frame. Selected payment frame has accusative pęning in 28.11, but selected reward frame dative edlēane in 39.6; show both with their task-specific limits. \textbf{underfōn} (v.). receive; take charge of.\par
% glossary-block:B01458:end
% glossary-block:B01459:begin
\needspace{3\baselineskip}
\textbf{recognize} — \textbf{oncnāwan} (v.). recognize; perceive. Supplied past source form. recognize; acknowledge. recognize; know by recognition.\par
% glossary-block:B01459:end
% glossary-block:B01460:begin
\needspace{3\baselineskip}
\textbf{recompense} — \textbf{edlēan} (n.). reward; recompense. \textbf{forgieldan} (v.). pay; repay; requite. Dative recipient, accusative thing repaid in the selected frame.\par
% glossary-block:B01460:end
% glossary-block:B01461:begin
\needspace{3\baselineskip}
\textbf{Redeemer} — \textbf{ālīesend} (n.). Redeemer; one who delivers.\par
% glossary-block:B01461:end
% glossary-block:B01462:begin
\needspace{3\baselineskip}
\textbf{reduce to misery} — \textbf{forierman} (v.). ruin; reduce to misery.\par
% glossary-block:B01462:end
% glossary-block:B01463:begin
\needspace{3\baselineskip}
\textbf{reduce to servitude} — \textbf{ðēowian} (v.). enslave; reduce to servitude.\par
% glossary-block:B01463:end
% glossary-block:B01464:begin
\needspace{3\baselineskip}
\textbf{reed} — \textbf{hrēod} (n.).\par
% glossary-block:B01464:end
% glossary-block:B01465:begin
\needspace{3\baselineskip}
\textbf{reeve} — \textbf{gerēfa} (n.). steward; reeve.\par
% glossary-block:B01465:end
% glossary-block:B01466:begin
\needspace{3\baselineskip}
\textbf{refuse} — \textbf{forwiernan} (v.). deny; refuse. Dative person denied, genitive thing denied. Dative person denied, genitive thing denied. \textbf{willan} (v.). be unwilling; refuse. With ne or a licensed negative contraction; followed by plain infinitive.\par
% glossary-block:B01466:end
% glossary-block:B01467:begin
\needspace{3\baselineskip}
\textbf{refuse to} — \textbf{willan} (v.). ne + willan, or the licensed contracted negative, with a plain infinitive.\par
% glossary-block:B01467:end
% glossary-block:B01468:begin
\needspace{3\baselineskip}
\textbf{regard} — \textbf{gesēon} (v.). regard; take notice of.\par
% glossary-block:B01468:end
% glossary-block:B01469:begin
\needspace{3\baselineskip}
\textbf{region} — \textbf{ðēod} (n.). region; country.\par
% glossary-block:B01469:end
% glossary-block:B01470:begin
\needspace{3\baselineskip}
\textbf{regularly} — \textbf{gelōmlice} (adv.). frequently; regularly.\par
% glossary-block:B01470:end
% glossary-block:B01471:begin
\needspace{3\baselineskip}
\textbf{rejoice} — \textbf{blissian} (v.). Source genitive plural present-participle use, of people rejoicing. Intransitive; mid + dative companion or ofer + accusative cause of joy in the assigned frames.\par
% glossary-block:B01471:end
% glossary-block:B01472:begin
\needspace{3\baselineskip}
\textbf{rejoice in} — \textbf{on} (prep.). in. Source/object of joy in the supplied frame.\par
% glossary-block:B01472:end
% glossary-block:B01473:begin
\needspace{3\baselineskip}
\textbf{rejoice over} — \textbf{ofer} (prep.). over; about. Object of rejoicing, with the assigned accusative frame.\par
% glossary-block:B01473:end
% glossary-block:B01474:begin
\needspace{3\baselineskip}
\textbf{relate to} — \textbf{belimpan} (v.). belong; relate. Supplied participial construction.\par
% glossary-block:B01474:end
% glossary-block:B01475:begin
\needspace{3\baselineskip}
\textbf{release} — \textbf{ālīesan} (v.). deliver; release. Accusative person delivered or released; fram + dative marks separation. \textbf{onlīesan} (v.). Accusative person or thing released.\par
% glossary-block:B01475:end
% glossary-block:B01476:begin
\needspace{3\baselineskip}
\textbf{remain} — \textbf{ðurhwunian} (v.). remain; endure; last. Intransitive continuing of a building; distinguish the following dative-governing wiðstandan. \textbf{wunian} (v.). dwell; remain.\par
% glossary-block:B01476:end
% glossary-block:B01477:begin
\needspace{3\baselineskip}
\textbf{remaining} — \textbf{ōðer} (adj.; ordinal). remaining; the rest. The rest of the forest/earth in the source.\par
% glossary-block:B01477:end
% glossary-block:B01478:begin
\needspace{3\baselineskip}
\textbf{remarkable} — \textbf{sellic} (adj.). strange; remarkable.\par
% glossary-block:B01478:end
% glossary-block:B01479:begin
\needspace{3\baselineskip}
\textbf{remember} — \textbf{gemunan} (v.).  \textbf{geðęncean} (v.). remember; bear in mind.\par
% glossary-block:B01479:end
% glossary-block:B01480:begin
\needspace{3\baselineskip}
\textbf{remove} — \textbf{āfierran} (v.). remove; take away.\par
% glossary-block:B01480:end
% glossary-block:B01481:begin
\needspace{3\baselineskip}
\textbf{render} — \textbf{dōn} (v.). make; render. Accusative object with agreeing object predicate. \textbf{geręccan} (v.). explain; render. \textbf{gewyrcean} (v.). make; render. Accusative object with an agreeing predicate noun/adjective.\par
% glossary-block:B01481:end
% glossary-block:B01482:begin
\needspace{3\baselineskip}
\textbf{renowned} — \textbf{mǣre} (adj.). great; renowned. Distinct from simple physical size where the task requires the renowned sense.\par
% glossary-block:B01482:end
% glossary-block:B01483:begin
\needspace{3\baselineskip}
\textbf{repay} — \textbf{āgyldan} (v.). reward; repay. C06 corrects Cook's āgyltan to strong III āgyldan; supplied contracted present āgylt and dative recipient. \textbf{forgieldan} (v.). pay; repay; requite. Dative recipient, accusative thing repaid in the selected frame.\par
% glossary-block:B01483:end
% glossary-block:B01484:begin
\needspace{3\baselineskip}
\textbf{reply} — \textbf{andswarian} (v.). answer; reply. Dative addressee; the reply may be given as direct speech.\par
% glossary-block:B01484:end
% glossary-block:B01485:begin
\needspace{3\baselineskip}
\textbf{report} — \textbf{hlīsa} (n.). report; account heard.\par
% glossary-block:B01485:end
% glossary-block:B01486:begin
\needspace{3\baselineskip}
\textbf{repose} — \textbf{geręstan} (v.). rest; repose. \textbf{ręst} (n.). rest; repose.\par
% glossary-block:B01486:end
% glossary-block:B01487:begin
\needspace{3\baselineskip}
\textbf{request} — \textbf{biddan} (v.). ask; request. Genitive thing requested; an accusative person may be added. Accusative person beseeched; ðæt + subjunctive gives the requested act. Prayer to God with gebiddan is a different entry. \textbf{giernan} (v.). desire; request; woo. Genitive desired thing/person in the selected frames, including the maiden wooed.\par
% glossary-block:B01487:end
% glossary-block:B01488:begin
\needspace{3\baselineskip}
\textbf{require} — \textbf{beðurfan} (v.). need; require. Genitive thing needed. Genitive thing needed; a relative particle may represent that governed relation.\par
% glossary-block:B01488:end
% glossary-block:B01489:begin
\needspace{3\baselineskip}
\textbf{requite} — \textbf{forgieldan} (v.). pay; repay; requite. Dative recipient, accusative thing repaid in the selected frame.\par
% glossary-block:B01489:end
% glossary-block:B01490:begin
\needspace{3\baselineskip}
\textbf{rescue} — \textbf{hǣlu} (n.). safety; rescue.\par
% glossary-block:B01490:end
% glossary-block:B01491:begin
\needspace{3\baselineskip}
\textbf{resist} — \textbf{wiðstandan} (v.). withstand; resist. Dative of what is resisted.\par
% glossary-block:B01491:end
% glossary-block:B01492:begin
\needspace{3\baselineskip}
\textbf{rest} — \textbf{geręstan} (v.). rest; repose. \textbf{ręst} (n.). rest; repose.\par
% glossary-block:B01492:end
% glossary-block:B01493:begin
\needspace{3\baselineskip}
\textbf{rest from} — Fram + dative expresses from the work or activity in the supplied rest-from construction; the verb of resting must also be supplied.\par
% glossary-block:B01493:end
% glossary-block:B01494:begin
\needspace{3\baselineskip}
\textbf{restore} — \textbf{āgiefan} (v.). give back; restore.\par
% glossary-block:B01494:end
% glossary-block:B01495:begin
\needspace{3\baselineskip}
\textbf{result} — \textbf{weorðan} (v.). come about; result. Supplied resultative expression.\par
% glossary-block:B01495:end
% glossary-block:B01496:begin
\needspace{3\baselineskip}
\textbf{retainer} — \textbf{ðegn} (n.). vassal; retainer; thane.\par
% glossary-block:B01496:end
% glossary-block:B01497:begin
\needspace{3\baselineskip}
\textbf{retreat} — \textbf{dīegelnes} (n.). retreat; seclusion. Supplied plural after of.\par
% glossary-block:B01497:end
% glossary-block:B01498:begin
\needspace{3\baselineskip}
\textbf{return} — \textbf{ęfthweorfan} (v.). return; turn back. \textbf{gewęndan} (v.). return; turn. \textbf{hweorfan} (v.). return; turn back. Supplied source preterite. \textbf{węndan} (v.). turn; return. Supplied source form.\par
% glossary-block:B01498:end
% glossary-block:B01499:begin
\needspace{3\baselineskip}
\textbf{reveal} — \textbf{ætīewan} (v.). show; display; reveal. \textbf{gesweotolian} (v.). reveal; make clear. Supplied source form. make manifest; reveal.\par
% glossary-block:B01499:end
% glossary-block:B01500:begin
\needspace{3\baselineskip}
\textbf{reward} — \textbf{āgyldan} (v.). reward; repay. C06 corrects Cook's āgyltan to strong III āgyldan; supplied contracted present āgylt and dative recipient. \textbf{edlēan} (n.). reward; recompense. \textbf{mēd} (n.).\par
% glossary-block:B01500:end
% glossary-block:B01501:begin
\needspace{3\baselineskip}
\textbf{rich} — \textbf{rīce} (adj.). rich; wealthy. Includes substantival rich people in the assigned frame. \textbf{welig} (adj.). rich; wealthy. \textbf{wlanc} (adj.). proud; wealthy. Contextually the rich man addressed. \textbf{welig} (adj.), \textbf{rīce} (adj.). Welig and rīce can mean rich/wealthy; rīce also occurs substantivally for rich people. Inflect the selected adjective for its role.\par
% glossary-block:B01501:end
% glossary-block:B01502:begin
\needspace{3\baselineskip}
\textbf{riches} — \textbf{wela} (n.). riches; wealth.\par
% glossary-block:B01502:end
% glossary-block:B01503:begin
\needspace{3\baselineskip}
\textbf{rich in} — \textbf{beran} (v.). rich in; productive of. Supplied berende on + dative ore expression; participial adjective. Berende on + dative is the supplied rich-in/productive-of expression. The meaning belongs to the participial expression, not to on alone. \textbf{welig} (adj.). rich in; abounding in. on + dative thing abounded in. welig on + dative of what is abundant.\par
% glossary-block:B01503:end
% glossary-block:B01504:begin
\needspace{3\baselineskip}
\textbf{rich people} — \textbf{rīce} (adj.). rich; wealthy. Includes substantival rich people in the assigned frame.\par
% glossary-block:B01504:end
% glossary-block:B01505:begin
\needspace{3\baselineskip}
\textbf{ride} — \textbf{rīdan} (v.). Intransitive in the introduced frame.\par
% glossary-block:B01505:end
% glossary-block:B01506:begin
\needspace{3\baselineskip}
\textbf{right} — \textbf{riht} (adj.). right; just. \textbf{swīðra} (adj.). Right-hand/side expression, distinct from stronger.\par
% glossary-block:B01506:end
% glossary-block:B01507:begin
\needspace{3\baselineskip}
\textbf{righteous} — \textbf{rihtwīs} (adj.).\par
% glossary-block:B01507:end
% glossary-block:B01508:begin
\needspace{3\baselineskip}
\textbf{righteousness} — \textbf{rihtwīsnes} (n.).\par
% glossary-block:B01508:end
% glossary-block:B01509:begin
\needspace{3\baselineskip}
\textbf{righteous people} — \textbf{rihtwīs} (adj.). Substantival adjective.\par
% glossary-block:B01509:end
% glossary-block:B01510:begin
\needspace{3\baselineskip}
\textbf{right-hand side} — \textbf{swīðra} (adj.). right. Right-hand/side expression, distinct from stronger.\par
% glossary-block:B01510:end
% glossary-block:B01511:begin
\needspace{3\baselineskip}
\textbf{ripe} — \textbf{rīpe} (adj.).\par
% glossary-block:B01511:end
% glossary-block:B01512:begin
\needspace{3\baselineskip}
\textbf{rise} — \textbf{rīsan} (v.). Earlier taught rīsan supplies this command; prefixed ārīsan is separately used later.\par
% glossary-block:B01512:end
% glossary-block:B01513:begin
\needspace{3\baselineskip}
\textbf{rise up} — \textbf{rīsan} (v.). Intransitive rīsan, optionally with ūp; ongēan + accusative marks opposition. \textbf{ūpāstīgan} (v.). rise up; ascend.\par
% glossary-block:B01513:end
% glossary-block:B01514:begin
\needspace{3\baselineskip}
\textbf{rising} — \textbf{ūpgang} (n.). rising; sunrise.\par
% glossary-block:B01514:end
% glossary-block:B01515:begin
\needspace{3\baselineskip}
\textbf{risk} — \textbf{pliht} (n.). peril; risk.\par
% glossary-block:B01515:end
% glossary-block:B01516:begin
\needspace{3\baselineskip}
\textbf{road} — \textbf{weg} (n.). Weg means road, path or way.\par
% glossary-block:B01516:end
% glossary-block:B01517:begin
\needspace{3\baselineskip}
\textbf{Roman} — \textbf{Rōmanisc} (adj.).\par
% glossary-block:B01517:end
% glossary-block:B01518:begin
\needspace{3\baselineskip}
\textbf{Romans} — \textbf{Rōmane} (ethnonym).\par
% glossary-block:B01518:end
% glossary-block:B01519:begin
\needspace{3\baselineskip}
\textbf{Rome} — \textbf{Rome} (proper name).\par
% glossary-block:B01519:end
% glossary-block:B01520:begin
\needspace{3\baselineskip}
\textbf{roof} — \textbf{hrōf} (n.).  \textbf{ðæce} (n.).\par
% glossary-block:B01520:end
% glossary-block:B01521:begin
\needspace{3\baselineskip}
\textbf{room} — \textbf{būr} (n.). chamber; room.\par
% glossary-block:B01521:end
% glossary-block:B01522:begin
\needspace{3\baselineskip}
\textbf{root} — \textbf{wyrttruma} (n.).\par
% glossary-block:B01522:end
% glossary-block:B01523:begin
\needspace{3\baselineskip}
\textbf{round about} — \textbf{ymbhwyrft} (n.). circuit. Whole on ymbhwyrfte expression means all around; supplied dependent landmark. Supplied on ymbhwyrfte + genitive landmark in this exercise.\par
% glossary-block:B01523:end
% glossary-block:B01524:begin
\needspace{3\baselineskip}
\textbf{row} — \textbf{rōwan} (v.). Intransitive; the assigned route is expressed with ofer.\par
% glossary-block:B01524:end
% glossary-block:B01525:begin
\needspace{3\baselineskip}
\textbf{royal seat} — \textbf{cynesętl} (n.). royal seat; throne.\par
% glossary-block:B01525:end
% glossary-block:B01526:begin
\needspace{3\baselineskip}
\textbf{ruin} — \textbf{forierman} (v.). ruin; reduce to misery. \textbf{lor} (n.). destruction; ruin.\par
% glossary-block:B01526:end
% glossary-block:B01527:begin
\needspace{3\baselineskip}
\textbf{rule} — \textbf{geweald} (n.). control; rule; dominion. Ordinary singular; on gewealde is the taught control phrase. \textbf{rīcsian} (v.). rule; have dominion. Used absolutely in the taught frame; no new object government. rīcsian can be used absolutely; English have/bear need no independent translation. \textbf{rīce} (n.). Genitive rīces after geswīcan in ceased to bear rule. Do not insert a literal bear verb.\par
% glossary-block:B01527:end
% glossary-block:B01528:begin
\needspace{3\baselineskip}
\textbf{rule over} — \textbf{gīeman} (v.). Genitive of what is ruled; no additional over preposition. gīeman governs genitive of what is ruled; English over supplies no additional preposition.\par
% glossary-block:B01528:end
% glossary-block:B01529:begin
\needspace{3\baselineskip}
\textbf{Ruler} — \textbf{wealdend} (n.).\par
% glossary-block:B01529:end
% glossary-block:B01530:begin
\needspace{3\baselineskip}
\textbf{run} — \textbf{iernan} (v.). Intransitive; geond + accusative path. A perception construction may make its actor accusative.\par
% glossary-block:B01530:end
\pdfbookmark[1]{S}{outline-B01531}
% glossary-block:B01531:begin
\needspace{5\baselineskip}
\subsection*{S}
% glossary-block:B01531:end
% glossary-block:B01532:begin
\textbf{sack} — \textbf{sacc} (n.).\par
% glossary-block:B01532:end
% glossary-block:B01533:begin
\needspace{3\baselineskip}
\textbf{safety} — \textbf{hǣlu} (n.). safety; rescue.\par
% glossary-block:B01533:end
% glossary-block:B01534:begin
\needspace{3\baselineskip}
\textbf{Saint} — \textbf{Sanctus} (foreign adjective). Saint; holy. Latin word in the source title.\par
% glossary-block:B01534:end
% glossary-block:B01535:begin
\needspace{3\baselineskip}
\textbf{saint} — \textbf{hālga} (n.). saint; holy one.\par
% glossary-block:B01535:end
% glossary-block:B01536:begin
\needspace{3\baselineskip}
\textbf{saints} — \textbf{hālig} (adj.). holy ones; saints. Substantival adjective in the source.\par
% glossary-block:B01536:end
% glossary-block:B01537:begin
\needspace{3\baselineskip}
\textbf{sake} — \textbf{intinga} (n.). cause; sake (in for … intingan). Within supplied for intingan expressions.\par
% glossary-block:B01537:end
% glossary-block:B01538:begin
\needspace{3\baselineskip}
\textbf{salute} — \textbf{grētan} (v.). greet; take leave of. Accusative person/place; greeting and farewell depend on context. Greeting with accusative person; the same verb also has farewell/leave-taking uses.\par
% glossary-block:B01538:end
% glossary-block:B01539:begin
\needspace{3\baselineskip}
\textbf{salvation} — \textbf{hǣlu} (n.).\par
% glossary-block:B01539:end
% glossary-block:B01540:begin
\needspace{3\baselineskip}
\textbf{same} — \textbf{ilca} (demonstrative adjective).\par
% glossary-block:B01540:end
% glossary-block:B01541:begin
\needspace{3\baselineskip}
\textbf{save} — \textbf{hǣlan} (v.).  \textbf{nęrian} (v.).\par
% glossary-block:B01541:end
% glossary-block:B01542:begin
\needspace{3\baselineskip}
\textbf{Savior} — \textbf{Hǣlend} (n.). Saviour; Jesus. Proper title.\par
% glossary-block:B01542:end
% glossary-block:B01543:begin
\needspace{3\baselineskip}
\textbf{Saviour} — \textbf{Hǣlend} (n.). Saviour; Jesus. Proper title.\par
% glossary-block:B01543:end
% glossary-block:B01544:begin
\needspace{3\baselineskip}
\textbf{say} — \textbf{cweðan} (v.). Source genitive plural present participle, of people saying. Supplied preterite said. \textbf{gecweðan} (v.). speak; say. \textbf{sęcgan} (v.). tell; say. Dative addressee and a reported statement or question as content. Supplied source preterite. say; tell. Person addressed is dative; direct or reported content is separate. \textbf{sprecan} (v.). This task explicitly selects sprecan for say; tō marks the addressee.\par
% glossary-block:B01544:end
% glossary-block:B01545:begin
\needspace{3\baselineskip}
\textbf{saying} — \textbf{sprǣc} (n.). saying; speech.\par
% glossary-block:B01545:end
% glossary-block:B01546:begin
\needspace{3\baselineskip}
\textbf{say to} — In supplied cwæð him tō, cwæð expresses said and postposed tō marks the addressee. Tō alone does not express say. Tō marks the dative addressee with the supplied verb of speaking; the verb supplies say, not tō alone.\par
% glossary-block:B01546:end
% glossary-block:B01547:begin
\needspace{3\baselineskip}
\textbf{scorch} — \textbf{forbærnan} (v.). parch; scorch.\par
% glossary-block:B01547:end
% glossary-block:B01548:begin
\needspace{3\baselineskip}
\textbf{Scots} — \textbf{Scottas} (ethnonym).\par
% glossary-block:B01548:end
% glossary-block:B01549:begin
\needspace{3\baselineskip}
\textbf{sea} — \textbf{sǣ} (n.).\par
% glossary-block:B01549:end
% glossary-block:B01550:begin
\needspace{3\baselineskip}
\textbf{seal} — \textbf{seolh} (n.).\par
% glossary-block:B01550:end
% glossary-block:B01551:begin
\needspace{3\baselineskip}
\textbf{sea-shore} — \textbf{strand} (n.). strand; sea-shore.\par
% glossary-block:B01551:end
% glossary-block:B01552:begin
\needspace{3\baselineskip}
\textbf{season} — \textbf{tīd} (n.). season; time.\par
% glossary-block:B01552:end
% glossary-block:B01553:begin
\needspace{3\baselineskip}
\textbf{seat} — \textbf{setl} (n.). seat; throne.\par
% glossary-block:B01553:end
% glossary-block:B01554:begin
\needspace{3\baselineskip}
\textbf{seclusion} — \textbf{dīegelnes} (n.). retreat; seclusion. Supplied plural after of.\par
% glossary-block:B01554:end
% glossary-block:B01555:begin
\needspace{3\baselineskip}
\textbf{second} — \textbf{ōðer} (adj.; ordinal).\par
% glossary-block:B01555:end
% glossary-block:B01556:begin
\needspace{3\baselineskip}
\textbf{second time} — Ōðre sīðe is the supplied instrumental singular phrase a second time, without a separate indefinite article.\par
% glossary-block:B01556:end
% glossary-block:B01557:begin
\needspace{3\baselineskip}
\textbf{see} — \textbf{gesēon} (v.). Absolute use, with no expressed object. Accusative perceived actor and plain infinitive. With a content clause stating the circumstance seen. With the assigned clause describing how a state of affairs stands. \textbf{sēon} (v.).\par
% glossary-block:B01557:end
% glossary-block:B01558:begin
\needspace{3\baselineskip}
\textbf{seed} — \textbf{sǣd} (n.). seed; descendants. Collective descendants, distinct from literal seed.\par
% glossary-block:B01558:end
% glossary-block:B01559:begin
\needspace{3\baselineskip}
\textbf{seedtime} — \textbf{sǣdtīma} (n.). seedtime; sowing season.\par
% glossary-block:B01559:end
% glossary-block:B01560:begin
\needspace{3\baselineskip}
\textbf{see how} — \textbf{gesēon} (v.). see. With the assigned clause describing how a state of affairs stands.\par
% glossary-block:B01560:end
% glossary-block:B01561:begin
\needspace{3\baselineskip}
\textbf{seek} — \textbf{sēcan} (v.). Accusative person or thing sought. sēcan with accusative sought thing, no extra word for English for/out.\par
% glossary-block:B01561:end
% glossary-block:B01562:begin
\needspace{3\baselineskip}
\textbf{seek in marriage} — \textbf{giernan} (v.). desire; request; woo. Genitive desired thing/person in the selected frames, including the maiden wooed. Genitive person wooed; keep the marriage-seeking sense distinct from a general want.\par
% glossary-block:B01562:end
% glossary-block:B01563:begin
\needspace{3\baselineskip}
\textbf{seek out} — \textbf{sēcan} (v.). sēcan with accusative sought thing, no extra word for English for/out.\par
% glossary-block:B01563:end
% glossary-block:B01564:begin
\needspace{3\baselineskip}
\textbf{seem} — \textbf{ðyncan} (v.). Nominative thing that seems, dative experiencer and agreeing predicate. Dative experiencer; the thing that seems is nominative. Supplied past source form. Wesan is rendered seem only within the supplied impersonal expression it seemed to the king. Retain its full dative-experiencer construction; ordinary wesan means be.\par
% glossary-block:B01564:end
% glossary-block:B01565:begin
\needspace{3\baselineskip}
\textbf{seem good to} — \textbf{ðyncan} (v.). seem. Nominative thing that seems, dative experiencer and agreeing predicate. Subject thing nominative, dative experiencer, agreeing predicate gōd.\par
% glossary-block:B01565:end
% glossary-block:B01566:begin
\needspace{3\baselineskip}
\textbf{see someone do something} — \textbf{gesēon} (v.). see. Accusative perceived actor and plain infinitive.\par
% glossary-block:B01566:end
% glossary-block:B01567:begin
\needspace{3\baselineskip}
\textbf{see that} — \textbf{gesēon} (v.). see. With a content clause stating the circumstance seen.\par
% glossary-block:B01567:end
% glossary-block:B01568:begin
\needspace{3\baselineskip}
\textbf{seize} — \textbf{fōn} (v.). seize; catch. Accusative patient in the selected frame. Accusative person overtaken in these tasks; the source's bracketed caught is the route.\par
% glossary-block:B01568:end
% glossary-block:B01569:begin
\needspace{3\baselineskip}
\textbf{select} — \textbf{cēosan} (v.). Accusative person selected.\par
% glossary-block:B01569:end
% glossary-block:B01570:begin
\needspace{3\baselineskip}
\textbf{self} — \textbf{self} (adj.). self; the same person or thing. Agreeing emphatic/reflexive self; ordinary reflexive personal pronouns do not always require it.\par
% glossary-block:B01570:end
% glossary-block:B01571:begin
\needspace{3\baselineskip}
\textbf{sell} — \textbf{cīepan} (v.). Accusative goods sold. \textbf{sęllan} (v.). Accusative thing sold. Sęllan means sell in the selected trading contexts; the thing sold is accusative.\par
% glossary-block:B01571:end
% glossary-block:B01572:begin
\needspace{3\baselineskip}
\textbf{sell into} — \textbf{on} (prep.). into. Destination of sale in the supplied expression.\par
% glossary-block:B01572:end
% glossary-block:B01573:begin
\needspace{3\baselineskip}
\textbf{send} — \textbf{āsęndan} (v.).  \textbf{onsęndan} (v.).  \textbf{sęndan} (v.). Accusative person or thing sent, dative recipient.\par
% glossary-block:B01573:end
% glossary-block:B01574:begin
\needspace{3\baselineskip}
\textbf{sense} — \textbf{andgiet} (n.). meaning; sense.\par
% glossary-block:B01574:end
% glossary-block:B01575:begin
\needspace{3\baselineskip}
\textbf{sense for sense} — \textbf{of} (prep.). Literal correspondence in the supplied translation expression.\par
% glossary-block:B01575:end
% glossary-block:B01576:begin
\needspace{3\baselineskip}
\textbf{sent back} — \textbf{gesęndan} (v.).\par
% glossary-block:B01576:end
% glossary-block:B01577:begin
\needspace{3\baselineskip}
\textbf{separate} — \textbf{gedǣlan} (v.). divide; separate. Preserve the assigned separation frame. \textbf{tōdǣlen / tōdǣlden (source forms)} (verb; supplied inflected forms). separate; distinguish. Source forms tōdǣlen/tōdǣlden are retained; no silently invented citation spelling. \textbf{tōsceadan} (v.). separate; divide. \textbf{tōtwǣman} (v.). divide; separate. Waters divided from waters in the supplied frame. Accusative beings separated; fram + dative other party.\par
% glossary-block:B01577:end
% glossary-block:B01578:begin
\needspace{3\baselineskip}
\textbf{serpent} — \textbf{nædre} (n.).\par
% glossary-block:B01578:end
% glossary-block:B01579:begin
\needspace{3\baselineskip}
\textbf{servant} — \textbf{cniht} (n.).  \textbf{ðēow} (n.). servant; enslaved person. \textbf{ðēowa} (n.).  \textbf{cniht} (n.), \textbf{ðēowa} (n.), \textbf{ðēow} (n.). Cniht, weak ðēowa and strong ðēow can denote a servant. Use the particular noun required by the exercise and its own declension.\par
% glossary-block:B01579:end
% glossary-block:B01580:begin
\needspace{3\baselineskip}
\textbf{servants} — \textbf{ðēowa manna} (noun phrase). enslaved men; servants. Whole supplied phrase ðēowa manna.\par
% glossary-block:B01580:end
% glossary-block:B01581:begin
\needspace{3\baselineskip}
\textbf{serve} — \textbf{gehīersumian} (v.). serve; obey. Supplied source form and government. \textbf{ðēnian} (v.). serve; attend. Supplied preterite source form.\par
% glossary-block:B01581:end
% glossary-block:B01582:begin
\needspace{3\baselineskip}
\textbf{service} — \textbf{ðēnung} (n.). service; attendance. \textbf{ðēowdōm} (n.).\par
% glossary-block:B01582:end
% glossary-block:B01583:begin
\needspace{3\baselineskip}
\textbf{set} — \textbf{āsęttan} (v.). set; place. \textbf{gelōgian} (v.). place; set. \textbf{gesęttan} (v.). set; place. Supplied gesętte. set; place; establish. \textbf{lęcgean} (v.). lay; place. Supplied source form. \textbf{sęttan} (v.). set; place. Place an accusative object; destination is a separate phrase.\par
% glossary-block:B01583:end
% glossary-block:B01584:begin
\needspace{3\baselineskip}
\textbf{set aside} — \textbf{forlǣtan} (v.). Supplied source sense.\par
% glossary-block:B01584:end
% glossary-block:B01585:begin
\needspace{3\baselineskip}
\textbf{set as lord} — \textbf{gesęttan} (v.). appoint as lord. Supplied frame with person appointed, tō + dative role, and dative people governed.\par
% glossary-block:B01585:end
% glossary-block:B01586:begin
\needspace{3\baselineskip}
\textbf{set forth} — \textbf{lęcgean} (v.). lęcgean + accusative thing + forð.\par
% glossary-block:B01586:end
% glossary-block:B01587:begin
\needspace{3\baselineskip}
\textbf{set in motion} — \textbf{astyrian} (v.). move; set in motion. Supplied source sense.\par
% glossary-block:B01587:end
% glossary-block:B01588:begin
\needspace{3\baselineskip}
\textbf{settle} — \textbf{gestillan} (v.). quell; settle. Accusative disturbance settled; distinguish intransitive become still.\par
% glossary-block:B01588:end
% glossary-block:B01589:begin
\needspace{3\baselineskip}
\textbf{seven} — \textbf{seofon} (cardinal numeral). Cardinal 7.\par
% glossary-block:B01589:end
% glossary-block:B01590:begin
\needspace{3\baselineskip}
\textbf{seventh} — \textbf{seofoða} (ordinal).\par
% glossary-block:B01590:end
% glossary-block:B01591:begin
\needspace{3\baselineskip}
\textbf{shake} — \textbf{geswęngan} (v.). toss; shake. Supplied participial source use. \textbf{gewǣgan} (v.). buffet; shake. Supplied participial source use.\par
% glossary-block:B01591:end
% glossary-block:B01592:begin
\needspace{3\baselineskip}
\textbf{shall} — Supplied gīet sculon fīf means that five are yet to come, with be expressly omitted. This bounded ellipsis is not a general rule that English shall always takes sculan.\par
% glossary-block:B01592:end
% glossary-block:B01593:begin
\needspace{3\baselineskip}
\textbf{shall (future reference)} — Present morphology can carry future reference; bēo is used in its bounded expected-state frame. Do not add willan unless volition or a stated source target requires it.\par
% glossary-block:B01593:end
% glossary-block:B01594:begin
\needspace{3\baselineskip}
\textbf{shape} — \textbf{scieppan} (v.). create; shape.\par
% glossary-block:B01594:end
% glossary-block:B01595:begin
\needspace{3\baselineskip}
\textbf{shared} — \textbf{gemǣne} (adj.). common; shared.\par
% glossary-block:B01595:end
% glossary-block:B01596:begin
\needspace{3\baselineskip}
\textbf{she} — Hēo is feminine singular nominative. Third-person pronouns follow the grammatical gender of their referent. Grammatical gender and antecedent determine reference; it may be represented by masculine/feminine forms for a gendered inanimate noun.\par
% glossary-block:B01596:end
% glossary-block:B01597:begin
\needspace{3\baselineskip}
\textbf{Sheba} — \textbf{Saba} (proper name).\par
% glossary-block:B01597:end
% glossary-block:B01598:begin
\needspace{3\baselineskip}
\textbf{sheep} — \textbf{scēap} (n.).\par
% glossary-block:B01598:end
% glossary-block:B01599:begin
\needspace{3\baselineskip}
\textbf{shepherd} — \textbf{scēaphierde} (n.).\par
% glossary-block:B01599:end
% glossary-block:B01600:begin
\needspace{3\baselineskip}
\textbf{shield} — \textbf{gescieldan} (v.). shield; protect.\par
% glossary-block:B01600:end
% glossary-block:B01601:begin
\needspace{3\baselineskip}
\textbf{shine} — \textbf{scīnan} (v.).\par
% glossary-block:B01601:end
% glossary-block:B01602:begin
\needspace{3\baselineskip}
\textbf{shining} — \textbf{scīnan} (v.).\par
% glossary-block:B01602:end
% glossary-block:B01603:begin
\needspace{3\baselineskip}
\textbf{ship} — \textbf{scip} (n.).\par
% glossary-block:B01603:end
% glossary-block:B01604:begin
\needspace{3\baselineskip}
\textbf{shipwreck} — \textbf{forlidennes} (n.).\par
% glossary-block:B01604:end
% glossary-block:B01605:begin
\needspace{3\baselineskip}
\textbf{shipwrecked} — \textbf{forliden} (adj.).\par
% glossary-block:B01605:end
% glossary-block:B01606:begin
\needspace{3\baselineskip}
\textbf{shoemaker} — \textbf{scēowyrhta} (n.).\par
% glossary-block:B01606:end
% glossary-block:B01607:begin
\needspace{3\baselineskip}
\textbf{shore} — \textbf{strand} (n.). Strand means shore or sea-shore in the assigned noun and translation tasks.\par
% glossary-block:B01607:end
% glossary-block:B01608:begin
\needspace{3\baselineskip}
\textbf{short} — \textbf{medmicel} (adj.). short; brief.\par
% glossary-block:B01608:end
% glossary-block:B01609:begin
\needspace{3\baselineskip}
\textbf{should} — In the supplied interpretive question, sculan expresses what events should/were to mean. Retain that reported or prospective interpretation rather than substituting the ordinary obligation sense. Select intention, obligation, reported prospect or counterfactual context; English spelling alone does not determine tense, mood or an OE modal.\par
% glossary-block:B01609:end
% glossary-block:B01610:begin
\needspace{3\baselineskip}
\textbf{show} — \textbf{ætīewan} (v.). show; display; reveal. \textbf{ge-cȳðan} (v.).  \textbf{getǣcan} (v.). show; direct. Supplied source form. \textbf{tǣcan} (v.). teach; show.\par
% glossary-block:B01610:end
% glossary-block:B01611:begin
\needspace{3\baselineskip}
\textbf{sick} — \textbf{untrum} (adj.).\par
% glossary-block:B01611:end
% glossary-block:B01612:begin
\needspace{3\baselineskip}
\textbf{sick people} — \textbf{untrum} (adj.). Substantival adjective.\par
% glossary-block:B01612:end
% glossary-block:B01613:begin
\needspace{3\baselineskip}
\textbf{side} — \textbf{hand} (n.). hand; side. Side in right-/left-hand expressions.\par
% glossary-block:B01613:end
% glossary-block:B01614:begin
\needspace{3\baselineskip}
\textbf{sight} — \textbf{gesihð} (n.). sight; gaze. Supplied accusative singular.\par
% glossary-block:B01614:end
% glossary-block:B01615:begin
\needspace{3\baselineskip}
\textbf{sign} — \textbf{tācen} (n.).\par
% glossary-block:B01615:end
% glossary-block:B01616:begin
\needspace{3\baselineskip}
\textbf{signal} — The supplied expression slōg tācen means made a signal; tācen is its object. Slēan alone has the separate strike sense.\par
% glossary-block:B01616:end
% glossary-block:B01617:begin
\needspace{3\baselineskip}
\textbf{silence} — \textbf{swīge} (n.).\par
% glossary-block:B01617:end
% glossary-block:B01618:begin
\needspace{3\baselineskip}
\textbf{silk} — \textbf{sīde} (n.).\par
% glossary-block:B01618:end
% glossary-block:B01619:begin
\needspace{3\baselineskip}
\textbf{silver} — \textbf{seolfor} (n.). Metal noun. \textbf{seolfren} (adj.). silver; made of silver. Material adjective, distinct from the noun seolfor.\par
% glossary-block:B01619:end
% glossary-block:B01620:begin
\needspace{3\baselineskip}
\textbf{similarly} — \textbf{gelīce} (adv.). likewise; similarly.\par
% glossary-block:B01620:end
% glossary-block:B01621:begin
\needspace{3\baselineskip}
\textbf{sin} — \textbf{synn} (n.).\par
% glossary-block:B01621:end
% glossary-block:B01622:begin
\needspace{3\baselineskip}
\textbf{since} — \textbf{siððan} (adv.; conj.). since; after; afterwards. Conjunction since/after versus adverb afterwards.\par
% glossary-block:B01622:end
% glossary-block:B01623:begin
\needspace{3\baselineskip}
\textbf{sinful} — \textbf{synfull} (adj.).\par
% glossary-block:B01623:end
% glossary-block:B01624:begin
\needspace{3\baselineskip}
\textbf{sing} — \textbf{singan} (v.). Intransitive in the introduced frame. Source genitive plural present participle, of people singing.\par
% glossary-block:B01624:end
% glossary-block:B01625:begin
\needspace{3\baselineskip}
\textbf{single} — \textbf{ān} (numeral; adj.). one; single. Cardinal one/single and the agreeing adjective alone are distinct uses of ān. Preserve the supplied form where the exercise provides one.\par
% glossary-block:B01625:end
% glossary-block:B01626:begin
\needspace{3\baselineskip}
\textbf{sir} — \textbf{lēof} (adj.). Supplied respectful address lēof. Supplied respectful lēof address; the older gloss dear alone does not give this lookup.\par
% glossary-block:B01626:end
% glossary-block:B01627:begin
\needspace{3\baselineskip}
\textbf{sister} — \textbf{sweostor} (n.).\par
% glossary-block:B01627:end
% glossary-block:B01628:begin
\needspace{3\baselineskip}
\textbf{sit} — \textbf{sittan} (v.).\par
% glossary-block:B01628:end
% glossary-block:B01629:begin
\needspace{3\baselineskip}
\textbf{sit around} — \textbf{ymbsittan} (v.). Supplied participle referring to people sitting around.\par
% glossary-block:B01629:end
% glossary-block:B01630:begin
\needspace{3\baselineskip}
\textbf{situated} — \textbf{gesęttan} (v.). placed; situated. Supplied participial complement of understood us; not the identical-looking finite preterite.\par
% glossary-block:B01630:end
% glossary-block:B01631:begin
\needspace{3\baselineskip}
\textbf{six} — \textbf{siex} (cardinal numeral). Cardinal 6.\par
% glossary-block:B01631:end
% glossary-block:B01632:begin
\needspace{3\baselineskip}
\textbf{sixth} — \textbf{siexta} (ordinal numeral).\par
% glossary-block:B01632:end
% glossary-block:B01633:begin
\needspace{3\baselineskip}
\textbf{sixty} — \textbf{siextig} (cardinal numeral). Within the supplied 365-days expression.\par
% glossary-block:B01633:end
% glossary-block:B01634:begin
\needspace{3\baselineskip}
\textbf{skill} — \textbf{cræft} (n.). skill; accomplishment. \textbf{snelnes} (n.). skill; agility.\par
% glossary-block:B01634:end
% glossary-block:B01635:begin
\needspace{3\baselineskip}
\textbf{skilled} — \textbf{gelǣred} (adj.). educated; trained; skilled.\par
% glossary-block:B01635:end
% glossary-block:B01636:begin
\needspace{3\baselineskip}
\textbf{sky} — \textbf{heofon} (n.). heaven; sky. \textbf{heofone} (n.). heaven; sky. Weak feminine lexical alternative in the source.\par
% glossary-block:B01636:end
% glossary-block:B01637:begin
\needspace{3\baselineskip}
\textbf{slaughter} — \textbf{wæl} (n.).\par
% glossary-block:B01637:end
% glossary-block:B01638:begin
\needspace{3\baselineskip}
\textbf{slay} — \textbf{ofslēan} (v.). slay; kill. \textbf{ongeslēan} (v.). inflict slaughter; slay. \textbf{slēan} (v.). slay; kill. Accusative person killed.\par
% glossary-block:B01638:end
% glossary-block:B01639:begin
\needspace{3\baselineskip}
\textbf{small} — \textbf{lȳtel} (adj.). little; small.\par
% glossary-block:B01639:end
% glossary-block:B01640:begin
\needspace{3\baselineskip}
\textbf{smaller} — \textbf{lǣssa} (adj.).\par
% glossary-block:B01640:end
% glossary-block:B01641:begin
\needspace{3\baselineskip}
\textbf{smallest} — \textbf{lǣst} (adj.).\par
% glossary-block:B01641:end
% glossary-block:B01642:begin
\needspace{3\baselineskip}
\textbf{smite} — \textbf{geslēan} (v.). strike; smite. Supplied past/participial source form. \textbf{slēan} (v.). strike; smite. strike; smite. Objectless striking in the assigned hammer sentences. \textbf{slēan} (v.), \textbf{bēatan} (v.). Slēan means strike/smite and can be objectless in the hammer sentences. Bēatan means beat/strike and takes an accusative object; the explicit bēatan target must be retained.\par
% glossary-block:B01642:end
% glossary-block:B01643:begin
\needspace{3\baselineskip}
\textbf{smithy} — \textbf{smiððe} (n.).\par
% glossary-block:B01643:end
% glossary-block:B01644:begin
\needspace{3\baselineskip}
\textbf{smooth} — \textbf{smylte} (adj.). calm; smooth.\par
% glossary-block:B01644:end
% glossary-block:B01645:begin
\needspace{3\baselineskip}
\textbf{snake} — \textbf{næddre} (n.). snake; adder. Supplied weak feminine plural.\par
% glossary-block:B01645:end
% glossary-block:B01646:begin
\needspace{3\baselineskip}
\textbf{snow} — \textbf{snāw} (n.).\par
% glossary-block:B01646:end
% glossary-block:B01647:begin
\needspace{3\baselineskip}
\textbf{so} — Tō ðan … ðæt is the supplied degree-and-result construction so/to that degree … that. Keep the whole construction. \textbf{swā} (adverb; conjunction). so; thus.\par
% glossary-block:B01647:end
% glossary-block:B01648:begin
\needspace{3\baselineskip}
\textbf{sojourn} — \textbf{wician} (v.). sojourn; lodge. Intransitive; mid + dative person with whom one lodges.\par
% glossary-block:B01648:end
% glossary-block:B01649:begin
\needspace{3\baselineskip}
\textbf{soldier} — \textbf{cęmpa} (n.).\par
% glossary-block:B01649:end
% glossary-block:B01650:begin
\needspace{3\baselineskip}
\textbf{solemnly charge} — \textbf{hālsian} (v.). conjure; solemnly charge. Accusative person; ðurh + accusative oath standard; a requested ðæt clause may follow.\par
% glossary-block:B01650:end
% glossary-block:B01651:begin
\needspace{3\baselineskip}
\textbf{Solomon} — \textbf{Salomon} (proper name).\par
% glossary-block:B01651:end
% glossary-block:B01652:begin
\needspace{3\baselineskip}
\textbf{some} — \textbf{hwilc} (interrogative/indefinite adj. or pron.). any; some. Indefinite adjective, distinct from an interrogative. \textbf{sum} (indefinite adjective/pronoun). Members selected from a group; genitive whole. a certain; some. \textbf{sum} (indefinite adjective/pronoun), \textbf{hwilc} (interrogative/indefinite adj. or pron.). Sum means some/a certain; partitive sum selects members from a genitive whole. Indefinite hwilc has the separate supplied some/any use.\par
% glossary-block:B01652:end
% glossary-block:B01653:begin
\needspace{3\baselineskip}
\textbf{some of} — Partitive sum selects some members from a genitive whole; English of is expressed by that genitive, without an extra preposition.\par
% glossary-block:B01653:end
% glossary-block:B01654:begin
\needspace{3\baselineskip}
\textbf{something} — \textbf{hwæthwega} (indefinite pronoun/substantive). Supplied local indefinite form; no unrestricted new paradigm. \textbf{hwæthwugu} (indefinite). Supplied hwæthwugu swilces.\par
% glossary-block:B01654:end
% glossary-block:B01655:begin
\needspace{3\baselineskip}
\textbf{something of that sort} — \textbf{hwæthwugu} (indefinite). something. Supplied hwæthwugu swilces.\par
% glossary-block:B01655:end
% glossary-block:B01656:begin
\needspace{3\baselineskip}
\textbf{something of this sort} — Hwæthwega swilces and the supplied variant hwæthwugu swilces mean something of this sort. Swilces is the genitive expression of quality; preserve the supplied indefinite variant.\par
% glossary-block:B01656:end
% glossary-block:B01657:begin
\needspace{3\baselineskip}
\textbf{sometimes} — \textbf{hwīlum} (adv.).\par
% glossary-block:B01657:end
% glossary-block:B01658:begin
\needspace{3\baselineskip}
\textbf{so much} — Swīðe expresses much as degree, as in love very much; swā swīðe expresses so much. This differs from micel expressing quantity.\par
% glossary-block:B01658:end
% glossary-block:B01659:begin
\needspace{3\baselineskip}
\textbf{son} — \textbf{bearn} (n.). son; child. \textbf{sunu} (n.).\par
% glossary-block:B01659:end
% glossary-block:B01660:begin
\needspace{3\baselineskip}
\textbf{song} — \textbf{sang} (n.).\par
% glossary-block:B01660:end
% glossary-block:B01661:begin
\needspace{3\baselineskip}
\textbf{song of joy} — \textbf{blissesang} (n.).\par
% glossary-block:B01661:end
% glossary-block:B01662:begin
\needspace{3\baselineskip}
\textbf{soon} — \textbf{sōna} (adv.). sōna carries soon in this future-reference prompt.\par
% glossary-block:B01662:end
% glossary-block:B01663:begin
\needspace{3\baselineskip}
\textbf{sore} — \textbf{wund} (n.). wound; sore.\par
% glossary-block:B01663:end
% glossary-block:B01664:begin
\needspace{3\baselineskip}
\textbf{sorrow} — \textbf{sorg} (n.). distress; anxiety; trouble; sorrow.\par
% glossary-block:B01664:end
% glossary-block:B01665:begin
\needspace{3\baselineskip}
\textbf{so that} — Tō ðon ðætte is the complete supplied result connection so that. The demonstrative component alone is not a conjunction meaning so that. \textbf{ðæt} (conj.). that; so that. Supplied expanded spelling; distinguish content from result by the source clause. so that; in order that. Purpose conjunction with the assigned subjunctive construction. that; so that. Asserted result, distinct from intended purpose. Purpose normally uses ðæt + subjunctive in the assigned frames. Do not add a permission modal merely for English may; retain mugan when the prompt explicitly requests it. Tō ðon ðætte is the complete supplied result connection so that; tō alone does not supply that connection.\par
% glossary-block:B01665:end
% glossary-block:B01666:begin
\needspace{3\baselineskip}
\textbf{so that not} — \textbf{ðȳ-lǣs-ðe} (conj.). lest; so that not. ðȳ-lǣs-ðe with the assigned subjunctive negative-purpose clause.\par
% glossary-block:B01666:end
% glossary-block:B01667:begin
\needspace{3\baselineskip}
\textbf{soul} — \textbf{sāwol} (n.).\par
% glossary-block:B01667:end
% glossary-block:B01668:begin
\needspace{3\baselineskip}
\textbf{sound} — \textbf{swēg} (n.). music; sound.\par
% glossary-block:B01668:end
% glossary-block:B01669:begin
\needspace{3\baselineskip}
\textbf{south} — The supplied whole phrase be sūðan means in the south. Be alone does not mean south. \textbf{sūðan} (adv.). in the south; from the south. Retain the supplied spatial expression.\par
% glossary-block:B01669:end
% glossary-block:B01670:begin
\needspace{3\baselineskip}
\textbf{southeastern part} — \textbf{ēastsūðdǣl} (n.).\par
% glossary-block:B01670:end
% glossary-block:B01671:begin
\needspace{3\baselineskip}
\textbf{southern} — \textbf{sūðerne} (adj.).\par
% glossary-block:B01671:end
% glossary-block:B01672:begin
\needspace{3\baselineskip}
\textbf{southwestern} — \textbf{sūðwesterne} (adj.).\par
% glossary-block:B01672:end
% glossary-block:B01673:begin
\needspace{3\baselineskip}
\textbf{sowing season} — \textbf{sǣdtīma} (n.). seedtime; sowing season.\par
% glossary-block:B01673:end
% glossary-block:B01674:begin
\needspace{3\baselineskip}
\textbf{spark} — \textbf{fȳrspearca} (n.).\par
% glossary-block:B01674:end
% glossary-block:B01675:begin
\needspace{3\baselineskip}
\textbf{speak} — \textbf{gecweðan} (v.). speak; say. \textbf{sprecan} (v.). speak with. mid + dative interlocutor in the taught production frame. This task explicitly selects sprecan for say; tō marks the addressee.\par
% glossary-block:B01675:end
% glossary-block:B01676:begin
\needspace{3\baselineskip}
\textbf{speak with} — \textbf{mid} (prep.; adv. in fixed expressions). with. Dative conversational partner. \textbf{sprecan} (v.). mid + dative interlocutor in the taught production frame.\par
% glossary-block:B01676:end
% glossary-block:B01677:begin
\needspace{3\baselineskip}
\textbf{speech} — \textbf{sprǣc} (n.). saying; speech. \textbf{word} (n.). speech; utterance.\par
% glossary-block:B01677:end
% glossary-block:B01678:begin
\needspace{3\baselineskip}
\textbf{speed} — \textbf{hrædnes} (n.). speed; quickness.\par
% glossary-block:B01678:end
% glossary-block:B01679:begin
\needspace{3\baselineskip}
\textbf{spice} — \textbf{wyrtgemang} (n.). spice; mixture of spices. \textbf{wyrt} (n.). Cook specifically routes these spices to wyrt; wyrtgemang is the separate earlier spice word.\par
% glossary-block:B01679:end
% glossary-block:B01680:begin
\needspace{3\baselineskip}
\textbf{spices} — \textbf{wyrt} (n.). Cook specifically routes these spices to wyrt; wyrtgemang is the separate earlier spice word.\par
% glossary-block:B01680:end
% glossary-block:B01681:begin
\needspace{3\baselineskip}
\textbf{spirit} — \textbf{gāst} (n.).\par
% glossary-block:B01681:end
% glossary-block:B01682:begin
\needspace{3\baselineskip}
\textbf{spit} — \textbf{spǣtan} (v.).\par
% glossary-block:B01682:end
% glossary-block:B01683:begin
\needspace{3\baselineskip}
\textbf{'s (possession)} — Possessive/relational English of or 's often becomes a dependent genitive, not preposition of. Partitive, source and lexical government require their own frames.\par
% glossary-block:B01683:end
% glossary-block:B01684:begin
\needspace{3\baselineskip}
\textbf{spread} — \textbf{brǣdan} (v.). spread; extend. Supplied source form spreads.\par
% glossary-block:B01684:end
% glossary-block:B01685:begin
\needspace{3\baselineskip}
\textbf{spy} — \textbf{scēawere} (n.).\par
% glossary-block:B01685:end
% glossary-block:B01686:begin
\needspace{3\baselineskip}
\textbf{stand} — \textbf{standan} (v.).\par
% glossary-block:B01686:end
% glossary-block:B01687:begin
\needspace{3\baselineskip}
\textbf{stand against} — \textbf{standan} (v.). oppose; stand against. Supplied standan ongēan frame, distinct from simple stand.\par
% glossary-block:B01687:end
% glossary-block:B01688:begin
\needspace{3\baselineskip}
\textbf{star} — \textbf{steorra} (n.).\par
% glossary-block:B01688:end
% glossary-block:B01689:begin
\needspace{3\baselineskip}
\textbf{steal} — \textbf{forstelan} (v.). steal away; steal. \textbf{stelan} (v.).\par
% glossary-block:B01689:end
% glossary-block:B01690:begin
\needspace{3\baselineskip}
\textbf{steal away} — \textbf{forstelan} (v.). steal away; steal.\par
% glossary-block:B01690:end
% glossary-block:B01691:begin
\needspace{3\baselineskip}
\textbf{steward} — \textbf{gerēfa} (n.). steward; reeve.\par
% glossary-block:B01691:end
% glossary-block:B01692:begin
\needspace{3\baselineskip}
\textbf{still} — \textbf{gīet} (adv.). still; further. \textbf{stille} (adj.). still; quiet.\par
% glossary-block:B01692:end
% glossary-block:B01693:begin
\needspace{3\baselineskip}
\textbf{stillness} — \textbf{stilnes} (n.). stillness; calm. Supplied genitive singular. stillness; quietness.\par
% glossary-block:B01693:end
% glossary-block:B01694:begin
\needspace{3\baselineskip}
\textbf{stone} — \textbf{stān} (n.).  \textbf{stānen} (adj.). stone; made of stone. Material adjective.\par
% glossary-block:B01694:end
% glossary-block:B01695:begin
\needspace{3\baselineskip}
\textbf{stone wall} — \textbf{stānweall} (n.).\par
% glossary-block:B01695:end
% glossary-block:B01696:begin
\needspace{3\baselineskip}
\textbf{storm} — \textbf{hrēohnes} (n.). tempest; storm. \textbf{storm} (n.).  \textbf{ȳst} (n.).\par
% glossary-block:B01696:end
% glossary-block:B01697:begin
\needspace{3\baselineskip}
\textbf{strand} — \textbf{strand} (n.). strand; sea-shore.\par
% glossary-block:B01697:end
% glossary-block:B01698:begin
\needspace{3\baselineskip}
\textbf{strange} — \textbf{sellic} (adj.). strange; remarkable.\par
% glossary-block:B01698:end
% glossary-block:B01699:begin
\needspace{3\baselineskip}
\textbf{stranger} — \textbf{cuma} (n.). stranger; guest.\par
% glossary-block:B01699:end
% glossary-block:B01700:begin
\needspace{3\baselineskip}
\textbf{strangers} — \textbf{ęlðēodig} (adj.). foreigners; strangers. Substantival adjective.\par
% glossary-block:B01700:end
% glossary-block:B01701:begin
\needspace{3\baselineskip}
\textbf{straw} — \textbf{strēaw} (n.).\par
% glossary-block:B01701:end
% glossary-block:B01702:begin
\needspace{3\baselineskip}
\textbf{stream of water} — \textbf{flōd} (n.). flood; stream of water.\par
% glossary-block:B01702:end
% glossary-block:B01703:begin
\needspace{3\baselineskip}
\textbf{street} — \textbf{strǣt} (n.).\par
% glossary-block:B01703:end
% glossary-block:B01704:begin
\needspace{3\baselineskip}
\textbf{strength} — \textbf{horn} (n.). horn; strength; dignity. Strength/dignity is the figurative force in 37.14, not a replacement for the literal source horn. \textbf{mægen} (n.). power; strength. Heavenly powers in 37.2 is the contextual plural use. \textbf{stręngðu} (n.).  \textbf{strengu} (n.). Recognition class-reference noun; preserve the source spelling and do not infer a plural count meaning.\par
% glossary-block:B01704:end
% glossary-block:B01705:begin
\needspace{3\baselineskip}
\textbf{strengthen} — \textbf{gestrangian} (v.).\par
% glossary-block:B01705:end
% glossary-block:B01706:begin
\needspace{3\baselineskip}
\textbf{strife} — \textbf{geflīt} (n.). dispute; strife.\par
% glossary-block:B01706:end
% glossary-block:B01707:begin
\needspace{3\baselineskip}
\textbf{strike} — \textbf{bēatan} (v.). beat; strike. Accusative person or thing struck. \textbf{geslēan} (v.). strike; smite. Supplied past/participial source form. Supplied source form struck. \textbf{slēan} (v.). strike; smite. strike; smite. Objectless striking in the assigned hammer sentences. Supplied past source form. \textbf{slēan} (v.), \textbf{bēatan} (v.). Slēan means strike/smite and can be objectless in the hammer sentences. Bēatan means beat/strike and takes an accusative object; the explicit bēatan target must be retained.\par
% glossary-block:B01707:end
% glossary-block:B01708:begin
\needspace{3\baselineskip}
\textbf{strong} — \textbf{strang} (adj.). Masculine singular predicate in the source. Attributive adjective with wind in the source. \textbf{swīð} (adjective citation).\par
% glossary-block:B01708:end
% glossary-block:B01709:begin
\needspace{3\baselineskip}
\textbf{stronger} — \textbf{swīðra} (adj.). stronger; fiercer. Comparative degree, distinct from right-side sense.\par
% glossary-block:B01709:end
% glossary-block:B01710:begin
\needspace{3\baselineskip}
\textbf{stronghold} — \textbf{fæsten} (n.). fortification; stronghold.\par
% glossary-block:B01710:end
% glossary-block:B01711:begin
\needspace{3\baselineskip}
\textbf{struck} — \textbf{ðrēan} (v.). chastised; struck.\par
% glossary-block:B01711:end
% glossary-block:B01712:begin
\needspace{3\baselineskip}
\textbf{structure built} — \textbf{getimbrung} (n.). building; structure built.\par
% glossary-block:B01712:end
% glossary-block:B01713:begin
\needspace{3\baselineskip}
\textbf{struggle} — \textbf{gewinn} (n.). distress; struggle. \textbf{winnan} (v.). struggle; toil.\par
% glossary-block:B01713:end
% glossary-block:B01714:begin
\needspace{3\baselineskip}
\textbf{study} — \textbf{geleornian} (v.). learn; study. \textbf{lār} (n.). learning; study. \textbf{leornung} (n.). study; learning.\par
% glossary-block:B01714:end
% glossary-block:B01715:begin
\needspace{3\baselineskip}
\textbf{subjection} — \textbf{underðēodnes} (n.). subjection; submission.\par
% glossary-block:B01715:end
% glossary-block:B01716:begin
\needspace{3\baselineskip}
\textbf{subjoin} — \textbf{underðēodan} (v.). add; subjoin. Supplied passive/participial addition to speech.\par
% glossary-block:B01716:end
% glossary-block:B01717:begin
\needspace{3\baselineskip}
\textbf{submission} — \textbf{underðēodnes} (n.). subjection; submission.\par
% glossary-block:B01717:end
% glossary-block:B01718:begin
\needspace{3\baselineskip}
\textbf{subside} — \textbf{gesmyltan} (v.). become calm. Supplied preterite of the sea becoming calm. \textbf{geswīcan} (v.). Absolute command to the sea; distinct from cessation with a genitive complement.\par
% glossary-block:B01718:end
% glossary-block:B01719:begin
\needspace{3\baselineskip}
\textbf{such} — \textbf{swilc} (indefinite adjective).\par
% glossary-block:B01719:end
% glossary-block:B01720:begin
\needspace{3\baselineskip}
\textbf{suddenly} — \textbf{fǣringa} (adv.).  \textbf{sęmninga} (adv.).\par
% glossary-block:B01720:end
% glossary-block:B01721:begin
\needspace{3\baselineskip}
\textbf{suffer} — \textbf{ðolian} (v.). suffer; endure. Accusative hardship endured.\par
% glossary-block:B01721:end
% glossary-block:B01722:begin
\needspace{3\baselineskip}
\textbf{suffice} — \textbf{genyhtsumian} (v.). suffice; be sufficient for. Complete supplied construction with dative beneficiary. Supplied complete construction with dative beneficiary; plural books require plural subjunctive.\par
% glossary-block:B01722:end
% glossary-block:B01723:begin
\needspace{3\baselineskip}
\textbf{summer} — \textbf{sumer} (n.).\par
% glossary-block:B01723:end
% glossary-block:B01724:begin
\needspace{3\baselineskip}
\textbf{summon} — \textbf{cleopian} (v.). call; summon. Accusative person summoned; distinguish intransitive cry aloud.\par
% glossary-block:B01724:end
% glossary-block:B01725:begin
\needspace{3\baselineskip}
\textbf{sun} — \textbf{sunne} (n.).\par
% glossary-block:B01725:end
% glossary-block:B01726:begin
\needspace{3\baselineskip}
\textbf{sunrise} — \textbf{ūpgang} (n.). rising; sunrise.\par
% glossary-block:B01726:end
% glossary-block:B01727:begin
\needspace{3\baselineskip}
\textbf{support} — \textbf{āfēdan} (v.). support; nourish. Accusative person supported. \textbf{fultum} (n.). help; support.\par
% glossary-block:B01727:end
% glossary-block:B01728:begin
\needspace{3\baselineskip}
\textbf{suppose} — \textbf{wēnan} (v.). suppose; expect; think. Genitive thing expected in the selected nominal frame; content-clause uses are separately taught.\par
% glossary-block:B01728:end
% glossary-block:B01729:begin
\needspace{3\baselineskip}
\textbf{surface} — \textbf{brādnes} (n.). surface; face. Of the deep in the assigned expression.\par
% glossary-block:B01729:end
% glossary-block:B01730:begin
\needspace{3\baselineskip}
\textbf{surge} — \textbf{weallan} (v.). surge; boil. Supplied present participle used adjectivally of the sea.\par
% glossary-block:B01730:end
% glossary-block:B01731:begin
\needspace{3\baselineskip}
\textbf{surround} — \textbf{ymbsęllan} (v.). surround; envelop. Supplied source form. surround; beset. Supplied source form.\par
% glossary-block:B01731:end
% glossary-block:B01732:begin
\needspace{3\baselineskip}
\textbf{sustain} — \textbf{healdan} (v.). maintain; sustain.\par
% glossary-block:B01732:end
% glossary-block:B01733:begin
\needspace{3\baselineskip}
\textbf{swear} — \textbf{swęrian} (v.).\par
% glossary-block:B01733:end
% glossary-block:B01734:begin
\needspace{3\baselineskip}
\textbf{swear by} — \textbf{ðurh, þurh} (prep.). by. Accusative oath standard.\par
% glossary-block:B01734:end
% glossary-block:B01735:begin
\needspace{3\baselineskip}
\textbf{sweet} — \textbf{swēte} (adj.).\par
% glossary-block:B01735:end
% glossary-block:B01736:begin
\needspace{3\baselineskip}
\textbf{sweeter} — \textbf{swētra} (adj.).\par
% glossary-block:B01736:end
% glossary-block:B01737:begin
\needspace{3\baselineskip}
\textbf{sweetest} — \textbf{swētest} (adj.).\par
% glossary-block:B01737:end
% glossary-block:B01738:begin
\needspace{3\baselineskip}
\textbf{sweetness} — \textbf{swētnes} (n.).\par
% glossary-block:B01738:end
% glossary-block:B01739:begin
\needspace{3\baselineskip}
\textbf{swift} — \textbf{swift} (adj.).\par
% glossary-block:B01739:end
% glossary-block:B01740:begin
\needspace{3\baselineskip}
\textbf{swiftness} — \textbf{swiftnes} (n.).\par
% glossary-block:B01740:end
% glossary-block:B01741:begin
\needspace{3\baselineskip}
\textbf{swim} — \textbf{swimman} (v.).\par
% glossary-block:B01741:end
% glossary-block:B01742:begin
\needspace{3\baselineskip}
\textbf{swimming} — \textbf{sund} (n.).\par
% glossary-block:B01742:end
% glossary-block:B01743:begin
\needspace{3\baselineskip}
\textbf{swing} — \textbf{swingan} (v.). swing; brandish. Supplied past source form.\par
% glossary-block:B01743:end
% glossary-block:B01744:begin
\needspace{3\baselineskip}
\textbf{swollen} — \textbf{āðindan} (v.). Supplied weak participial adjective describing the sea.\par
% glossary-block:B01744:end
\pdfbookmark[1]{T}{outline-B01745}
% glossary-block:B01745:begin
\needspace{5\baselineskip}
\subsection*{T}
% glossary-block:B01745:end
% glossary-block:B01746:begin
\textbf{table of the law} — \textbf{tabele} (n.). tablet; table. Tables of the law, not an unrestricted furniture meaning.\par
% glossary-block:B01746:end
% glossary-block:B01747:begin
\needspace{3\baselineskip}
\textbf{tablet} — \textbf{tabele} (n.). tablet; table. Tables of the law, not an unrestricted furniture meaning.\par
% glossary-block:B01747:end
% glossary-block:B01748:begin
\needspace{3\baselineskip}
\textbf{tailor} — \textbf{sēamere} (n.).\par
% glossary-block:B01748:end
% glossary-block:B01749:begin
\needspace{3\baselineskip}
\textbf{take} — \textbf{geniman} (v.).  \textbf{lǣdan} (v.). lead; bring; take. A person is the object in the assigned frame. carry; take. A person is conveyed in this supplied frame. take; carry along. \textbf{niman} (v.). Accusative person or thing taken, including food.\par
% glossary-block:B01749:end
% glossary-block:B01750:begin
\needspace{3\baselineskip}
\textbf{take away} — \textbf{āfierran} (v.). remove; take away.\par
% glossary-block:B01750:end
% glossary-block:B01751:begin
\needspace{3\baselineskip}
\textbf{take charge of} — \textbf{underfōn} (v.). receive; take charge of.\par
% glossary-block:B01751:end
% glossary-block:B01752:begin
\needspace{3\baselineskip}
\textbf{take leave of} — \textbf{grētan} (v.). greet; take leave of. Accusative person/place; greeting and farewell depend on context. grētan + accusative person/place; no separate take or bid is needed.\par
% glossary-block:B01752:end
% glossary-block:B01753:begin
\needspace{3\baselineskip}
\textbf{take notice of} — \textbf{gesēon} (v.). regard; take notice of.\par
% glossary-block:B01753:end
% glossary-block:B01754:begin
\needspace{3\baselineskip}
\textbf{take off} — \textbf{unscrȳdan} (v.). undress; take off. Supplied source form; retain its object/reflexive roles.\par
% glossary-block:B01754:end
% glossary-block:B01755:begin
\needspace{3\baselineskip}
\textbf{teach} — \textbf{lǣran} (v.). teach; instruct. \textbf{tǣcan} (v.). teach; show.\par
% glossary-block:B01755:end
% glossary-block:B01756:begin
\needspace{3\baselineskip}
\textbf{teacher} — \textbf{lārēow} (n.). master; teacher.\par
% glossary-block:B01756:end
% glossary-block:B01757:begin
\needspace{3\baselineskip}
\textbf{tear} — \textbf{tēar} (n.).  \textbf{tō-slītan} (v.). Supplied past source form.\par
% glossary-block:B01757:end
% glossary-block:B01758:begin
\needspace{3\baselineskip}
\textbf{tear apart} — \textbf{tō-slītan} (v.). tear. Supplied past source form.\par
% glossary-block:B01758:end
% glossary-block:B01759:begin
\needspace{3\baselineskip}
\textbf{tearful} — \textbf{wēpan} (v.). weeping; tearful. Supplied participial/adjectival use.\par
% glossary-block:B01759:end
% glossary-block:B01760:begin
\needspace{3\baselineskip}
\textbf{tell} — \textbf{sęcgan} (v.). tell; say. Dative addressee and a reported statement or question as content. say; tell. Person addressed is dative; direct or reported content is separate.\par
% glossary-block:B01760:end
% glossary-block:B01761:begin
\needspace{3\baselineskip}
\textbf{tell of beforehand} — \textbf{foresęcgean} (v.). mention before; tell of beforehand.\par
% glossary-block:B01761:end
% glossary-block:B01762:begin
\needspace{3\baselineskip}
\textbf{tell someone to} — \textbf{hātan} (v.). hātan + accusative actor + plain infinitive.\par
% glossary-block:B01762:end
% glossary-block:B01763:begin
\needspace{3\baselineskip}
\textbf{tempest} — \textbf{hrēohnes} (n.). tempest; storm.\par
% glossary-block:B01763:end
% glossary-block:B01764:begin
\needspace{3\baselineskip}
\textbf{temple} — \textbf{tempel} (n.).\par
% glossary-block:B01764:end
% glossary-block:B01765:begin
\needspace{3\baselineskip}
\textbf{temptation} — \textbf{costnung} (n.).\par
% glossary-block:B01765:end
% glossary-block:B01766:begin
\needspace{3\baselineskip}
\textbf{ten} — \textbf{tīen} (cardinal numeral). Cardinal 10.\par
% glossary-block:B01766:end
% glossary-block:B01767:begin
\needspace{3\baselineskip}
\textbf{tenacious} — \textbf{fæsthafol} (adj.). tenacious; withholding. Genitive of what is held back in the source.\par
% glossary-block:B01767:end
% glossary-block:B01768:begin
\needspace{3\baselineskip}
\textbf{terrible} — \textbf{ęgeslic} (adj.). dreadful; terrible.\par
% glossary-block:B01768:end
% glossary-block:B01769:begin
\needspace{3\baselineskip}
\textbf{territory} — \textbf{landgemǣre} (n.). territory; borders.\par
% glossary-block:B01769:end
% glossary-block:B01770:begin
\needspace{3\baselineskip}
\textbf{than} — \textbf{ðonne} (conj.; adv.). Comparison standard follows ðonne; do not confuse temporal ðonne then/when.\par
% glossary-block:B01770:end
% glossary-block:B01771:begin
\needspace{3\baselineskip}
\textbf{thane} — \textbf{ðegn} (n.). vassal; retainer; thane.\par
% glossary-block:B01771:end
% glossary-block:B01772:begin
\needspace{3\baselineskip}
\textbf{thanksgiving} — \textbf{ðancung} (n.).\par
% glossary-block:B01772:end
% glossary-block:B01773:begin
\needspace{3\baselineskip}
\textbf{that} — \textbf{se, sēo, ðæt} (article; demonstrative; relative pron.). who; which; that. Inflected relative; case follows its clause role. that; those. Ðes supplies this/these; se supplies demonstrative that/those. Select the form by case, number and gender. Conjunction and relative that have separate entries. \textbf{se-ðe} (relative). which; that. Supplied se-ðe relative complex referring to the storm. \textbf{ðæt} (conj.). Content conjunction, distinct from demonstrative/relative that. that; so that. Supplied expanded spelling; distinguish content from result by the source clause. that; so that. Asserted result, distinct from intended purpose. Cook leaves the source linkage uncertain; do not impose one content/result classification. \textbf{ðe} (relative particle; clause marker). which; that. Subject-relative referring to the waters. that; which; who. \textbf{ðe} (relative particle; clause marker), \textbf{se, sēo, ðæt} (article; demonstrative; relative pron.). Relative ðe is invariant; the inflected relative/demonstrative se takes the form required by its role in the relative clause. The antecedent’s case does not automatically determine that form.\par
% glossary-block:B01773:end
% glossary-block:B01774:begin
\needspace{3\baselineskip}
\textbf{that … may} — \textbf{ðæt} (conj.). Purpose normally uses ðæt + subjunctive in the assigned frames. Do not add a permission modal merely for English may; retain mugan when the prompt explicitly requests it.\par
% glossary-block:B01774:end
% glossary-block:B01775:begin
\needspace{3\baselineskip}
\textbf{that … might} — \textbf{ðæt} (conj.). Purpose normally uses ðæt + subjunctive in the assigned frames. Do not add a permission modal merely for English may; retain mugan when the prompt explicitly requests it.\par
% glossary-block:B01775:end
% glossary-block:B01776:begin
\needspace{3\baselineskip}
\textbf{that (reported content)} — \textbf{ðæt} (conj.). Content ðæt is separate from demonstrative and relative that; choose mood/time from the reporting construction.\par
% glossary-block:B01776:end
% glossary-block:B01777:begin
\needspace{3\baselineskip}
\textbf{that (result)} — \textbf{ðæt} (conj.). Result clause, often with degree swā/tō ðan; distinguish intended purpose.\par
% glossary-block:B01777:end
% glossary-block:B01778:begin
\needspace{3\baselineskip}
\textbf{that which} — \textbf{se, sēo, ðæt} (article; demonstrative; relative pron.). what; that which. \textbf{ðæt-ðæt} (pron.). that which; what.\par
% glossary-block:B01778:end
% glossary-block:B01779:begin
\needspace{3\baselineskip}
\textbf{the} — The definite article/demonstrative has se, sēo, ðæt as its nominative singular forms; choose its case, gender and number from the noun phrase. Preserve the exercises’ explicit article omissions. Select article agreement only where required; preserve all explicit Cook omission controls and bare noun predicates.\par
% glossary-block:B01779:end
% glossary-block:B01780:begin
\needspace{3\baselineskip}
\textbf{the Danes} — \textbf{Dęne} (n.). Danes. Plural people-name.\par
% glossary-block:B01780:end
% glossary-block:B01781:begin
\needspace{3\baselineskip}
\textbf{the dead} — \textbf{dēad} (adj.). dead people. Substantival adjective. Substantival adjectives take their own agreement/case; do not add man/people/ones where Cook omits it.\par
% glossary-block:B01781:end
% glossary-block:B01782:begin
\needspace{3\baselineskip}
\textbf{the deep} — \textbf{nēowolnes} (n.). abyss; the deep.\par
% glossary-block:B01782:end
% glossary-block:B01783:begin
\needspace{3\baselineskip}
\textbf{thee} — Ðē is singular dative or accusative thee/you; the governed case follows the verb or preposition. Singular ðū, dual git and plural gē systems; thou/thee mark singular, ye plural. Unspecified you allows appropriate choices unless the task specifies number.\par
% glossary-block:B01783:end
% glossary-block:B01784:begin
\needspace{3\baselineskip}
\textbf{the elderly} — \textbf{eald} (adj.). old people; the elderly. Substantival adjective.\par
% glossary-block:B01784:end
% glossary-block:B01785:begin
\needspace{3\baselineskip}
\textbf{the English} — \textbf{Engle} (n.). People, distinct from the language/adjective.\par
% glossary-block:B01785:end
% glossary-block:B01786:begin
\needspace{3\baselineskip}
\textbf{the evil} — \textbf{yfel} (adj.). evil; evil people. Substantival adjective in the assigned use.\par
% glossary-block:B01786:end
% glossary-block:B01787:begin
\needspace{3\baselineskip}
\textbf{the Highest} — \textbf{hēah} (adj.). Divine title.\par
% glossary-block:B01787:end
% glossary-block:B01788:begin
\needspace{3\baselineskip}
\textbf{the hungry} — \textbf{hungrig} (adj.). hungry people. Substantival adjective. Substantival adjectives take their own agreement/case; do not add man/people/ones where Cook omits it.\par
% glossary-block:B01788:end
% glossary-block:B01789:begin
\needspace{3\baselineskip}
\textbf{their} — \textbf{hē} (personal pron.). his; her; its; their. Invariant genitive possessors, unlike agreeing possessive adjectives. Source genitive possessor. Invariant personal genitive possessor; distinguish the same surface used as a governed personal case.\par
% glossary-block:B01789:end
% glossary-block:B01790:begin
\needspace{3\baselineskip}
\textbf{the living} — \textbf{cwic} (adj.). living people. Substantival adjective. Substantival adjectives take their own agreement/case; do not add man/people/ones where Cook omits it.\par
% glossary-block:B01790:end
% glossary-block:B01791:begin
\needspace{3\baselineskip}
\textbf{the Lord} — \textbf{dryhten} (n.). lord; the Lord.\par
% glossary-block:B01791:end
% glossary-block:B01792:begin
\needspace{3\baselineskip}
\textbf{the lowly} — \textbf{ēaðmōd} (adj.). humble people; the lowly. Substantival adjective.\par
% glossary-block:B01792:end
% glossary-block:B01793:begin
\needspace{3\baselineskip}
\textbf{them} — Hīe is plural accusative; him is plural dative. These plural forms serve all genders. Grammatical gender and antecedent determine reference; it may be represented by masculine/feminine forms for a gendered inanimate noun.\par
% glossary-block:B01793:end
% glossary-block:B01794:begin
\needspace{3\baselineskip}
\textbf{the mighty} — \textbf{mihtig} (adj.). the powerful; the mighty. Substantival adjective.\par
% glossary-block:B01794:end
% glossary-block:B01795:begin
\needspace{3\baselineskip}
\textbf{the more … the more} — \textbf{swā} (adverb; conjunction). the … the. Proportional comparison, not the definite article.\par
% glossary-block:B01795:end
% glossary-block:B01796:begin
\needspace{3\baselineskip}
\textbf{themselves} — An ordinary reflexive uses the personal pronoun in the case governed by its verb or preposition. Agreeing self adds the taught emphasis; English -self does not automatically require a separate self. Ordinary reflexive personal pronouns take the governed case; agreeing self adds the taught emphasis. English -self does not automatically require self.\par
% glossary-block:B01796:end
% glossary-block:B01797:begin
\needspace{3\baselineskip}
\textbf{then} — \textbf{ðā} (adv.).  \textbf{ðonne} (conj.; adv.). Sequence of future judgment events. then; at that point.\par
% glossary-block:B01797:end
% glossary-block:B01798:begin
\needspace{3\baselineskip}
\textbf{the naked} — \textbf{nacod} (adj.). naked people. Substantival adjective.\par
% glossary-block:B01798:end
% glossary-block:B01799:begin
\needspace{3\baselineskip}
\textbf{thence} — \textbf{ðanan} (adv.). thence; from there.\par
% glossary-block:B01799:end
% glossary-block:B01800:begin
\needspace{3\baselineskip}
\textbf{the noble} — \textbf{æðele} (adj.). noble people. Substantival adjective in the source.\par
% glossary-block:B01800:end
% glossary-block:B01801:begin
\needspace{3\baselineskip}
\textbf{the poor} — \textbf{earm} (adj.). the poor; the unfortunate. Substantival adjective.\par
% glossary-block:B01801:end
% glossary-block:B01802:begin
\needspace{3\baselineskip}
\textbf{the powerful} — \textbf{mihtig} (adj.). the powerful; the mighty. Substantival adjective.\par
% glossary-block:B01802:end
% glossary-block:B01803:begin
\needspace{3\baselineskip}
\textbf{there} — Ðǣr is locative there, at that place. Presentational there in there is/there was has no obligatory matching word and follows the assigned whole construction.\par
% glossary-block:B01803:end
% glossary-block:B01804:begin
\needspace{3\baselineskip}
\textbf{therefore} — \textbf{for-ðon-ðe} (conj.; adv. in for-ðon). therefore; for that reason. Adverbial consequence, distinct from because. \textbf{for-ðon-ðe} (conj.; adv. in for-ðon), \textbf{eornostlice} (adv.). The supplied consequence use of for-ðon-ðe means therefore/for that reason; eornostlice can mean accordingly/thus. Do not substitute the causal because sense.\par
% glossary-block:B01804:end
% glossary-block:B01805:begin
\needspace{3\baselineskip}
\textbf{therein} — \textbf{ðǣron} (adverb or particle). therein; in it. Supplied ðǣron points back to the sacks' location.\par
% glossary-block:B01805:end
% glossary-block:B01806:begin
\needspace{3\baselineskip}
\textbf{there is} — Presentational there has no obligatory OE word. Cook may supply ðā, use weorðan, or explicitly omit there; locative ðǣr is a distinct retained word.\par
% glossary-block:B01806:end
% glossary-block:B01807:begin
\needspace{3\baselineskip}
\textbf{the rest} — \textbf{ōðer} (adj.; ordinal). remaining; the rest. The rest of the forest/earth in the source.\par
% glossary-block:B01807:end
% glossary-block:B01808:begin
\needspace{3\baselineskip}
\textbf{there was} — Presentational there has no obligatory OE word. Cook may supply ðā, use weorðan, or explicitly omit there; locative ðǣr is a distinct retained word.\par
% glossary-block:B01808:end
% glossary-block:B01809:begin
\needspace{3\baselineskip}
\textbf{the rich} — \textbf{rīce} (adj.). rich; wealthy. Includes substantival rich people in the assigned frame.\par
% glossary-block:B01809:end
% glossary-block:B01810:begin
\needspace{3\baselineskip}
\textbf{the righteous} — \textbf{rihtwīs} (adj.). righteous people. Substantival adjective. Substantival adjectives take their own agreement/case; do not add man/people/ones where Cook omits it.\par
% glossary-block:B01810:end
% glossary-block:B01811:begin
\needspace{3\baselineskip}
\textbf{these} — \textbf{ðes, ðēos, ðis} (demonstrative). this; these. Ðes supplies this/these; se supplies demonstrative that/those. Select the form by case, number and gender. Conjunction and relative that have separate entries.\par
% glossary-block:B01811:end
% glossary-block:B01812:begin
\needspace{3\baselineskip}
\textbf{the sick} — \textbf{untrum} (adj.). sick people. Substantival adjective. Substantival adjectives take their own agreement/case; do not add man/people/ones where Cook omits it.\par
% glossary-block:B01812:end
% glossary-block:B01813:begin
\needspace{3\baselineskip}
\textbf{the … the} — \textbf{swā} (adverb; conjunction). Proportional comparison, not the definite article.\par
% glossary-block:B01813:end
% glossary-block:B01814:begin
\needspace{3\baselineskip}
\textbf{the … the (proportional comparison)} — \textbf{swā} (adverb; conjunction). swā … swā with comparatives; English the is not the definite article here.\par
% glossary-block:B01814:end
% glossary-block:B01815:begin
\needspace{3\baselineskip}
\textbf{the unnoble} — \textbf{unæðele} (adj.). people not noble. Substantival adjective in the supplied contrast.\par
% glossary-block:B01815:end
% glossary-block:B01816:begin
\needspace{3\baselineskip}
\textbf{the unrighteous} — \textbf{unrihtwīs} (adj.). unrighteous people. Substantival adjective.\par
% glossary-block:B01816:end
% glossary-block:B01817:begin
\needspace{3\baselineskip}
\textbf{they} — Hīe is plural nominative they for all genders. The plural dative is him, accusative hīe, genitive hiera. Grammatical gender and antecedent determine reference; it may be represented by masculine/feminine forms for a gendered inanimate noun.\par
% glossary-block:B01817:end
% glossary-block:B01818:begin
\needspace{3\baselineskip}
\textbf{thine} — \textbf{ðīn} (possessive adjective). your; yours. One possessor; agrees with the possessed noun. ðīn identifies one possessor and agrees with the possessed noun. ēower and incer mark plural and dual possessors separately.\par
% glossary-block:B01818:end
% glossary-block:B01819:begin
\needspace{3\baselineskip}
\textbf{thing} — \textbf{ðing} (n.). Ordinary singular/plural noun. thing; fact; matter. Matter known/discussed, distinct from a generic physical object.\par
% glossary-block:B01819:end
% glossary-block:B01820:begin
\needspace{3\baselineskip}
\textbf{things} — \textbf{ðing} (n.). thing. Ordinary singular/plural noun.\par
% glossary-block:B01820:end
% glossary-block:B01821:begin
\needspace{3\baselineskip}
\textbf{think} — \textbf{hycgan} (v.).  \textbf{ðęncan} (v.). Supplied source preterite. \textbf{wēnan} (v.). suppose; expect; think. Genitive thing expected in the selected nominal frame; content-clause uses are separately taught.\par
% glossary-block:B01821:end
% glossary-block:B01822:begin
\needspace{3\baselineskip}
\textbf{third} — \textbf{ðridda} (ordinal numeral).\par
% glossary-block:B01822:end
% glossary-block:B01823:begin
\needspace{3\baselineskip}
\textbf{thirsty} — \textbf{ðurstig} (adj.).\par
% glossary-block:B01823:end
% glossary-block:B01824:begin
\needspace{3\baselineskip}
\textbf{thirty} — \textbf{ðrītig} (numeral).\par
% glossary-block:B01824:end
% glossary-block:B01825:begin
\needspace{3\baselineskip}
\textbf{this} — \textbf{ðes, ðēos, ðis} (demonstrative). this; these. Ðes supplies this/these; se supplies demonstrative that/those. Select the form by case, number and gender. Conjunction and relative that have separate entries.\par
% glossary-block:B01825:end
% glossary-block:B01826:begin
\needspace{3\baselineskip}
\textbf{thither} — \textbf{ðider} (adv.). thither; to that place.\par
% glossary-block:B01826:end
% glossary-block:B01827:begin
\needspace{3\baselineskip}
\textbf{thorn} — \textbf{ðorn} (n.).\par
% glossary-block:B01827:end
% glossary-block:B01828:begin
\needspace{3\baselineskip}
\textbf{those} — \textbf{se, sēo, ðæt} (article; demonstrative; relative pron.). that; those. Supplied substantival dative plural referring to waters. Ðes supplies this/these; se supplies demonstrative that/those. Select the form by case, number and gender. Conjunction and relative that have separate entries.\par
% glossary-block:B01828:end
% glossary-block:B01829:begin
\needspace{3\baselineskip}
\textbf{thou} — Ðū is the singular nominative subject thou/you. Singular ðū, dual git and plural gē systems; thou/thee mark singular, ye plural. Unspecified you allows appropriate choices unless the task specifies number.\par
% glossary-block:B01829:end
% glossary-block:B01830:begin
\needspace{3\baselineskip}
\textbf{though} — \textbf{ðēah ðe}. Supplied whole concessive conjunction in EX-XII-XIII-43-10. \textbf{ðēah} (conj.; adv.). ðēah with the taught concessive clause; a subjunctive need not make an admitted fact unreal.\par
% glossary-block:B01830:end
% glossary-block:B01831:begin
\needspace{3\baselineskip}
\textbf{thought} — \textbf{geðanc} (n.). thought; mind.\par
% glossary-block:B01831:end
% glossary-block:B01832:begin
\needspace{3\baselineskip}
\textbf{thoughtful} — \textbf{behygdig} (adj.). attentive; thoughtful.\par
% glossary-block:B01832:end
% glossary-block:B01833:begin
\needspace{3\baselineskip}
\textbf{thousand} — \textbf{ðūsend} (numeral). Within the complete supplied high-number expression. High quantities retain the complete local numerical expressions; a glossary headword does not replace the supplied genitive/counting construction.\par
% glossary-block:B01833:end
% glossary-block:B01834:begin
\needspace{3\baselineskip}
\textbf{three} — \textbf{ðrīe} (cardinal numeral).\par
% glossary-block:B01834:end
% glossary-block:B01835:begin
\needspace{3\baselineskip}
\textbf{three times} — \textbf{ðrīwa} (adv.). thrice; three times. ðrīwa is an adverb; not a cardinal with a silently invented time noun.\par
% glossary-block:B01835:end
% glossary-block:B01836:begin
\needspace{3\baselineskip}
\textbf{thrice} — \textbf{ðrīwa} (adv.). thrice; three times. ðrīwa is an adverb; not a cardinal with a silently invented time noun.\par
% glossary-block:B01836:end
% glossary-block:B01837:begin
\needspace{3\baselineskip}
\textbf{throne} — \textbf{cynesętl} (n.). royal seat; throne. \textbf{setl} (n.). seat; throne. \textbf{ðrymsetl} (n.). throne; throne of glory.\par
% glossary-block:B01837:end
% glossary-block:B01838:begin
\needspace{3\baselineskip}
\textbf{throne of glory} — \textbf{ðrymsetl} (n.). throne; throne of glory.\par
% glossary-block:B01838:end
% glossary-block:B01839:begin
\needspace{3\baselineskip}
\textbf{through} — \textbf{geond} (prep.). along; through. With accusative. \textbf{ðurh, þurh} (prep.). Accusative landmark or medium. through; by. Mediated action through a person in the supplied frame. through; by means of. Supplied service/evidence expressions, including Ðurh hwæt.\par
% glossary-block:B01839:end
% glossary-block:B01840:begin
\needspace{3\baselineskip}
\textbf{throw} — \textbf{weorpan} (v.). Accusative thing thrown. Accusative thing cast; a following destination has its own government.\par
% glossary-block:B01840:end
% glossary-block:B01841:begin
\needspace{3\baselineskip}
\textbf{thunderbolt} — \textbf{līeg} (n.). thunderbolt; lightning. Source levin sense.\par
% glossary-block:B01841:end
% glossary-block:B01842:begin
\needspace{3\baselineskip}
\textbf{thus} — \textbf{eornostlice} (adv.). accordingly; thus. \textbf{swā} (adverb; conjunction). so; thus. \textbf{ðus} (adv.).\par
% glossary-block:B01842:end
% glossary-block:B01843:begin
\needspace{3\baselineskip}
\textbf{thy} — \textbf{ðīn} (possessive adjective). your; yours. One possessor; agrees with the possessed noun. ðīn identifies one possessor and agrees with the possessed noun. ēower and incer mark plural and dual possessors separately.\par
% glossary-block:B01843:end
% glossary-block:B01844:begin
\needspace{3\baselineskip}
\textbf{tie} — \textbf{bindan} (v.). bind; tie. Accusative person or thing bound.\par
% glossary-block:B01844:end
% glossary-block:B01845:begin
\needspace{3\baselineskip}
\textbf{time} — \textbf{sīð} (n.). time; occasion. \textbf{tīd} (n.). season; time.\par
% glossary-block:B01845:end
% glossary-block:B01846:begin
\needspace{3\baselineskip}
\textbf{tin} — \textbf{tin} (n.).\par
% glossary-block:B01846:end
% glossary-block:B01847:begin
\needspace{3\baselineskip}
\textbf{to} — On introduces the succeeding generation in the supplied from-generation-to-generation range. Retain the full supplied temporal phrase; on alone does not mean from generation to generation. \textbf{tō} (prep.). Endpoint/result death. Dative destination. to; toward. Person to whom knees/prayer are directed in the supplied phrase. Postposed in the supplied cwæð him tō speech phrase. at; to. Bounded supplied root expression with genitive dependent and dative root. to; into. Result in supplied tō lore, tō sorge or tō nāhte. The person addressed in speech. into; to. Resulting language/literary form. Complement of belonging in the supplied frame. Resulting condition of youth in the supplied comparison. Tō + inflected infinitive expresses the taught purpose, to/in order to. Plain infinitives after modals do not take an inserted tō; recipient to may instead be expressed by a dative. A dative noun/pronoun often expresses the recipient/beneficiary without an added preposition. Distinguish destinations and lexical government.\par
% glossary-block:B01847:end
% glossary-block:B01848:begin
\needspace{3\baselineskip}
\textbf{to-day} — \textbf{tō-dæg} (adv.), \textbf{tōdæg} (adv.). Tō-dæg and tōdæg are the retained source/course spellings of today.\par
% glossary-block:B01848:end
% glossary-block:B01849:begin
\needspace{3\baselineskip}
\textbf{today} — \textbf{tō-dæg} (adv.).  \textbf{tōdæg} (adv.).  \textbf{tō-dæg} (adv.), \textbf{tōdæg} (adv.). Tō-dæg and tōdæg are the retained source/course spellings of today.\par
% glossary-block:B01849:end
% glossary-block:B01850:begin
\needspace{3\baselineskip}
\textbf{to death} — Tō introduces death as the endpoint or result; use the dative form of the death noun in the complete supplied expression.\par
% glossary-block:B01850:end
% glossary-block:B01851:begin
\needspace{3\baselineskip}
\textbf{together} — \textbf{ætgædere} (adv.). together; at once. Paired with samod in the supplied source. \textbf{samod} (adv.).  \textbf{tōgædere} (adv.).\par
% glossary-block:B01851:end
% glossary-block:B01852:begin
\needspace{3\baselineskip}
\textbf{toil} — \textbf{gedeorf} (n.). labor; toil. \textbf{winnan} (v.). struggle; toil. Verb; the subject labours with the hands.\par
% glossary-block:B01852:end
% glossary-block:B01853:begin
\needspace{3\baselineskip}
\textbf{to meet} — \textbf{tōgēanes} (postposition). towards; to meet.\par
% glossary-block:B01853:end
% glossary-block:B01854:begin
\needspace{3\baselineskip}
\textbf{to-morrow} — \textbf{tō-morgen} (adv.). Adverb tō-morgen, not separately to plus morning.\par
% glossary-block:B01854:end
% glossary-block:B01855:begin
\needspace{3\baselineskip}
\textbf{tomorrow} — \textbf{tō-morgen} (adv.). Adverb tō-morgen, not separately to plus morning.\par
% glossary-block:B01855:end
% glossary-block:B01856:begin
\needspace{3\baselineskip}
\textbf{to one another} — \textbf{betwēonan} (adverb or particle). between; to one another. Supplied reciprocal betwēonan him in 45.6. Supplied reciprocal betwēonan him; do not translate one and another independently.\par
% glossary-block:B01856:end
% glossary-block:B01857:begin
\needspace{3\baselineskip}
\textbf{tooth} — \textbf{tōð} (n.).\par
% glossary-block:B01857:end
% glossary-block:B01858:begin
\needspace{3\baselineskip}
\textbf{to (purpose)} — \textbf{tō} (prep.). tō + inflected infinitive expresses the assigned elementary purpose; modal complements use plain infinitives.\par
% glossary-block:B01858:end
% glossary-block:B01859:begin
\needspace{3\baselineskip}
\textbf{toss} — \textbf{geswęngan} (v.). toss; shake. Supplied participial source use.\par
% glossary-block:B01859:end
% glossary-block:B01860:begin
\needspace{3\baselineskip}
\textbf{to that degree} — Tō ðan … ðæt is the supplied degree-and-result construction so/to that degree … that. Keep the whole construction.\par
% glossary-block:B01860:end
% glossary-block:B01861:begin
\needspace{3\baselineskip}
\textbf{to that place} — \textbf{ðider} (adv.). thither; to that place.\par
% glossary-block:B01861:end
% glossary-block:B01862:begin
\needspace{3\baselineskip}
\textbf{to the benefit of} — \textbf{for} (prep.). for; to the benefit of. \textbf{tō} (prep.). for; to the benefit of.\par
% glossary-block:B01862:end
% glossary-block:B01863:begin
\needspace{3\baselineskip}
\textbf{toward} — \textbf{in} (prep.). toward; onto. Supplied phrase directing the mind toward these things. \textbf{ofer} (prep.). In gōd ofer + accusative people. \textbf{tō} (prep.). to; toward. Person to whom knees/prayer are directed in the supplied phrase. Conduct toward a person in the supplied expression. Target of a look in the supplied frame. \textbf{wið} (prep.). Conduct toward a person in the supplied local frame.\par
% glossary-block:B01863:end
% glossary-block:B01864:begin
\needspace{3\baselineskip}
\textbf{towards} — \textbf{in} (prep.). toward; onto. Supplied phrase directing the mind toward these things. \textbf{tō} (prep.). toward. Conduct toward a person in the supplied expression. toward. Target of a look in the supplied frame. \textbf{tōgēanes} (postposition). towards; to meet. \textbf{wið} (prep.). toward. Conduct toward a person in the supplied local frame.\par
% glossary-block:B01864:end
% glossary-block:B01865:begin
\needspace{3\baselineskip}
\textbf{tower} — \textbf{hlīfian} (v.). Verb meaning stand/tower high in the source; distinct from noun tower. \textbf{torr} (n.).\par
% glossary-block:B01865:end
% glossary-block:B01866:begin
\needspace{3\baselineskip}
\textbf{to where} — \textbf{hwider} (adv.). whither; to where.\par
% glossary-block:B01866:end
% glossary-block:B01867:begin
\needspace{3\baselineskip}
\textbf{to wife} — Tō wīfe is the supplied role phrase as wife/to wife.\par
% glossary-block:B01867:end
% glossary-block:B01868:begin
\needspace{3\baselineskip}
\textbf{town} — \textbf{burg} (n.). town; city. \textbf{ceaster} (n.). town; city.\par
% glossary-block:B01868:end
% glossary-block:B01869:begin
\needspace{3\baselineskip}
\textbf{townspeople} — \textbf{ceasterwaran} (n.). townspeople; citizens.\par
% glossary-block:B01869:end
% glossary-block:B01870:begin
\needspace{3\baselineskip}
\textbf{to youth} — Tō introduces the resulting youthful condition in the supplied comparison; the accompanying youth noun supplies that meaning.\par
% glossary-block:B01870:end
% glossary-block:B01871:begin
\needspace{3\baselineskip}
\textbf{trade} — \textbf{cēapian} (v.). trade; traffic. Traffic means trade; retain the supplied trading frame. Cook's traffic means trade/buy through trade, not modern road traffic. \textbf{cræft} (n.). trade; occupation; craft.\par
% glossary-block:B01871:end
% glossary-block:B01872:begin
\needspace{3\baselineskip}
\textbf{traffic} — \textbf{cēapian} (v.). trade; traffic. Traffic means trade; retain the supplied trading frame. Cook's traffic means trade/buy through trade, not modern road traffic.\par
% glossary-block:B01872:end
% glossary-block:B01873:begin
\needspace{3\baselineskip}
\textbf{train} — \textbf{getēon} (v.). train; educate.\par
% glossary-block:B01873:end
% glossary-block:B01874:begin
\needspace{3\baselineskip}
\textbf{trained} — \textbf{gelǣred} (adj.). educated; trained; skilled.\par
% glossary-block:B01874:end
% glossary-block:B01875:begin
\needspace{3\baselineskip}
\textbf{trample down} — \textbf{fortrēdan} (v.).\par
% glossary-block:B01875:end
% glossary-block:B01876:begin
\needspace{3\baselineskip}
\textbf{translate} — \textbf{gewęndan} (v.). translate; turn from one language into another. \textbf{węndan} (v.).\par
% glossary-block:B01876:end
% glossary-block:B01877:begin
\needspace{3\baselineskip}
\textbf{translate into} — Tō introduces the resulting language or literary form with the supplied translation verb; tō alone does not express translate.\par
% glossary-block:B01877:end
% glossary-block:B01878:begin
\needspace{3\baselineskip}
\textbf{translation} — \textbf{āwęndednes} (n.).\par
% glossary-block:B01878:end
% glossary-block:B01879:begin
\needspace{3\baselineskip}
\textbf{translator} — \textbf{wealhstōd} (n.). interpreter; translator.\par
% glossary-block:B01879:end
% glossary-block:B01880:begin
\needspace{3\baselineskip}
\textbf{travel} — \textbf{faran} (v.). go; travel. Verb faran; distinct from noun journey sīðfæt.\par
% glossary-block:B01880:end
% glossary-block:B01881:begin
\needspace{3\baselineskip}
\textbf{treasure} — \textbf{goldhord} (n.). treasure; hoard of gold. \textbf{hord} (n.). hoard; treasure.\par
% glossary-block:B01881:end
% glossary-block:B01882:begin
\needspace{3\baselineskip}
\textbf{tree} — \textbf{trēo} (n.).\par
% glossary-block:B01882:end
% glossary-block:B01883:begin
\needspace{3\baselineskip}
\textbf{tribe} — \textbf{mǣgð} (n.). tribe; nation.\par
% glossary-block:B01883:end
% glossary-block:B01884:begin
\needspace{3\baselineskip}
\textbf{trouble} — \textbf{bisgian} (v.). occupy; trouble. \textbf{sorg} (n.). distress; anxiety; trouble; sorrow.\par
% glossary-block:B01884:end
% glossary-block:B01885:begin
\needspace{3\baselineskip}
\textbf{true} — \textbf{sōð} (adj.).\par
% glossary-block:B01885:end
% glossary-block:B01886:begin
\needspace{3\baselineskip}
\textbf{truly} — \textbf{sōðlice} (adv.). truly; indeed. Adverbial confirmation; distinct from discourse connective sōðlice where the key uses but/now.\par
% glossary-block:B01886:end
% glossary-block:B01887:begin
\needspace{3\baselineskip}
\textbf{trumpet} — \textbf{biema} (n.). trumpet; clarion.\par
% glossary-block:B01887:end
% glossary-block:B01888:begin
\needspace{3\baselineskip}
\textbf{trust} — \textbf{trēownes} (n.). trust; confidence.\par
% glossary-block:B01888:end
% glossary-block:B01889:begin
\needspace{3\baselineskip}
\textbf{trust in} — \textbf{on} (prep.). in. Accusative person believed/trusted in. \textbf{gelīefan} (v.). gelīefan on + accusative person in this assigned frame.\par
% glossary-block:B01889:end
% glossary-block:B01890:begin
\needspace{3\baselineskip}
\textbf{truth} — \textbf{sōðfæstnes} (n.).\par
% glossary-block:B01890:end
% glossary-block:B01891:begin
\needspace{3\baselineskip}
\textbf{turn} — \textbf{bewęndan} (v.). The assigned reflexive frame uses dative him and tō + dative goal. \textbf{gewęndan} (v.). return; turn. turn; change back. Supplied source form. \textbf{węndan} (v.). turn; return. Supplied source form.\par
% glossary-block:B01891:end
% glossary-block:B01892:begin
\needspace{3\baselineskip}
\textbf{turn back} — \textbf{ęfthweorfan} (v.). return; turn back. \textbf{hweorfan} (v.). return; turn back. Supplied source preterite.\par
% glossary-block:B01892:end
% glossary-block:B01893:begin
\needspace{3\baselineskip}
\textbf{turn from one language into another} — \textbf{gewęndan} (v.). translate; turn from one language into another.\par
% glossary-block:B01893:end
% glossary-block:B01894:begin
\needspace{3\baselineskip}
\textbf{turn one's attention to} — \textbf{āðęnian} (v.). direct; stretch toward. Directing the mind toward what is heard; supplied figurative phrase.\par
% glossary-block:B01894:end
% glossary-block:B01895:begin
\needspace{3\baselineskip}
\textbf{turn one's gaze} — \textbf{besēon} (v.). look; turn one's gaze. Supplied source form and frame.\par
% glossary-block:B01895:end
% glossary-block:B01896:begin
\needspace{3\baselineskip}
\textbf{twelve} — \textbf{twęlf} (cardinal numeral). Cardinal 12.\par
% glossary-block:B01896:end
% glossary-block:B01897:begin
\needspace{3\baselineskip}
\textbf{twenty} — \textbf{twēntig} (numeral).\par
% glossary-block:B01897:end
% glossary-block:B01898:begin
\needspace{3\baselineskip}
\textbf{two} — \textbf{twēgen} (cardinal numeral). two; in two. Two parts in the supplied division expression.\par
% glossary-block:B01898:end
% glossary-block:B01899:begin
\needspace{3\baselineskip}
\textbf{twofold} — \textbf{twīfeald} (adj.). double; twofold.\par
% glossary-block:B01899:end
\pdfbookmark[1]{U}{outline-B01900}
% glossary-block:B01900:begin
\needspace{5\baselineskip}
\subsection*{U}
% glossary-block:B01900:end
% glossary-block:B01901:begin
\textbf{uncounted quantity} — \textbf{unrīm} (n.). abundance; uncounted quantity. Supplied quantity with genitive of the substance.\par
% glossary-block:B01901:end
% glossary-block:B01902:begin
\needspace{3\baselineskip}
\textbf{under} — \textbf{under} (prep.). Dative landmark in the taught static frame. Directional accusative in the assigned entry-under-the-roof frame.\par
% glossary-block:B01902:end
% glossary-block:B01903:begin
\needspace{3\baselineskip}
\textbf{under one's control} — \textbf{on} (prep.). in one's control. on + dative possessed geweald. On + dative possessed geweald expresses in/under someone’s control. In the assigned habban construction, the things controlled are its accusative object; the controller is expressed within the geweald phrase.\par
% glossary-block:B01903:end
% glossary-block:B01904:begin
\needspace{3\baselineskip}
\textbf{understand} — \textbf{ongietan} (v.). understand; perceive. Supplied infinitive after a modal. understand; perceive.\par
% glossary-block:B01904:end
% glossary-block:B01905:begin
\needspace{3\baselineskip}
\textbf{understanding} — \textbf{andgiet} (n.). andgiet denotes the knowledge possessed and desired in this exercise.\par
% glossary-block:B01905:end
% glossary-block:B01906:begin
\needspace{3\baselineskip}
\textbf{undo} — \textbf{undōn} (v.). undo; open.\par
% glossary-block:B01906:end
% glossary-block:B01907:begin
\needspace{3\baselineskip}
\textbf{undress} — \textbf{unscrȳdan} (v.). undress; take off. Supplied source form; retain its object/reflexive roles.\par
% glossary-block:B01907:end
% glossary-block:B01908:begin
\needspace{3\baselineskip}
\textbf{unfamiliar} — \textbf{ungewunelic} (adj.). unusual; unfamiliar.\par
% glossary-block:B01908:end
% glossary-block:B01909:begin
\needspace{3\baselineskip}
\textbf{unfortunate people} — \textbf{earm} (adj.). the poor; the unfortunate. Substantival adjective.\par
% glossary-block:B01909:end
% glossary-block:B01910:begin
\needspace{3\baselineskip}
\textbf{unfree} — \textbf{ðēow} (adj.). the unfree; enslaved people. Substantival adjective in the source comparison.\par
% glossary-block:B01910:end
% glossary-block:B01911:begin
\needspace{3\baselineskip}
\textbf{unknown} — \textbf{ungecnāwen} (past participle; supplied predicate).\par
% glossary-block:B01911:end
% glossary-block:B01912:begin
\needspace{3\baselineskip}
\textbf{unless} — \textbf{būtan} (conj.). unless; except if. Būtan introducing the supplied conditional exception; distinguish a prepositional complement.\par
% glossary-block:B01912:end
% glossary-block:B01913:begin
\needspace{3\baselineskip}
\textbf{unrest} — \textbf{unstilnes} (n.). unrest; disorder.\par
% glossary-block:B01913:end
% glossary-block:B01914:begin
\needspace{3\baselineskip}
\textbf{unrighteous} — \textbf{unrihtwīs} (adj.). Adjective, distinct from the substantive unrighteous people.\par
% glossary-block:B01914:end
% glossary-block:B01915:begin
\needspace{3\baselineskip}
\textbf{unrighteousness} — \textbf{unrihtwīsnes} (n.).\par
% glossary-block:B01915:end
% glossary-block:B01916:begin
\needspace{3\baselineskip}
\textbf{unrighteous people} — \textbf{unrihtwīs} (adj.). Substantival adjective.\par
% glossary-block:B01916:end
% glossary-block:B01917:begin
\needspace{3\baselineskip}
\textbf{unthankful} — \textbf{ungeðancfull} (adj.).\par
% glossary-block:B01917:end
% glossary-block:B01918:begin
\needspace{3\baselineskip}
\textbf{until} — \textbf{oð-ðæt} (conj.). Oð-ðæt introduces until; an anticipated endpoint takes the taught subjunctive construction. \textbf{oð} (prep.). up to; until. With accusative.\par
% glossary-block:B01918:end
% glossary-block:B01919:begin
\needspace{3\baselineskip}
\textbf{unusual} — \textbf{ungewunelic} (adj.). unusual; unfamiliar.\par
% glossary-block:B01919:end
% glossary-block:B01920:begin
\needspace{3\baselineskip}
\textbf{up} — \textbf{ūp} (adv.).\par
% glossary-block:B01920:end
% glossary-block:B01921:begin
\needspace{3\baselineskip}
\textbf{upon} — \textbf{on} (prep.). upon; into. Supplied chancing-upon idiom. Accusative beloved in the supplied love-falls-upon idiom. \textbf{uppan} (prep.). upon; on top of.\par
% glossary-block:B01921:end
% glossary-block:B01922:begin
\needspace{3\baselineskip}
\textbf{up to} — \textbf{oð} (prep.). up to; until. With accusative. \textbf{ōð} (prep.). as far as; up to.\par
% glossary-block:B01922:end
% glossary-block:B01923:begin
\needspace{3\baselineskip}
\textbf{us} — Ūs is first-person plural dative or accusative us; unc is the explicitly dual us two. Select case from the verb/preposition and preserve dual versus plural; dual subjects use plural verb forms.\par
% glossary-block:B01923:end
% glossary-block:B01924:begin
\needspace{3\baselineskip}
\textbf{useful} — \textbf{behēfe} (adj.). Dative person to whom something is useful.\par
% glossary-block:B01924:end
% glossary-block:B01925:begin
\needspace{3\baselineskip}
\textbf{us two} — Unc is first-person dual dative or accusative us two. Select case from the verb/preposition and preserve dual versus plural; dual subjects use plural verb forms.\par
% glossary-block:B01925:end
% glossary-block:B01926:begin
\needspace{3\baselineskip}
\textbf{utterance} — \textbf{word} (n.). speech; utterance.\par
% glossary-block:B01926:end
\pdfbookmark[1]{V}{outline-B01927}
% glossary-block:B01927:begin
\needspace{5\baselineskip}
\subsection*{V}
% glossary-block:B01927:end
% glossary-block:B01928:begin
\textbf{valuable} — \textbf{dēorwierðe} (adj.). precious; valuable.\par
% glossary-block:B01928:end
% glossary-block:B01929:begin
\needspace{3\baselineskip}
\textbf{value} — \textbf{weorð} (n.). price; value.\par
% glossary-block:B01929:end
% glossary-block:B01930:begin
\needspace{3\baselineskip}
\textbf{various} — \textbf{mislic} (adj.).\par
% glossary-block:B01930:end
% glossary-block:B01931:begin
\needspace{3\baselineskip}
\textbf{vassal} — \textbf{ðegn} (n.). vassal; retainer; thane.\par
% glossary-block:B01931:end
% glossary-block:B01932:begin
\needspace{3\baselineskip}
\textbf{vein of metal} — \textbf{ōra} (n.). vein of metal; ore (sense uncertain). Supplied dative plural after on.\par
% glossary-block:B01932:end
% glossary-block:B01933:begin
\needspace{3\baselineskip}
\textbf{verily} — \textbf{sōðlice} (adv.). Truly; indeed.\par
% glossary-block:B01933:end
% glossary-block:B01934:begin
\needspace{3\baselineskip}
\textbf{verse} — \textbf{lēoð} (n.). poem; poetry; verse.\par
% glossary-block:B01934:end
% glossary-block:B01935:begin
\needspace{3\baselineskip}
\textbf{very} — \textbf{full} (adv.).  \textbf{swīðe} (adv.). Degree adverb swīðe; retain the separate great/large adjective micel.\par
% glossary-block:B01935:end
% glossary-block:B01936:begin
\needspace{3\baselineskip}
\textbf{very calm} — \textbf{smyltest} (adj.). very calm; calmest. Supplied superlative adjective describing the sea.\par
% glossary-block:B01936:end
% glossary-block:B01937:begin
\needspace{3\baselineskip}
\textbf{very greatly} — Swīðost means most/very greatly. It is the recognition-only superlative reference form of swīðe.\par
% glossary-block:B01937:end
% glossary-block:B01938:begin
\needspace{3\baselineskip}
\textbf{vessel} — \textbf{fæt} (n.). vessel; container.\par
% glossary-block:B01938:end
% glossary-block:B01939:begin
\needspace{3\baselineskip}
\textbf{victory} — \textbf{sige} (n.).\par
% glossary-block:B01939:end
% glossary-block:B01940:begin
\needspace{3\baselineskip}
\textbf{vineyard} — \textbf{wīngeard} (n.).\par
% glossary-block:B01940:end
% glossary-block:B01941:begin
\needspace{3\baselineskip}
\textbf{violently} — \textbf{ðearlice} (adv.). violently; fiercely.\par
% glossary-block:B01941:end
% glossary-block:B01942:begin
\needspace{3\baselineskip}
\textbf{virginity} — \textbf{mægðhād} (n.). maidenhood; virginity.\par
% glossary-block:B01942:end
% glossary-block:B01943:begin
\needspace{3\baselineskip}
\textbf{visible} — \textbf{gesewenlic} (adj.).\par
% glossary-block:B01943:end
% glossary-block:B01944:begin
\needspace{3\baselineskip}
\textbf{visit} — \textbf{genēosian} (v.). Accusative person in the selected prefixed-verb frame.\par
% glossary-block:B01944:end
% glossary-block:B01945:begin
\needspace{3\baselineskip}
\textbf{voice} — \textbf{stefn} (n.).\par
% glossary-block:B01945:end
% glossary-block:B01946:begin
\needspace{3\baselineskip}
\textbf{void} — \textbf{ǣmtig} (adj.). empty; void. \textbf{īdel} (adj.). empty; void.\par
% glossary-block:B01946:end
\pdfbookmark[1]{W}{outline-B01947}
% glossary-block:B01947:begin
\needspace{5\baselineskip}
\subsection*{W}
% glossary-block:B01947:end
% glossary-block:B01948:begin
\textbf{wages paid for work} — \textbf{mēd} (n.). hire; wages paid for work.\par
% glossary-block:B01948:end
% glossary-block:B01949:begin
\needspace{3\baselineskip}
\textbf{walk} — \textbf{gān} (v.). go; walk.\par
% glossary-block:B01949:end
% glossary-block:B01950:begin
\needspace{3\baselineskip}
\textbf{wall} — \textbf{weall} (n.).\par
% glossary-block:B01950:end
% glossary-block:B01951:begin
\needspace{3\baselineskip}
\textbf{want to} — \textbf{willan} (v.). wish; want to. willan + plain infinitive denotes volition; intention does not assert the action occurred.\par
% glossary-block:B01951:end
% glossary-block:B01952:begin
\needspace{3\baselineskip}
\textbf{wares} — \textbf{hlæst} (n.). wares; merchandise; cargo. Selected plural source use.\par
% glossary-block:B01952:end
% glossary-block:B01953:begin
\needspace{3\baselineskip}
\textbf{warrior} — \textbf{beorn} (n.).\par
% glossary-block:B01953:end
% glossary-block:B01954:begin
\needspace{3\baselineskip}
\textbf{was} — \textbf{wesan} (v.). be. Copular present/past identification or state; other auxiliary constructions are separate senses. Copular eom/wæs and bēo indicative series; later passive and compound uses require their own construction.\par
% glossary-block:B01954:end
% glossary-block:B01955:begin
\needspace{3\baselineskip}
\textbf{wash} — \textbf{ðwēan} (v.). Accusative body part, mid + dative liquid.\par
% glossary-block:B01955:end
% glossary-block:B01956:begin
\needspace{3\baselineskip}
\textbf{was … -ing} — Simple OE finite forms can render English progressive aspect; wesan + present participle is only one bounded source construction, not a required default.\par
% glossary-block:B01956:end
% glossary-block:B01957:begin
\needspace{3\baselineskip}
\textbf{watch} — \textbf{scēawian} (v.). watch; observe. Accusative person or event watched, including a journey.\par
% glossary-block:B01957:end
% glossary-block:B01958:begin
\needspace{3\baselineskip}
\textbf{water} — \textbf{wæter} (n.).\par
% glossary-block:B01958:end
% glossary-block:B01959:begin
\needspace{3\baselineskip}
\textbf{water for the feet} — The supplied phrase combines water with tō and the dative feet expression; tō supplies the relation for, not the meanings water or feet.\par
% glossary-block:B01959:end
% glossary-block:B01960:begin
\needspace{3\baselineskip}
\textbf{wave} — \textbf{ȳð} (n.). wave; billow. Supplied dative plural after fram.\par
% glossary-block:B01960:end
% glossary-block:B01961:begin
\needspace{3\baselineskip}
\textbf{way} — \textbf{weg} (n.). way; path. way; means of access. Figurative access to the Father in the declaration; distinct from physical road. Weg means road, path or way.\par
% glossary-block:B01961:end
% glossary-block:B01962:begin
\needspace{3\baselineskip}
\textbf{we} — Wē is nominative plural we; wit is the explicitly dual we two. Both take a plural verb form. Select case from the verb/preposition and preserve dual versus plural; dual subjects use plural verb forms.\par
% glossary-block:B01962:end
% glossary-block:B01963:begin
\needspace{3\baselineskip}
\textbf{wealth} — \textbf{wela} (n.). riches; wealth.\par
% glossary-block:B01963:end
% glossary-block:B01964:begin
\needspace{3\baselineskip}
\textbf{wealthy} — \textbf{rīce} (adj.). rich; wealthy. Includes substantival rich people in the assigned frame. \textbf{welig} (adj.). rich; wealthy. \textbf{wlanc} (adj.). proud; wealthy. Contextually the rich man addressed. \textbf{welig} (adj.), \textbf{rīce} (adj.). Welig and rīce can mean rich/wealthy; rīce also occurs substantivally for rich people. Inflect the selected adjective for its role.\par
% glossary-block:B01964:end
% glossary-block:B01965:begin
\needspace{3\baselineskip}
\textbf{weapon} — \textbf{wǣpen} (n.).\par
% glossary-block:B01965:end
% glossary-block:B01966:begin
\needspace{3\baselineskip}
\textbf{wearing} — \textbf{on} (prep.). in; wearing. Supplied on scēapa gegierelum, in sheep's garments.\par
% glossary-block:B01966:end
% glossary-block:B01967:begin
\needspace{3\baselineskip}
\textbf{weep} — \textbf{wēpan} (v.).\par
% glossary-block:B01967:end
% glossary-block:B01968:begin
\needspace{3\baselineskip}
\textbf{weeping} — \textbf{wēpan} (v.). weeping; tearful. Supplied participial/adjectival use.\par
% glossary-block:B01968:end
% glossary-block:B01969:begin
\needspace{3\baselineskip}
\textbf{weight} — \textbf{gewiht} (n.).\par
% glossary-block:B01969:end
% glossary-block:B01970:begin
\needspace{3\baselineskip}
\textbf{well} — \textbf{gesund} (adj.). well; in good health. \textbf{hāl} (adj.). hale; well. Supplied greeting. \textbf{wel} (adv.).\par
% glossary-block:B01970:end
% glossary-block:B01971:begin
\needspace{3\baselineskip}
\textbf{were} — \textbf{wesan} (v.). be. Copular present/past identification or state; other auxiliary constructions are separate senses. Copular eom/wæs and bēo indicative series; later passive and compound uses require their own construction.\par
% glossary-block:B01971:end
% glossary-block:B01972:begin
\needspace{3\baselineskip}
\textbf{were it not for} — \textbf{wesan} (v.). Whole supplied past-subjunctive conditional and wesan + cumen result; preserve its counterfactual anterior sense.\par
% glossary-block:B01972:end
% glossary-block:B01973:begin
\needspace{3\baselineskip}
\textbf{we two} — Wit is nominative dual we two and takes a plural verb form. Select case from the verb/preposition and preserve dual versus plural; dual subjects use plural verb forms.\par
% glossary-block:B01973:end
% glossary-block:B01974:begin
\needspace{3\baselineskip}
\textbf{whale} — \textbf{hwæl} (n.).\par
% glossary-block:B01974:end
% glossary-block:B01975:begin
\needspace{3\baselineskip}
\textbf{what} — \textbf{hwā, hwæt} (interrogative/indefinite pron.). who? what?. Case follows role/government; distinguish personal and neuter reference. Inflect the interrogative by its role/government; dependent questions and relatives are separate constructions. \textbf{se, sēo, ðæt} (article; demonstrative; relative pron.). what; that which. \textbf{ðæt-ðæt} (pron.). that which; what.\par
% glossary-block:B01975:end
% glossary-block:B01976:begin
\needspace{3\baselineskip}
\textbf{whatever} — \textbf{swā-hwæt-swā} (pron.). whatsoever; whatever. Swā hwā swā means whoever, swā hwæt swā whatever, and swā hwæðer swā whichever of two. Inflect the internal pronoun for its clause role; swā alone does not carry the whole free-choice meaning. Swā hwā swā means whoever; swā hwæt swā means whatever. Inflect the internal pronoun for its clause role, so whomsoever requires the appropriate oblique form. Retain any explicitly supplied attraction.\par
% glossary-block:B01976:end
% glossary-block:B01977:begin
\needspace{3\baselineskip}
\textbf{what reason} — \textbf{hwā, hwæt} (interrogative/indefinite pron.). Instrumental hwon in the supplied why expressions.\par
% glossary-block:B01977:end
% glossary-block:B01978:begin
\needspace{3\baselineskip}
\textbf{whatsoever} — \textbf{swā-hwæt-swā} (pron.). whatsoever; whatever. Swā hwā swā means whoever, swā hwæt swā whatever, and swā hwæðer swā whichever of two. Inflect the internal pronoun for its clause role; swā alone does not carry the whole free-choice meaning. Swā hwā swā means whoever; swā hwæt swā means whatever. Inflect the internal pronoun for its clause role, so whomsoever requires the appropriate oblique form. Retain any explicitly supplied attraction.\par
% glossary-block:B01978:end
% glossary-block:B01979:begin
\needspace{3\baselineskip}
\textbf{what sort of} — \textbf{hūlic} (interrogative adjective). Interrogative hūlic asks what sort/kind of; inflect it for the noun phrase required by the task.\par
% glossary-block:B01979:end
% glossary-block:B01980:begin
\needspace{3\baselineskip}
\textbf{wheat} — \textbf{hwǣte} (n.).\par
% glossary-block:B01980:end
% glossary-block:B01981:begin
\needspace{3\baselineskip}
\textbf{when} — \textbf{hwonne} (adv.). when?. Direct interrogative. \textbf{mid-ðȳ} (conj.). when; while. \textbf{ðā} (conj.). when; factual past time. \textbf{ðonne} (conj.). when; recurring or future time. \textbf{ðā} (conj.), \textbf{ðonne} (conj.), \textbf{hwonne} (adv.). Ðā introduces factual past when; ðonne recurring/future when; hwonne asks when? Clause mood follows the particular temporal construction.\par
% glossary-block:B01981:end
% glossary-block:B01982:begin
\needspace{3\baselineskip}
\textbf{whence} — \textbf{hwanon} (adv.). whence; from where.\par
% glossary-block:B01982:end
% glossary-block:B01983:begin
\needspace{3\baselineskip}
\textbf{where} — \textbf{hwǣr} (adv.).\par
% glossary-block:B01983:end
% glossary-block:B01984:begin
\needspace{3\baselineskip}
\textbf{whether} — \textbf{hwæðer} (conj.).\par
% glossary-block:B01984:end
% glossary-block:B01985:begin
\needspace{3\baselineskip}
\textbf{which} — \textbf{hwilc} (interrogative/indefinite adj. or pron.). Selection from a group expressed in genitive. \textbf{se, sēo, ðæt} (article; demonstrative; relative pron.). who; which; that. Inflected relative; case follows its clause role. \textbf{se-ðe} (relative). which; that. Supplied se-ðe relative complex referring to the storm. \textbf{ðe} (relative particle; clause marker). which; that. Subject-relative referring to the waters. that; which; who.\par
% glossary-block:B01985:end
% glossary-block:B01986:begin
\needspace{3\baselineskip}
\textbf{whichever} — Swā hwā swā means whoever, swā hwæt swā whatever, and swā hwæðer swā whichever of two. Inflect the internal pronoun for its clause role; swā alone does not carry the whole free-choice meaning.\par
% glossary-block:B01986:end
% glossary-block:B01987:begin
\needspace{3\baselineskip}
\textbf{whichever of two} — Swā hwæðer swā is the whole free-choice expression whichever of two; the supplied attraction example has a special case arrangement. \textbf{hwæðer} (interrogative pronoun/adjective). Use the whole free-choice expression swā hwæðer swā, with the pronoun inflected for its role. The supplied attraction example has a special case arrangement.\par
% glossary-block:B01987:end
% glossary-block:B01988:begin
\needspace{3\baselineskip}
\textbf{which of} — \textbf{hwilc} (interrogative/indefinite adj. or pron.). which. Selection from a group expressed in genitive.\par
% glossary-block:B01988:end
% glossary-block:B01989:begin
\needspace{3\baselineskip}
\textbf{which of two} — \textbf{hwæðer} (interrogative pronoun/adjective). which of two?. Interrogative hwæðer asks which of two; its case follows its role.\par
% glossary-block:B01989:end
% glossary-block:B01990:begin
\needspace{3\baselineskip}
\textbf{which? of what kind} — \textbf{hwilc} (interrogative/indefinite adj. or pron.). which? of what kind?.\par
% glossary-block:B01990:end
% glossary-block:B01991:begin
\needspace{3\baselineskip}
\textbf{which (relative)} — \textbf{ðe} (relative particle; clause marker), \textbf{se, sēo, ðæt} (article; demonstrative; relative pron.). Relative ðe is invariant; the inflected relative/demonstrative se takes the form required by its role in the relative clause. The antecedent’s case does not automatically determine that form.\par
% glossary-block:B01991:end
% glossary-block:B01992:begin
\needspace{3\baselineskip}
\textbf{while} — \textbf{fæc} (n.). interval; period of time; while. Noun in the supplied time expression. \textbf{mid-ðȳ} (conj.). when; while. \textbf{ðā-hwīle-ðe} (conj.).  \textbf{mid-ðȳ} (conj.), \textbf{ðā-hwīle-ðe} (conj.). Mid-ðȳ and ðā-hwīle-ðe are complete conjunctions meaning while in their assigned temporal frames.\par
% glossary-block:B01992:end
% glossary-block:B01993:begin
\needspace{3\baselineskip}
\textbf{white} — \textbf{hwīt} (adj.).\par
% glossary-block:B01993:end
% glossary-block:B01994:begin
\needspace{3\baselineskip}
\textbf{whither} — \textbf{hwider} (adv.). whither; to where.\par
% glossary-block:B01994:end
% glossary-block:B01995:begin
\needspace{3\baselineskip}
\textbf{who} — \textbf{hwā, hwæt} (interrogative/indefinite pron.). who? what?. Case follows role/government; distinguish personal and neuter reference. Inflect the interrogative by its role/government; dependent questions and relatives are separate constructions. \textbf{se, sēo, ðæt} (article; demonstrative; relative pron.). who; which; that. Inflected relative; case follows its clause role. \textbf{ðe} (relative particle; clause marker). that; which; who. \textbf{ðe} (relative particle; clause marker), \textbf{se, sēo, ðæt} (article; demonstrative; relative pron.). Relative ðe is invariant; the inflected relative/demonstrative se takes the form required by its role in the relative clause. The antecedent’s case does not automatically determine that form.\par
% glossary-block:B01995:end
% glossary-block:B01996:begin
\needspace{3\baselineskip}
\textbf{whoever} — Swā hwā swā means whoever, swā hwæt swā whatever, and swā hwæðer swā whichever of two. Inflect the internal pronoun for its clause role; swā alone does not carry the whole free-choice meaning. Swā hwā swā means whoever; swā hwæt swā means whatever. Inflect the internal pronoun for its clause role, so whomsoever requires the appropriate oblique form. Retain any explicitly supplied attraction.\par
% glossary-block:B01996:end
% glossary-block:B01997:begin
\needspace{3\baselineskip}
\textbf{whole} — \textbf{eall} (adj.; substantive; adv.). all; whole. Eall agrees with a noun in all/whole expressions. Substantival plural refers to all people; neuter singular can mean everything.\par
% glossary-block:B01997:end
% glossary-block:B01998:begin
\needspace{3\baselineskip}
\textbf{wholeheartedly} — \textbf{inweardlīce} (adv.). fervently; wholeheartedly.\par
% glossary-block:B01998:end
% glossary-block:B01999:begin
\needspace{3\baselineskip}
\textbf{whom} — \textbf{hwā, hwæt} (interrogative/indefinite pron.). who? what?. Case follows role/government; distinguish personal and neuter reference. Inflect the interrogative by its role/government; dependent questions and relatives are separate constructions. \textbf{se, sēo, ðæt} (article; demonstrative; relative pron.). who; which; that. Inflected relative; case follows its clause role.\par
% glossary-block:B01999:end
% glossary-block:B02000:begin
\needspace{3\baselineskip}
\textbf{whomsoever} — Swā hwā swā means whoever, swā hwæt swā whatever, and swā hwæðer swā whichever of two. Inflect the internal pronoun for its clause role; swā alone does not carry the whole free-choice meaning. Swā hwā swā means whoever; swā hwæt swā means whatever. Inflect the internal pronoun for its clause role, so whomsoever requires the appropriate oblique form. Retain any explicitly supplied attraction.\par
% glossary-block:B02000:end
% glossary-block:B02001:begin
\needspace{3\baselineskip}
\textbf{whose} — \textbf{hwā, hwæt} (interrogative/indefinite pron.). who? what?. Case follows role/government; distinguish personal and neuter reference. Inflect the interrogative by its role/government; dependent questions and relatives are separate constructions. \textbf{se, sēo, ðæt} (article; demonstrative; relative pron.). who; which; that. Inflected relative; case follows its clause role.\par
% glossary-block:B02001:end
% glossary-block:B02002:begin
\needspace{3\baselineskip}
\textbf{whosoever} — Swā hwā swā means whoever, swā hwæt swā whatever, and swā hwæðer swā whichever of two. Inflect the internal pronoun for its clause role; swā alone does not carry the whole free-choice meaning. Swā hwā swā means whoever; swā hwæt swā means whatever. Inflect the internal pronoun for its clause role, so whomsoever requires the appropriate oblique form. Retain any explicitly supplied attraction.\par
% glossary-block:B02002:end
% glossary-block:B02003:begin
\needspace{3\baselineskip}
\textbf{why} — \textbf{for} (prep.). why; for what reason. In for hwȳ or for hwon. \textbf{hwā, hwæt} (interrogative/indefinite pron.). what reason. Instrumental hwon in the supplied why expressions. \textbf{hwȳ} (adv.). why?. Tō hwon is the whole expression why/for what purpose.\par
% glossary-block:B02003:end
% glossary-block:B02004:begin
\needspace{3\baselineskip}
\textbf{wide} — \textbf{wīd} (adjective citation).\par
% glossary-block:B02004:end
% glossary-block:B02005:begin
\needspace{3\baselineskip}
\textbf{widely} — \textbf{wīde} (adv.). widely; far and wide.\par
% glossary-block:B02005:end
% glossary-block:B02006:begin
\needspace{3\baselineskip}
\textbf{wife} — \textbf{wīf} (n.). Wīf is grammatically neuter. Choose the woman or wife sense from the context.\par
% glossary-block:B02006:end
% glossary-block:B02007:begin
\needspace{3\baselineskip}
\textbf{will} — \textbf{willa} (n.). will; disposition. Present morphology can carry future reference; bēo is used in its bounded expected-state frame. Do not add willan unless volition or a stated source target requires it.\par
% glossary-block:B02007:end
% glossary-block:B02008:begin
\needspace{3\baselineskip}
\textbf{wind} — \textbf{wind} (n.).  \textbf{windan} (v.). wind; plait.\par
% glossary-block:B02008:end
% glossary-block:B02009:begin
\needspace{3\baselineskip}
\textbf{window} — \textbf{ēagðȳrel} (n.).\par
% glossary-block:B02009:end
% glossary-block:B02010:begin
\needspace{3\baselineskip}
\textbf{wine} — \textbf{wīn} (n.).\par
% glossary-block:B02010:end
% glossary-block:B02011:begin
\needspace{3\baselineskip}
\textbf{winged} — \textbf{gefiðred} (adj.). Adjective agreeing with the fowl.\par
% glossary-block:B02011:end
% glossary-block:B02012:begin
\needspace{3\baselineskip}
\textbf{winsome} — \textbf{wynsum} (adj.). pleasant; winsome.\par
% glossary-block:B02012:end
% glossary-block:B02013:begin
\needspace{3\baselineskip}
\textbf{winter} — \textbf{winter} (n.).  \textbf{winterlic} (adj.). winter; of winter.\par
% glossary-block:B02013:end
% glossary-block:B02014:begin
\needspace{3\baselineskip}
\textbf{wisdom} — \textbf{wīsdōm} (n.).\par
% glossary-block:B02014:end
% glossary-block:B02015:begin
\needspace{3\baselineskip}
\textbf{wise} — \textbf{wīs} (adj.).\par
% glossary-block:B02015:end
% glossary-block:B02016:begin
\needspace{3\baselineskip}
\textbf{wise man} — \textbf{wīs} (adj.). wise man; wise one. Substantival direct address.\par
% glossary-block:B02016:end
% glossary-block:B02017:begin
\needspace{3\baselineskip}
\textbf{wise one} — \textbf{wīs} (adj.). wise man; wise one. Substantival direct address.\par
% glossary-block:B02017:end
% glossary-block:B02018:begin
\needspace{3\baselineskip}
\textbf{wiser} — \textbf{wīsra} (adj.).\par
% glossary-block:B02018:end
% glossary-block:B02019:begin
\needspace{3\baselineskip}
\textbf{wish} — \textbf{gewilnian} (v.). desire; wish. \textbf{willan} (v.). wish; want to.\par
% glossary-block:B02019:end
% glossary-block:B02020:begin
\needspace{3\baselineskip}
\textbf{wish to} — \textbf{willan} (v.). willan + plain infinitive denotes volition; intention does not assert the action occurred.\par
% glossary-block:B02020:end
% glossary-block:B02021:begin
\needspace{3\baselineskip}
\textbf{wish well} — \textbf{biddan} (v.). wish well; bid farewell. Supplied farewell formula.\par
% glossary-block:B02021:end
% glossary-block:B02022:begin
\needspace{3\baselineskip}
\textbf{with} — \textbf{mid} (prep.; adv. in fixed expressions). with; accompanying. Dative person or cargo in the ordinary accompaniment frame. with; by means of. Dative instrument or means. Circumstances of risk or loss in the supplied phrase. Manner in the supplied compassion phrase. with; along with. Supplied mid wes, be with them. Dative garment in the assigned clothing frame. Dative conversational partner. Accompanying quality or mental disposition in the source. Material or clothing in the supplied adornment phrase. Exact cloth/cloak complements of begyrd and bewǣfed. with; by means of. Manner/circumstance of victory. \textbf{wið} (prep.). with; against. Dative fellow players in the exact supplied playing frame. Supplied conversational counterpart; do not infer hostility.\par
% glossary-block:B02022:end
% glossary-block:B02023:begin
\needspace{3\baselineskip}
\textbf{with compassion} — Mid + the supplied dative noun of compassion forms the whole manner expression with compassion; mid alone means with.\par
% glossary-block:B02023:end
% glossary-block:B02024:begin
\needspace{3\baselineskip}
\textbf{with difficulty} — \textbf{unēaðe} (adv.). with difficulty; barely.\par
% glossary-block:B02024:end
% glossary-block:B02025:begin
\needspace{3\baselineskip}
\textbf{withholding} — \textbf{fæsthafol} (adj.). tenacious; withholding. Genitive of what is held back in the source.\par
% glossary-block:B02025:end
% glossary-block:B02026:begin
\needspace{3\baselineskip}
\textbf{within} — \textbf{binnan} (prep.).  \textbf{innan} (adv.). within; inside.\par
% glossary-block:B02026:end
% glossary-block:B02027:begin
\needspace{3\baselineskip}
\textbf{without} — \textbf{būtan} (prep.). With dative. Būtan + dative means without a companion or outside a city in the respective sense. Conditional unless is separate.\par
% glossary-block:B02027:end
% glossary-block:B02028:begin
\needspace{3\baselineskip}
\textbf{with permission} — \textbf{be} (prep.). by permission. Supplied permission expression.\par
% glossary-block:B02028:end
% glossary-block:B02029:begin
\needspace{3\baselineskip}
\textbf{with respect to} — \textbf{on} (prep.). in; with respect to. Quality or domain of description.\par
% glossary-block:B02029:end
% glossary-block:B02030:begin
\needspace{3\baselineskip}
\textbf{withstand} — \textbf{wiðstandan} (v.). withstand; resist. Dative of what is resisted.\par
% glossary-block:B02030:end
% glossary-block:B02031:begin
\needspace{3\baselineskip}
\textbf{with victory} — \textbf{mid} (prep.; adv. in fixed expressions). with. Manner/circumstance of victory.\par
% glossary-block:B02031:end
% glossary-block:B02032:begin
\needspace{3\baselineskip}
\textbf{wolf} — \textbf{wulf} (n.).\par
% glossary-block:B02032:end
% glossary-block:B02033:begin
\needspace{3\baselineskip}
\textbf{woman} — \textbf{wīf} (n.). Wīf is grammatically neuter. Choose the woman or wife sense from the context.\par
% glossary-block:B02033:end
% glossary-block:B02034:begin
\needspace{3\baselineskip}
\textbf{wonder} — \textbf{wundor} (n.).  \textbf{wundrian} (v.). Supplied source form.\par
% glossary-block:B02034:end
% glossary-block:B02035:begin
\needspace{3\baselineskip}
\textbf{woo} — \textbf{giernan} (v.). desire; request; woo. Genitive desired thing/person in the selected frames, including the maiden wooed. Genitive person wooed; keep the marriage-seeking sense distinct from a general want.\par
% glossary-block:B02035:end
% glossary-block:B02036:begin
\needspace{3\baselineskip}
\textbf{wood} — \textbf{wudu} (n.). wood; forest. Collective vegetation, not a new material sense.\par
% glossary-block:B02036:end
% glossary-block:B02037:begin
\needspace{3\baselineskip}
\textbf{word} — \textbf{word} (n.).\par
% glossary-block:B02037:end
% glossary-block:B02038:begin
\needspace{3\baselineskip}
\textbf{word by word} — \textbf{be} (prep.). In word be worde; literal correspondence.\par
% glossary-block:B02038:end
% glossary-block:B02039:begin
\needspace{3\baselineskip}
\textbf{work} — \textbf{gewinn} (n.). labour; work. Noun, distinct from labour as a verb. Noun in the expression from his labour. \textbf{weorc} (n.). weorc includes deed/action in these prompts; dǣd is a separate lexical option elsewhere. \textbf{wyrcean} (v.). work; perform. Accusative work or deed.\par
% glossary-block:B02039:end
% glossary-block:B02040:begin
\needspace{3\baselineskip}
\textbf{workman} — \textbf{wyrhta} (n.). workman; craftsman.\par
% glossary-block:B02040:end
% glossary-block:B02041:begin
\needspace{3\baselineskip}
\textbf{world} — \textbf{middangeard} (n.). world; earth. \textbf{woruld} (n.).\par
% glossary-block:B02041:end
% glossary-block:B02042:begin
\needspace{3\baselineskip}
\textbf{worldly dignity} — \textbf{woruldgeðyngð} (n.). worldly honour; worldly dignity.\par
% glossary-block:B02042:end
% glossary-block:B02043:begin
\needspace{3\baselineskip}
\textbf{worldly honour} — \textbf{woruldgeðyngð} (n.). worldly honour; worldly dignity.\par
% glossary-block:B02043:end
% glossary-block:B02044:begin
\needspace{3\baselineskip}
\textbf{worldly life} — \textbf{woruldlīf} (n.).\par
% glossary-block:B02044:end
% glossary-block:B02045:begin
\needspace{3\baselineskip}
\textbf{worldly occupation} — \textbf{woruldcræft} (n.).\par
% glossary-block:B02045:end
% glossary-block:B02046:begin
\needspace{3\baselineskip}
\textbf{worldly occupations} — \textbf{woruldbisgu} (n.). worldly occupations; cares.\par
% glossary-block:B02046:end
% glossary-block:B02047:begin
\needspace{3\baselineskip}
\textbf{worse} — \textbf{wiersa} (adj.).\par
% glossary-block:B02047:end
% glossary-block:B02048:begin
\needspace{3\baselineskip}
\textbf{worst} — \textbf{wierrest} (adj.).\par
% glossary-block:B02048:end
% glossary-block:B02049:begin
\needspace{3\baselineskip}
\textbf{worthy} — \textbf{wierðe} (adj.).\par
% glossary-block:B02049:end
% glossary-block:B02050:begin
\needspace{3\baselineskip}
\textbf{would} — Select intention, obligation, reported prospect or counterfactual context; English spelling alone does not determine tense, mood or an OE modal.\par
% glossary-block:B02050:end
% glossary-block:B02051:begin
\needspace{3\baselineskip}
\textbf{would have come} — The supplied past-subjunctive wesan + cumen construction expresses would have come in the counterfactual result. It does not require separate Old English words for would and have. \textbf{wesan} (v.). Whole supplied past-subjunctive conditional and wesan + cumen result; preserve its counterfactual anterior sense.\par
% glossary-block:B02051:end
% glossary-block:B02052:begin
\needspace{3\baselineskip}
\textbf{would rather} — \textbf{lēofra} (adj.). preferable; dearer. Dative experiencer + is lēofre + infinitive in the taught preference frame. The preference construction is dative experiencer + is lēofre + plain infinitive, is preferable to someone. It renders had/would rather without requiring past time. Wician in the supplied example is the action preferred, sojourn/live with, not a word meaning prefer.\par
% glossary-block:B02052:end
% glossary-block:B02053:begin
\needspace{3\baselineskip}
\textbf{would that} — A present-unreal wish introduced by English would that uses the taught preterite subjunctive. The lexical verb expresses the desired state; no mandatory willan translates would.\par
% glossary-block:B02053:end
% glossary-block:B02054:begin
\needspace{3\baselineskip}
\textbf{wound} — \textbf{wund} (n.). wound; sore.\par
% glossary-block:B02054:end
% glossary-block:B02055:begin
\needspace{3\baselineskip}
\textbf{wrap} — \textbf{bewǣfan} (v.). wrap; clothe. Supplied participial clothing expression.\par
% glossary-block:B02055:end
% glossary-block:B02056:begin
\needspace{3\baselineskip}
\textbf{wrapped with} — \textbf{mid} (prep.; adv. in fixed expressions). with. Exact cloth/cloak complements of begyrd and bewǣfed.\par
% glossary-block:B02056:end
% glossary-block:B02057:begin
\needspace{3\baselineskip}
\textbf{write} — \textbf{āwrītan} (v.).  \textbf{wrītan} (v.). Accusative text written.\par
% glossary-block:B02057:end
% glossary-block:B02058:begin
\needspace{3\baselineskip}
\textbf{written document} — \textbf{gewrit} (n.). letter; written document.\par
% glossary-block:B02058:end
% glossary-block:B02059:begin
\needspace{3\baselineskip}
\textbf{wrongly} — \textbf{yfele} (adv.). badly; wrongly.\par
% glossary-block:B02059:end
\pdfbookmark[1]{Y}{outline-B02060}
% glossary-block:B02060:begin
\needspace{5\baselineskip}
\subsection*{Y}
% glossary-block:B02060:end
% glossary-block:B02061:begin
\textbf{ye} — Gē is plural nominative ye/you; dative/accusative plural is ēow. Singular ðū, dual git and plural gē systems; thou/thee mark singular, ye plural. Unspecified you allows appropriate choices unless the task specifies number.\par
% glossary-block:B02061:end
% glossary-block:B02062:begin
\needspace{3\baselineskip}
\textbf{year} — \textbf{gēar} (n.).  \textbf{winter} (n.). Supplied age expression uses wintra; gēar is the ordinary year lemma elsewhere.\par
% glossary-block:B02062:end
% glossary-block:B02063:begin
\needspace{3\baselineskip}
\textbf{years of age} — \textbf{winter} (n.). Supplied age expression uses wintra; gēar is the ordinary year lemma elsewhere.\par
% glossary-block:B02063:end
% glossary-block:B02064:begin
\needspace{3\baselineskip}
\textbf{yet} — \textbf{ðā-hwæðre} (conj.). nevertheless; yet. Whole supplied connective. \textbf{ðēah} (conj.; adv.). nevertheless; yet. Independent adverb, distinct from concessive conjunction although. \textbf{gīet} (adv.). Still; continuing until now, distinct from adversative nevertheless. Retain exact nū gīet where it is supplied.\par
% glossary-block:B02064:end
% glossary-block:B02065:begin
\needspace{3\baselineskip}
\textbf{yield} — \textbf{gieldan} (v.).  \textbf{wyrcean} (v.). yield; produce. Accusative fruit in the taught/source expression. Accusative fruit; birð wæstm is a recognized lexical alternative in the key.\par
% glossary-block:B02065:end
% glossary-block:B02066:begin
\needspace{3\baselineskip}
\textbf{you} — Singular: nominative ðū, dative/accusative ðē; dual: nominative git, dative/accusative inc; plural: nominative gē, dative/accusative ēow. Dual subjects take plural verb forms. Singular ðū, dual git and plural gē systems; thou/thee mark singular, ye plural. Unspecified you allows appropriate choices unless the task specifies number.\par
% glossary-block:B02066:end
% glossary-block:B02067:begin
\needspace{3\baselineskip}
\textbf{young} — \textbf{geong} (adj.).\par
% glossary-block:B02067:end
% glossary-block:B02068:begin
\needspace{3\baselineskip}
\textbf{youngest} — Supplied weak superlative gingstan means youngest in the assigned accusative/genitive masculine singular contexts; positive geong means young.\par
% glossary-block:B02068:end
% glossary-block:B02069:begin
\needspace{3\baselineskip}
\textbf{young man} — \textbf{cniht} (n.). young man; youth.\par
% glossary-block:B02069:end
% glossary-block:B02070:begin
\needspace{3\baselineskip}
\textbf{young woman} — \textbf{fǣmne} (n.). maiden; young woman.\par
% glossary-block:B02070:end
% glossary-block:B02071:begin
\needspace{3\baselineskip}
\textbf{your} — \textbf{ēower} (possessive adjective). your; yours. Plural possessors; inflection agrees with the possessed noun. \textbf{incer} (possessive adjective). your; yours. Two possessors; agreeing possessive. \textbf{ðīn} (possessive adjective). your; yours. One possessor; agrees with the possessed noun. ðīn identifies one possessor and agrees with the possessed noun. ēower and incer mark plural and dual possessors separately. \textbf{ðīn} (possessive adjective), \textbf{incer} (possessive adjective), \textbf{ēower} (possessive adjective). Ðīn is for one second-person possessor, incer for two, ēower for plural possessors. The possessive agrees with the possessed noun.\par
% glossary-block:B02071:end
% glossary-block:B02072:begin
\needspace{3\baselineskip}
\textbf{yours} — \textbf{ēower} (possessive adjective). your; yours. Plural possessors; inflection agrees with the possessed noun. \textbf{incer} (possessive adjective). your; yours. Two possessors; agreeing possessive. \textbf{ðīn} (possessive adjective). your; yours. One possessor; agrees with the possessed noun. ðīn identifies one possessor and agrees with the possessed noun. ēower and incer mark plural and dual possessors separately. \textbf{ðīn} (possessive adjective), \textbf{incer} (possessive adjective), \textbf{ēower} (possessive adjective). Ðīn is for one second-person possessor, incer for two, ēower for plural possessors. The possessive agrees with the possessed noun.\par
% glossary-block:B02072:end
% glossary-block:B02073:begin
\needspace{3\baselineskip}
\textbf{yourself} — An ordinary reflexive uses the personal pronoun in the case governed by its verb or preposition. Agreeing self adds the taught emphasis; English -self does not automatically require a separate self. Ordinary reflexive personal pronouns take the governed case; agreeing self adds the taught emphasis. English -self does not automatically require self.\par
% glossary-block:B02073:end
% glossary-block:B02074:begin
\needspace{3\baselineskip}
\textbf{yourselves} — An ordinary reflexive uses the personal pronoun in the case governed by its verb or preposition. Agreeing self adds the taught emphasis; English -self does not automatically require a separate self. Ordinary reflexive personal pronouns take the governed case; agreeing self adds the taught emphasis. English -self does not automatically require self.\par
% glossary-block:B02074:end
% glossary-block:B02075:begin
\needspace{3\baselineskip}
\textbf{youth} — \textbf{cniht} (n.). young man; youth. \textbf{geoguð} (n.).\par
% glossary-block:B02075:end
% glossary-block:B02076:begin
\needspace{3\baselineskip}
\textbf{you two} — Git is nominative dual you two; inc is dative/accusative dual. A dual subject takes a plural verb form.\par
% glossary-block:B02076:end
\pdfbookmark[1]{Æ}{outline-B02077}
% glossary-block:B02077:begin
\needspace{5\baselineskip}
\subsection*{Æ}
% glossary-block:B02077:end
% glossary-block:B02078:begin
\textbf{Æthelthryth} — \textbf{Æðelðrȳð} (proper name). Etheldred; Æthelthryth. Proper name.\par
% glossary-block:B02078:end
% glossary-block:B02079:begin
\needspace{3\baselineskip}
\textbf{Æthelwald} — \textbf{Æðelwald} (proper name). Proper name in apposition to father.\par
% glossary-block:B02079:end
% glossary-block:B02080:begin
\needspace{3\baselineskip}
\textbf{Æthelwulf} — \textbf{Æthelwulf} (proper name).\par
% glossary-block:B02080:end
\end{multicols}
\end{document}
